%versi 3 (18-12-2016)
\chapter{Kode dan Model Machine Learning}
\label{lamp:A}

\begin{lstlisting}[language=Python, caption={Kode 1}]
# ============================================================
# versi rentals.csv + rental_logs.csv
# ============================================================
import pandas as pd
import numpy as np
import joblib
import seaborn as sns
import matplotlib.pyplot as plt

from sklearn.model_selection import train_test_split, GridSearchCV
from sklearn.experimental import enable_halving_search_cv
from sklearn.model_selection import HalvingGridSearchCV
from sklearn.preprocessing import LabelEncoder, StandardScaler, MinMaxScaler
from sklearn.metrics import classification_report, confusion_matrix
from sklearn.feature_selection import chi2, mutual_info_classif
from scipy.stats import pearsonr

# Models & Pipelines
from sklearn.pipeline import Pipeline, make_pipeline
from imblearn.over_sampling import SMOTE
from sklearn.tree import DecisionTreeClassifier
from sklearn.ensemble import RandomForestClassifier
from sklearn.neighbors import KNeighborsClassifier
from sklearn.svm import SVC
from xgboost import XGBClassifier
from sklearn.utils.class_weight import compute_class_weight

# ------------------------------------------------------------
# 1. LOAD DATA
# ------------------------------------------------------------
rentals = pd.read_csv('rentals.csv')
logs = pd.read_csv('rental_logs.csv')
# Standarisasi nama kolom ke huruf kecil
rentals.columns = rentals.columns.str.lower()
logs.columns = logs.columns.str.lower()

if 'komentar' not in rentals.columns:
    raise ValueError("Kolom 'komentar' tidak ditemukan di rentals.csv")

print("rentals.csv")
print(rentals.head())
print("\nlogs.csv")
print(logs.head())

# ------------------------------------------------------------
# 2. LABELING & LOG MERGING
# ------------------------------------------------------------
# Bersihkan format teks (ubah ke huruf kecil dan hapus spasi berlebih)
rentals['komentar'] = rentals['komentar'].astype(str).str.lower().str.strip()

# Definisikan nilai-nilai yang dianggap "kosong" / tidak ada komentar
empty_values = ["", "nan", "na", "null", "-", "0", "tidak ada", "none"]

# Jika isi komentar TIDAK ada di dalam empty_values, berarti ada indikasi pelanggaran (label 1)
rentals['label'] = (~rentals['komentar'].isin(empty_values)).astype(int)

# Menggabungkan logs ke rentals (membuat fitur: jumlah edit & editor unik)
logs['rental_id'] = logs['rental_id'].astype(int)
edit_count = logs.groupby('rental_id').size().reset_index(name='edit_count')
unique_editors = logs.groupby('rental_id')['email'].nunique().reset_index(name='unique_editors')

df = rentals.merge(edit_count, left_on='id', right_on='rental_id', how='left')
df = df.merge(unique_editors, left_on='id', right_on='rental_id', how='left')

# Isi nilai kosong dengan 0 untuk penyewaan yang tidak memiliki log
df['edit_count'] = df['edit_count'].fillna(0)
df['unique_editors'] = df['unique_editors'].fillna(0)

# ------------------------------------------------------------
# 3. FEATURE ENGINEERING & PREPARATION
# ------------------------------------------------------------
# Mengisi nilai kosong pada fitur kategorikal dan numerik
df['metode_pembayaran'] = df['metode_pembayaran'].fillna('Others')
if 'metode_lain' in df.columns:
    df['metode_lain'] = df['metode_lain'].fillna('')

df['lama_menginap'] = df['lama_menginap'].fillna(0)

cat_cols = ['tower', 'status_kewarganegaraan', 'metode_pembayaran', 'jenis_sewa']
numeric_cols = ['lama_menginap', 'edit_count', 'unique_editors']

# Daftar seluruh fitur relevan yang akan digunakan
feature_cols = cat_cols + numeric_cols

# Pastikan data kategorikal berupa string
for c in cat_cols:
    df[c] = df[c].astype(str)

X = df[feature_cols].copy()
y = df['label']

# ------------------------------------------------------------
# 4. TRAIN-TEST SPLIT
# ------------------------------------------------------------
X_train, X_test, y_train, y_test = train_test_split(
    X, y, stratify=y, test_size=0.2, random_state=42
)

# Menghitung bobot kelas (class weights) karena data tidak seimbang
classes = np.array([0, 1])
weights = compute_class_weight(class_weight='balanced', classes=classes, y=y_train)
class_weights = {0: weights[0], 1: weights[1]}
print(f"\nBobot Kelas (Class weights): {class_weights}")

# ------------------------------------------------------------
# 4B. ENCODING MENCEGAH DATA LEAKAGE
# ------------------------------------------------------------
label_encoders = {}

for c in cat_cols:
    le = LabelEncoder()
    # Fit & Transform pada data latih
    X_train[c] = le.fit_transform(X_train[c])

    # Tangani unseen labels pada data uji (cegah error saat deployment)
    X_test[c] = X_test[c].map(lambda s: s if s in le.classes_ else 'Unknown')
    if 'Unknown' not in le.classes_:
        le.classes_ = np.append(le.classes_, 'Unknown')

    X_test[c] = le.transform(X_test[c])
    label_encoders[c] = le

# ------------------------------------------------------------
# 5. FEATURE EVALUATION
# ------------------------------------------------------------
print("\n=== Evaluasi Korelasi Fitur ===")
# Catatan: Pemodelan menggunakan seluruh kolom relevan yang tersedia dengan
# spesifik encoding (tidak otomatis menyeleksi sebagian/membuang fitur).
# Pengujian ini sekadar untuk analisis signifikansi.

chi_vals, _ = chi2(X_train, y_train)
mi_vals = mutual_info_classif(X_train, y_train)
pearson_vals = [abs(pearsonr(X_train[col], y_train)[0]) for col in X_train.columns]

fs_df = pd.DataFrame({
    'Fitur': X_train.columns,
    'Chi2': chi_vals,
    'MutualInfo': mi_vals,
    'Pearson': pearson_vals
}).sort_values(by='MutualInfo', ascending=False)

print(fs_df)

# ------------------------------------------------------------
# 6. Fungsi Evaluasi
# ------------------------------------------------------------
def evaluate(name, model):
    y_pred = model.predict(X_test)
    print(f"\n===== {name} =====")
    print("Classification Report:")
    print(classification_report(y_test, y_pred, digits=4))

    cm = confusion_matrix(y_test, y_pred)

    # --------------------------------------------------------
    # Confusion Matrix (Text Base)
    # --------------------------------------------------------
    print("Confusion Matrix (Text):")
    print("                Predicted")
    print("               0       1")
    print("          -----------------")
    print(f"True   0 | {cm[0][0]:6d}  {cm[0][1]:6d}")
    print(f"       1 | {cm[1][0]:6d}  {cm[1][1]:6d}")
    print("          -----------------\n")

    # --------------------------------------------------------
    # Confusion Matrix (Gambar / Visualisasi)
    # --------------------------------------------------------
    # plt.figure(figsize=(5, 4))
    # sns.heatmap(cm, annot=True, fmt='d', cmap='Blues', cbar=False,
    #             xticklabels=['0 (Aman)', '1 (Melanggar)'],
    #             yticklabels=['0 (Aman)', '1 (Melanggar)'])
    #
    # plt.title(f'Confusion Matrix - {name}')
    # plt.xlabel('Predicted Label')
    # plt.ylabel('True Label')
    # plt.tight_layout()
    # plt.show()

# ------------------------------------------------------------
# 7. LATIH SEMUA MODEL & HYPERPARAMETER TUNING
# ------------------------------------------------------------
models = {}
print("\nMemulai proses Hyperparameter Tuning untuk semua model. Ini mungkin memakan waktu...\n")

# Penanganan SMOTE khusus untuk KNN
sm = SMOTE(random_state=42)
X_train_res, y_train_res = sm.fit_resample(X_train, y_train)

# --- 1. Decision Tree ---
print(">> Tuning Decision Tree...")
param_dt = {
    'max_depth': [5, 10, None],
    'min_samples_split': [2, 5],
    'min_samples_leaf': [1, 2],
    'class_weight': ['balanced']
}
grid_dt = GridSearchCV(DecisionTreeClassifier(random_state=42), param_grid=param_dt, scoring='recall', cv=3, n_jobs=-1, verbose=0)
grid_dt.fit(X_train, y_train)
models["Decision Tree (Tuned)"] = grid_dt.best_estimator_
print("Parameter Terbaik (DT):", grid_dt.best_params_)

# --- 2. Random Forest ---
print("\n>> Tuning Random Forest...")
param_rf = {
    'n_estimators': [100, 200],
    'max_depth': [5, 10, None],
    'min_samples_split': [2, 5],
    'min_samples_leaf': [1, 2],
    'class_weight': ['balanced']
}
grid_rf = GridSearchCV(RandomForestClassifier(random_state=42), param_grid=param_rf, scoring='recall', cv=3, n_jobs=-1, verbose=0)
grid_rf.fit(X_train, y_train)
models["Random Forest (Tuned)"] = grid_rf.best_estimator_
print("Parameter Terbaik (RF):", grid_rf.best_params_)

# --- 3. KNN (dengan SMOTE) ---
print("\n>> Tuning KNN (SMOTE)...")
param_knn = {
    'n_neighbors': [3, 5, 7],
    'weights': ['uniform', 'distance']
}
grid_knn = GridSearchCV(KNeighborsClassifier(), param_grid=param_knn, scoring='recall', cv=3, n_jobs=-1, verbose=0)
grid_knn.fit(X_train_res, y_train_res) # Latih dengan data yang sudah di-resample
models["KNN (SMOTE Tuned)"] = grid_knn.best_estimator_
print("Parameter Terbaik (KNN):", grid_knn.best_params_)

# --- 4. SVM (HalvingGridSearchCV) ---
print("\n>> Tuning SVM Optimasi...")
pipe_svm = Pipeline([
    ('scaler', StandardScaler()),
    ('svm', SVC(random_state=42, class_weight='balanced'))
])
param_svm = {
    'svm__C': [0.1, 1, 10],
    'svm__gamma': ['scale', 0.01, 0.1],
    'svm__kernel': ['rbf']
}
halving_svm = HalvingGridSearchCV(pipe_svm, param_grid=param_svm, scoring='recall', cv=2, factor=2, n_jobs=-1, verbose=0)
halving_svm.fit(X_train, y_train)
models["SVM Optimasi (Tuned)"] = halving_svm.best_estimator_
print("Parameter Terbaik (SVM):", halving_svm.best_params_)

# --- 5. XGBoost ---
print("\n>> Tuning XGBoost...")
pos_weight = (y_train == 0).sum() / (y_train == 1).sum()
param_xgb = {
    'n_estimators': [100, 200],
    'max_depth': [3, 6],
    'learning_rate': [0.01, 0.1, 0.3],
    'subsample': [0.8, 1.0],
    'scale_pos_weight': [pos_weight]
}
grid_xgb = GridSearchCV(XGBClassifier(eval_metric='logloss'), param_grid=param_xgb, scoring='recall', cv=3, n_jobs=-1, verbose=0)
grid_xgb.fit(X_train, y_train)
models["XGBoost (Tuned)"] = grid_xgb.best_estimator_
print("Parameter Terbaik (XGB):", grid_xgb.best_params_)

# --- Evaluasi Semua Model ---
for name, model in models.items():
    evaluate(name, model)
\end{lstlisting}

\begin{lstlisting}[caption={Output 1}]
rentals.csv
   id           nama               nik status_pasutri tower  lantai    unit  \
0   1   V*** P******   2**************             ya     C       1  C-0111   
1   2  I**** S******  7***************  belum menikah     C      15  C-1524   
2   3  I**** H******  2***************             ya     B       8  B-0836   
3   4   W*** P******  5***************             ya     C      17  C-1712   
4   5   R**** W*****  5***************  belum menikah     B       9  B-0934   

  status_kewarganegaraan jenis_sewa metode_pembayaran metode_lain  \
0                    WNI   mingguan       Kartu Debit         NaN   
1                    WNI     harian       Kartu Debit         NaN   
2                    WNA   mingguan              Cash     Voucher   
3                    WNI   mingguan              Cash         NaN   
4                    WNI    bulanan              QRIS         NaN   

  tanggal_checkin waktu_checkin tanggal_checkout waktu_checkout  \
0      2024-07-23         06:48       2024-10-11          10:06   
1      2024-03-05         08:44       2025-02-12          21:41   
2      2025-03-28         07:08       2025-10-02          04:51   
3      2024-09-03         16:38       2025-07-28          08:57   
4      2024-03-09         23:44       2024-12-11          06:45   

   lama_menginap                                komentar        email_agent  \
0             80    Ditemukan kondom dalam jumlah banyak  a******@srusun.id   
1            344          Kerusakan parah pada fasilitas  a******@srusun.id   
2            188  Ditemukan alat suntik di tempat sampah    s****@gmail.com   
3            328                                     NaN  a******@srusun.id   
4            277                                     NaN  a******@srusun.id   

         diedit_oleh           created_at  
0  a******@srusun.id  2024-11-10 05:34:02  
1  a******@srusun.id  2024-07-23 08:46:51  
2    s****@gmail.com  2025-05-22 03:08:25  
3  a******@srusun.id  2025-06-25 09:19:29  
4  a******@srusun.id  2024-09-19 00:09:23  

logs.csv
   id  rental_id  action field_changed old_value  \
0   1       8866  insert           all       NaN   
1   2       5194  insert           all       NaN   
2   3       6383  insert           all       NaN   
3   4       2573  insert           all       NaN   
4   5       6042  insert           all       NaN   

                                           new_value              email  \
0  {"nama": "T**** P******", "nik": "5*****************************************"nama": "R**** N******", "nik": "1*****************************************"nama": "R**** H***", "nik": "9********************************************"nama": "N**** P******", "nik": "4*****************************************"nama": "E** H****", "nik": "5***************...  a******@srusun.id   

            time_stamp  
0  2023-01-05 23:30:37  
1  2023-01-06 19:55:07  
2  2023-01-07 12:38:50  
3  2023-01-08 19:27:08  
4  2023-01-10 05:05:24  

Bobot Kelas (Class weights): {0: np.float64(0.5***************), 1: np.float64(3.9**************)}

=== Evaluasi Korelasi Fitur ===
                    Fitur       Chi2  MutualInfo   Pearson
2       metode_pembayaran   0.000006    0.004583  0.000028
6          unique_editors   7.911098    0.004272  0.069049
1  status_kewarganegaraan   0.043046    0.003034  0.005025
0                   tower   0.003847    0.001495  0.000754
5              edit_count  14.396443    0.000033  0.058444
4           lama_menginap   0.785871    0.000000  0.001272
3              jenis_sewa   0.085542    0.000000  0.004456

Memulai proses Hyperparameter Tuning untuk semua model. Ini mungkin memakan waktu...

>> Tuning Decision Tree...
Parameter Terbaik (DT): {'class_weight': 'balanced', 'max_depth': 5, 'min_samples_leaf': 1, 'min_samples_split': 2}

>> Tuning Random Forest...
Parameter Terbaik (RF): {'class_weight': 'balanced', 'max_depth': 5, 'min_samples_leaf': 1, 'min_samples_split': 5, 'n_estimators': 200}

>> Tuning KNN (SMOTE)...
Parameter Terbaik (KNN): {'n_neighbors': 7, 'weights': 'distance'}

>> Tuning SVM Optimasi...
Parameter Terbaik (SVM): {'svm__C': 0.1, 'svm__gamma': 'scale', 'svm__kernel': 'rbf'}

>> Tuning XGBoost...
Parameter Terbaik (XGB): {'learning_rate': 0.01, 'max_depth': 3, 'n_estimators': 200, 'scale_pos_weight': np.float64(6.9**************), 'subsample': 1.0}

===== Decision Tree (Tuned) =====
Classification Report:
              precision    recall  f1-score   support

           0     0.8841    0.4231    0.5723      1749
           1     0.1324    0.6135    0.2178       251

    accuracy                         0.4470      2000
   macro avg     0.5083    0.5183    0.3951      2000
weighted avg     0.7898    0.4470    0.5278      2000

Confusion Matrix (Text):
                Predicted
               0       1
          -----------------
True   0 |    740    1009
       1 |     97     154
          -----------------


===== Random Forest (Tuned) =====
Classification Report:
              precision    recall  f1-score   support

           0     0.8830    0.6861    0.7722      1749
           1     0.1435    0.3665    0.2063       251

    accuracy                         0.6460      2000
   macro avg     0.5133    0.5263    0.4892      2000
weighted avg     0.7902    0.6460    0.7012      2000

Confusion Matrix (Text):
                Predicted
               0       1
          -----------------
True   0 |   1200     549
       1 |    159      92
          -----------------


===== KNN (SMOTE Tuned) =====
Classification Report:
              precision    recall  f1-score   support

           0     0.8719    0.6575    0.7497      1749
           1     0.1204    0.3267    0.1760       251

    accuracy                         0.6160      2000
   macro avg     0.4961    0.4921    0.4628      2000
weighted avg     0.7776    0.6160    0.6777      2000

Confusion Matrix (Text):
                Predicted
               0       1
          -----------------
True   0 |   1150     599
       1 |    169      82
          -----------------


===== SVM Optimasi (Tuned) =====
Classification Report:
              precision    recall  f1-score   support

           0     0.8762    0.6352    0.7365      1749
           1     0.1284    0.3745    0.1913       251

    accuracy                         0.6025      2000
   macro avg     0.5023    0.5049    0.4639      2000
weighted avg     0.7823    0.6025    0.6681      2000

Confusion Matrix (Text):
                Predicted
               0       1
          -----------------
True   0 |   1111     638
       1 |    157      94
          -----------------


===== XGBoost (Tuned) =====
Classification Report:
              precision    recall  f1-score   support

           0     0.8872    0.5843    0.7046      1749
           1     0.1427    0.4821    0.2202       251

    accuracy                         0.5715      2000
   macro avg     0.5149    0.5332    0.4624      2000
weighted avg     0.7937    0.5715    0.6438      2000

Confusion Matrix (Text):
                Predicted
               0       1
          -----------------
True   0 |   1022     727
       1 |    130     121
          -----------------
\end{lstlisting}

\begin{lstlisting}[language=Python, caption={Kode 2}]
import pandas as pd
import numpy as np
import joblib
import seaborn as sns
import matplotlib.pyplot as plt

from sklearn.model_selection import train_test_split, GridSearchCV
from sklearn.experimental import enable_halving_search_cv  # Dibutuhkan untuk HalvingGridSearchCV
from sklearn.model_selection import HalvingGridSearchCV
from sklearn.preprocessing import LabelEncoder, StandardScaler
from sklearn.metrics import classification_report, confusion_matrix
from sklearn.feature_selection import chi2, mutual_info_classif
from scipy.stats import pearsonr

# Models & Pipelines
from sklearn.pipeline import Pipeline, make_pipeline
from imblearn.over_sampling import SMOTE
from sklearn.tree import DecisionTreeClassifier
from sklearn.ensemble import RandomForestClassifier
from sklearn.neighbors import KNeighborsClassifier
from sklearn.svm import SVC
from xgboost import XGBClassifier
from sklearn.utils.class_weight import compute_class_weight

# ------------------------------------------------------------
# 1. LOAD DATA
# ------------------------------------------------------------
rentals = pd.read_csv('rentals.csv')
logs = pd.read_csv('rental_logs.csv')
agents = pd.read_csv('agents.csv')
units = pd.read_csv('units.csv')
users = pd.read_csv('users.csv')
violations = pd.read_csv('violations.csv')

# Standarisasi nama kolom ke huruf kecil
rentals.columns = rentals.columns.str.lower()
logs.columns = logs.columns.str.lower()

if 'komentar' not in rentals.columns:
    raise ValueError("Kolom 'komentar' tidak ditemukan di rentals.csv")

print("rentals.csv")
print(rentals.head())
print("\nlogs.csv")
print(logs.head())

# ------------------------------------------------------------
# 2. LABELING & LOG MERGING
# ------------------------------------------------------------
# Merge logs into rentals (feature: edit_count & unique editors)
logs['rental_id'] = logs['rental_id'].astype(int)
edit_count = logs.groupby('rental_id').size().reset_index(name='edit_count')
unique_editors = logs.groupby('rental_id')['email'].nunique().reset_index(name='unique_editors')

df = rentals.merge(edit_count, left_on='id', right_on='rental_id', how='left')
df = df.merge(unique_editors, left_on='id', right_on='rental_id', how='left')

df['edit_count'] = df['edit_count'].fillna(0)
df['unique_editors'] = df['unique_editors'].fillna(0)

# Labeling Target: has_violation -> label
violation_keywords = [
    'alat suntik', 'kondom', 'kerusakan', 'kebisingan', 'alkohol', 'narkoba',
    'kekerasan', 'ancaman', 'merokok', 'kebersihan', 'tamu', 'menyewakan',
    'transit', 'lalai'
]

df['komentar'] = df['komentar'].astype(str).str.lower()
df['label'] = df['komentar'].apply(lambda x: int(any(k in x for k in violation_keywords)))

# ------------------------------------------------------------
# 3. FEATURE ENGINEERING & PREPARATION
# ------------------------------------------------------------
# Mengisi nilai kosong pada fitur kategorikal dan numerik
df['metode_pembayaran'] = df['metode_pembayaran'].fillna('Others')
if 'metode_lain' in df.columns:
    df['metode_lain'] = df['metode_lain'].fillna('')

df['lama_menginap'] = df['lama_menginap'].fillna(0)

cat_cols = ['tower', 'status_kewarganegaraan', 'metode_pembayaran', 'jenis_sewa']
numeric_cols = ['lantai', 'lama_menginap', 'edit_count', 'unique_editors']

# Daftar seluruh fitur relevan yang akan digunakan
feature_cols = cat_cols + numeric_cols

# Pastikan data kategorikal berupa string
for c in cat_cols:
    df[c] = df[c].astype(str)

X = df[feature_cols].copy()
y = df['label']

# ------------------------------------------------------------
# 4. TRAIN-TEST SPLIT
# ------------------------------------------------------------
X_train, X_test, y_train, y_test = train_test_split(
    X, y, stratify=y, test_size=0.2, random_state=42
)

# Menghitung bobot kelas (class weights) karena data tidak seimbang
classes = np.array([0, 1])
weights = compute_class_weight(class_weight='balanced', classes=classes, y=y_train)
class_weights = {0: weights[0], 1: weights[1]}
print(f"\nBobot Kelas (Class weights): {class_weights}")

# ------------------------------------------------------------
# 4B. ENCODING MENCEGAH DATA LEAKAGE
# ------------------------------------------------------------
label_encoders = {}

for c in cat_cols:
    le = LabelEncoder()
    # Fit & Transform pada data latih
    X_train[c] = le.fit_transform(X_train[c])

    # Tangani unseen labels pada data uji (cegah error saat deployment)
    X_test[c] = X_test[c].map(lambda s: s if s in le.classes_ else 'Unknown')
    if 'Unknown' not in le.classes_:
        le.classes_ = np.append(le.classes_, 'Unknown')

    X_test[c] = le.transform(X_test[c])
    label_encoders[c] = le

# ------------------------------------------------------------
# 5. FEATURE EVALUATION
# ------------------------------------------------------------
print("\n=== Evaluasi Korelasi Fitur ===")
# Catatan: Pemodelan menggunakan seluruh kolom relevan yang tersedia dengan
# spesifik encoding (tidak otomatis menyeleksi sebagian/membuang fitur).
# Pengujian ini sekadar untuk analisis signifikansi.

chi_vals, _ = chi2(X_train, y_train)
mi_vals = mutual_info_classif(X_train, y_train)
pearson_vals = [abs(pearsonr(X_train[col], y_train)[0]) for col in X_train.columns]

fs_df = pd.DataFrame({
    'Fitur': X_train.columns,
    'Chi2': chi_vals,
    'MutualInfo': mi_vals,
    'Pearson': pearson_vals
}).sort_values(by='MutualInfo', ascending=False)

print(fs_df)

# ------------------------------------------------------------
# 6. Fungsi Evaluasi
# ------------------------------------------------------------
def evaluate(name, model):
    y_pred = model.predict(X_test)
    print(f"\n===== {name} =====")
    print("Classification Report:")
    print(classification_report(y_test, y_pred, digits=4))

    cm = confusion_matrix(y_test, y_pred)

    # --------------------------------------------------------
    # Confusion Matrix (Text Base)
    # --------------------------------------------------------
    print("Confusion Matrix (Text):")
    print("                Predicted")
    print("               0       1")
    print("          -----------------")
    print(f"True   0 | {cm[0][0]:6d}  {cm[0][1]:6d}")
    print(f"       1 | {cm[1][0]:6d}  {cm[1][1]:6d}")
    print("          -----------------\n")

    # --------------------------------------------------------
    # Confusion Matrix (Gambar / Visualisasi)
    # --------------------------------------------------------
    # plt.figure(figsize=(5, 4))
    # sns.heatmap(cm, annot=True, fmt='d', cmap='Blues', cbar=False,
    #             xticklabels=['0 (Aman)', '1 (Melanggar)'],
    #             yticklabels=['0 (Aman)', '1 (Melanggar)'])
    #
    # plt.title(f'Confusion Matrix - {name}')
    # plt.xlabel('Predicted Label')
    # plt.ylabel('True Label')
    # plt.tight_layout()
    # plt.show()

# ------------------------------------------------------------
# 7. LATIH SEMUA MODEL & HYPERPARAMETER TUNING
# ------------------------------------------------------------
models = {}
print("\nMemulai proses Hyperparameter Tuning untuk semua model. Ini mungkin memakan waktu...\n")

# Penanganan SMOTE khusus untuk KNN
sm = SMOTE(random_state=42)
X_train_res, y_train_res = sm.fit_resample(X_train, y_train)

# --- 1. Decision Tree ---
print(">> Tuning Decision Tree...")
param_dt = {
    'max_depth': [5, 10, None],
    'min_samples_split': [2, 5],
    'min_samples_leaf': [1, 2],
    'class_weight': ['balanced']
}
grid_dt = GridSearchCV(DecisionTreeClassifier(random_state=42), param_grid=param_dt, scoring='recall', cv=3, n_jobs=-1, verbose=0)
grid_dt.fit(X_train, y_train)
models["Decision Tree (Tuned)"] = grid_dt.best_estimator_
print("Parameter Terbaik (DT):", grid_dt.best_params_)

# --- 2. Random Forest ---
print("\n>> Tuning Random Forest...")
param_rf = {
    'n_estimators': [100, 200],
    'max_depth': [5, 10, None],
    'min_samples_split': [2, 5],
    'min_samples_leaf': [1, 2],
    'class_weight': ['balanced']
}
grid_rf = GridSearchCV(RandomForestClassifier(random_state=42), param_grid=param_rf, scoring='recall', cv=3, n_jobs=-1, verbose=0)
grid_rf.fit(X_train, y_train)
models["Random Forest (Tuned)"] = grid_rf.best_estimator_
print("Parameter Terbaik (RF):", grid_rf.best_params_)

# --- 3. KNN (dengan SMOTE) ---
print("\n>> Tuning KNN (SMOTE)...")
param_knn = {
    'n_neighbors': [3, 5, 7],
    'weights': ['uniform', 'distance']
}
grid_knn = GridSearchCV(KNeighborsClassifier(), param_grid=param_knn, scoring='recall', cv=3, n_jobs=-1, verbose=0)
grid_knn.fit(X_train_res, y_train_res)  # Latih dengan data yang sudah di-resample
models["KNN (SMOTE Tuned)"] = grid_knn.best_estimator_
print("Parameter Terbaik (KNN):", grid_knn.best_params_)

# --- 4. SVM (HalvingGridSearchCV) ---
print("\n>> Tuning SVM Optimasi...")
pipe_svm = Pipeline([
    ('scaler', StandardScaler()),
    ('svm', SVC(random_state=42, class_weight='balanced'))
])
param_svm = {
    'svm__C': [0.1, 1, 10],
    'svm__gamma': ['scale', 0.01, 0.1],
    'svm__kernel': ['rbf']
}
halving_svm = HalvingGridSearchCV(pipe_svm, param_grid=param_svm, scoring='recall', cv=2, factor=2, n_jobs=-1, verbose=0)
halving_svm.fit(X_train, y_train)
models["SVM Optimasi (Tuned)"] = halving_svm.best_estimator_
print("Parameter Terbaik (SVM):", halving_svm.best_params_)

# --- 5. XGBoost ---
print("\n>> Tuning XGBoost...")
pos_weight = (y_train == 0).sum() / (y_train == 1).sum()
param_xgb = {
    'n_estimators': [100, 200],
    'max_depth': [3, 6],
    'learning_rate': [0.01, 0.1, 0.3],
    'subsample': [0.8, 1.0],
    'scale_pos_weight': [pos_weight]
}
grid_xgb = GridSearchCV(XGBClassifier(use_label_encoder=False, eval_metric='logloss'), param_grid=param_xgb, scoring='recall', cv=3, n_jobs=-1, verbose=0)
grid_xgb.fit(X_train, y_train)
models["XGBoost (Tuned)"] = grid_xgb.best_estimator_
print("Parameter Terbaik (XGB):", grid_xgb.best_params_)

# --- Evaluasi Semua Model ---
for name, model in models.items():
    evaluate(name, model)
\end{lstlisting}

\begin{lstlisting}[caption={Output 2}]
rentals.csv
   id            nama               nik status_pasutri tower  lantai  unit  \
0   1  N*** P****** 1  3***************  Belum Menikah     A       1    13   
1   2  N*** P****** 2  3***************        Menikah     D       3     4   
2   3  N*** P****** 3  3***************  Belum Menikah     B      18     2   
3   4  N*** P****** 4  3***************  Belum Menikah     B      11     7   
4   5  N*** P****** 5  3***************        Menikah     A      14     6   

  status_kewarganegaraan jenis_sewa metode_pembayaran  metode_lain  \
0                    WNI    Bulanan              Cash          NaN   
1                    WNA    Bulanan              Cash          NaN   
2                    WNA   Mingguan              QRIS          NaN   
3                    WNI     Harian              QRIS          NaN   
4                    WNA    Bulanan          Transfer          NaN   

  tanggal_checkin waktu_checkin tanggal_checkout waktu_checkout  \
0      2024-01-22      18:00:00       2024-01-31       06:00:00   
1      2023-06-16      18:00:00       2023-06-29       10:00:00   
2      2023-08-01      15:00:00       2023-08-02       09:00:00   
3      2023-03-30      10:00:00       2023-04-05       12:00:00   
4      2024-10-29      11:00:00       2024-11-02       08:00:00   

   lama_menginap                             komentar       email_agent  \
0              9                                  NaN  a*****@gmail.com   
1             13  Kebisingan berlebihan di malam hari  a*****@gmail.com   
2              1                                  NaN  a*****@gmail.com   
3              6                                  NaN  a*****@gmail.com   
4              4                                  NaN   a****@gmail.com   

        diedit_oleh           created_at  
0  a*****@gmail.com  2025-11-17 09:13:08  
1  a*****@gmail.com  2025-11-17 09:13:08  
2  a*****@gmail.com  2025-11-17 09:13:08  
3  a*****@gmail.com  2025-11-17 09:13:08  
4   a****@gmail.com  2025-11-17 09:13:08  

logs.csv
   id  rental_id  action  field_changed  old_value  new_value  \
0   1          1  CREATE            NaN        NaN        NaN   
1   2          2  CREATE            NaN        NaN        NaN   
2   3          3  CREATE            NaN        NaN        NaN   
3   4          4  CREATE            NaN        NaN        NaN   
4   5          5  CREATE            NaN        NaN        NaN   

              email           time_stamp  
0  a*****@gmail.com  2025-11-17 09:13:08  
1  a*****@gmail.com  2025-11-17 09:13:08  
2  a*****@gmail.com  2025-11-17 09:13:08  
3  a*****@gmail.com  2025-11-17 09:13:08  
4   a****@gmail.com  2025-11-17 09:13:08  

Bobot Kelas (Class weights): {0: np.float64(0.5***************), 1: np.float64(5.0**************)}

=== Evaluasi Korelasi Fitur ===
                    Fitur      Chi2  MutualInfo   Pearson
3              jenis_sewa  0.015141    0.004646  0.001676
2       metode_pembayaran  0.432687    0.004396  0.008965
7          unique_editors  0.000000    0.000856       NaN
6              edit_count  0.000000    0.000121       NaN
0                   tower  2.706459    0.000000  0.020519
1  status_kewarganegaraan  0.009222    0.000000  0.001511
5           lama_menginap  1.717225    0.000000  0.009953
4                  lantai  3.590630    0.000000  0.011865

Memulai proses Hyperparameter Tuning untuk semua model. Ini mungkin memakan waktu...

>> Tuning Decision Tree...

/tmp/ipykernel_5100/1987556958.py:135: ConstantInputWarning: An input array is constant; the correlation coefficient is not defined.
  pearson_vals = [abs(pearsonr(X_train[col], y_train)[0]) for col in X_train.columns]

Parameter Terbaik (DT): {'class_weight': 'balanced', 'max_depth': 5, 'min_samples_leaf': 1, 'min_samples_split': 2}

>> Tuning Random Forest...
Parameter Terbaik (RF): {'class_weight': 'balanced', 'max_depth': 5, 'min_samples_leaf': 1, 'min_samples_split': 2, 'n_estimators': 100}

>> Tuning KNN (SMOTE)...
Parameter Terbaik (KNN): {'n_neighbors': 7, 'weights': 'distance'}

>> Tuning SVM Optimasi...
Parameter Terbaik (SVM): {'svm__C': 1, 'svm__gamma': 0.01, 'svm__kernel': 'rbf'}

>> Tuning XGBoost...

/usr/local/lib/python3.12/dist-packages/xgboost/training.py:200: UserWarning: [06:42:40] WARNING: /__w/xgboost/xgboost/src/learner.cc:782: 
Parameters: { "use_label_encoder" } are not used.

  bst.update(dtrain, iteration=i, fobj=obj)

Parameter Terbaik (XGB): {'learning_rate': 0.01, 'max_depth': 3, 'n_estimators': 100, 'scale_pos_weight': np.float64(9.0**************), 'subsample': 1.0}

===== Decision Tree (Tuned) =====
Classification Report:
              precision    recall  f1-score   support

           0     0.9054    0.5475    0.6824      1801
           1     0.1054    0.4824    0.1730       199

    accuracy                         0.5410      2000
   macro avg     0.5054    0.5149    0.4277      2000
weighted avg     0.8258    0.5410    0.6317      2000

Confusion Matrix (Text):
                Predicted
               0       1
          -----------------
True   0 |    986     815
       1 |    103      96
          -----------------


===== Random Forest (Tuned) =====
Classification Report:
              precision    recall  f1-score   support

           0     0.8985    0.6635    0.7633      1801
           1     0.0955    0.3216    0.1473       199

    accuracy                         0.6295      2000
   macro avg     0.4970    0.4926    0.4553      2000
weighted avg     0.8186    0.6295    0.7020      2000

Confusion Matrix (Text):
                Predicted
               0       1
          -----------------
True   0 |   1195     606
       1 |    135      64
          -----------------


===== KNN (SMOTE Tuned) =====
Classification Report:
              precision    recall  f1-score   support

           0     0.9007    0.7501    0.8185      1801
           1     0.1000    0.2513    0.1431       199

    accuracy                         0.7005      2000
   macro avg     0.5003    0.5007    0.4808      2000
weighted avg     0.8210    0.7005    0.7513      2000

Confusion Matrix (Text):
                Predicted
               0       1
          -----------------
True   0 |   1351     450
       1 |    149      50
          -----------------


===== SVM Optimasi (Tuned) =====
Classification Report:
              precision    recall  f1-score   support

           0     0.9086    0.5136    0.6563      1801
           1     0.1079    0.5327    0.1795       199

    accuracy                         0.5155      2000
   macro avg     0.5083    0.5231    0.4179      2000
weighted avg     0.8290    0.5155    0.6088      2000

Confusion Matrix (Text):
                Predicted
               0       1
          -----------------
True   0 |    925     876
       1 |     93     106
          -----------------


===== XGBoost (Tuned) =====
Classification Report:
              precision    recall  f1-score   support

           0     0.9163    0.4986    0.6458      1801
           1     0.1147    0.5879    0.1920       199

    accuracy                         0.5075      2000
   macro avg     0.5155    0.5433    0.4189      2000
weighted avg     0.8366    0.5075    0.6007      2000

Confusion Matrix (Text):
                Predicted
               0       1
          -----------------
True   0 |    898     903
       1 |     82     117
          -----------------
\end{lstlisting}

\begin{lstlisting}[language=Python, caption={Kode 3}]

import pandas as pd
import numpy as np
import joblib
import seaborn as sns
import matplotlib.pyplot as plt

from sklearn.model_selection import train_test_split
from sklearn.preprocessing import LabelEncoder
from sklearn.metrics import classification_report, confusion_matrix, precision_recall_curve

from imblearn.combine import SMOTETomek
from xgboost import XGBClassifier

# ------------------------------------------------------------
# 1. LOAD DATA
# ------------------------------------------------------------
rentals    = pd.read_csv('rentals.csv')
logs       = pd.read_csv('rental_logs.csv')
agents     = pd.read_csv('agents.csv')
units      = pd.read_csv('units.csv')
users      = pd.read_csv('users.csv')
violations = pd.read_csv('violations.csv')

# Standarisasi nama kolom ke huruf kecil
rentals.columns = rentals.columns.str.lower()
logs.columns = logs.columns.str.lower()

if 'komentar' not in rentals.columns:
    raise ValueError("Kolom 'komentar' tidak ditemukan di rentals.csv")

print("rentals.csv")
print(rentals.head())
print("\nlogs.csv")
print(logs.head())

# ------------------------------------------------------------
# 2. LABELING & LOG MERGING
# ------------------------------------------------------------
# Menggabungkan logs ke rentals (membuat fitur: jumlah edit & editor unik)
logs['rental_id'] = logs['rental_id'].astype(int)
edit_count = logs.groupby('rental_id').size().reset_index(name='edit_count')
unique_editors = logs.groupby('rental_id')['email'].nunique().reset_index(name='unique_editors')

df = rentals.merge(edit_count, left_on='id', right_on='rental_id', how='left')
df = df.merge(unique_editors, left_on='id', right_on='rental_id', how='left')

# Isi nilai kosong dengan 0 untuk penyewaan yang tidak memiliki log
df['edit_count'] = df['edit_count'].fillna(0)
df['unique_editors'] = df['unique_editors'].fillna(0)

# Labeling Target: has_violation -> label
violation_keywords = [
    'alat suntik', 'kondom', 'kerusakan', 'kebisingan', 'alkohol', 'narkoba',
    'kekerasan', 'ancaman', 'merokok', 'kebersihan', 'tamu', 'menyewakan',
    'transit', 'lalai'
]

df['komentar'] = df['komentar'].astype(str).str.lower()
df['label'] = df['komentar'].apply(lambda x: int(any(k in x for k in violation_keywords)))

# ------------------------------------------------------------
# 3. FEATURE ENGINEERING & PREPARATION
# ------------------------------------------------------------
# Mengisi nilai kosong pada fitur kategorikal dan numerik
df['metode_pembayaran'] = df['metode_pembayaran'].fillna('Others')
if 'metode_lain' in df.columns:
    df['metode_lain'] = df['metode_lain'].fillna('')

df['lama_menginap'] = df['lama_menginap'].fillna(0)

cat_cols = ['tower', 'status_kewarganegaraan', 'metode_pembayaran', 'jenis_sewa']
numeric_cols = ['lantai', 'lama_menginap', 'edit_count', 'unique_editors']

# Daftar seluruh fitur relevan yang akan digunakan
feature_cols = cat_cols + numeric_cols

# Pastikan data kategorikal berupa string
for c in cat_cols:
    df[c] = df[c].astype(str)

X = df[feature_cols].copy()
y = df['label']

# ------------------------------------------------------------
# 4. TRAIN-TEST SPLIT
# ------------------------------------------------------------
X_train, X_test, y_train, y_test = train_test_split(
    X, y, stratify=y, test_size=0.2, random_state=42
)

# ------------------------------------------------------------
# 4B. ENCODING MENCEGAH DATA LEAKAGE
# ------------------------------------------------------------
# Simpan semua LabelEncoder ke dalam dictionary agar bisa di-export
label_encoders = {}

for c in cat_cols:
    le = LabelEncoder()
    # Fit & Transform pada data latih
    X_train[c] = le.fit_transform(X_train[c])

    # Tangani unseen labels pada data uji (ubah jadi 'Unknown' jika tidak pernah dilihat di X_train)
    X_test[c] = X_test[c].map(lambda s: s if s in le.classes_ else 'Unknown')
    if 'Unknown' not in le.classes_:
        le.classes_ = np.append(le.classes_, 'Unknown')

    X_test[c] = le.transform(X_test[c])
    label_encoders[c] = le

# ------------------------------------------------------------
# 5. RESAMPLING (SMOTETomek)
# ------------------------------------------------------------
print("\n=== Proses Resampling (SMOTETomek) ===")
smk = SMOTETomek(random_state=42)
X_res, y_res = smk.fit_resample(X_train, y_train)
print(f"Distribusi kelas sebelum resampling: \n{y_train.value_counts()}")
print(f"Distribusi kelas setelah resampling: \n{y_res.value_counts()}")

# ------------------------------------------------------------
# 6. TRAIN XGBOOST MODEL
# ------------------------------------------------------------
print("\n>> Melatih Final Model XGBoost...")
pos_weight = (y_train == 0).sum() / (y_train == 1).sum()

model = XGBClassifier(
    n_estimators=300,
    max_depth=4,
    learning_rate=0.05,
    subsample=0.9,
    colsample_bytree=0.8,
    scale_pos_weight=pos_weight,
    eval_metric='logloss',
    random_state=42
)

model.fit(X_res, y_res)

# ------------------------------------------------------------
# 7. THRESHOLD TUNING (F1-Score)
# ------------------------------------------------------------
print("\n=== Proses Threshold Tuning ===")
y_prob = model.predict_proba(X_test)[:, 1]

# Cari threshold terbaik berdasarkan metrik F1
prec, rec, thresh = precision_recall_curve(y_test, y_prob)
f1 = 2 * prec * rec / (prec + rec + 1e-9)
best_t = thresh[np.argmax(f1)]

print(f"-> Threshold Terbaik (Berdasarkan F1): {best_t:.4f}")

# Terapkan threshold terbaik untuk prediksi akhir
y_pred = (y_prob >= best_t).astype(int)

# ------------------------------------------------------------
# 8. EVALUATION
# ------------------------------------------------------------
print("\n===== FINAL MODEL CLASSIFICATION REPORT =====")
print(classification_report(y_test, y_pred, digits=4))

cm = confusion_matrix(y_test, y_pred)
print("Confusion Matrix (Text):")
print("                Predicted")
print("               0       1")
print("          -----------------")
print(f"True   0 | {cm[0][0]:6d}  {cm[0][1]:6d}")
print(f"       1 | {cm[1][0]:6d}  {cm[1][1]:6d}")
print("          -----------------\n")

# Output Confusion Matrix berupa visualisasi
# plt.figure(figsize=(4, 3))
# sns.heatmap(cm, annot=True, fmt='d', cmap='Blues')
# plt.title('Balanced XGBoost Confusion Matrix')
# plt.xlabel("Label Prediksi")
# plt.ylabel("Label Asli")
# plt.show()
\end{lstlisting}

\begin{lstlisting}[caption={Output 3}]
rentals.csv
   id            nama               nik status_pasutri tower  lantai  unit  \
0   1  N*** P****** 1  3***************  Belum Menikah     A       1    13   
1   2  N*** P****** 2  3***************        Menikah     D       3     4   
2   3  N*** P****** 3  3***************  Belum Menikah     B      18     2   
3   4  N*** P****** 4  3***************  Belum Menikah     B      11     7   
4   5  N*** P****** 5  3***************        Menikah     A      14     6   

  status_kewarganegaraan jenis_sewa metode_pembayaran  metode_lain  \
0                    WNI    Bulanan              Cash          NaN   
1                    WNA    Bulanan              Cash          NaN   
2                    WNA   Mingguan              QRIS          NaN   
3                    WNI     Harian              QRIS          NaN   
4                    WNA    Bulanan          Transfer          NaN   

  tanggal_checkin waktu_checkin tanggal_checkout waktu_checkout  \
0      2024-01-22      18:00:00       2024-01-31       06:00:00   
1      2023-06-16      18:00:00       2023-06-29       10:00:00   
2      2023-08-01      15:00:00       2023-08-02       09:00:00   
3      2023-03-30      10:00:00       2023-04-05       12:00:00   
4      2024-10-29      11:00:00       2024-11-02       08:00:00   

   lama_menginap                             komentar       email_agent  \
0              9                                  NaN  a*****@gmail.com   
1             13  Kebisingan berlebihan di malam hari  a*****@gmail.com   
2              1                                  NaN  a*****@gmail.com   
3              6                                  NaN  a*****@gmail.com   
4              4                                  NaN   a****@gmail.com   

        diedit_oleh           created_at  
0  a*****@gmail.com  2025-11-17 09:13:08  
1  a*****@gmail.com  2025-11-17 09:13:08  
2  a*****@gmail.com  2025-11-17 09:13:08  
3  a*****@gmail.com  2025-11-17 09:13:08  
4   a****@gmail.com  2025-11-17 09:13:08  

logs.csv
   id  rental_id  action  field_changed  old_value  new_value  \
0   1          1  CREATE            NaN        NaN        NaN   
1   2          2  CREATE            NaN        NaN        NaN   
2   3          3  CREATE            NaN        NaN        NaN   
3   4          4  CREATE            NaN        NaN        NaN   
4   5          5  CREATE            NaN        NaN        NaN   

              email           time_stamp  
0  a*****@gmail.com  2025-11-17 09:13:08  
1  a*****@gmail.com  2025-11-17 09:13:08  
2  a*****@gmail.com  2025-11-17 09:13:08  
3  a*****@gmail.com  2025-11-17 09:13:08  
4   a****@gmail.com  2025-11-17 09:13:08  

=== Proses Resampling (SMOTETomek) ===
Distribusi kelas sebelum resampling: 
label
0    7202
1     798
Name: count, dtype: int64
Distribusi kelas setelah resampling: 
label
0    7161
1    7161
Name: count, dtype: int64

>> Melatih Final Model XGBoost...

=== Proses Threshold Tuning ===
-> Threshold Terbaik (Berdasarkan F1): 0.3928

===== FINAL MODEL CLASSIFICATION REPORT =====
              precision    recall  f1-score   support

           0     0.9398    0.0433    0.0828      1801
           1     0.1012    0.9749    0.1834       199

    accuracy                         0.1360      2000
   macro avg     0.5205    0.5091    0.1331      2000
weighted avg     0.8563    0.1360    0.0928      2000

Confusion Matrix (Text):
                Predicted
               0       1
          -----------------
True   0 |     78    1723
       1 |      5     194
          -----------------
\end{lstlisting}

\begin{lstlisting}[language=Python, caption={Kode 4}]
import pandas as pd
import numpy as np
import joblib
import seaborn as sns
import matplotlib.pyplot as plt

from sklearn.model_selection import train_test_split
from sklearn.preprocessing import LabelEncoder, StandardScaler
from sklearn.metrics import classification_report, confusion_matrix, precision_recall_curve, precision_score
from imblearn.combine import SMOTETomek

# Models & Pipelines
from sklearn.pipeline import make_pipeline
from sklearn.tree import DecisionTreeClassifier
from sklearn.ensemble import RandomForestClassifier
from sklearn.neighbors import KNeighborsClassifier
from sklearn.svm import SVC
from xgboost import XGBClassifier

# ------------------------------------------------------------
# 1. LOAD DATA
# ------------------------------------------------------------
rentals = pd.read_csv('rentals.csv')
logs = pd.read_csv('rental_logs.csv')

# Standarisasi nama kolom ke huruf kecil
rentals.columns = rentals.columns.str.lower()
logs.columns = logs.columns.str.lower()

if 'komentar' not in rentals.columns:
    raise ValueError("Kolom 'komentar' tidak ditemukan di rentals.csv")

print("rentals.csv")
print(rentals.head())
print("\nlogs.csv")
print(logs.head())

# ------------------------------------------------------------
# 2. LABELING & LOG MERGING
# ------------------------------------------------------------
# Menggabungkan logs ke rentals (membuat fitur: jumlah edit & editor unik)
logs['rental_id'] = logs['rental_id'].astype(int)
edit_count = logs.groupby('rental_id').size().reset_index(name='edit_count')
unique_editors = logs.groupby('rental_id')['email'].nunique().reset_index(name='unique_editors')

df = rentals.merge(edit_count, left_on='id', right_on='rental_id', how='left')
df = df.merge(unique_editors, left_on='id', right_on='rental_id', how='left')

# Isi nilai kosong dengan 0 untuk penyewaan yang tidak memiliki log
df['edit_count'] = df['edit_count'].fillna(0)
df['unique_editors'] = df['unique_editors'].fillna(0)

# Labeling Target: has_violation -> label
violation_keywords = [
    'alat suntik', 'kondom', 'kerusakan', 'kebisingan', 'alkohol', 'narkoba',
    'kekerasan', 'ancaman', 'merokok', 'kebersihan', 'tamu', 'menyewakan',
    'transit', 'lalai'
]

# Bersihkan format teks dan berikan label
df['komentar'] = df['komentar'].astype(str).str.lower().str.strip()
df['label'] = df['komentar'].apply(lambda x: int(any(k in x for k in violation_keywords)))

# ------------------------------------------------------------
# 3. FEATURE ENGINEERING & PREPARATION
# ------------------------------------------------------------
# Mengisi nilai kosong pada fitur kategorikal dan numerik
df['metode_pembayaran'] = df['metode_pembayaran'].fillna('Others')
df['lama_menginap'] = df['lama_menginap'].fillna(0)

cat_cols = ['tower', 'status_kewarganegaraan', 'metode_pembayaran', 'jenis_sewa']
numeric_cols = ['lantai', 'lama_menginap', 'edit_count', 'unique_editors']

# Daftar seluruh fitur relevan yang akan digunakan (Tanpa membuang fitur)
feature_cols = cat_cols + numeric_cols

# Pastikan data kategorikal berupa string
for c in cat_cols:
    df[c] = df[c].astype(str)

X = df[feature_cols].copy()
y = df['label']

# ------------------------------------------------------------
# 4. TRAIN-TEST SPLIT
# ------------------------------------------------------------
X_train, X_test, y_train, y_test = train_test_split(
    X, y, stratify=y, test_size=0.2, random_state=42
)

# ------------------------------------------------------------
# 4B. ENCODING MENCEGAH DATA LEAKAGE
# ------------------------------------------------------------
# Simpan semua LabelEncoder ke dalam dictionary agar bisa di-export
label_encoders = {}

for c in cat_cols:
    le = LabelEncoder()
    # Fit & Transform hanya pada data latih
    X_train[c] = le.fit_transform(X_train[c])

    # Tangani unseen labels pada data uji (ubah jadi 'Unknown' jika tidak pernah dilihat di X_train)
    X_test[c] = X_test[c].map(lambda s: s if s in le.classes_ else 'Unknown')
    if 'Unknown' not in le.classes_:
        le.classes_ = np.append(le.classes_, 'Unknown')

    X_test[c] = le.transform(X_test[c])
    label_encoders[c] = le

# ------------------------------------------------------------
# 5. RESAMPLING (SMOTETomek)
# ------------------------------------------------------------
print("\n=== Proses Resampling (SMOTETomek) ===")
smk = SMOTETomek(random_state=42)
X_res, y_res = smk.fit_resample(X_train, y_train)
print(f"Distribusi kelas sebelum resampling: \n{y_train.value_counts()}")
print(f"Distribusi kelas setelah resampling: \n{y_res.value_counts()}")

# ------------------------------------------------------------
# 6. FUNGSI EVALUASI & THRESHOLD TUNING (BALANCED PRECISION)
# ------------------------------------------------------------
def best_threshold(y_true, y_prob):
    """Mencari threshold probabilitas terbaik berdasarkan F1/Balanced Precision."""
    prec, rec, thresh = precision_recall_curve(y_true, y_prob)

    best_bal = 0
    best_t = 0.5

    for t in thresh:
        y_pred = (y_prob >= t).astype(int)
        p0 = precision_score(y_true, y_pred, pos_label=0, zero_division=0)
        p1 = precision_score(y_true, y_pred, pos_label=1, zero_division=0)
        bal = (p0 + p1) / 2

        if bal > best_bal:
            best_bal = bal
            best_t = t

    return best_t


def evaluate(name, model):
    y_pred = model.predict(X_test)
    print(f"\n===== {name} =====")
    print("Classification Report:")
    print(classification_report(y_test, y_pred, digits=4))

    cm = confusion_matrix(y_test, y_pred)

    # --------------------------------------------------------
    # Confusion Matrix (Text Base)
    # --------------------------------------------------------
    print("Confusion Matrix (Text):")
    print("                Predicted")
    print("               0       1")
    print("          -----------------")
    print(f"True   0 | {cm[0][0]:6d}  {cm[0][1]:6d}")
    print(f"       1 | {cm[1][0]:6d}  {cm[1][1]:6d}")
    print("          -----------------\n")

    # Hitung nilai threshold jika model mendukung predict_proba
    if hasattr(model, "predict_proba"):
        y_prob = model.predict_proba(X_test)[:, 1]
        t = best_threshold(y_test, y_prob)
        print(f"-> Best Threshold Calculated: {t:.4f}")
    else:
        t = 0.5 # Default fallback
        print(f"-> Default Threshold Used: {t:.4f}")

    return t

# ------------------------------------------------------------
# 7. LATIH SEMUA MODEL
# ------------------------------------------------------------
models = {}
best_thresholds = {}

print("\nMemulai proses pelatihan untuk semua model...\n")

# --- 1. Decision Tree ---
print(">> Melatih Decision Tree...")
dt = DecisionTreeClassifier(class_weight='balanced', max_depth=5, random_state=42)
dt.fit(X_res, y_res)
models["Decision Tree"] = dt
best_thresholds["Decision Tree"] = evaluate("Decision Tree", dt) # Posisi Argumen Diperbaiki

# --- 2. Random Forest ---
print(">> Melatih Random Forest...")
rf = RandomForestClassifier(class_weight='balanced', n_estimators=200, max_depth=8, random_state=42)
rf.fit(X_res, y_res)
models["Random Forest"] = rf
best_thresholds["Random Forest"] = evaluate("Random Forest", rf) # Posisi Argumen Diperbaiki

# --- 3. KNN (SMOTETomek) ---
print(">> Melatih KNN...")
knn = KNeighborsClassifier(n_neighbors=7, weights='distance')
knn.fit(X_res, y_res)
models["KNN (SMOTETomek)"] = knn
best_thresholds["KNN (SMOTETomek)"] = evaluate("KNN (SMOTETomek)", knn) # Posisi Argumen Diperbaiki

# --- 4. SVM (RBF) ---
print(">> Melatih SVM (RBF)...")
svm = make_pipeline(
    StandardScaler(),
    SVC(kernel='rbf', probability=True, class_weight='balanced', random_state=42)
)
svm.fit(X_res, y_res)
models["SVM (RBF)"] = svm
best_thresholds["SVM (RBF)"] = evaluate("SVM (RBF)", svm) # Posisi Argumen Diperbaiki

# --- 5. XGBoost ---
print(">> Melatih XGBoost...")
pos_weight = (y_train == 0).sum() / (y_train == 1).sum()
xgb = XGBClassifier(
    n_estimators=300,
    max_depth=4,
    learning_rate=0.05,
    subsample=0.9,
    colsample_bytree=0.8,
    scale_pos_weight=pos_weight,
    eval_metric='logloss',
    random_state=42
)
xgb.fit(X_res, y_res)
models["XGBoost"] = xgb
best_thresholds["XGBoost"] = evaluate("XGBoost", xgb) # Posisi Argumen Diperbaiki
\end{lstlisting}

\begin{lstlisting}[caption={Output 4}]
rentals.csv
   id            nama               nik status_pasutri tower  lantai  unit  \
0   1  N*** P****** 1  3***************  Belum Menikah     A       1    13   
1   2  N*** P****** 2  3***************        Menikah     D       3     4   
2   3  N*** P****** 3  3***************  Belum Menikah     B      18     2   
3   4  N*** P****** 4  3***************  Belum Menikah     B      11     7   
4   5  N*** P****** 5  3***************        Menikah     A      14     6   

  status_kewarganegaraan jenis_sewa metode_pembayaran  metode_lain  \
0                    WNI    Bulanan              Cash          NaN   
1                    WNA    Bulanan              Cash          NaN   
2                    WNA   Mingguan              QRIS          NaN   
3                    WNI     Harian              QRIS          NaN   
4                    WNA    Bulanan          Transfer          NaN   

  tanggal_checkin waktu_checkin tanggal_checkout waktu_checkout  \
0      2024-01-22      18:00:00       2024-01-31       06:00:00   
1      2023-06-16      18:00:00       2023-06-29       10:00:00   
2      2023-08-01      15:00:00       2023-08-02       09:00:00   
3      2023-03-30      10:00:00       2023-04-05       12:00:00   
4      2024-10-29      11:00:00       2024-11-02       08:00:00   

   lama_menginap                             komentar       email_agent  \
0              9                                  NaN  a*****@gmail.com   
1             13  Kebisingan berlebihan di malam hari  a*****@gmail.com   
2              1                                  NaN  a*****@gmail.com   
3              6                                  NaN  a*****@gmail.com   
4              4                                  NaN   a****@gmail.com   

        diedit_oleh           created_at  
0  a*****@gmail.com  2025-11-17 09:13:08  
1  a*****@gmail.com  2025-11-17 09:13:08  
2  a*****@gmail.com  2025-11-17 09:13:08  
3  a*****@gmail.com  2025-11-17 09:13:08  
4   a****@gmail.com  2025-11-17 09:13:08  

logs.csv
   id  rental_id  action  field_changed  old_value  new_value  \
0   1          1  CREATE            NaN        NaN        NaN   
1   2          2  CREATE            NaN        NaN        NaN   
2   3          3  CREATE            NaN        NaN        NaN   
3   4          4  CREATE            NaN        NaN        NaN   
4   5          5  CREATE            NaN        NaN        NaN   

              email           time_stamp  
0  a*****@gmail.com  2025-11-17 09:13:08  
1  a*****@gmail.com  2025-11-17 09:13:08  
2  a*****@gmail.com  2025-11-17 09:13:08  
3  a*****@gmail.com  2025-11-17 09:13:08  
4   a****@gmail.com  2025-11-17 09:13:08  

=== Proses Resampling (SMOTETomek) ===
Distribusi kelas sebelum resampling: 
label
0    7202
1     798
Name: count, dtype: int64
Distribusi kelas setelah resampling: 
label
0    7161
1    7161
Name: count, dtype: int64

Memulai proses pelatihan untuk semua model...

>> Melatih Decision Tree...

===== Decision Tree =====
Classification Report:
              precision    recall  f1-score   support

           0     0.8894    0.6430    0.7464      1801
           1     0.0788    0.2764    0.1226       199

    accuracy                         0.6065      2000
   macro avg     0.4841    0.4597    0.4345      2000
weighted avg     0.8087    0.6065    0.6843      2000

Confusion Matrix (Text):
                Predicted
               0       1
          -----------------
True   0 |   1158     643
       1 |    144      55
          -----------------

-> Best Threshold Calculated: 0.8939
>> Melatih Random Forest...

===== Random Forest =====
Classification Report:
              precision    recall  f1-score   support

           0     0.8971    0.6052    0.7228      1801
           1     0.0943    0.3719    0.1504       199

    accuracy                         0.5820      2000
   macro avg     0.4957    0.4885    0.4366      2000
weighted avg     0.8172    0.5820    0.6659      2000

Confusion Matrix (Text):
                Predicted
               0       1
          -----------------
True   0 |   1090     711
       1 |    125      74
          -----------------

-> Best Threshold Calculated: 0.0821
>> Melatih KNN...

===== KNN (SMOTETomek) =====
Classification Report:
              precision    recall  f1-score   support

           0     0.9015    0.7568    0.8228      1801
           1     0.1025    0.2513    0.1456       199

    accuracy                         0.7065      2000
   macro avg     0.5020    0.5040    0.4842      2000
weighted avg     0.8220    0.7065    0.7554      2000

Confusion Matrix (Text):
                Predicted
               0       1
          -----------------
True   0 |   1363     438
       1 |    149      50
          -----------------

-> Best Threshold Calculated: 0.4361
>> Melatih SVM (RBF)...

===== SVM (RBF) =====
Classification Report:
              precision    recall  f1-score   support

           0     0.9005    0.5480    0.6814      1801
           1     0.0996    0.4523    0.1632       199

    accuracy                         0.5385      2000
   macro avg     0.5001    0.5001    0.4223      2000
weighted avg     0.8208    0.5385    0.6298      2000

Confusion Matrix (Text):
                Predicted
               0       1
          -----------------
True   0 |    987     814
       1 |    109      90
          -----------------

-> Best Threshold Calculated: 0.1798
>> Melatih XGBoost...

===== XGBoost =====
Classification Report:
              precision    recall  f1-score   support

           0     0.8865    0.0694    0.1287      1801
           1     0.0984    0.9196    0.1778       199

    accuracy                         0.1540      2000
   macro avg     0.4925    0.4945    0.1533      2000
weighted avg     0.8081    0.1540    0.1336      2000

Confusion Matrix (Text):
                Predicted
               0       1
          -----------------
True   0 |    125    1676
       1 |     16     183
          -----------------

-> Best Threshold Calculated: 0.9754
\end{lstlisting}

\begin{lstlisting}[language=Python, caption={Kode 5}]
# ============================================================
# DETEKSI PENYEWA BERPOTENSI MELANGGAR (All Features)
# Balanced Precision & Data Leakage Prevention
# ============================================================

import pandas as pd
import numpy as np
import joblib
import seaborn as sns
import matplotlib.pyplot as plt
import warnings

from sklearn.model_selection import train_test_split
from sklearn.preprocessing import LabelEncoder, StandardScaler
from sklearn.metrics import (
    precision_score,
    classification_report,
    confusion_matrix,
    precision_recall_curve
)
from imblearn.combine import SMOTETomek

# Models & Pipelines
from sklearn.pipeline import make_pipeline
from sklearn.tree import DecisionTreeClassifier
from sklearn.ensemble import RandomForestClassifier
from sklearn.neighbors import KNeighborsClassifier
from sklearn.svm import SVC
from xgboost import XGBClassifier

warnings.filterwarnings("ignore")

# ------------------------------------------------------------
# 1. LOAD DATA
# ------------------------------------------------------------
rentals    = pd.read_csv('rentals.csv')
logs       = pd.read_csv('rental_logs.csv')
agents     = pd.read_csv('agents.csv')
units      = pd.read_csv('units.csv')
users      = pd.read_csv('users.csv')
violations = pd.read_csv('violations.csv')

# Standarisasi nama kolom ke huruf kecil untuk semua dataframe
for frame in [rentals, logs, agents, units, users, violations]:
    frame.columns = frame.columns.str.lower()

if 'komentar' not in rentals.columns:
    rentals['komentar'] = ''

print("rentals.csv")
print(rentals.head())

# ------------------------------------------------------------
# 2. LABELING & LOG MERGING
# ------------------------------------------------------------
# Menggabungkan logs ke rentals (membuat fitur: jumlah edit & editor unik)
if 'rental_id' in logs.columns:
    logs['rental_id'] = logs['rental_id'].astype(int)
    edit_count = logs.groupby('rental_id').size().reset_index(name='edit_count')
    unique_editors = logs.groupby('rental_id')['email'].nunique().reset_index(name='unique_editors')

    df = rentals.merge(edit_count, left_on='id', right_on='rental_id', how='left')
    df = df.merge(unique_editors, left_on='id', right_on='rental_id', how='left')

    df['edit_count'] = df['edit_count'].fillna(0)
    df['unique_editors'] = df['unique_editors'].fillna(0)
else:
    raise ValueError("Kolom 'rental_id' tidak ditemukan di rental_logs.csv")

# Bersihkan format teks komentar
df['komentar'] = df['komentar'].astype(str).str.lower().str.strip()
empty_values = ["", "nan", "na", "null", "-", "0", "tidak ada", "none", " "]

# Definisi keyword pelanggaran
violation_keywords = [
    'alat suntik', 'kondom', 'kerusakan', 'kebisingan', 'alkohol', 'narkoba',
    'kekerasan', 'ancaman', 'merokok', 'kebersihan', 'tamu', 'menyewakan',
    'transit', 'lalai'
]

# a. Flag pelanggaran dari teks komentar
df['has_violation_text'] = df['komentar'].apply(
    lambda x: int(any(k in x for k in violation_keywords)) if x not in empty_values else 0
)

# b. Flag pelanggaran dari tabel violations
if 'rental_id' in violations.columns:
    violation_ids = set(violations['rental_id'].astype(int).tolist())
    df['has_violation_flag'] = df['id'].astype(int).isin(violation_ids).astype(int)
else:
    df['has_violation_flag'] = 0

# Final Target (Label): 1 jika ada indikasi di komentar ATAU terdata di tabel violations
df['label'] = ((df['has_violation_text'] == 1) | (df['has_violation_flag'] == 1)).astype(int)

# ------------------------------------------------------------
# 3. FEATURE ENGINEERING & PREPARATION
# ------------------------------------------------------------
# 3a. Join users untuk menentukan peran admin (is_admin_agent)
if 'email_agent' in df.columns and 'email' in users.columns:
    users_small = users[['email', 'role']].drop_duplicates()
    df = df.merge(users_small, left_on='email_agent', right_on='email', how='left', suffixes=('', '_user'))
    df['role'] = df['role'].fillna('unknown')

    admin_roles = ['ketua agen', 'p3srs', 'pkj']
    df['is_admin_agent'] = df['role'].str.lower().isin(admin_roles).astype(int)
else:
    df['is_admin_agent'] = 0

# 3b. Feature risk per agent
if 'email_agent' in df.columns:
    agent_risk = df.groupby('email_agent')['label'].mean().reset_index()
    agent_risk.rename(columns={'label': 'agent_risk'}, inplace=True)
    df = df.merge(agent_risk, on='email_agent', how='left')
    df['agent_risk'] = df['agent_risk'].fillna(0)
else:
    df['agent_risk'] = 0.0

# 3c. Feature risk per unit
if all(col in df.columns for col in ['tower', 'lantai', 'unit']):
    df['unit_key'] = df['tower'].astype(str) + '-' + df['lantai'].astype(str) + '-' + df['unit'].astype(str)
    unit_risk = df.groupby('unit_key')['label'].mean().reset_index()
    unit_risk.rename(columns={'label': 'unit_risk'}, inplace=True)
    df = df.merge(unit_risk, on='unit_key', how='left')
    df['unit_risk'] = df['unit_risk'].fillna(0)
else:
    df['unit_risk'] = 0.0

# 3d. Parse waktu checkout
if 'waktu_checkout' in df.columns:
    df['waktu_checkout_dt'] = pd.to_datetime(df['waktu_checkout'], errors='coerce')
    df['checkout_hour'] = df['waktu_checkout_dt'].dt.hour.fillna(0).astype(int)
else:
    df['checkout_hour'] = 0

# 3e. Pengisian nilai kosong dasar
df['metode_pembayaran'] = df.get('metode_pembayaran', 'Others').fillna('Others')
df['jenis_sewa'] = df.get('jenis_sewa', 'Unknown').fillna('Unknown')
df['status_kewarganegaraan'] = df.get('status_kewarganegaraan', 'Unknown').fillna('Unknown')
df['lama_menginap'] = df.get('lama_menginap', 0).fillna(0)

# Daftar kategori dan numerik relevan
cat_cols = [c for c in ['tower', 'status_kewarganegaraan', 'metode_pembayaran', 'jenis_sewa'] if c in df.columns]
numeric_cols = [c for c in ['lantai', 'lama_menginap', 'edit_count', 'unique_editors', 'is_admin_agent', 'agent_risk', 'unit_risk', 'checkout_hour'] if c in df.columns]

feature_cols = cat_cols + numeric_cols

for c in cat_cols:
    df[c] = df[c].astype(str)

X = df[feature_cols].copy()
y = df['label']

print("\n=== Distribusi Final Label ===")
print(y.value_counts())
print(f"Proporsi pelanggaran: {y.mean():.4f}")

if len(y.unique()) < 2:
    raise ValueError("Data masih hanya 1 kelas! Periksa kembali sumber labeling.")

# ------------------------------------------------------------
# 4. TRAIN-TEST SPLIT
# ------------------------------------------------------------
X_train, X_test, y_train, y_test = train_test_split(
    X, y, stratify=y, test_size=0.2, random_state=42
)

# ------------------------------------------------------------
# 4B. ENCODING MENCEGAH DATA LEAKAGE
# ------------------------------------------------------------
label_encoders = {}

for c in cat_cols:
    le = LabelEncoder()
    # Fit & Transform hanya pada data latih
    X_train[c] = le.fit_transform(X_train[c])

    # Tangani unseen labels pada data uji
    X_test[c] = X_test[c].map(lambda s: s if s in le.classes_ else 'Unknown')
    if 'Unknown' not in le.classes_:
        le.classes_ = np.append(le.classes_, 'Unknown')

    X_test[c] = le.transform(X_test[c])
    label_encoders[c] = le

# ------------------------------------------------------------
# 5. RESAMPLING (SMOTETomek)
# ------------------------------------------------------------
print("\n=== Proses Resampling (SMOTETomek) ===")
smk = SMOTETomek(random_state=42)
X_train_res, y_train_res = smk.fit_resample(X_train, y_train)

print(f"Distribusi kelas setelah resampling: \n{y_train_res.value_counts()}")

# ------------------------------------------------------------
# 6. FUNGSI EVALUASI & THRESHOLD TUNING (BALANCED PRECISION)
# ------------------------------------------------------------
def best_threshold(y_true, y_prob):
    """Mencari threshold probabilitas terbaik berdasarkan F1/Balanced Precision."""
    prec, rec, thresh = precision_recall_curve(y_true, y_prob)

    best_bal = 0
    best_t = 0.5

    for t in thresh:
        y_pred = (y_prob >= t).astype(int)
        p0 = precision_score(y_true, y_pred, pos_label=0, zero_division=0)
        p1 = precision_score(y_true, y_pred, pos_label=1, zero_division=0)
        bal = (p0 + p1) / 2

        if bal > best_bal:
            best_bal = bal
            best_t = t

    return best_t

def evaluate(name, model):
    y_pred = model.predict(X_test)
    print(f"\n===== {name} =====")
    print("Classification Report:")
    print(classification_report(y_test, y_pred, digits=4))

    cm = confusion_matrix(y_test, y_pred)

    # --------------------------------------------------------
    # Confusion Matrix (Text Base)
    # --------------------------------------------------------
    print("Confusion Matrix (Text):")
    print("                Predicted")
    print("               0       1")
    print("          -----------------")
    print(f"True   0 | {cm[0][0]:6d}  {cm[0][1]:6d}")
    print(f"       1 | {cm[1][0]:6d}  {cm[1][1]:6d}")
    print("          -----------------\n")

    # Hitung nilai threshold jika model mendukung predict_proba
    if hasattr(model, "predict_proba"):
        y_prob = model.predict_proba(X_test)[:, 1]
        t = best_threshold(y_test, y_prob)
        print(f"-> Best Threshold Calculated: {t:.4f}")
    else:
        t = 0.5 # Default fallback
        print(f"-> Default Threshold Used: {t:.4f}")

    return t

# ------------------------------------------------------------
# 7. LATIH SEMUA MODEL
# ------------------------------------------------------------
models = {}
best_thresholds = {}

print("\nMemulai proses pelatihan untuk semua model...\n")

# --- 1. Decision Tree ---
print(">> Melatih Decision Tree...")
dt = DecisionTreeClassifier(class_weight='balanced', max_depth=5, random_state=42)
dt.fit(X_train_res, y_train_res)
models["Decision Tree"] = dt
best_thresholds["Decision Tree"] = evaluate("Decision Tree", dt) # Urutan Diperbaiki

# --- 2. Random Forest ---
print(">> Melatih Random Forest...")
rf = RandomForestClassifier(class_weight='balanced', n_estimators=200, max_depth=8, random_state=42)
rf.fit(X_train_res, y_train_res)
models["Random Forest"] = rf
best_thresholds["Random Forest"] = evaluate("Random Forest", rf) # Urutan Diperbaiki

# --- 3. KNN (SMOTETomek) ---
print(">> Melatih KNN...")
knn = make_pipeline(StandardScaler(), KNeighborsClassifier(n_neighbors=7, weights='distance'))
knn.fit(X_train_res, y_train_res)
models["KNN (SMOTETomek)"] = knn
best_thresholds["KNN (SMOTETomek)"] = evaluate("KNN (SMOTETomek)", knn) # Urutan Diperbaiki

# --- 4. SVM (RBF) ---
print(">> Melatih SVM (RBF)...")
svm = make_pipeline(
    StandardScaler(),
    SVC(kernel='rbf', probability=True, class_weight='balanced', random_state=42)
)
svm.fit(X_train_res, y_train_res)
models["SVM (RBF)"] = svm
best_thresholds["SVM (RBF)"] = evaluate("SVM (RBF)", svm) # Urutan Diperbaiki

# --- 5. XGBoost ---
print(">> Melatih XGBoost...")
pos_weight = (y_train == 0).sum() / max(1, (y_train == 1).sum())
xgb = XGBClassifier(
    n_estimators=300,
    max_depth=4,
    learning_rate=0.05,
    subsample=0.9,
    colsample_bytree=0.8,
    scale_pos_weight=pos_weight,
    eval_metric='logloss',
    random_state=42
)
xgb.fit(X_train_res, y_train_res)
models["XGBoost"] = xgb
best_thresholds["XGBoost"] = evaluate("XGBoost", xgb)
\end{lstlisting}

\begin{lstlisting}[caption={Output 5}]
rentals.csv
   id            nama               nik status_pasutri tower  lantai  unit  \
0   1  N*** P****** 1  3***************  Belum Menikah     A       1    13   
1   2  N*** P****** 2  3***************        Menikah     D       3     4   
2   3  N*** P****** 3  3***************  Belum Menikah     B      18     2   
3   4  N*** P****** 4  3***************  Belum Menikah     B      11     7   
4   5  N*** P****** 5  3***************        Menikah     A      14     6   

  status_kewarganegaraan jenis_sewa metode_pembayaran  metode_lain  \
0                    WNI    Bulanan              Cash          NaN   
1                    WNA    Bulanan              Cash          NaN   
2                    WNA   Mingguan              QRIS          NaN   
3                    WNI     Harian              QRIS          NaN   
4                    WNA    Bulanan          Transfer          NaN   

  tanggal_checkin waktu_checkin tanggal_checkout waktu_checkout  \
0      2024-01-22      18:00:00       2024-01-31       06:00:00   
1      2023-06-16      18:00:00       2023-06-29       10:00:00   
2      2023-08-01      15:00:00       2023-08-02       09:00:00   
3      2023-03-30      10:00:00       2023-04-05       12:00:00   
4      2024-10-29      11:00:00       2024-11-02       08:00:00   

   lama_menginap                             komentar       email_agent  \
0              9                                  NaN  a*****@gmail.com   
1             13  Kebisingan berlebihan di malam hari  a*****@gmail.com   
2              1                                  NaN  a*****@gmail.com   
3              6                                  NaN  a*****@gmail.com   
4              4                                  NaN   a****@gmail.com   

        diedit_oleh           created_at  
0  a*****@gmail.com  2025-11-17 09:13:08  
1  a*****@gmail.com  2025-11-17 09:13:08  
2  a*****@gmail.com  2025-11-17 09:13:08  
3  a*****@gmail.com  2025-11-17 09:13:08  
4   a****@gmail.com  2025-11-17 09:13:08  

=== Distribusi Final Label ===
label
0    9003
1     997
Name: count, dtype: int64
Proporsi pelanggaran: 0.0997

=== Proses Resampling (SMOTETomek) ===
Distribusi kelas setelah resampling: 
label
0    7146
1    7146
Name: count, dtype: int64

Memulai proses pelatihan untuk semua model...

>> Melatih Decision Tree...

===== Decision Tree =====
Classification Report:
              precision    recall  f1-score   support

           0     0.9163    0.6807    0.7811      1801
           1     0.1314    0.4372    0.2021       199

    accuracy                         0.6565      2000
   macro avg     0.5239    0.5590    0.4916      2000
weighted avg     0.8382    0.6565    0.7235      2000

Confusion Matrix (Text):
                Predicted
               0       1
          -----------------
True   0 |   1226     575
       1 |    112      87
          -----------------

-> Best Threshold Calculated: 0.0201
>> Melatih Random Forest...

===== Random Forest =====
Classification Report:
              precision    recall  f1-score   support

           0     0.9148    0.6852    0.7835      1801
           1     0.1290    0.4221    0.1976       199

    accuracy                         0.6590      2000
   macro avg     0.5219    0.5536    0.4906      2000
weighted avg     0.8366    0.6590    0.7252      2000

Confusion Matrix (Text):
                Predicted
               0       1
          -----------------
True   0 |   1234     567
       1 |    115      84
          -----------------

-> Best Threshold Calculated: 0.8091
>> Melatih KNN...

===== KNN (SMOTETomek) =====
Classification Report:
              precision    recall  f1-score   support

           0     0.9070    0.6607    0.7645      1801
           1     0.1119    0.3869    0.1736       199

    accuracy                         0.6335      2000
   macro avg     0.5095    0.5238    0.4691      2000
weighted avg     0.8279    0.6335    0.7057      2000

Confusion Matrix (Text):
                Predicted
               0       1
          -----------------
True   0 |   1190     611
       1 |    122      77
          -----------------

-> Best Threshold Calculated: 0.9073
>> Melatih SVM (RBF)...

===== SVM (RBF) =====
Classification Report:
              precision    recall  f1-score   support

           0     0.9132    0.6663    0.7705      1801
           1     0.1239    0.4271    0.1921       199

    accuracy                         0.6425      2000
   macro avg     0.5186    0.5467    0.4813      2000
weighted avg     0.8347    0.6425    0.7129      2000

Confusion Matrix (Text):
                Predicted
               0       1
          -----------------
True   0 |   1200     601
       1 |    114      85
          -----------------

-> Best Threshold Calculated: 0.0165
>> Melatih XGBoost...

===== XGBoost =====
Classification Report:
              precision    recall  f1-score   support

           0     0.9484    0.3470    0.5081      1801
           1     0.1230    0.8291    0.2143       199

    accuracy                         0.3950      2000
   macro avg     0.5357    0.5881    0.3612      2000
weighted avg     0.8663    0.3950    0.4789      2000

Confusion Matrix (Text):
                Predicted
               0       1
          -----------------
True   0 |    625    1176
       1 |     34     165
          -----------------

-> Best Threshold Calculated: 0.9655
\end{lstlisting}

\begin{lstlisting}[language=Python, caption={Kode 6}]
# ============================================================
# DATA SEED
# ============================================================

import pandas as pd
import numpy as np
import joblib
import seaborn as sns
import matplotlib.pyplot as plt
import warnings

from sklearn.model_selection import train_test_split
from sklearn.preprocessing import LabelEncoder, StandardScaler, MinMaxScaler
from sklearn.metrics import classification_report, confusion_matrix
from sklearn.feature_selection import SelectKBest, chi2, mutual_info_classif
from scipy.stats import pearsonr

from imblearn.over_sampling import RandomOverSampler

# Models & Pipelines
from sklearn.pipeline import make_pipeline
from sklearn.tree import DecisionTreeClassifier
from sklearn.ensemble import RandomForestClassifier
from sklearn.neighbors import KNeighborsClassifier
from sklearn.svm import SVC
from xgboost import XGBClassifier

warnings.filterwarnings("ignore")

# ------------------------------------------------------------
# 1. LOAD DATA
# ------------------------------------------------------------
rentals    = pd.read_csv('rentals.csv')
logs       = pd.read_csv('rental_logs.csv')
units      = pd.read_csv('units.csv')
users      = pd.read_csv('users.csv')
violations = pd.read_csv('violations.csv')

# Standarisasi nama kolom ke huruf kecil
rentals.columns = rentals.columns.str.lower()

if 'komentar' not in rentals.columns:
    raise ValueError("Kolom 'komentar' tidak ditemukan di rentals.csv")

print("rentals.csv")
print(rentals.head())

# ------------------------------------------------------------
# 2. LABELING
# ------------------------------------------------------------
rentals['komentar'] = rentals['komentar'].astype(str).str.lower().str.strip()

empty_values = ["", "nan", "na", "null", "-", "0", "tidak ada", "none", " "]

rentals['label'] = (~rentals['komentar'].isin(empty_values)).astype(int)

print("\n=== Distribusi Label Setelah FIX ===")
print(rentals['label'].value_counts())
print(f"Proporsi pelanggaran: {rentals['label'].mean():.4f}")

# ------------------------------------------------------------
# 3. FEATURE ENGINEERING & PREPARATION
# ------------------------------------------------------------
required_cols = ['jenis_sewa', 'metode_pembayaran', 'waktu_checkout']
for col in required_cols:
    if col not in rentals.columns:
        raise ValueError(f"Kolom '{col}' tidak ditemukan")

# Parse waktu checkout
rentals['waktu_checkout_dt'] = pd.to_datetime(rentals['waktu_checkout'], errors='coerce')
rentals['checkout_hour'] = rentals['waktu_checkout_dt'].dt.hour.fillna(0).astype(int)

cat_cols = ['jenis_sewa', 'metode_pembayaran']
numeric_cols = ['checkout_hour']

# Daftar seluruh fitur relevan yang akan digunakan (Tanpa membuang fitur)
feature_cols = cat_cols + numeric_cols

# Pastikan data kategorikal berupa string
for c in cat_cols:
    rentals[c] = rentals[c].astype(str)

X = rentals[feature_cols].copy()
y = rentals['label']

# ------------------------------------------------------------
# 4. TRAIN-TEST SPLIT
# ------------------------------------------------------------
if len(y.unique()) < 2:
    raise ValueError("Data masih hanya 1 kelas! Periksa kembali kolom komentar.")

X_train, X_test, y_train, y_test = train_test_split(
    X, y, stratify=y, test_size=0.2, random_state=42
)

print("\nDistribusi TRAIN:")
print(y_train.value_counts())
print("\nDistribusi TEST:")
print(y_test.value_counts())

# ------------------------------------------------------------
# 4B. ENCODING MENCEGAH DATA LEAKAGE
# ------------------------------------------------------------
# Simpan semua LabelEncoder ke dalam dictionary agar bisa di-export
label_encoders = {}

for c in cat_cols:
    le = LabelEncoder()
    # Fit & Transform hanya pada data latih
    X_train[c] = le.fit_transform(X_train[c])

    # Tangani unseen labels pada data uji (ubah jadi 'Unknown' jika tidak pernah dilihat di X_train)
    X_test[c] = X_test[c].map(lambda s: s if s in le.classes_ else 'Unknown')
    if 'Unknown' not in le.classes_:
        le.classes_ = np.append(le.classes_, 'Unknown')

    X_test[c] = le.transform(X_test[c])
    label_encoders[c] = le

# ------------------------------------------------------------
# 5. FEATURE EVALUATION (Chi-square, Pearson, MI)
# ------------------------------------------------------------
print("\n=== Evaluasi Korelasi Fitur ===")
# Catatan: Pemodelan menggunakan seluruh kolom relevan yang tersedia dengan
# spesifik encoding (tidak otomatis menyeleksi sebagian/membuang fitur).
# Pengujian ini sekadar untuk analisis signifikansi fitur.

chi2_selector = SelectKBest(score_func=chi2, k='all')
chi2_selector.fit(X_train, y_train)
chi2_scores = chi2_selector.scores_

train_df = X_train.copy()
train_df['label'] = y_train
pearson_scores = train_df.corr(method='pearson')['label'].drop('label').abs().values

mi_selector = SelectKBest(score_func=mutual_info_classif, k='all')
mi_selector.fit(X_train, y_train)
mi_scores = mi_selector.scores_

fs_df = pd.DataFrame({
    'Feature': X_train.columns,
    'Chi-Square': chi2_scores,
    'Pearson (Abs)': pearson_scores,
    'Mutual Info': mi_scores
})

# Normalisasi skor
scaler = MinMaxScaler()
fs_df[['Chi2_Norm', 'Pearson_Norm', 'MI_Norm']] = scaler.fit_transform(
    fs_df[['Chi-Square', 'Pearson (Abs)', 'Mutual Info']]
)
fs_df['Total_Score'] = fs_df[['Chi2_Norm', 'Pearson_Norm', 'MI_Norm']].mean(axis=1)
fs_df = fs_df.sort_values(by='Total_Score', ascending=False).reset_index(drop=True)

print(fs_df[['Feature', 'Chi-Square', 'Pearson (Abs)', 'Mutual Info', 'Total_Score']])

# ------------------------------------------------------------
# 6. RESAMPLING (RandomOverSampler)
# ------------------------------------------------------------
print("\n=== Proses Resampling (Oversampling) ===")
ros = RandomOverSampler(random_state=42)
X_train_res, y_train_res = ros.fit_resample(X_train, y_train)

print(f"Distribusi kelas setelah Oversampling:\n{y_train_res.value_counts()}")

# ------------------------------------------------------------
# 7. FUNGSI EVALUASI
# ------------------------------------------------------------
def evaluate(name, model, X_eval=X_test):
    y_pred = model.predict(X_eval)

    print(f"\n===== {name} =====")
    print("Classification Report:")
    print(classification_report(y_test, y_pred, digits=4))

    cm = confusion_matrix(y_test, y_pred)

    print("Confusion Matrix (Text):")
    print("                Predicted")
    print("               0       1")
    print("          -----------------")
    print(f"True   0 | {cm[0][0]:6d}  {cm[0][1]:6d}")
    print(f"       1 | {cm[1][0]:6d}  {cm[1][1]:6d}")
    print("          -----------------\n")

    # # Output Confusion Matrix berupa visualisasi
    # plt.figure(figsize=(4, 3))
    # sns.heatmap(cm, annot=True, fmt='d', cmap='Blues')
    # plt.title(f'Confusion Matrix: {name}')
    # plt.xlabel('Label Prediksi')
    # plt.ylabel('Label Asli')
    # plt.show()

# ------------------------------------------------------------
# 8. LATIH SEMUA MODEL
# ------------------------------------------------------------
models = {}

print("\nMemulai proses pelatihan untuk semua model...\n")

# --- 1. Decision Tree ---
print(">> Melatih Decision Tree...")
dt = DecisionTreeClassifier(max_depth=5, class_weight='balanced', random_state=42)
dt.fit(X_train_res, y_train_res)
models["Decision Tree"] = dt

# --- 2. Random Forest ---
print(">> Melatih Random Forest...")
rf = RandomForestClassifier(n_estimators=200, max_depth=8, class_weight='balanced', random_state=42)
rf.fit(X_train_res, y_train_res)
models["Random Forest"] = rf

# --- 3. KNN ---
print(">> Melatih KNN...")
knn = make_pipeline(StandardScaler(), KNeighborsClassifier(n_neighbors=7))
knn.fit(X_train_res, y_train_res)
models["KNN"] = knn

# --- 4. SVM (RBF) ---
print(">> Melatih SVM (RBF)...")
svm = make_pipeline(
    StandardScaler(),
    SVC(kernel='rbf', probability=True, class_weight='balanced', random_state=42)
)
svm.fit(X_train_res, y_train_res)
models["SVM RBF"] = svm

# --- 5. XGBoost ---
print(">> Melatih XGBoost...")
pos_weight = (y_train == 0).sum() / max(1, (y_train == 1).sum())
xgb = XGBClassifier(
    n_estimators=300,
    max_depth=4,
    learning_rate=0.05,
    scale_pos_weight=pos_weight,
    eval_metric='logloss',
    random_state=42
)
xgb.fit(X_train_res, y_train_res)
models["XGBoost"] = xgb

# ------------------------------------------------------------
# 9. EVALUASI SEMUA MODEL
# ------------------------------------------------------------
for name, model in models.items():
    evaluate(name, model)

# ------------------------------------------------------------
# 10. EKSPORT MODEL TERBAIK & LABEL ENCODERS
# ------------------------------------------------------------
print("\n=== Proses Export Model ===")

# Anda dapat menyesuaikan indeks string ini ke model yang paling stabil performanya
target_model_name = "XGBoost"

export_package = {
    'model': models[target_model_name],
    'label_encoders': label_encoders,
    'feature_cols': feature_cols
}

filename = 'final_model_advanced_features.pkl'
joblib.dump(export_package, filename)
print(f"Sukses! Final Model ({target_model_name}) beserta komponennya telah disimpan sebagai file: {filename}")
\end{lstlisting}

\begin{lstlisting}[caption={Output 6}]
rentals.csv
   id                nama               nik status_pasutri  \
0   1        F**** B*****  5***************        Menikah   
1   2      E***** P******  3***************        Menikah   
2   3        O**** H*****  3***************  Belum Menikah   
3   4  N***** B*****-B***  3***************  Belum Menikah   
4   5       R****** S****  3***************        Menikah   

  status_kewarganegaraan jenis_sewa unit_number metode_pembayaran metode_lain  \
0                    WNI   Mingguan      C-1-13              Cash         NaN   
1                    WNA     Harian       D-5-2      Kartu Kredit         NaN   
2                    WNI    Bulanan     B-25-17              QRIS         NaN   
3                    WNI    Bulanan       A-2-8            Others       gopay   
4                    WNI     Harian      A-21-9       Kartu Debit         NaN   

  tanggal_checkin waktu_checkin tanggal_checkout waktu_checkout  \
0      2024-06-20      20:22:00       2024-07-04       07:46:00   
1      2023-02-19      05:33:00       2023-02-20       12:33:00   
2      2023-12-17      17:48:00       2024-03-16       09:02:00   
3      2018-05-18      13:15:00       2018-10-15       12:45:00   
4      2025-01-05      17:28:00       2025-01-05       22:28:00   

   lama_menginap                                komentar  \
0             14                                     NaN   
1              1                                     NaN   
2             90                                     NaN   
3            150                                     NaN   
4              0  Ditemukan alat suntik di tempat sampah   

                     user_email diedit_oleh           created_at  \
0   O**************@hotmail.com         NaN  2025-12-08 19:13:06   
1  D***************@hotmail.com         NaN  2025-12-08 19:13:06   
2              K*****@gmail.com         NaN  2025-12-08 19:13:06   
3              T*****@gmail.com         NaN  2025-12-08 19:13:06   
4      K*************@yahoo.com         NaN  2025-12-08 19:13:06   

  classification  
0         normal  
1         normal  
2         normal  
3         normal  
4   tidak normal  

=== Distribusi Label Setelah FIX ===
label
0    9243
1     746
Name: count, dtype: int64
Proporsi pelanggaran: 0.0747

Distribusi TRAIN:
label
0    7394
1     597
Name: count, dtype: int64

Distribusi TEST:
label
0    1849
1     149
Name: count, dtype: int64

=== Evaluasi Korelasi Fitur ===
             Feature  Chi-Square  Pearson (Abs)  Mutual Info  Total_Score
0      checkout_hour   46.241510       0.103399     0.059024     0.876096
1         jenis_sewa    0.080931       0.003885     0.093944     0.333333
2  metode_pembayaran    0.242216       0.005513     0.000000     0.006619

=== Proses Resampling (Oversampling) ===
Distribusi kelas setelah Oversampling:
label
0    7394
1    7394
Name: count, dtype: int64

Memulai proses pelatihan untuk semua model...

>> Melatih Decision Tree...
>> Melatih Random Forest...
>> Melatih KNN...
>> Melatih SVM (RBF)...
>> Melatih XGBoost...

===== Decision Tree =====
Classification Report:
              precision    recall  f1-score   support

           0     1.0000    0.7301    0.8440      1849
           1     0.2299    1.0000    0.3739       149

    accuracy                         0.7503      1998
   macro avg     0.6150    0.8651    0.6090      1998
weighted avg     0.9426    0.7503    0.8090      1998

Confusion Matrix (Text):
                Predicted
               0       1
          -----------------
True   0 |   1350     499
       1 |      0     149
          -----------------


===== Random Forest =====
Classification Report:
              precision    recall  f1-score   support

           0     1.0000    0.7307    0.8444      1849
           1     0.2303    1.0000    0.3744       149

    accuracy                         0.7508      1998
   macro avg     0.6151    0.8653    0.6094      1998
weighted avg     0.9426    0.7508    0.8093      1998

Confusion Matrix (Text):
                Predicted
               0       1
          -----------------
True   0 |   1351     498
       1 |      0     149
          -----------------


===== KNN =====
Classification Report:
              precision    recall  f1-score   support

           0     0.9524    0.9946    0.9730      1849
           1     0.8507    0.3826    0.5278       149

    accuracy                         0.9489      1998
   macro avg     0.9016    0.6886    0.7504      1998
weighted avg     0.9448    0.9489    0.9398      1998

Confusion Matrix (Text):
                Predicted
               0       1
          -----------------
True   0 |   1839      10
       1 |     92      57
          -----------------


===== SVM RBF =====
Classification Report:
              precision    recall  f1-score   support

           0     1.0000    0.7301    0.8440      1849
           1     0.2299    1.0000    0.3739       149

    accuracy                         0.7503      1998
   macro avg     0.6150    0.8651    0.6090      1998
weighted avg     0.9426    0.7503    0.8090      1998

Confusion Matrix (Text):
                Predicted
               0       1
          -----------------
True   0 |   1350     499
       1 |      0     149
          -----------------


===== XGBoost =====
Classification Report:
              precision    recall  f1-score   support

           0     1.0000    0.7301    0.8440      1849
           1     0.2299    1.0000    0.3739       149

    accuracy                         0.7503      1998
   macro avg     0.6150    0.8651    0.6090      1998
weighted avg     0.9426    0.7503    0.8090      1998

Confusion Matrix (Text):
                Predicted
               0       1
          -----------------
True   0 |   1350     499
       1 |      0     149
          -----------------


=== Proses Export Model ===
Sukses! Final Model (XGBoost) beserta komponennya telah disimpan sebagai file: final_model_advanced_features.pkl
\end{lstlisting}

\begin{lstlisting}[language=Python, caption={Kode 7}]
# ============================================================
# DATA SEED (Cost-Sensitive Learning / Imbalanced Handling)
# ============================================================

import pandas as pd
import numpy as np
import joblib
import seaborn as sns
import matplotlib.pyplot as plt
import warnings

from sklearn.model_selection import train_test_split
from sklearn.preprocessing import LabelEncoder, StandardScaler, MinMaxScaler
from sklearn.metrics import classification_report, confusion_matrix
from sklearn.feature_selection import SelectKBest, chi2, mutual_info_classif

# Models & Pipelines
from sklearn.pipeline import make_pipeline
from sklearn.tree import DecisionTreeClassifier
from sklearn.ensemble import RandomForestClassifier
from sklearn.neighbors import KNeighborsClassifier
from sklearn.svm import SVC
from xgboost import XGBClassifier

warnings.filterwarnings("ignore")

# ------------------------------------------------------------
# 1. LOAD DATA
# ------------------------------------------------------------
rentals    = pd.read_csv('rentals.csv')
logs       = pd.read_csv('rental_logs.csv')
units      = pd.read_csv('units.csv')
users      = pd.read_csv('users.csv')
violations = pd.read_csv('violations.csv')

# Standarisasi nama kolom ke huruf kecil
rentals.columns = rentals.columns.str.lower()

if 'komentar' not in rentals.columns:
    raise ValueError("Kolom 'komentar' tidak ditemukan di rentals.csv")

print("rentals.csv")
print(rentals.head())

# ------------------------------------------------------------
# 2. LABELING
# ------------------------------------------------------------
# Bersihkan format teks (ubah ke huruf kecil dan hapus spasi berlebih)
rentals['komentar'] = rentals['komentar'].astype(str).str.lower().str.strip()

empty_values = ["", "nan", "na", "null", "-", "0", "tidak ada", "none", " "]

# Jika isi komentar TIDAK ada di dalam empty_values, berarti ada indikasi pelanggaran (label 1)
rentals['label'] = (~rentals['komentar'].isin(empty_values)).astype(int)

print("\n=== Distribusi Label Setelah FIX ===")
print(rentals['label'].value_counts())
print(f"Proporsi pelanggaran: {rentals['label'].mean():.4f}")

# ------------------------------------------------------------
# 3. FEATURE ENGINEERING & PREPARATION
# ------------------------------------------------------------
required_cols = ['jenis_sewa', 'metode_pembayaran', 'waktu_checkout']
for col in required_cols:
    if col not in rentals.columns:
        raise ValueError(f"Kolom '{col}' tidak ditemukan")

# Parse waktu checkout
rentals['waktu_checkout_dt'] = pd.to_datetime(rentals['waktu_checkout'], errors='coerce')
rentals['checkout_hour'] = rentals['waktu_checkout_dt'].dt.hour.fillna(0).astype(int)

cat_cols = ['jenis_sewa', 'metode_pembayaran']
numeric_cols = ['checkout_hour']

# Daftar seluruh fitur relevan yang akan digunakan (Tanpa membuang fitur)
feature_cols = cat_cols + numeric_cols

# Pastikan data kategorikal berupa string sebelum di-encode
for c in cat_cols:
    rentals[c] = rentals[c].astype(str)

X = rentals[feature_cols].copy()
y = rentals['label']

# ------------------------------------------------------------
# 4. TRAIN-TEST SPLIT
# ------------------------------------------------------------
if len(y.unique()) < 2:
    raise ValueError("Data masih hanya 1 kelas! Periksa kembali kolom komentar.")

X_train, X_test, y_train, y_test = train_test_split(
    X, y, stratify=y, test_size=0.2, random_state=42
)

print("\nDistribusi TRAIN:")
print(y_train.value_counts())
print("\nDistribusi TEST:")
print(y_test.value_counts())

# ------------------------------------------------------------
# 4B. ENCODING MENCEGAH DATA LEAKAGE
# ------------------------------------------------------------
# Simpan semua LabelEncoder ke dalam dictionary agar bisa di-export
label_encoders = {}

for c in cat_cols:
    le = LabelEncoder()
    # Fit & Transform hanya pada data latih
    X_train[c] = le.fit_transform(X_train[c])

    # Tangani unseen labels pada data uji (ubah jadi 'Unknown' jika tidak pernah dilihat di X_train)
    X_test[c] = X_test[c].map(lambda s: s if s in le.classes_ else 'Unknown')
    if 'Unknown' not in le.classes_:
        le.classes_ = np.append(le.classes_, 'Unknown')

    X_test[c] = le.transform(X_test[c])
    label_encoders[c] = le

# ------------------------------------------------------------
# 5. FEATURE EVALUATION (Chi-square, Pearson, MI)
# ------------------------------------------------------------
print("\n=== Evaluasi Signifikansi Fitur ===")
# Catatan: Pemodelan menggunakan seluruh kolom relevan yang tersedia dengan spesifik
# encoding (tidak otomatis membuang fitur). Pengujian ini untuk analisis signifikansi.

chi2_selector = SelectKBest(score_func=chi2, k='all')
chi2_selector.fit(X_train, y_train)
chi2_scores = chi2_selector.scores_

train_df = X_train.copy()
train_df['label'] = y_train
pearson_scores = train_df.corr(method='pearson')['label'].drop('label').abs().values

mi_selector = SelectKBest(score_func=mutual_info_classif, k='all')
mi_selector.fit(X_train, y_train)
mi_scores = mi_selector.scores_

fs_df = pd.DataFrame({
    'Feature': X_train.columns,
    'Chi-Square': chi2_scores,
    'Pearson (Abs)': pearson_scores,
    'Mutual Info': mi_scores
})

# Normalisasi skor agar bisa dibandingkan
scaler = MinMaxScaler()
fs_df[['Chi2_Norm', 'Pearson_Norm', 'MI_Norm']] = scaler.fit_transform(
    fs_df[['Chi-Square', 'Pearson (Abs)', 'Mutual Info']]
)
fs_df['Total_Score'] = fs_df[['Chi2_Norm', 'Pearson_Norm', 'MI_Norm']].mean(axis=1)
fs_df = fs_df.sort_values(by='Total_Score', ascending=False).reset_index(drop=True)

print(fs_df[['Feature', 'Chi-Square', 'Pearson (Abs)', 'Mutual Info', 'Total_Score']])

# ------------------------------------------------------------
# 6. FUNGSI EVALUASI
# ------------------------------------------------------------
def evaluate(name, model, X_eval=X_test):
    y_pred = model.predict(X_eval)

    print(f"\n===== {name} =====")
    print("Classification Report:")
    print(classification_report(y_test, y_pred, digits=4))

    cm = confusion_matrix(y_test, y_pred)

    print("Confusion Matrix (Text):")
    print("                Predicted")
    print("               0       1")
    print("          -----------------")
    print(f"True   0 | {cm[0][0]:6d}  {cm[0][1]:6d}")
    print(f"       1 | {cm[1][0]:6d}  {cm[1][1]:6d}")
    print("          -----------------\n")

    # # Output Confusion Matrix berupa visualisasi
    # plt.figure(figsize=(4, 3))
    # sns.heatmap(cm, annot=True, fmt='d', cmap='Blues')
    # plt.title(f'Confusion Matrix: {name}')
    # plt.xlabel('Label Prediksi')
    # plt.ylabel('Label Asli')
    # plt.show()

# ------------------------------------------------------------
# 7. LATIH SEMUA MODEL (DENGAN COST-SENSITIVE LEARNING)
# ------------------------------------------------------------
models = {}

print("\nMemulai proses pelatihan untuk semua model...\n")

# --- 1. Decision Tree ---
print(">> Melatih Decision Tree...")
dt = DecisionTreeClassifier(max_depth=5, class_weight='balanced', random_state=42)
dt.fit(X_train, y_train)
models["Decision Tree"] = dt

# --- 2. Random Forest ---
print(">> Melatih Random Forest...")
rf = RandomForestClassifier(n_estimators=200, max_depth=6, class_weight='balanced_subsample', random_state=42)
rf.fit(X_train, y_train)
models["Random Forest"] = rf

# --- 3. KNN ---
# KNN tidak memiliki fitur class_weight bawaan, kita gunakan pembobotan jarak
print(">> Melatih KNN...")
knn = make_pipeline(StandardScaler(), KNeighborsClassifier(n_neighbors=5, weights='distance'))
knn.fit(X_train, y_train)
models["KNN"] = knn

# --- 4. SVM (RBF) ---
print(">> Melatih SVM (RBF)...")
svm = make_pipeline(
    StandardScaler(),
    SVC(kernel='rbf', probability=True, class_weight='balanced', random_state=42)
)
svm.fit(X_train, y_train)
models["SVM RBF"] = svm

# --- 5. XGBoost ---
print(">> Melatih XGBoost...")
pos_weight = (y_train == 0).sum() / max(1, (y_train == 1).sum())
xgb = XGBClassifier(
    n_estimators=200,
    max_depth=4,
    learning_rate=0.05,
    scale_pos_weight=pos_weight,
    eval_metric='logloss',
    random_state=42
)
xgb.fit(X_train, y_train)
models["XGBoost"] = xgb

# ------------------------------------------------------------
# 8. EVALUASI SEMUA MODEL
# ------------------------------------------------------------
for name, model in models.items():
    evaluate(name, model)

# ------------------------------------------------------------
# 9. EKSPORT MODEL TERBAIK & LABEL ENCODERS
# ------------------------------------------------------------
print("\n=== Proses Export Model ===")

# Anda dapat menyesuaikan indeks string ini dengan model yang performanya paling seimbang
target_model_name = "XGBoost"

export_package = {
    'model': models[target_model_name],
    'label_encoders': label_encoders,
    'feature_cols': feature_cols
}

filename = 'final_model_cost_sensitive.pkl'
joblib.dump(export_package, filename)
print(f"Sukses! Final Model ({target_model_name}) beserta komponennya telah disimpan sebagai file: {filename}")
\end{lstlisting}

\begin{lstlisting}[caption={Output 7}]
rentals.csv
   id                nama               nik status_pasutri  \
0   1        F**** B*****  5***************        Menikah   
1   2      E***** P******  3***************        Menikah   
2   3        O**** H*****  3***************  Belum Menikah   
3   4  N***** B*****-B***  3***************  Belum Menikah   
4   5       R****** S****  3***************        Menikah   

  status_kewarganegaraan jenis_sewa unit_number metode_pembayaran metode_lain  \
0                    WNI   Mingguan      C-1-13              Cash         NaN   
1                    WNA     Harian       D-5-2      Kartu Kredit         NaN   
2                    WNI    Bulanan     B-25-17              QRIS         NaN   
3                    WNI    Bulanan       A-2-8            Others       gopay   
4                    WNI     Harian      A-21-9       Kartu Debit         NaN   

  tanggal_checkin waktu_checkin tanggal_checkout waktu_checkout  \
0      2024-06-20      20:22:00       2024-07-04       07:46:00   
1      2023-02-19      05:33:00       2023-02-20       12:33:00   
2      2023-12-17      17:48:00       2024-03-16       09:02:00   
3      2018-05-18      13:15:00       2018-10-15       12:45:00   
4      2025-01-05      17:28:00       2025-01-05       22:28:00   

   lama_menginap                                komentar  \
0             14                                     NaN   
1              1                                     NaN   
2             90                                     NaN   
3            150                                     NaN   
4              0  Ditemukan alat suntik di tempat sampah   

                     user_email diedit_oleh           created_at  \
0   O**************@hotmail.com         NaN  2025-12-08 19:13:06   
1  D***************@hotmail.com         NaN  2025-12-08 19:13:06   
2              K*****@gmail.com         NaN  2025-12-08 19:13:06   
3              T*****@gmail.com         NaN  2025-12-08 19:13:06   
4      K*************@yahoo.com         NaN  2025-12-08 19:13:06   

  classification  
0         normal  
1         normal  
2         normal  
3         normal  
4   tidak normal  

=== Distribusi Label Setelah FIX ===
label
0    9243
1     746
Name: count, dtype: int64
Proporsi pelanggaran: 0.0747

Distribusi TRAIN:
label
0    7394
1     597
Name: count, dtype: int64

Distribusi TEST:
label
0    1849
1     149
Name: count, dtype: int64

=== Evaluasi Signifikansi Fitur ===
             Feature  Chi-Square  Pearson (Abs)  Mutual Info  Total_Score
0      checkout_hour   46.241510       0.103399     0.061000     0.889383
1         jenis_sewa    0.080931       0.003885     0.089653     0.333333
2  metode_pembayaran    0.242216       0.005513     0.003309     0.006619

Memulai proses pelatihan untuk semua model...

>> Melatih Decision Tree...
>> Melatih Random Forest...
>> Melatih KNN...
>> Melatih SVM (RBF)...
>> Melatih XGBoost...

===== Decision Tree =====
Classification Report:
              precision    recall  f1-score   support

           0     1.0000    0.7301    0.8440      1849
           1     0.2299    1.0000    0.3739       149

    accuracy                         0.7503      1998
   macro avg     0.6150    0.8651    0.6090      1998
weighted avg     0.9426    0.7503    0.8090      1998

Confusion Matrix (Text):
                Predicted
               0       1
          -----------------
True   0 |   1350     499
       1 |      0     149
          -----------------


===== Random Forest =====
Classification Report:
              precision    recall  f1-score   support

           0     1.0000    0.7301    0.8440      1849
           1     0.2299    1.0000    0.3739       149

    accuracy                         0.7503      1998
   macro avg     0.6150    0.8651    0.6090      1998
weighted avg     0.9426    0.7503    0.8090      1998

Confusion Matrix (Text):
                Predicted
               0       1
          -----------------
True   0 |   1350     499
       1 |      0     149
          -----------------


===== KNN =====
Classification Report:
              precision    recall  f1-score   support

           0     0.9498    0.9816    0.9654      1849
           1     0.6092    0.3557    0.4492       149

    accuracy                         0.9349      1998
   macro avg     0.7795    0.6687    0.7073      1998
weighted avg     0.9244    0.9349    0.9269      1998

Confusion Matrix (Text):
                Predicted
               0       1
          -----------------
True   0 |   1815      34
       1 |     96      53
          -----------------


===== SVM RBF =====
Classification Report:
              precision    recall  f1-score   support

           0     1.0000    0.7301    0.8440      1849
           1     0.2299    1.0000    0.3739       149

    accuracy                         0.7503      1998
   macro avg     0.6150    0.8651    0.6090      1998
weighted avg     0.9426    0.7503    0.8090      1998

Confusion Matrix (Text):
                Predicted
               0       1
          -----------------
True   0 |   1350     499
       1 |      0     149
          -----------------


===== XGBoost =====
Classification Report:
              precision    recall  f1-score   support

           0     1.0000    0.7301    0.8440      1849
           1     0.2299    1.0000    0.3739       149

    accuracy                         0.7503      1998
   macro avg     0.6150    0.8651    0.6090      1998
weighted avg     0.9426    0.7503    0.8090      1998

Confusion Matrix (Text):
                Predicted
               0       1
          -----------------
True   0 |   1350     499
       1 |      0     149
          -----------------


=== Proses Export Model ===
Sukses! Final Model (XGBoost) beserta komponennya telah disimpan sebagai file: final_model_cost_sensitive.pkl
\end{lstlisting}

\begin{lstlisting}[language=Python, caption={Kode 8}]
# ============================================================
# DATA SEED (REVISI IMBALANCED HANDLING - Custom Weights)
# ============================================================

import pandas as pd
import numpy as np
import joblib
import seaborn as sns
import matplotlib.pyplot as plt
import warnings

from sklearn.model_selection import train_test_split
from sklearn.preprocessing import LabelEncoder, StandardScaler, MinMaxScaler
from sklearn.metrics import classification_report, confusion_matrix
from sklearn.feature_selection import SelectKBest, chi2, mutual_info_classif

# Models & Pipelines
from sklearn.pipeline import make_pipeline
from sklearn.tree import DecisionTreeClassifier
from sklearn.ensemble import RandomForestClassifier
from sklearn.neighbors import KNeighborsClassifier
from sklearn.svm import SVC
from xgboost import XGBClassifier

warnings.filterwarnings("ignore")

# ------------------------------------------------------------
# 1. LOAD DATA
# ------------------------------------------------------------
rentals    = pd.read_csv('rentals.csv')
logs       = pd.read_csv('rental_logs.csv')
units      = pd.read_csv('units.csv')
users      = pd.read_csv('users.csv')
violations = pd.read_csv('violations.csv')

# Standarisasi nama kolom ke huruf kecil
rentals.columns = rentals.columns.str.lower()

if 'komentar' not in rentals.columns:
    raise ValueError("Kolom 'komentar' tidak ditemukan di rentals.csv")

print("rentals.csv")
print(rentals.head())
print("\nlogs.csv")
print(logs.head())
print("\nunits.csv")
print(units.head())
print("\nusers.csv")
print(users.head())
print("\nviolations.csv")
print(violations.head())

# ------------------------------------------------------------
# 2. LABELING
# ------------------------------------------------------------
# Bersihkan format teks (ubah ke huruf kecil dan hapus spasi berlebih)
rentals['komentar'] = rentals['komentar'].astype(str).str.lower().str.strip()

empty_values = ["", "nan", "na", "null", "-", "0", "tidak ada", "none", " "]

# Jika isi komentar TIDAK ada di dalam empty_values, berarti ada indikasi pelanggaran (label 1)
rentals['label'] = (~rentals['komentar'].isin(empty_values)).astype(int)

print("\n=== Distribusi Label Setelah FIX ===")
print(rentals['label'].value_counts())
print(f"Proporsi pelanggaran: {rentals['label'].mean():.4f}")

# ------------------------------------------------------------
# 3. FEATURE ENGINEERING & PREPARATION
# ------------------------------------------------------------
required_cols = ['jenis_sewa', 'metode_pembayaran', 'waktu_checkout']
for col in required_cols:
    if col not in rentals.columns:
        raise ValueError(f"Kolom '{col}' tidak ditemukan")

# Parse waktu checkout
rentals['waktu_checkout_dt'] = pd.to_datetime(rentals['waktu_checkout'], errors='coerce')
rentals['checkout_hour'] = rentals['waktu_checkout_dt'].dt.hour.fillna(0).astype(int)

cat_cols = ['jenis_sewa', 'metode_pembayaran']
numeric_cols = ['checkout_hour']

# Daftar seluruh fitur relevan yang akan digunakan (Tanpa membuang fitur)
feature_cols = cat_cols + numeric_cols

# Pastikan data kategorikal berupa string sebelum di-encode
for c in cat_cols:
    rentals[c] = rentals[c].astype(str)

X = rentals[feature_cols].copy()
y = rentals['label']

# ------------------------------------------------------------
# 4. TRAIN-TEST SPLIT
# ------------------------------------------------------------
if len(y.unique()) < 2:
    raise ValueError("Data masih hanya 1 kelas! Periksa kembali kolom komentar.")

X_train, X_test, y_train, y_test = train_test_split(
    X, y, stratify=y, test_size=0.2, random_state=42
)

print("\nDistribusi TRAIN:")
print(y_train.value_counts())

# ------------------------------------------------------------
# 4B. ENCODING MENCEGAH DATA LEAKAGE
# ------------------------------------------------------------
# Simpan semua LabelEncoder ke dalam dictionary agar bisa di-export
label_encoders = {}

for c in cat_cols:
    le = LabelEncoder()
    # Fit & Transform hanya pada data latih
    X_train[c] = le.fit_transform(X_train[c])

    # Tangani unseen labels pada data uji (ubah jadi 'Unknown' jika tidak pernah dilihat di X_train)
    X_test[c] = X_test[c].map(lambda s: s if s in le.classes_ else 'Unknown')
    if 'Unknown' not in le.classes_:
        le.classes_ = np.append(le.classes_, 'Unknown')

    X_test[c] = le.transform(X_test[c])
    label_encoders[c] = le

# ------------------------------------------------------------
# 5. FEATURE EVALUATION (Chi-square, Pearson, MI)
# ------------------------------------------------------------
print("\n=== Evaluasi Signifikansi Fitur ===")
# Catatan: Pemodelan menggunakan seluruh kolom relevan yang tersedia dengan spesifik
# encoding (tidak otomatis membuang fitur). Pengujian ini untuk keperluan analisis data.

chi2_selector = SelectKBest(score_func=chi2, k='all')
chi2_selector.fit(X_train, y_train)
chi2_scores = chi2_selector.scores_

train_df = X_train.copy()
train_df['label'] = y_train
pearson_scores = train_df.corr(method='pearson')['label'].drop('label').abs().values

mi_selector = SelectKBest(score_func=mutual_info_classif, k='all')
mi_selector.fit(X_train, y_train)
mi_scores = mi_selector.scores_

fs_df = pd.DataFrame({
    'Feature': X_train.columns,
    'Chi-Square': chi2_scores,
    'Pearson (Abs)': pearson_scores,
    'Mutual Info': mi_scores
})

# Normalisasi skor agar bisa dibandingkan
scaler = MinMaxScaler()
fs_df[['Chi2_Norm', 'Pearson_Norm', 'MI_Norm']] = scaler.fit_transform(
    fs_df[['Chi-Square', 'Pearson (Abs)', 'Mutual Info']]
)
fs_df['Total_Score'] = fs_df[['Chi2_Norm', 'Pearson_Norm', 'MI_Norm']].mean(axis=1)
fs_df = fs_df.sort_values(by='Total_Score', ascending=False).reset_index(drop=True)

print(fs_df[['Feature', 'Chi-Square', 'Pearson (Abs)', 'Mutual Info', 'Total_Score']])

# ------------------------------------------------------------
# 6. FUNGSI EVALUASI
# ------------------------------------------------------------
def evaluate(name, model, X_eval=X_test):
    y_pred = model.predict(X_eval)

    print(f"\n===== {name} =====")
    print("Classification Report:")
    print(classification_report(y_test, y_pred, digits=4))

    cm = confusion_matrix(y_test, y_pred)

    print("Confusion Matrix (Text):")
    print("                Predicted")
    print("               0       1")
    print("          -----------------")
    print(f"True   0 | {cm[0][0]:6d}  {cm[0][1]:6d}")
    print(f"       1 | {cm[1][0]:6d}  {cm[1][1]:6d}")
    print("          -----------------\n")

    # # Output Confusion Matrix berupa visualisasi
    # plt.figure(figsize=(4, 3))
    # sns.heatmap(cm, annot=True, fmt='d', cmap='Blues')
    # plt.title(f'Confusion Matrix: {name}')
    # plt.xlabel('Label Prediksi')
    # plt.ylabel('Label Asli')
    # plt.show()

# ------------------------------------------------------------
# 7. LATIH SEMUA MODEL (REVISI: CUSTOM WEIGHTS)
# ------------------------------------------------------------
models = {}

print("\nMemulai proses pelatihan untuk semua model dengan Custom Weights...\n")

# Kita atur bobot manual: Kelas 0 bobotnya 1, Kelas 1 bobotnya 4
# (Artinya AI diminta 4x lebih waspada terhadap pelanggaran)
custom_weight = {0: 1, 1: 4}

# --- 1. Decision Tree ---
print(">> Melatih Decision Tree...")
dt = DecisionTreeClassifier(max_depth=5, class_weight=custom_weight, random_state=42)
dt.fit(X_train, y_train)
models["Decision Tree"] = dt

# --- 2. Random Forest ---
print(">> Melatih Random Forest...")
rf = RandomForestClassifier(n_estimators=200, max_depth=6, class_weight=custom_weight, random_state=42)
rf.fit(X_train, y_train)
models["Random Forest"] = rf

# --- 3. KNN ---
print(">> Melatih KNN...")
# KNN tidak memiliki fitur class_weight, kita gunakan pembobotan jarak
knn = make_pipeline(StandardScaler(), KNeighborsClassifier(n_neighbors=5, weights='distance'))
knn.fit(X_train, y_train)
models["KNN"] = knn

# --- 4. SVM (RBF) ---
print(">> Melatih SVM (RBF)...")
svm = make_pipeline(
    StandardScaler(),
    SVC(kernel='rbf', probability=True, class_weight=custom_weight, random_state=42)
)
svm.fit(X_train, y_train)
models["SVM RBF"] = svm

# --- 5. XGBoost ---
print(">> Melatih XGBoost...")
xgb = XGBClassifier(
    n_estimators=200,
    max_depth=4,
    learning_rate=0.05,
    scale_pos_weight=4,
    eval_metric='logloss',
    random_state=42
)
xgb.fit(X_train, y_train)
models["XGBoost"] = xgb

# ------------------------------------------------------------
# 8. EVALUASI SEMUA MODEL
# ------------------------------------------------------------
for name, model in models.items():
    evaluate(name, model)

# ------------------------------------------------------------
# 9. EKSPORT MODEL TERBAIK & LABEL ENCODERS
# ------------------------------------------------------------
print("\n=== Proses Export Model ===")

# Anda dapat mengganti indeks string di bawah ini sesuai dengan model yang paling seimbang hasilnya
target_model_name = "XGBoost"

export_package = {
    'model': models[target_model_name],
    'label_encoders': label_encoders,
    'feature_cols': feature_cols
}

filename = 'final_model_custom_weights.pkl'
joblib.dump(export_package, filename)
print(f"Sukses! Final Model ({target_model_name}) beserta komponennya telah disimpan sebagai file: {filename}")
\end{lstlisting}

\begin{lstlisting}[caption={Output 8}]
rentals.csv
   id                nama               nik status_pasutri  \
0   1        F**** B*****  5***************        Menikah   
1   2      E***** P******  3***************        Menikah   
2   3        O**** H*****  3***************  Belum Menikah   
3   4  N***** B*****-B***  3***************  Belum Menikah   
4   5       R****** S****  3***************        Menikah   

  status_kewarganegaraan jenis_sewa unit_number metode_pembayaran metode_lain  \
0                    WNI   Mingguan      C-1-13              Cash         NaN   
1                    WNA     Harian       D-5-2      Kartu Kredit         NaN   
2                    WNI    Bulanan     B-25-17              QRIS         NaN   
3                    WNI    Bulanan       A-2-8            Others       gopay   
4                    WNI     Harian      A-21-9       Kartu Debit         NaN   

  tanggal_checkin waktu_checkin tanggal_checkout waktu_checkout  \
0      2024-06-20      20:22:00       2024-07-04       07:46:00   
1      2023-02-19      05:33:00       2023-02-20       12:33:00   
2      2023-12-17      17:48:00       2024-03-16       09:02:00   
3      2018-05-18      13:15:00       2018-10-15       12:45:00   
4      2025-01-05      17:28:00       2025-01-05       22:28:00   

   lama_menginap                                komentar  \
0             14                                     NaN   
1              1                                     NaN   
2             90                                     NaN   
3            150                                     NaN   
4              0  Ditemukan alat suntik di tempat sampah   

                     user_email diedit_oleh           created_at  \
0   O**************@hotmail.com         NaN  2025-12-08 19:13:06   
1  D***************@hotmail.com         NaN  2025-12-08 19:13:06   
2              K*****@gmail.com         NaN  2025-12-08 19:13:06   
3              T*****@gmail.com         NaN  2025-12-08 19:13:06   
4      K*************@yahoo.com         NaN  2025-12-08 19:13:06   

  classification  
0         normal  
1         normal  
2         normal  
3         normal  
4   tidak normal  

logs.csv
   id  rental_id  action field_changed old_value  \
0   1          1  create           all       NaN   
1   2          2  create           all       NaN   
2   3          3  create           all       NaN   
3   4          4  create           all       NaN   
4   5          5  create           all       NaN   

                                           new_value  \
0  {"id":1,"nama":"F**** B*****","nik":"5*******************"id":2,"nama":"E***** P******","nik":"3*****************"id":3,"nama":"O**** H*****","nik":"3*******************"id":4,"nama":"N***** B*****-B***","nik":"3*************"id":5,"nama":"R****** S****","nik":"33966831...   

                          email            timestamp  
0   O**************@hotmail.com  2025-12-08 19:13:06  
1  D***************@hotmail.com  2025-12-08 19:13:06  
2              K*****@gmail.com  2025-12-08 19:13:06  
3              T*****@gmail.com  2025-12-08 19:13:06  
4      K*************@yahoo.com  2025-12-08 19:13:06  

units.csv
  unit_number                   user_email tower  lantai  unit
0      A-1-14     K*************@yahoo.com     A       1    14
1      A-1-18         A*******@hotmail.com     A       1    18
2      A-1-22  O**************@hotmail.com     A       1    22
3      A-1-24                p**@gmail.com     A       1    24
4      A-1-28  O**************@hotmail.com     A       1    28

users.csv
   id                         email         password           nib  role  \
0   1    Y***************@yahoo.com  oe4eaGaBLpQv76y  7.265157e+12  agen   
1   2           R******@hotmail.com  UhoFvefkUMB44Cj  2.780507e+12  agen   
2   3  F*****************@yahoo.com  Jf5D9lZJIGVf011  1.335353e+12  agen   
3   4       M************@yahoo.com  DYPYEZs_Zh17IjQ  8.262239e+12  agen   
4   5         E**********@gmail.com  1gcfQuFUGjUATHy  8.127913e+11  agen   

               name           created_at  
0        K** M*** I  2025-12-08 19:13:05  
1    D*. J**** L***  2025-12-08 19:13:05  
2  R******** G*****  2025-12-08 19:13:05  
3  B******* K******  2025-12-08 19:13:05  
4      M****** D***  2025-12-08 19:13:05  

violations.csv
   id  rental_id                                       photo_url  \
0   1        233      https://picsum.photos/seed/u717oX/983/2412   
1   2       7182      https://picsum.photos/seed/0cZJtXQ/243/249   
2   3       9117     https://picsum.photos/seed/Pl0nan/3034/2819   
3   4       8582  https://picsum.photos/seed/LKFORW1qWE/3127/725   
4   5       4165    https://picsum.photos/seed/dkpBkZPg/336/1118   

                                         description  \
0              Crepusculum torrens vomica atque rem.   
1  Inventore articulus corrumpo pax clamo absterg...   
2  Nostrum desino comitatus debilito perferendis ...   
3  Tremo accedo aegre votum adhuc aveho vilicus v...   
4                         Comminor vis thorax absum.   

                         uploaded_by           created_at  
0     L*******************@yahoo.com  2025-12-08 19:13:07  
1        O**************@hotmail.com  2025-12-08 19:13:07  
2       A*****************@gmail.com  2025-12-08 19:13:07  
3                  S******@gmail.com  2025-12-08 19:13:07  
4  I********************@hotmail.com  2025-12-08 19:13:07  

=== Distribusi Label Setelah FIX ===
label
0    9243
1     746
Name: count, dtype: int64
Proporsi pelanggaran: 0.0747

Distribusi TRAIN:
label
0    7394
1     597
Name: count, dtype: int64

=== Evaluasi Signifikansi Fitur ===
             Feature  Chi-Square  Pearson (Abs)  Mutual Info  Total_Score
0      checkout_hour   46.241510       0.103399     0.058345     0.876701
1         jenis_sewa    0.080931       0.003885     0.092596     0.333333
2  metode_pembayaran    0.242216       0.005513     0.000000     0.006619

Memulai proses pelatihan untuk semua model dengan Custom Weights...

>> Melatih Decision Tree...
>> Melatih Random Forest...
>> Melatih KNN...
>> Melatih SVM (RBF)...
>> Melatih XGBoost...

===== Decision Tree =====
Classification Report:
              precision    recall  f1-score   support

           0     0.9524    0.9946    0.9730      1849
           1     0.8507    0.3826    0.5278       149

    accuracy                         0.9489      1998
   macro avg     0.9016    0.6886    0.7504      1998
weighted avg     0.9448    0.9489    0.9398      1998

Confusion Matrix (Text):
                Predicted
               0       1
          -----------------
True   0 |   1839      10
       1 |     92      57
          -----------------


===== Random Forest =====
Classification Report:
              precision    recall  f1-score   support

           0     0.9524    0.9946    0.9730      1849
           1     0.8507    0.3826    0.5278       149

    accuracy                         0.9489      1998
   macro avg     0.9016    0.6886    0.7504      1998
weighted avg     0.9448    0.9489    0.9398      1998

Confusion Matrix (Text):
                Predicted
               0       1
          -----------------
True   0 |   1839      10
       1 |     92      57
          -----------------


===== KNN =====
Classification Report:
              precision    recall  f1-score   support

           0     0.9498    0.9816    0.9654      1849
           1     0.6092    0.3557    0.4492       149

    accuracy                         0.9349      1998
   macro avg     0.7795    0.6687    0.7073      1998
weighted avg     0.9244    0.9349    0.9269      1998

Confusion Matrix (Text):
                Predicted
               0       1
          -----------------
True   0 |   1815      34
       1 |     96      53
          -----------------


===== SVM RBF =====
Classification Report:
              precision    recall  f1-score   support

           0     0.9524    0.9946    0.9730      1849
           1     0.8507    0.3826    0.5278       149

    accuracy                         0.9489      1998
   macro avg     0.9016    0.6886    0.7504      1998
weighted avg     0.9448    0.9489    0.9398      1998

Confusion Matrix (Text):
                Predicted
               0       1
          -----------------
True   0 |   1839      10
       1 |     92      57
          -----------------


===== XGBoost =====
Classification Report:
              precision    recall  f1-score   support

           0     0.9544    0.9724    0.9633      1849
           1     0.5526    0.4228    0.4791       149

    accuracy                         0.9314      1998
   macro avg     0.7535    0.6976    0.7212      1998
weighted avg     0.9244    0.9314    0.9272      1998

Confusion Matrix (Text):
                Predicted
               0       1
          -----------------
True   0 |   1798      51
       1 |     86      63
          -----------------


=== Proses Export Model ===
Sukses! Final Model (XGBoost) beserta komponennya telah disimpan sebagai file: final_model_custom_weights.pkl
\end{lstlisting}

\begin{lstlisting}[language=Python, caption={Kode 9}]
# ============================================================
# DATA SEED (REVISI PENAMBAHAN FITUR BARU & FIX ERROR)
# ============================================================

import pandas as pd
import numpy as np
import joblib
import seaborn as sns
import matplotlib.pyplot as plt
import warnings

from sklearn.model_selection import train_test_split
from sklearn.preprocessing import LabelEncoder, StandardScaler, MinMaxScaler
from sklearn.metrics import classification_report, confusion_matrix
from sklearn.feature_selection import SelectKBest, chi2, mutual_info_classif

# Models & Pipelines
from sklearn.pipeline import make_pipeline
from sklearn.tree import DecisionTreeClassifier
from sklearn.ensemble import RandomForestClassifier
from sklearn.neighbors import KNeighborsClassifier
from sklearn.svm import SVC
from xgboost import XGBClassifier

warnings.filterwarnings("ignore")

# ------------------------------------------------------------
# 1. LOAD DATA
# ------------------------------------------------------------
rentals    = pd.read_csv('rentals.csv')
# logs       = pd.read_csv('rental_logs.csv')
# units      = pd.read_csv('units.csv')
# users      = pd.read_csv('users.csv')
# violations = pd.read_csv('violations.csv')

# Standarisasi nama kolom ke huruf kecil
rentals.columns = rentals.columns.str.lower()
logs.columns = logs.columns.str.lower()
units.columns = units.columns.str.lower()
users.columns = users.columns.str.lower()
violations.columns = violations.columns.str.lower()

if 'komentar' not in rentals.columns:
    raise ValueError("Kolom 'komentar' tidak ditemukan di rentals.csv")

print("rentals.csv")
print(rentals.head())
# print("\nlogs.csv")
# print(logs.head())
# print("\nunits.csv")
# print(units.head())
# print("\nusers.csv")
# print(users.head())
# print("\nviolations.csv")
# print(violations.head())

# ------------------------------------------------------------
# 2. LABELING
# ------------------------------------------------------------
# Bersihkan format teks (ubah ke huruf kecil dan hapus spasi berlebih)
rentals['komentar'] = rentals['komentar'].astype(str).str.lower().str.strip()

empty_values = ["", "nan", "na", "null", "-", "0", "tidak ada", "none", " "]

# Jika isi komentar TIDAK ada di dalam empty_values, berarti ada indikasi pelanggaran (label 1)
rentals['label'] = (~rentals['komentar'].isin(empty_values)).astype(int)

print("\n=== Distribusi Label Setelah FIX ===")
print(rentals['label'].value_counts())
print(f"Proporsi pelanggaran: {rentals['label'].mean():.4f}")

# ------------------------------------------------------------
# 3. FEATURE ENGINEERING & PREPARATION
# ------------------------------------------------------------
# Parse waktu checkout
rentals['waktu_checkout_dt'] = pd.to_datetime(rentals['waktu_checkout'], errors='coerce')
rentals['checkout_hour'] = rentals['waktu_checkout_dt'].dt.hour.fillna(0).astype(int)

# Parse waktu checkin
rentals['waktu_checkin_dt'] = pd.to_datetime(rentals['waktu_checkin'], errors='coerce')
rentals['checkin_hour'] = rentals['waktu_checkin_dt'].dt.hour.fillna(0).astype(int)

# Pastikan lama_menginap berupa angka (Fitur Baru)
rentals['lama_menginap'] = pd.to_numeric(rentals['lama_menginap'], errors='coerce').fillna(0)

# Kategori yang akan di-encode (Tambah status_pasutri & status_kewarganegaraan)
cat_cols = ['jenis_sewa', 'metode_pembayaran', 'status_pasutri', 'status_kewarganegaraan']
numeric_cols = ['checkout_hour', 'checkin_hour', 'lama_menginap']

# Daftar seluruh fitur relevan yang akan digunakan (Tanpa membuang fitur)
feature_cols = cat_cols + numeric_cols

# Pastikan data kategorikal berupa string sebelum di-encode
for c in cat_cols:
    rentals[c] = rentals[c].fillna('Unknown').astype(str)

X = rentals[feature_cols].copy()
y = rentals['label']

# ------------------------------------------------------------
# 4. TRAIN-TEST SPLIT
# ------------------------------------------------------------
if len(y.unique()) < 2:
    raise ValueError("Data masih hanya 1 kelas! Periksa kembali kolom komentar.")

X_train, X_test, y_train, y_test = train_test_split(
    X, y, stratify=y, test_size=0.2, random_state=42
)

print("\nDistribusi TRAIN:")
print(y_train.value_counts())

# ------------------------------------------------------------
# 4B. ENCODING MENCEGAH DATA LEAKAGE
# ------------------------------------------------------------
# Simpan semua LabelEncoder ke dalam dictionary agar bisa di-export
label_encoders = {}

for c in cat_cols:
    le = LabelEncoder()
    # Fit & Transform hanya pada data latih
    X_train[c] = le.fit_transform(X_train[c])

    # Tangani unseen labels pada data uji (ubah jadi 'Unknown' jika tidak pernah dilihat di X_train)
    X_test[c] = X_test[c].map(lambda s: s if s in le.classes_ else 'Unknown')
    if 'Unknown' not in le.classes_:
        le.classes_ = np.append(le.classes_, 'Unknown')

    X_test[c] = le.transform(X_test[c])
    label_encoders[c] = le

# ------------------------------------------------------------
# 5. FEATURE EVALUATION (Chi-square, Pearson, MI)
# ------------------------------------------------------------
print("\n=== Evaluasi Signifikansi Fitur ===")
# Catatan: Pemodelan menggunakan seluruh kolom relevan yang tersedia dengan spesifik
# encoding (tidak otomatis membuang fitur). Pengujian ini untuk keperluan analisis.

chi2_selector = SelectKBest(score_func=chi2, k='all')
chi2_selector.fit(X_train, y_train)
chi2_scores = chi2_selector.scores_

train_df = X_train.copy()
train_df['label'] = y_train
pearson_scores = train_df.corr(method='pearson')['label'].drop('label').abs().values

mi_selector = SelectKBest(score_func=mutual_info_classif, k='all')
mi_selector.fit(X_train, y_train)
mi_scores = mi_selector.scores_

fs_df = pd.DataFrame({
    'Feature': X_train.columns,
    'Chi-Square': chi2_scores,
    'Pearson (Abs)': pearson_scores,
    'Mutual Info': mi_scores
})

# Normalisasi skor agar bisa dibandingkan
scaler = MinMaxScaler()
fs_df[['Chi2_Norm', 'Pearson_Norm', 'MI_Norm']] = scaler.fit_transform(
    fs_df[['Chi-Square', 'Pearson (Abs)', 'Mutual Info']]
)
fs_df['Total_Score'] = fs_df[['Chi2_Norm', 'Pearson_Norm', 'MI_Norm']].mean(axis=1)
fs_df = fs_df.sort_values(by='Total_Score', ascending=False).reset_index(drop=True)

print(fs_df[['Feature', 'Chi-Square', 'Pearson (Abs)', 'Mutual Info', 'Total_Score']])

# ------------------------------------------------------------
# 6. FUNGSI EVALUASI
# ------------------------------------------------------------
def evaluate(name, model, X_eval=X_test):
    y_pred = model.predict(X_eval)

    print(f"\n===== {name} =====")
    print("Classification Report:")
    print(classification_report(y_test, y_pred, digits=4))

    cm = confusion_matrix(y_test, y_pred)

    print("Confusion Matrix (Text):")
    print("                Predicted")
    print("               0       1")
    print("          -----------------")
    print(f"True   0 | {cm[0][0]:6d}  {cm[0][1]:6d}")
    print(f"       1 | {cm[1][0]:6d}  {cm[1][1]:6d}")
    print("          -----------------\n")

    # # Output Confusion Matrix berupa visualisasi
    # plt.figure(figsize=(4, 3))
    # sns.heatmap(cm, annot=True, fmt='d', cmap='Blues')
    # plt.title(f'Confusion Matrix: {name}')
    # plt.xlabel('Label Prediksi')
    # plt.ylabel('Label Asli')
    # plt.show()

# ------------------------------------------------------------
# 7. LATIH SEMUA MODEL (DENGAN CUSTOM WEIGHTS)
# ------------------------------------------------------------
models = {}

print("\nMemulai proses pelatihan untuk semua model dengan Custom Weights...\n")

# Kita atur bobot manual: Kelas 0 bobotnya 1, Kelas 1 bobotnya 4
# (Artinya AI diminta 4x lebih waspada terhadap pelanggaran)
custom_weight = {0: 1, 1: 4}

# --- 1. Decision Tree ---
print(">> Melatih Decision Tree...")
models["Decision Tree"] = DecisionTreeClassifier(max_depth=6, class_weight=custom_weight, random_state=42)

# --- 2. Random Forest ---
print(">> Melatih Random Forest...")
models["Random Forest"] = RandomForestClassifier(n_estimators=200, max_depth=7, class_weight=custom_weight, random_state=42)

# --- 3. KNN ---
print(">> Melatih KNN...")
# KNN tidak memiliki fitur class_weight, menggunakan pembobotan jarak
models["KNN"] = make_pipeline(StandardScaler(), KNeighborsClassifier(n_neighbors=5, weights='distance'))

# --- 4. SVM (RBF) ---
print(">> Melatih SVM (RBF)...")
models["SVM RBF"] = make_pipeline(StandardScaler(), SVC(kernel='rbf', probability=True, class_weight=custom_weight, random_state=42))

# --- 5. XGBoost ---
print(">> Melatih XGBoost...")
models["XGBoost"] = XGBClassifier(
    n_estimators=200,
    max_depth=5,
    learning_rate=0.05,
    scale_pos_weight=4,
    eval_metric='logloss',
    random_state=42
)

# ------------------------------------------------------------
# 8. EVALUASI SEMUA MODEL
# ------------------------------------------------------------
for name, model in models.items():
    model.fit(X_train, y_train)
    evaluate(name, model)
\end{lstlisting}

\begin{lstlisting}[caption={Output 9}]
rentals.csv
   id                nama               nik status_pasutri  \
0   1        F**** B*****  5***************        Menikah   
1   2      E***** P******  3***************        Menikah   
2   3        O**** H*****  3***************  Belum Menikah   
3   4  N***** B*****-B***  3***************  Belum Menikah   
4   5       R****** S****  3***************        Menikah   

  status_kewarganegaraan jenis_sewa unit_number metode_pembayaran metode_lain  \
0                    WNI   Mingguan      C-1-13              Cash         NaN   
1                    WNA     Harian       D-5-2      Kartu Kredit         NaN   
2                    WNI    Bulanan     B-25-17              QRIS         NaN   
3                    WNI    Bulanan       A-2-8            Others       gopay   
4                    WNI     Harian      A-21-9       Kartu Debit         NaN   

  tanggal_checkin waktu_checkin tanggal_checkout waktu_checkout  \
0      2024-06-20      20:22:00       2024-07-04       07:46:00   
1      2023-02-19      05:33:00       2023-02-20       12:33:00   
2      2023-12-17      17:48:00       2024-03-16       09:02:00   
3      2018-05-18      13:15:00       2018-10-15       12:45:00   
4      2025-01-05      17:28:00       2025-01-05       22:28:00   

   lama_menginap                                komentar  \
0             14                                     NaN   
1              1                                     NaN   
2             90                                     NaN   
3            150                                     NaN   
4              0  Ditemukan alat suntik di tempat sampah   

                     user_email diedit_oleh           created_at  \
0   O**************@hotmail.com         NaN  2025-12-08 19:13:06   
1  D***************@hotmail.com         NaN  2025-12-08 19:13:06   
2              K*****@gmail.com         NaN  2025-12-08 19:13:06   
3              T*****@gmail.com         NaN  2025-12-08 19:13:06   
4      K*************@yahoo.com         NaN  2025-12-08 19:13:06   

  classification  
0         normal  
1         normal  
2         normal  
3         normal  
4   tidak normal  

=== Distribusi Label Setelah FIX ===
label
0    9243
1     746
Name: count, dtype: int64
Proporsi pelanggaran: 0.0747

Distribusi TRAIN:
label
0    7394
1     597
Name: count, dtype: int64

=== Evaluasi Signifikansi Fitur ===
                  Feature    Chi-Square  Pearson (Abs)  Mutual Info  \
0           lama_menginap  22182.967372       0.217366     0.217217   
1            checkin_hour    262.418775       0.249692     0.051909   
2           checkout_hour     46.241510       0.103399     0.059958   
3              jenis_sewa      0.080931       0.003885     0.091195   
4  status_kewarganegaraan      0.147592       0.006097     0.001381   
5       metode_pembayaran      0.242216       0.005513     0.000436   
6          status_pasutri      0.036919       0.003023     0.002246   

   Total_Score  
0     0.956317  
1     0.416422  
2     0.227861  
3     0.140721  
4     0.005608  
5     0.003368  
6     0.002782  

Memulai proses pelatihan untuk semua model dengan Custom Weights...

>> Melatih Decision Tree...
>> Melatih Random Forest...
>> Melatih KNN...
>> Melatih SVM (RBF)...
>> Melatih XGBoost...

===== Decision Tree =====
Classification Report:
              precision    recall  f1-score   support

           0     1.0000    0.9822    0.9910      1849
           1     0.8187    1.0000    0.9003       149

    accuracy                         0.9835      1998
   macro avg     0.9093    0.9911    0.9456      1998
weighted avg     0.9865    0.9835    0.9842      1998

Confusion Matrix (Text):
                Predicted
               0       1
          -----------------
True   0 |   1816      33
       1 |      0     149
          -----------------


===== Random Forest =====
Classification Report:
              precision    recall  f1-score   support

           0     1.0000    0.9816    0.9907      1849
           1     0.8142    1.0000    0.8976       149

    accuracy                         0.9830      1998
   macro avg     0.9071    0.9908    0.9442      1998
weighted avg     0.9861    0.9830    0.9838      1998

Confusion Matrix (Text):
                Predicted
               0       1
          -----------------
True   0 |   1815      34
       1 |      0     149
          -----------------


===== KNN =====
Classification Report:
              precision    recall  f1-score   support

           0     0.9674    0.9794    0.9734      1849
           1     0.6984    0.5906    0.6400       149

    accuracy                         0.9505      1998
   macro avg     0.8329    0.7850    0.8067      1998
weighted avg     0.9474    0.9505    0.9485      1998

Confusion Matrix (Text):
                Predicted
               0       1
          -----------------
True   0 |   1811      38
       1 |     61      88
          -----------------


===== SVM RBF =====
Classification Report:
              precision    recall  f1-score   support

           0     0.9653    0.9789    0.9721      1849
           1     0.6829    0.5638    0.6176       149

    accuracy                         0.9479      1998
   macro avg     0.8241    0.7713    0.7949      1998
weighted avg     0.9443    0.9479    0.9456      1998

Confusion Matrix (Text):
                Predicted
               0       1
          -----------------
True   0 |   1810      39
       1 |     65      84
          -----------------


===== XGBoost =====
Classification Report:
              precision    recall  f1-score   support

           0     1.0000    0.9816    0.9907      1849
           1     0.8142    1.0000    0.8976       149

    accuracy                         0.9830      1998
   macro avg     0.9071    0.9908    0.9442      1998
weighted avg     0.9861    0.9830    0.9838      1998

Confusion Matrix (Text):
                Predicted
               0       1
          -----------------
True   0 |   1815      34
       1 |      0     149
          -----------------
\end{lstlisting}

\begin{lstlisting}[language=Python, caption={Kode 10}]
# ============================================================
# DATA SEED (VERSI FINAL: MERGE UNIT + N_SELECT EKSPLISIT)
# ============================================================

import pandas as pd
import numpy as np

from sklearn.model_selection import train_test_split
from sklearn.preprocessing import LabelEncoder, StandardScaler, MinMaxScaler
from sklearn.metrics import classification_report, confusion_matrix
from sklearn.pipeline import make_pipeline
from sklearn.feature_selection import SelectKBest, chi2, mutual_info_classif

from sklearn.tree import DecisionTreeClassifier
from sklearn.ensemble import RandomForestClassifier
from sklearn.neighbors import KNeighborsClassifier
from sklearn.svm import SVC
from xgboost import XGBClassifier

# ------------------------------------------------------------
# 1. LOAD DATA & PENGGABUNGAN
# ------------------------------------------------------------
rentals    = pd.read_csv('rentals.csv')
logs       = pd.read_csv('rental_logs.csv')
units      = pd.read_csv('units.csv')
users      = pd.read_csv('users.csv')
violations = pd.read_csv('violations.csv')

rentals.columns = rentals.columns.str.lower()
units.columns = units.columns.str.lower()

if 'komentar' not in rentals.columns:
    raise ValueError("Kolom 'komentar' tidak ditemukan di rentals.csv")

print("rentals.csv")
print(rentals.head())
print("logs.csv")
print(logs.head())
print("units.csv")
print(units.head())
print("users.csv")
print(users.head())
print("violations.csv")
print(violations.head())

# === PROSES VLOOKUP / MERGE ===
# Kita timpa kembali ke variabel 'rentals' agar penamaan konsisten dengan kode lama
rentals = pd.merge(rentals, units[['unit_number', 'tower', 'lantai']], on='unit_number', how='left')

# ------------------------------------------------------------
# 2. LABELING
# ------------------------------------------------------------
rentals['komentar'] = rentals['komentar'].astype(str).str.lower().str.strip()
empty_values = ["", "nan", "na", "null", "-", "0", "tidak ada", "none", " "]
rentals['label'] = (~rentals['komentar'].isin(empty_values)).astype(int)

# ------------------------------------------------------------
# 3. FEATURE ENGINEERING
# ------------------------------------------------------------
# Parse waktu checkout
rentals['waktu_checkout_dt'] = pd.to_datetime(rentals['waktu_checkout'], errors='coerce')
rentals['checkout_hour'] = rentals['waktu_checkout_dt'].dt.hour.fillna(0).astype(int)

# Parse waktu checkin
rentals['waktu_checkin_dt'] = pd.to_datetime(rentals['waktu_checkin'], errors='coerce')
rentals['checkin_hour'] = rentals['waktu_checkin_dt'].dt.hour.fillna(0).astype(int)

# Pastikan lama_menginap & lantai berupa angka (Fitur Baru)
rentals['lama_menginap'] = pd.to_numeric(rentals['lama_menginap'], errors='coerce').fillna(0)
rentals['lantai'] = pd.to_numeric(rentals['lantai'], errors='coerce').fillna(0)

# Kategori yang akan di-encode (Tambah status_pasutri, status_kewarganegaraan, tower)
cat_cols = ['jenis_sewa', 'metode_pembayaran', 'status_pasutri', 'status_kewarganegaraan', 'tower']
for c in cat_cols:
    # Jika ada data kosong, isi dengan 'Unknown'
    rentals[c] = rentals[c].fillna('Unknown').astype(str)
    le = LabelEncoder()
    rentals[c] = le.fit_transform(rentals[c])

# Gabungkan semua fitur
numeric_cols = ['checkout_hour', 'checkin_hour', 'lama_menginap', 'lantai']
feature_cols = cat_cols + numeric_cols

X = rentals[feature_cols].copy()
y = rentals['label']

# ------------------------------------------------------------
# 4. TRAIN-TEST SPLIT
# ------------------------------------------------------------
X_train, X_test, y_train, y_test = train_test_split(
    X, y, stratify=y, test_size=0.2, random_state=42
)

# ------------------------------------------------------------
# 5. FEATURE SELECTION
# ------------------------------------------------------------
print("\n=== Proses Feature Selection ===")
chi2_selector = SelectKBest(score_func=chi2, k='all')
chi2_selector.fit(X_train, y_train)

train_df = X_train.copy()
train_df['label'] = y_train
pearson_scores = train_df.corr(method='pearson')['label'].drop('label').abs().values

mi_selector = SelectKBest(score_func=mutual_info_classif, k='all')
mi_selector.fit(X_train, y_train)

fs_df = pd.DataFrame({
    'Feature': X_train.columns,
    'Chi-Square': chi2_selector.scores_,
    'Pearson (Abs)': pearson_scores,
    'Mutual Info': mi_selector.scores_
})

scaler = MinMaxScaler()
fs_df[['Chi2_Norm', 'Pearson_Norm', 'MI_Norm']] = scaler.fit_transform(
    fs_df[['Chi-Square', 'Pearson (Abs)', 'Mutual Info']]
)
fs_df['Total_Score'] = fs_df[['Chi2_Norm', 'Pearson_Norm', 'MI_Norm']].mean(axis=1)
fs_df = fs_df.sort_values(by='Total_Score', ascending=False).reset_index(drop=True)

# JUMLAH FITUR
n_select = 7
selected_features = fs_df['Feature'].head(n_select).tolist()
print(f"\n-> Fitur yang dipilih ({n_select} terbaik): {selected_features}")

X_train = X_train[selected_features]
X_test = X_test[selected_features]

# ------------------------------------------------------------
# 6. Fungsi Evaluasi
# ------------------------------------------------------------
def evaluate(name, model):
    y_pred = model.predict(X_test)
    print(f"\n===== {name} =====")
    print("Classification Report:")
    print(classification_report(y_test, y_pred, digits=4))
    cm = confusion_matrix(y_test, y_pred)
    print("Confusion Matrix (Text):")
    print("                Predicted")
    print("               0       1")
    print("          -----------------")
    print(f"True   0 | {cm[0][0]:6d}  {cm[0][1]:6d}")
    print(f"       1 | {cm[1][0]:6d}  {cm[1][1]:6d}")
    print("          -----------------\n")

# ------------------------------------------------------------
# 7. Latih Semua Model (DENGAN CUSTOM WEIGHTS)
# ------------------------------------------------------------
models = {}
custom_weight = {0: 1, 1: 4}

models["Decision Tree"] = DecisionTreeClassifier(max_depth=6, class_weight=custom_weight, random_state=42)
models["Random Forest"] = RandomForestClassifier(n_estimators=200, max_depth=7, class_weight=custom_weight, random_state=42)
models["KNN"] = make_pipeline(StandardScaler(), KNeighborsClassifier(n_neighbors=5, weights='distance'))
models["SVM RBF"] = make_pipeline(StandardScaler(), SVC(kernel='rbf', probability=True, class_weight=custom_weight, random_state=42))
models["XGBoost"] = XGBClassifier(n_estimators=200, max_depth=5, learning_rate=0.05, scale_pos_weight=4, eval_metric='logloss', random_state=42)

# Latih dan Evaluasi
for name, model in models.items():
    model.fit(X_train, y_train)
    evaluate(name, model)
\end{lstlisting}

\begin{lstlisting}[caption={Output 10}]
rentals.csv
   id                nama               nik status_pasutri  \
0   1        F**** B*****  5***************        Menikah   
1   2      E***** P******  3***************        Menikah   
2   3        O**** H*****  3***************  Belum Menikah   
3   4  N***** B*****-B***  3***************  Belum Menikah   
4   5       R****** S****  3***************        Menikah   

  status_kewarganegaraan jenis_sewa unit_number metode_pembayaran metode_lain  \
0                    WNI   Mingguan      C-1-13              Cash         NaN   
1                    WNA     Harian       D-5-2      Kartu Kredit         NaN   
2                    WNI    Bulanan     B-25-17              QRIS         NaN   
3                    WNI    Bulanan       A-2-8            Others       gopay   
4                    WNI     Harian      A-21-9       Kartu Debit         NaN   

  tanggal_checkin waktu_checkin tanggal_checkout waktu_checkout  \
0      2024-06-20      20:22:00       2024-07-04       07:46:00   
1      2023-02-19      05:33:00       2023-02-20       12:33:00   
2      2023-12-17      17:48:00       2024-03-16       09:02:00   
3      2018-05-18      13:15:00       2018-10-15       12:45:00   
4      2025-01-05      17:28:00       2025-01-05       22:28:00   

   lama_menginap                                komentar  \
0             14                                     NaN   
1              1                                     NaN   
2             90                                     NaN   
3            150                                     NaN   
4              0  Ditemukan alat suntik di tempat sampah   

                     user_email diedit_oleh           created_at  \
0   O**************@hotmail.com         NaN  2025-12-08 19:13:06   
1  D***************@hotmail.com         NaN  2025-12-08 19:13:06   
2              K*****@gmail.com         NaN  2025-12-08 19:13:06   
3              T*****@gmail.com         NaN  2025-12-08 19:13:06   
4      K*************@yahoo.com         NaN  2025-12-08 19:13:06   

  classification  
0         normal  
1         normal  
2         normal  
3         normal  
4   tidak normal  
logs.csv
   id  rental_id  action field_changed old_value  \
0   1          1  create           all       NaN   
1   2          2  create           all       NaN   
2   3          3  create           all       NaN   
3   4          4  create           all       NaN   
4   5          5  create           all       NaN   

                                           new_value  \
0  {"id":1,"nama":"F**** B*****","nik":"5*******************"id":2,"nama":"E***** P******","nik":"3*****************"id":3,"nama":"O**** H*****","nik":"3*******************"id":4,"nama":"N***** B*****-B***","nik":"3*************"id":5,"nama":"R****** S****","nik":"33966831...   

                          email            timestamp  
0   O**************@hotmail.com  2025-12-08 19:13:06  
1  D***************@hotmail.com  2025-12-08 19:13:06  
2              K*****@gmail.com  2025-12-08 19:13:06  
3              T*****@gmail.com  2025-12-08 19:13:06  
4      K*************@yahoo.com  2025-12-08 19:13:06  
units.csv
  unit_number                   user_email tower  lantai  unit
0      A-1-14     K*************@yahoo.com     A       1    14
1      A-1-18         A*******@hotmail.com     A       1    18
2      A-1-22  O**************@hotmail.com     A       1    22
3      A-1-24                p**@gmail.com     A       1    24
4      A-1-28  O**************@hotmail.com     A       1    28
users.csv
   id                         email         password           nib  role  \
0   1    Y***************@yahoo.com  oe4eaGaBLpQv76y  7.265157e+12  agen   
1   2           R******@hotmail.com  UhoFvefkUMB44Cj  2.780507e+12  agen   
2   3  F*****************@yahoo.com  Jf5D9lZJIGVf011  1.335353e+12  agen   
3   4       M************@yahoo.com  DYPYEZs_Zh17IjQ  8.262239e+12  agen   
4   5         E**********@gmail.com  1gcfQuFUGjUATHy  8.127913e+11  agen   

               name           created_at  
0        K** M*** I  2025-12-08 19:13:05  
1    D*. J**** L***  2025-12-08 19:13:05  
2  R******** G*****  2025-12-08 19:13:05  
3  B******* K******  2025-12-08 19:13:05  
4      M****** D***  2025-12-08 19:13:05  
violations.csv
   id  rental_id                                       photo_url  \
0   1        233      https://picsum.photos/seed/u717oX/983/2412   
1   2       7182      https://picsum.photos/seed/0cZJtXQ/243/249   
2   3       9117     https://picsum.photos/seed/Pl0nan/3034/2819   
3   4       8582  https://picsum.photos/seed/LKFORW1qWE/3127/725   
4   5       4165    https://picsum.photos/seed/dkpBkZPg/336/1118   

                                         description  \
0              Crepusculum torrens vomica atque rem.   
1  Inventore articulus corrumpo pax clamo absterg...   
2  Nostrum desino comitatus debilito perferendis ...   
3  Tremo accedo aegre votum adhuc aveho vilicus v...   
4                         Comminor vis thorax absum.   

                         uploaded_by           created_at  
0     L*******************@yahoo.com  2025-12-08 19:13:07  
1        O**************@hotmail.com  2025-12-08 19:13:07  
2       A*****************@gmail.com  2025-12-08 19:13:07  
3                  S******@gmail.com  2025-12-08 19:13:07  
4  I********************@hotmail.com  2025-12-08 19:13:07  

=== Proses Feature Selection ===

/tmp/ipykernel_3999/3749411096.py:61: UserWarning: Could not infer format, so each element will be parsed individually, falling back to `dateutil`. To ensure parsing is consistent and as-expected, please specify a format.
  rentals['waktu_checkout_dt'] = pd.to_datetime(rentals['waktu_checkout'], errors='coerce')
/tmp/ipykernel_3999/3749411096.py:65: UserWarning: Could not infer format, so each element will be parsed individually, falling back to `dateutil`. To ensure parsing is consistent and as-expected, please specify a format.
  rentals['waktu_checkin_dt'] = pd.to_datetime(rentals['waktu_checkin'], errors='coerce')


-> Fitur yang dipilih (7 terbaik): ['lama_menginap', 'checkin_hour', 'checkout_hour', 'jenis_sewa', 'tower', 'status_kewarganegaraan', 'lantai']

===== Decision Tree =====
Classification Report:
              precision    recall  f1-score   support

           0     0.9983    0.9816    0.9899      1849
           1     0.8111    0.9799    0.8875       149

    accuracy                         0.9815      1998
   macro avg     0.9047    0.9807    0.9387      1998
weighted avg     0.9844    0.9815    0.9823      1998

Confusion Matrix (Text):
                Predicted
               0       1
          -----------------
True   0 |   1815      34
       1 |      3     146
          -----------------


===== Random Forest =====
Classification Report:
              precision    recall  f1-score   support

           0     1.0000    0.9816    0.9907      1849
           1     0.8142    1.0000    0.8976       149

    accuracy                         0.9830      1998
   macro avg     0.9071    0.9908    0.9442      1998
weighted avg     0.9861    0.9830    0.9838      1998

Confusion Matrix (Text):
                Predicted
               0       1
          -----------------
True   0 |   1815      34
       1 |      0     149
          -----------------


===== KNN =====
Classification Report:
              precision    recall  f1-score   support

           0     0.9568    0.9816    0.9690      1849
           1     0.6634    0.4497    0.5360       149

    accuracy                         0.9419      1998
   macro avg     0.8101    0.7156    0.7525      1998
weighted avg     0.9349    0.9419    0.9367      1998

Confusion Matrix (Text):
                Predicted
               0       1
          -----------------
True   0 |   1815      34
       1 |     82      67
          -----------------


===== SVM RBF =====
Classification Report:
              precision    recall  f1-score   support

           0     0.9671    0.9843    0.9756      1849
           1     0.7500    0.5839    0.6566       149

    accuracy                         0.9545      1998
   macro avg     0.8585    0.7841    0.8161      1998
weighted avg     0.9509    0.9545    0.9518      1998

Confusion Matrix (Text):
                Predicted
               0       1
          -----------------
True   0 |   1820      29
       1 |     62      87
          -----------------


===== XGBoost =====
Classification Report:
              precision    recall  f1-score   support

           0     0.9994    0.9816    0.9905      1849
           1     0.8132    0.9933    0.8943       149

    accuracy                         0.9825      1998
   macro avg     0.9063    0.9875    0.9424      1998
weighted avg     0.9856    0.9825    0.9833      1998

Confusion Matrix (Text):
                Predicted
               0       1
          -----------------
True   0 |   1815      34
       1 |      1     148
          -----------------
\end{lstlisting}

\begin{lstlisting}[language=Python, caption={Kode 11}]
# ============================================================
# DATA SEED (TANPA TUNING)
# ============================================================

import pandas as pd
import numpy as np
import joblib

from sklearn.model_selection import train_test_split
from sklearn.preprocessing import LabelEncoder, StandardScaler, MinMaxScaler
from sklearn.metrics import classification_report, confusion_matrix
from sklearn.pipeline import make_pipeline
from sklearn.feature_selection import SelectKBest, chi2, mutual_info_classif

from sklearn.tree import DecisionTreeClassifier
from sklearn.ensemble import RandomForestClassifier
from sklearn.neighbors import KNeighborsClassifier
from sklearn.svm import SVC
from xgboost import XGBClassifier

import matplotlib.pyplot as plt
import seaborn as sns

# ------------------------------------------------------------
# 1. LOAD DATA
# ------------------------------------------------------------
rentals    = pd.read_csv('rentals.csv')
logs       = pd.read_csv('rental_logs.csv')
units      = pd.read_csv('units.csv')
users      = pd.read_csv('users.csv')
violations = pd.read_csv('violations.csv')

rentals.columns = rentals.columns.str.lower()

if 'komentar' not in rentals.columns:
    raise ValueError("Kolom 'komentar' tidak ditemukan di rentals.csv")

print("rentals.csv")
print(rentals.head())
print("logs.csv")
print(logs.head())
print("units.csv")
print(units.head())
print("users.csv")
print(users.head())
print("violations.csv")
print(violations.head())

# ------------------------------------------------------------
# 2. LABELING
# ------------------------------------------------------------
rentals['komentar'] = rentals['komentar'].astype(str).str.lower().str.strip()
empty_values = ["", "nan", "na", "null", "-", "0", "tidak ada", "none", " "]
rentals['label'] = (~rentals['komentar'].isin(empty_values)).astype(int)


# ============================================================
# EXPLORATORY DATA ANALYSIS (EDA)
# ============================================================
print("\n" + "="*50)
print("=== EXPLORATORY DATA ANALYSIS (EDA) ===")
print("="*50)

df_eda = rentals.copy()

print("\n[1] Statistik Deskriptif:")
print(df_eda.describe(include='all'))

print("\n[2] Distribusi Kelas Target (0: Aman, 1: Melanggar):")
print(df_eda['label'].value_counts())

fig, ax = plt.subplots(figsize=(6, 4))
sns.countplot(data=df_eda, x='label', palette='Set2', ax=ax)
plt.title('Distribusi Kelas Klasifikasi (0: Aman, 1: Melanggar)')
plt.xlabel('Klasifikasi')
plt.ylabel('Jumlah Penyewaan')

y_max = df_eda['label'].value_counts().max()
plt.ylim(0, y_max * 1.15)

for p in ax.patches:
    if p.get_height() > 0:
        ax.annotate(f'{int(p.get_height())}', (p.get_x() + p.get_width() / 2., p.get_height()),
                    ha='center', va='baseline', fontsize=10, color='black', xytext=(0, 5),
                    textcoords='offset points')
plt.show()

kolom_dikecualikan = ['id', 'nama', 'nik', 'unit_number', 'tanggal_checkin', 'waktu_checkin',
                      'tanggal_checkout', 'waktu_checkout', 'komentar', 'user_email', 'diedit_oleh', 'created_at', 'metode_lain', 'classification']

semua_kategori = [col for col in df_eda.select_dtypes(include=['object', 'category']).columns if col not in kolom_dikecualikan]

print("\n[3] Rata-rata Pelanggaran Berdasarkan Seluruh Kategori:")
for col in semua_kategori:
    if col in df_eda.columns:
        print(f"\n--- Tingkat Pelanggaran Berdasarkan {col} ---")
        print(df_eda.groupby(col)['label'].mean().sort_values(ascending=False).round(3))

df_eda['lama_menginap_num'] = pd.to_numeric(df_eda['lama_menginap'], errors='coerce').fillna(0)
df_eda['checkin_hour_num'] = pd.to_datetime(df_eda['waktu_checkin'], errors='coerce').dt.hour.fillna(0)
df_eda['checkout_hour_num'] = pd.to_datetime(df_eda['waktu_checkout'], errors='coerce').dt.hour.fillna(0)

le_eda = LabelEncoder()
for c in semua_kategori:
    if c in df_eda.columns:
        df_eda[c + '_encoded'] = le_eda.fit_transform(df_eda[c].astype(str))

print("\n[4] Korelasi Antar Fitur:")
kolom_korelasi = ['lama_menginap_num', 'checkin_hour_num', 'checkout_hour_num', 'label']
kolom_korelasi += [c + '_encoded' for c in semua_kategori if c + '_encoded' in df_eda.columns]

corr_matrix = df_eda[kolom_korelasi].corr()

plt.figure(figsize=(10, 8))
sns.heatmap(corr_matrix, annot=True, cmap='coolwarm', fmt='.2f', linewidths=0.5)
plt.title('Heatmap Korelasi Antar Fitur dan Target (Klasifikasi)')
plt.tight_layout()
plt.show()

print("\n[5] Visualisasi Pola Penyewaan (Dikaitkan dengan Klasifikasi):")

custom_line_palette = {0: 'blue', 1: 'red'}
custom_labels = ['0: Aman', '1: Melanggar']

fig, ax1 = plt.subplots(figsize=(12, 5))
waktu_df = df_eda.groupby(['checkin_hour_num', 'label']).size().reset_index(name='jumlah')

sns.lineplot(data=waktu_df, x='checkin_hour_num', y='jumlah', hue='label', marker='o', palette=custom_line_palette, ax=ax1)

plt.title('Pola Waktu Kedatangan Berdasarkan Klasifikasi')
plt.xlabel('Jam Check-in (0-23)')
plt.ylabel('Jumlah Penyewaan')
plt.xticks(range(0, 24))
plt.grid(True, alpha=0.3)

y_max_line = waktu_df['jumlah'].max()
plt.ylim(-y_max_line * 0.20, y_max_line * 1.15)

handles, _ = ax1.get_legend_handles_labels()
plt.legend(handles=handles, labels=custom_labels, title='Klasifikasi')

for l in waktu_df['label'].unique():
    sub = waktu_df[waktu_df['label'] == l]
    warna_teks = custom_line_palette[l]

    for x, y in zip(sub['checkin_hour_num'], sub['jumlah']):
        y_offset = 8 if l == 0 else -12
        va_align = 'bottom' if l == 0 else 'top'
        plt.annotate(f'{int(y)}', (x, y), textcoords="offset points", xytext=(0, y_offset),
                     ha='center', va=va_align, fontsize=9, color=warna_teks, fontweight='bold')
plt.show()

fig, ax2 = plt.subplots(figsize=(12, 5))
lama_df = df_eda.groupby(['lama_menginap_num', 'label']).size().reset_index(name='jumlah')

sns.lineplot(data=lama_df, x='lama_menginap_num', y='jumlah', hue='label', marker='o', palette=custom_line_palette, ax=ax2)

plt.title('Tren Lama Menginap Berdasarkan Klasifikasi')
plt.xlabel('Lama Menginap (Hari)')
plt.ylabel('Jumlah Penyewaan')
plt.grid(True, alpha=0.3)

handles2, _ = ax2.get_legend_handles_labels()
plt.legend(handles=handles2, labels=custom_labels, title='Klasifikasi')
plt.show()

for col in semua_kategori:
    if col in df_eda.columns:
        fig, ax = plt.subplots(figsize=(8, 5))
        sns.countplot(data=df_eda, x=col, hue='label', palette='magma', ax=ax)

        judul_kolom = col.replace('_', ' ').title()
        plt.title(f'Perbandingan Jumlah Pelanggaran Berdasarkan {judul_kolom}')
        plt.xlabel(judul_kolom)
        plt.ylabel('Jumlah Penyewaan')
        plt.xticks(rotation=0)

        handles_bar, _ = ax.get_legend_handles_labels()
        plt.legend(handles=handles_bar, labels=custom_labels, title='Klasifikasi')

        y_max_bar = df_eda.groupby([col, 'label']).size().max()
        plt.ylim(0, y_max_bar * 1.15)

        for p in ax.patches:
            if p.get_height() > 0:
                ax.annotate(f'{int(p.get_height())}', (p.get_x() + p.get_width() / 2., p.get_height()),
                            ha='center', va='baseline', fontsize=9, color='black', xytext=(0, 5),
                            textcoords='offset points')
        plt.tight_layout()
        plt.show()


print("\n[5b] Visualisasi Metode Pembayaran 'Others' (metode_lain):")

# Memastikan kolom 'metode_lain' berbentuk string dan membersihkan spasi/nilai kosong
df_eda['metode_lain'] = df_eda['metode_lain'].astype(str).str.strip().str.title()
empty_lain = ["", "Nan", "Na", "Null", "-", "0", "None"]

# Filter hanya baris yang memiliki nilai di 'metode_lain'
df_others = df_eda[~df_eda['metode_lain'].isin(empty_lain)]

if not df_others.empty:
    fig, ax = plt.subplots(figsize=(10, 6))

    # Menghitung frekuensi masing-masing metode pembayaran lain
    order_others = df_others['metode_lain'].value_counts().index

    sns.countplot(
        data=df_others,
        y='metode_lain',
        order=order_others,
        palette='viridis',
        ax=ax
    )

    plt.title('Detail Metode Pembayaran "Others" (metode_lain)')
    plt.xlabel('Jumlah Penyewaan')
    plt.ylabel('Jenis Pembayaran Lainnya')

    # Menambahkan anotasi angka di sebelah setiap bar
    for p in ax.patches:
        width = p.get_width()
        if width > 0:
            ax.annotate(f'{int(width)}',
                        (width, p.get_y() + p.get_height() / 2.),
                        ha='left', va='center', fontsize=10, color='black',
                        xytext=(5, 0), textcoords='offset points')

    # Perlebar batas x-axis sedikit agar angka anotasi tidak terpotong
    x_max = df_others['metode_lain'].value_counts().max()
    plt.xlim(0, x_max * 1.15)

    plt.tight_layout()
    plt.show()
else:
    print("Tidak ada data metode pembayaran 'Others' yang valid untuk ditampilkan.")


# ------------------------------------------------------------
# 3. FEATURE ENGINEERING & PREPARATION
# ------------------------------------------------------------
rentals['waktu_checkout_dt'] = pd.to_datetime(rentals['waktu_checkout'], errors='coerce')
rentals['checkout_hour'] = rentals['waktu_checkout_dt'].dt.hour.fillna(0).astype(int)

rentals['waktu_checkin_dt'] = pd.to_datetime(rentals['waktu_checkin'], errors='coerce')
rentals['checkin_hour'] = rentals['waktu_checkin_dt'].dt.hour.fillna(0).astype(int)

rentals['lama_menginap'] = pd.to_numeric(rentals['lama_menginap'], errors='coerce').fillna(0)

cat_cols = ['jenis_sewa', 'metode_pembayaran', 'status_pasutri', 'status_kewarganegaraan']
for c in cat_cols:
    rentals[c] = rentals[c].fillna('Unknown').astype(str)

numeric_cols = ['checkout_hour', 'checkin_hour', 'lama_menginap']
feature_cols = cat_cols + numeric_cols

X = rentals[feature_cols].copy()
y = rentals['label']

# ------------------------------------------------------------
# 4. TRAIN-TEST SPLIT
# ------------------------------------------------------------
X_train, X_test, y_train, y_test = train_test_split(
    X, y, stratify=y, test_size=0.2, random_state=42
)

# ------------------------------------------------------------
# 4B. ENCODING MENCEGAH DATA LEAKAGE
# ------------------------------------------------------------
label_encoders = {}

for c in cat_cols:
    le = LabelEncoder()
    X_train[c] = le.fit_transform(X_train[c])

    X_test[c] = X_test[c].map(lambda s: s if s in le.classes_ else 'Unknown')
    if 'Unknown' not in le.classes_:
        le.classes_ = np.append(le.classes_, 'Unknown')

    X_test[c] = le.transform(X_test[c])

    label_encoders[c] = le

# ------------------------------------------------------------
# 5. FEATURE SELECTION
# ------------------------------------------------------------
print("\n=== Proses Feature Selection ===")
chi2_selector = SelectKBest(score_func=chi2, k='all')
chi2_selector.fit(X_train, y_train)

train_df = X_train.copy()
train_df['label'] = y_train
pearson_scores = train_df.corr(method='pearson')['label'].drop('label').abs().values

mi_selector = SelectKBest(score_func=mutual_info_classif, k='all')
mi_selector.fit(X_train, y_train)

fs_df = pd.DataFrame({
    'Feature': X_train.columns,
    'Chi-Square': chi2_selector.scores_,
    'Pearson (Abs)': pearson_scores,
    'Mutual Info': mi_selector.scores_
})

scaler = MinMaxScaler()
fs_df[['Chi2_Norm', 'Pearson_Norm', 'MI_Norm']] = scaler.fit_transform(
    fs_df[['Chi-Square', 'Pearson (Abs)', 'Mutual Info']]
)
fs_df['Total_Score'] = fs_df[['Chi2_Norm', 'Pearson_Norm', 'MI_Norm']].mean(axis=1)
fs_df = fs_df.sort_values(by='Total_Score', ascending=False).reset_index(drop=True)

# Jumlah Fitur
n_select = 4
selected_features = fs_df['Feature'].head(n_select).tolist()
print(f"\n-> Fitur yang dipilih ({n_select} terbaik): {selected_features}")

X_train = X_train[selected_features]
X_test = X_test[selected_features]

# ------------------------------------------------------------
# 6. Fungsi Evaluasi
# ------------------------------------------------------------
def evaluate(name, model):
    y_pred = model.predict(X_test)
    print(f"\n===== {name} =====")
    print("Classification Report:")
    print(classification_report(y_test, y_pred, digits=4))
    cm = confusion_matrix(y_test, y_pred)
    print("Confusion Matrix (Text):")
    print("                Predicted")
    print("               0       1")
    print("          -----------------")
    print(f"True   0 | {cm[0][0]:6d}  {cm[0][1]:6d}")
    print(f"       1 | {cm[1][0]:6d}  {cm[1][1]:6d}")
    print("          -----------------\n")

# ------------------------------------------------------------
# 7. Latih Semua Model (TANPA TUNING)
# ------------------------------------------------------------
models = {}
custom_weight = {0: 1, 1: 4}

# Model menggunakan pengaturan tetap
models["Decision Tree"] = DecisionTreeClassifier(max_depth=6, class_weight=custom_weight, random_state=42)
models["Random Forest"] = RandomForestClassifier(n_estimators=200, max_depth=7, class_weight=custom_weight, random_state=42)
models["KNN"] = make_pipeline(StandardScaler(), KNeighborsClassifier(n_neighbors=5, weights='distance'))
models["SVM RBF"] = make_pipeline(StandardScaler(), SVC(kernel='rbf', probability=True, class_weight=custom_weight, random_state=42))
models["XGBoost"] = XGBClassifier(n_estimators=200, max_depth=5, learning_rate=0.05, scale_pos_weight=4, eval_metric='logloss', random_state=42)

for name, model in models.items():
    model.fit(X_train, y_train)
    evaluate(name, model)

# ------------------------------------------------------------
# 8. Eksport Model Terbaik & Label Encoders
# ------------------------------------------------------------
print("\n=== Proses Export Model ===")

export_package = {
    'model': models["Random Forest"],
    'label_encoders': label_encoders,
    'selected_features': selected_features
}

filename = 'model_random_forest.pkl'
joblib.dump(export_package, filename)
print(f"Sukses! Model dan Label Encoders telah disimpan sebagai file: {filename}")
\end{lstlisting}

\begin{lstlisting}[caption={Output 11}]
rentals.csv
   id                nama               nik status_pasutri  \
0   1        F**** B*****  5***************        Menikah   
1   2      E***** P******  3***************        Menikah   
2   3        O**** H*****  3***************  Belum Menikah   
3   4  N***** B*****-B***  3***************  Belum Menikah   
4   5       R****** S****  3***************        Menikah   

  status_kewarganegaraan jenis_sewa unit_number metode_pembayaran metode_lain  \
0                    WNI   Mingguan      C-1-13              Cash         NaN   
1                    WNA     Harian       D-5-2      Kartu Kredit         NaN   
2                    WNI    Bulanan     B-25-17              QRIS         NaN   
3                    WNI    Bulanan       A-2-8            Others       gopay   
4                    WNI     Harian      A-21-9       Kartu Debit         NaN   

  tanggal_checkin waktu_checkin tanggal_checkout waktu_checkout  \
0      2024-06-20      20:22:00       2024-07-04       07:46:00   
1      2023-02-19      05:33:00       2023-02-20       12:33:00   
2      2023-12-17      17:48:00       2024-03-16       09:02:00   
3      2018-05-18      13:15:00       2018-10-15       12:45:00   
4      2025-01-05      17:28:00       2025-01-05       22:28:00   

   lama_menginap                                komentar  \
0             14                                     NaN   
1              1                                     NaN   
2             90                                     NaN   
3            150                                     NaN   
4              0  Ditemukan alat suntik di tempat sampah   

                     user_email diedit_oleh           created_at  \
0   O**************@hotmail.com         NaN  2025-12-08 19:13:06   
1  D***************@hotmail.com         NaN  2025-12-08 19:13:06   
2              K*****@gmail.com         NaN  2025-12-08 19:13:06   
3              T*****@gmail.com         NaN  2025-12-08 19:13:06   
4      K*************@yahoo.com         NaN  2025-12-08 19:13:06   

  classification  
0         normal  
1         normal  
2         normal  
3         normal  
4   tidak normal  
logs.csv
   id  rental_id  action field_changed old_value  \
0   1          1  create           all       NaN   
1   2          2  create           all       NaN   
2   3          3  create           all       NaN   
3   4          4  create           all       NaN   
4   5          5  create           all       NaN   

                                           new_value  \
0  {"id":1,"nama":"F**** B*****","nik":"5*******************"id":2,"nama":"E***** P******","nik":"3*****************"id":3,"nama":"O**** H*****","nik":"3*******************"id":4,"nama":"N***** B*****-B***","nik":"3*************"id":5,"nama":"R****** S****","nik":"33966831...   

                          email            timestamp  
0   O**************@hotmail.com  2025-12-08 19:13:06  
1  D***************@hotmail.com  2025-12-08 19:13:06  
2              K*****@gmail.com  2025-12-08 19:13:06  
3              T*****@gmail.com  2025-12-08 19:13:06  
4      K*************@yahoo.com  2025-12-08 19:13:06  
units.csv
  unit_number                   user_email tower  lantai  unit
0      A-1-14     K*************@yahoo.com     A       1    14
1      A-1-18         A*******@hotmail.com     A       1    18
2      A-1-22  O**************@hotmail.com     A       1    22
3      A-1-24                p**@gmail.com     A       1    24
4      A-1-28  O**************@hotmail.com     A       1    28
users.csv
   id                         email         password           nib  role  \
0   1    Y***************@yahoo.com  oe4eaGaBLpQv76y  7.265157e+12  agen   
1   2           R******@hotmail.com  UhoFvefkUMB44Cj  2.780507e+12  agen   
2   3  F*****************@yahoo.com  Jf5D9lZJIGVf011  1.335353e+12  agen   
3   4       M************@yahoo.com  DYPYEZs_Zh17IjQ  8.262239e+12  agen   
4   5         E**********@gmail.com  1gcfQuFUGjUATHy  8.127913e+11  agen   

               name           created_at  
0        K** M*** I  2025-12-08 19:13:05  
1    D*. J**** L***  2025-12-08 19:13:05  
2  R******** G*****  2025-12-08 19:13:05  
3  B******* K******  2025-12-08 19:13:05  
4      M****** D***  2025-12-08 19:13:05  
violations.csv
   id  rental_id                                       photo_url  \
0   1        233      https://picsum.photos/seed/u717oX/983/2412   
1   2       7182      https://picsum.photos/seed/0cZJtXQ/243/249   
2   3       9117     https://picsum.photos/seed/Pl0nan/3034/2819   
3   4       8582  https://picsum.photos/seed/LKFORW1qWE/3127/725   
4   5       4165    https://picsum.photos/seed/dkpBkZPg/336/1118   

                                         description  \
0              Crepusculum torrens vomica atque rem.   
1  Inventore articulus corrumpo pax clamo absterg...   
2  Nostrum desino comitatus debilito perferendis ...   
3  Tremo accedo aegre votum adhuc aveho vilicus v...   
4                         Comminor vis thorax absum.   

                         uploaded_by           created_at  
0     L*******************@yahoo.com  2025-12-08 19:13:07  
1        O**************@hotmail.com  2025-12-08 19:13:07  
2       A*****************@gmail.com  2025-12-08 19:13:07  
3                  S******@gmail.com  2025-12-08 19:13:07  
4  I********************@hotmail.com  2025-12-08 19:13:07  

==================================================
=== EXPLORATORY DATA ANALYSIS (EDA) ===
==================================================

[1] Statistik Deskriptif:
                  id           nama           nik status_pasutri  \
count    9989.000000           9989  9.989000e+03           9989   
unique           NaN           9894           NaN              2   
top              NaN  R*** D*******           NaN        Menikah   
freq             NaN              3           NaN           5004   
mean     5000.496746            NaN  3.649718e+15            NaN   
std      2887.229484            NaN  6.341279e+14            NaN   
min         1.000000            NaN  3.100037e+15            NaN   
25%      2501.000000            NaN  3.274772e+15            NaN   
50%      5001.000000            NaN  3.448126e+15            NaN   
75%      7500.000000            NaN  3.624302e+15            NaN   
max     10000.000000            NaN  5.199981e+15            NaN   

       status_kewarganegaraan jenis_sewa unit_number metode_pembayaran  \
count                    9989       9989        9989              9989   
unique                      2          3         794                 5   
top                       WNI   Mingguan     B-13-18      Kartu Kredit   
freq                     5011       3466          25              2045   
mean                      NaN        NaN         NaN               NaN   
std                       NaN        NaN         NaN               NaN   
min                       NaN        NaN         NaN               NaN   
25%                       NaN        NaN         NaN               NaN   
50%                       NaN        NaN         NaN               NaN   
75%                       NaN        NaN         NaN               NaN   
max                       NaN        NaN         NaN               NaN   

       metode_lain tanggal_checkin waktu_checkin tanggal_checkout  \
count         2000            9989          9989             9989   
unique           7            4742           754             4791   
top       paylater      2021-11-18      16:20:00       2019-01-29   
freq           297               8            34                9   
mean           NaN             NaN           NaN              NaN   
std            NaN             NaN           NaN              NaN   
min            NaN             NaN           NaN              NaN   
25%            NaN             NaN           NaN              NaN   
50%            NaN             NaN           NaN              NaN   
75%            NaN             NaN           NaN              NaN   
max            NaN             NaN           NaN              NaN   

       waktu_checkout  lama_menginap komentar        user_email  \
count            9989    9989.000000     9989              9989   
unique            659            NaN       11                53   
top          07:21:00            NaN      nan  K*****@gmail.com   
freq               43            NaN     9243               223   
mean              NaN      35.655321      NaN               NaN   
std               NaN      45.875125      NaN               NaN   
min               NaN       0.000000      NaN               NaN   
25%               NaN       5.000000      NaN               NaN   
50%               NaN      14.000000      NaN               NaN   
75%               NaN      60.000000      NaN               NaN   
max               NaN     150.000000      NaN               NaN   

                        diedit_oleh           created_at classification  \
count                            51                 9989           9989   
unique                           35                    1              2   
top     J**************@hotmail.com  2025-12-08 19:13:06         normal   
freq                              3                 9989           9243   
mean                            NaN                  NaN            NaN   
std                             NaN                  NaN            NaN   
min                             NaN                  NaN            NaN   
25%                             NaN                  NaN            NaN   
50%                             NaN                  NaN            NaN   
75%                             NaN                  NaN            NaN   
max                             NaN                  NaN            NaN   

              label  
count   9989.000000  
unique          NaN  
top             NaN  
freq            NaN  
mean       0.074682  
std        0.262891  
min        0.000000  
25%        0.000000  
50%        0.000000  
75%        0.000000  
max        1.000000  

[2] Distribusi Kelas Target (0: Aman, 1: Melanggar):
label
0    9243
1     746
Name: count, dtype: int64

/tmp/ipykernel_3382/994572986.py:73: FutureWarning: 

Passing `palette` without assigning `hue` is deprecated and will be removed in v0.14.0. Assign the `x` variable to `hue` and set `legend=False` for the same effect.

  sns.countplot(data=df_eda, x='label', palette='Set2', ax=ax)

<Figure size 600x400 with 1 Axes>
/tmp/ipykernel_3382/994572986.py:100: UserWarning: Could not infer format, so each element will be parsed individually, falling back to `dateutil`. To ensure parsing is consistent and as-expected, please specify a format.
  df_eda['checkin_hour_num'] = pd.to_datetime(df_eda['waktu_checkin'], errors='coerce').dt.hour.fillna(0)
/tmp/ipykernel_3382/994572986.py:101: UserWarning: Could not infer format, so each element will be parsed individually, falling back to `dateutil`. To ensure parsing is consistent and as-expected, please specify a format.
  df_eda['checkout_hour_num'] = pd.to_datetime(df_eda['waktu_checkout'], errors='coerce').dt.hour.fillna(0)


[3] Rata-rata Pelanggaran Berdasarkan Seluruh Kategori:

--- Tingkat Pelanggaran Berdasarkan status_pasutri ---
status_pasutri
Menikah          0.076
Belum Menikah    0.073
Name: label, dtype: float64

--- Tingkat Pelanggaran Berdasarkan status_kewarganegaraan ---
status_kewarganegaraan
WNA    0.078
WNI    0.072
Name: label, dtype: float64

--- Tingkat Pelanggaran Berdasarkan jenis_sewa ---
jenis_sewa
Harian      0.232
Bulanan     0.000
Mingguan    0.000
Name: label, dtype: float64

--- Tingkat Pelanggaran Berdasarkan metode_pembayaran ---
metode_pembayaran
Others          0.081
QRIS            0.075
Cash            0.075
Kartu Debit     0.074
Kartu Kredit    0.069
Name: label, dtype: float64

[4] Korelasi Antar Fitur:

<Figure size 1000x800 with 2 Axes>

[5] Visualisasi Pola Penyewaan (Dikaitkan dengan Klasifikasi):

<Figure size 1200x500 with 1 Axes>
<Figure size 1200x500 with 1 Axes>
<Figure size 800x500 with 1 Axes>
<Figure size 800x500 with 1 Axes>
<Figure size 800x500 with 1 Axes>
<Figure size 800x500 with 1 Axes>

[5b] Visualisasi Metode Pembayaran 'Others' (metode_lain):

/tmp/ipykernel_3382/994572986.py:208: FutureWarning: 

Passing `palette` without assigning `hue` is deprecated and will be removed in v0.14.0. Assign the `y` variable to `hue` and set `legend=False` for the same effect.

  sns.countplot(

<Figure size 1000x600 with 1 Axes>
/tmp/ipykernel_3382/994572986.py:242: UserWarning: Could not infer format, so each element will be parsed individually, falling back to `dateutil`. To ensure parsing is consistent and as-expected, please specify a format.
  rentals['waktu_checkout_dt'] = pd.to_datetime(rentals['waktu_checkout'], errors='coerce')
/tmp/ipykernel_3382/994572986.py:245: UserWarning: Could not infer format, so each element will be parsed individually, falling back to `dateutil`. To ensure parsing is consistent and as-expected, please specify a format.
  rentals['waktu_checkin_dt'] = pd.to_datetime(rentals['waktu_checkin'], errors='coerce')


=== Proses Feature Selection ===

-> Fitur yang dipilih (4 terbaik): ['lama_menginap', 'checkin_hour', 'checkout_hour', 'jenis_sewa']

===== Decision Tree =====
Classification Report:
              precision    recall  f1-score   support

           0     1.0000    0.9822    0.9910      1849
           1     0.8187    1.0000    0.9003       149

    accuracy                         0.9835      1998
   macro avg     0.9093    0.9911    0.9456      1998
weighted avg     0.9865    0.9835    0.9842      1998

Confusion Matrix (Text):
                Predicted
               0       1
          -----------------
True   0 |   1816      33
       1 |      0     149
          -----------------


===== Random Forest =====
Classification Report:
              precision    recall  f1-score   support

           0     1.0000    0.9822    0.9910      1849
           1     0.8187    1.0000    0.9003       149

    accuracy                         0.9835      1998
   macro avg     0.9093    0.9911    0.9456      1998
weighted avg     0.9865    0.9835    0.9842      1998

Confusion Matrix (Text):
                Predicted
               0       1
          -----------------
True   0 |   1816      33
       1 |      0     149
          -----------------


===== KNN =====
Classification Report:
              precision    recall  f1-score   support

           0     0.9891    0.9843    0.9867      1849
           1     0.8165    0.8658    0.8404       149

    accuracy                         0.9755      1998
   macro avg     0.9028    0.9250    0.9136      1998
weighted avg     0.9763    0.9755    0.9758      1998

Confusion Matrix (Text):
                Predicted
               0       1
          -----------------
True   0 |   1820      29
       1 |     20     129
          -----------------


===== SVM RBF =====
Classification Report:
              precision    recall  f1-score   support

           0     0.9988    0.9356    0.9662      1849
           1     0.5526    0.9866    0.7084       149

    accuracy                         0.9394      1998
   macro avg     0.7757    0.9611    0.8373      1998
weighted avg     0.9656    0.9394    0.9470      1998

Confusion Matrix (Text):
                Predicted
               0       1
          -----------------
True   0 |   1730     119
       1 |      2     147
          -----------------


===== XGBoost =====
Classification Report:
              precision    recall  f1-score   support

           0     0.9994    0.9822    0.9907      1849
           1     0.8177    0.9933    0.8970       149

    accuracy                         0.9830      1998
   macro avg     0.9086    0.9877    0.9438      1998
weighted avg     0.9859    0.9830    0.9837      1998

Confusion Matrix (Text):
                Predicted
               0       1
          -----------------
True   0 |   1816      33
       1 |      1     148
          -----------------


=== Proses Export Model ===
Sukses! Model dan Label Encoders telah disimpan sebagai file: model_random_forest.pkl
\end{lstlisting}

\begin{lstlisting}[language=Python, caption={Kode 12}]
# ============================================================
# DATA SEED + Tunning
# ============================================================

import pandas as pd
import numpy as np
import joblib

from sklearn.model_selection import train_test_split, RandomizedSearchCV
from sklearn.preprocessing import LabelEncoder, StandardScaler, MinMaxScaler
from sklearn.metrics import classification_report, confusion_matrix
from sklearn.pipeline import make_pipeline
from sklearn.feature_selection import SelectKBest, chi2, mutual_info_classif

from sklearn.tree import DecisionTreeClassifier
from sklearn.ensemble import RandomForestClassifier
from sklearn.neighbors import KNeighborsClassifier
from sklearn.svm import SVC
from xgboost import XGBClassifier

# ------------------------------------------------------------
# 1. LOAD DATA
# ------------------------------------------------------------
rentals   = pd.read_csv('rentals.csv')
logs      = pd.read_csv('rental_logs.csv')
units     = pd.read_csv('units.csv')
users     = pd.read_csv('users.csv')
violations= pd.read_csv('violations.csv')

rentals.columns = rentals.columns.str.lower()

if 'komentar' not in rentals.columns:
    raise ValueError("Kolom 'komentar' tidak ditemukan di rentals.csv")

print("renatls.csv")
print(rentals.head())
print("logs.csv")
print(logs.head())
print("units.csv")
print(units.head())
print("users.csv")
print(users.head())
print("violations.csv")
print(violations.head())

# ------------------------------------------------------------
# 2. LABELING
# ------------------------------------------------------------
rentals['komentar'] = rentals['komentar'].astype(str).str.lower().str.strip()
empty_values = ["", "nan", "na", "null", "-", "0", "tidak ada", "none", " "]
rentals['label'] = (~rentals['komentar'].isin(empty_values)).astype(int)

# ============================================================
# EXPLORATORY DATA ANALYSIS (EDA)
# ============================================================
print("\n" + "="*50)
print("=== EXPLORATORY DATA ANALYSIS (EDA) ===")
print("="*50)

df_eda = rentals.copy()

print("\n[1] Statistik Deskriptif:")
print(df_eda.describe(include='all'))

print("\n[2] Distribusi Kelas Target (0: Aman, 1: Melanggar):")
print(df_eda['label'].value_counts())

fig, ax = plt.subplots(figsize=(6, 4))
sns.countplot(data=df_eda, x='label', palette='Set2', ax=ax)
plt.title('Distribusi Kelas Klasifikasi (0: Aman, 1: Melanggar)')
plt.xlabel('Klasifikasi')
plt.ylabel('Jumlah Penyewaan')

y_max = df_eda['label'].value_counts().max()
plt.ylim(0, y_max * 1.15)

for p in ax.patches:
    if p.get_height() > 0:
        ax.annotate(f'{int(p.get_height())}', (p.get_x() + p.get_width() / 2., p.get_height()),
                    ha='center', va='baseline', fontsize=10, color='black', xytext=(0, 5),
                    textcoords='offset points')
plt.show()

kolom_dikecualikan = ['id', 'nama', 'nik', 'unit_number', 'tanggal_checkin', 'waktu_checkin',
                      'tanggal_checkout', 'waktu_checkout', 'komentar', 'user_email', 'diedit_oleh', 'created_at', 'metode_lain', 'classification']

semua_kategori = [col for col in df_eda.select_dtypes(include=['object', 'category']).columns if col not in kolom_dikecualikan]

print("\n[3] Rata-rata Pelanggaran Berdasarkan Seluruh Kategori:")
for col in semua_kategori:
    if col in df_eda.columns:
        print(f"\n--- Tingkat Pelanggaran Berdasarkan {col} ---")
        print(df_eda.groupby(col)['label'].mean().sort_values(ascending=False).round(3))

df_eda['lama_menginap_num'] = pd.to_numeric(df_eda['lama_menginap'], errors='coerce').fillna(0)
df_eda['checkin_hour_num'] = pd.to_datetime(df_eda['waktu_checkin'], errors='coerce').dt.hour.fillna(0)
df_eda['checkout_hour_num'] = pd.to_datetime(df_eda['waktu_checkout'], errors='coerce').dt.hour.fillna(0)

le_eda = LabelEncoder()
for c in semua_kategori:
    if c in df_eda.columns:
        df_eda[c + '_encoded'] = le_eda.fit_transform(df_eda[c].astype(str))

print("\n[4] Korelasi Antar Fitur:")
kolom_korelasi = ['lama_menginap_num', 'checkin_hour_num', 'checkout_hour_num', 'label']
kolom_korelasi += [c + '_encoded' for c in semua_kategori if c + '_encoded' in df_eda.columns]

corr_matrix = df_eda[kolom_korelasi].corr()

plt.figure(figsize=(10, 8))
sns.heatmap(corr_matrix, annot=True, cmap='coolwarm', fmt='.2f', linewidths=0.5)
plt.title('Heatmap Korelasi Antar Fitur dan Target (Klasifikasi)')
plt.tight_layout()
plt.show()

print("\n[5] Visualisasi Pola Penyewaan (Dikaitkan dengan Klasifikasi):")

custom_line_palette = {0: 'blue', 1: 'red'}
custom_labels = ['0: Aman', '1: Melanggar']

fig, ax1 = plt.subplots(figsize=(12, 5))
waktu_df = df_eda.groupby(['checkin_hour_num', 'label']).size().reset_index(name='jumlah')

sns.lineplot(data=waktu_df, x='checkin_hour_num', y='jumlah', hue='label', marker='o', palette=custom_line_palette, ax=ax1)

plt.title('Pola Waktu Kedatangan Berdasarkan Klasifikasi')
plt.xlabel('Jam Check-in (0-23)')
plt.ylabel('Jumlah Penyewaan')
plt.xticks(range(0, 24))
plt.grid(True, alpha=0.3)

y_max_line = waktu_df['jumlah'].max()
plt.ylim(-y_max_line * 0.20, y_max_line * 1.15)

handles, _ = ax1.get_legend_handles_labels()
plt.legend(handles=handles, labels=custom_labels, title='Klasifikasi')

for l in waktu_df['label'].unique():
    sub = waktu_df[waktu_df['label'] == l]
    warna_teks = custom_line_palette[l]

    for x, y in zip(sub['checkin_hour_num'], sub['jumlah']):
        y_offset = 8 if l == 0 else -12
        va_align = 'bottom' if l == 0 else 'top'
        plt.annotate(f'{int(y)}', (x, y), textcoords="offset points", xytext=(0, y_offset),
                     ha='center', va=va_align, fontsize=9, color=warna_teks, fontweight='bold')
plt.show()

fig, ax2 = plt.subplots(figsize=(12, 5))
lama_df = df_eda.groupby(['lama_menginap_num', 'label']).size().reset_index(name='jumlah')

sns.lineplot(data=lama_df, x='lama_menginap_num', y='jumlah', hue='label', marker='o', palette=custom_line_palette, ax=ax2)

plt.title('Tren Lama Menginap Berdasarkan Klasifikasi')
plt.xlabel('Lama Menginap (Hari)')
plt.ylabel('Jumlah Penyewaan')
plt.grid(True, alpha=0.3)

handles2, _ = ax2.get_legend_handles_labels()
plt.legend(handles=handles2, labels=custom_labels, title='Klasifikasi')
plt.show()

for col in semua_kategori:
    if col in df_eda.columns:
        fig, ax = plt.subplots(figsize=(8, 5))
        sns.countplot(data=df_eda, x=col, hue='label', palette='magma', ax=ax)

        judul_kolom = col.replace('_', ' ').title()
        plt.title(f'Perbandingan Jumlah Pelanggaran Berdasarkan {judul_kolom}')
        plt.xlabel(judul_kolom)
        plt.ylabel('Jumlah Penyewaan')
        plt.xticks(rotation=0)

        handles_bar, _ = ax.get_legend_handles_labels()
        plt.legend(handles=handles_bar, labels=custom_labels, title='Klasifikasi')

        y_max_bar = df_eda.groupby([col, 'label']).size().max()
        plt.ylim(0, y_max_bar * 1.15)

        for p in ax.patches:
            if p.get_height() > 0:
                ax.annotate(f'{int(p.get_height())}', (p.get_x() + p.get_width() / 2., p.get_height()),
                            ha='center', va='baseline', fontsize=9, color='black', xytext=(0, 5),
                            textcoords='offset points')
        plt.tight_layout()
        plt.show()


print("\n[5b] Visualisasi Metode Pembayaran 'Others' (metode_lain):")

# Memastikan kolom 'metode_lain' berbentuk string dan membersihkan spasi/nilai kosong
df_eda['metode_lain'] = df_eda['metode_lain'].astype(str).str.strip().str.title()
empty_lain = ["", "Nan", "Na", "Null", "-", "0", "None"]

# Filter hanya baris yang memiliki nilai di 'metode_lain'
df_others = df_eda[~df_eda['metode_lain'].isin(empty_lain)]

if not df_others.empty:
    fig, ax = plt.subplots(figsize=(10, 6))

    # Menghitung frekuensi masing-masing metode pembayaran lain
    order_others = df_others['metode_lain'].value_counts().index

    sns.countplot(
        data=df_others,
        y='metode_lain',
        order=order_others,
        palette='viridis',
        ax=ax
    )

    plt.title('Detail Metode Pembayaran "Others" (metode_lain)')
    plt.xlabel('Jumlah Penyewaan')
    plt.ylabel('Jenis Pembayaran Lainnya')

    # Menambahkan anotasi angka di sebelah setiap bar
    for p in ax.patches:
        width = p.get_width()
        if width > 0:
            ax.annotate(f'{int(width)}',
                        (width, p.get_y() + p.get_height() / 2.),
                        ha='left', va='center', fontsize=10, color='black',
                        xytext=(5, 0), textcoords='offset points')

    # Perlebar batas x-axis sedikit agar angka anotasi tidak terpotong
    x_max = df_others['metode_lain'].value_counts().max()
    plt.xlim(0, x_max * 1.15)

    plt.tight_layout()
    plt.show()
else:
    print("Tidak ada data metode pembayaran 'Others' yang valid untuk ditampilkan.")

# ------------------------------------------------------------
# 3. FEATURE ENGINEERING & PREPARATION
# ------------------------------------------------------------
# Parse waktu checkout
rentals['waktu_checkout_dt'] = pd.to_datetime(rentals['waktu_checkout'], errors='coerce')
rentals['checkout_hour'] = rentals['waktu_checkout_dt'].dt.hour.fillna(0).astype(int)

# Parse waktu checkin
rentals['waktu_checkin_dt'] = pd.to_datetime(rentals['waktu_checkin'], errors='coerce')
rentals['checkin_hour'] = rentals['waktu_checkin_dt'].dt.hour.fillna(0).astype(int)

# Pastikan lama_menginap berupa angka
rentals['lama_menginap'] = pd.to_numeric(rentals['lama_menginap'], errors='coerce').fillna(0)

# Isi data kosong untuk kolom kategorikal sebelum split
cat_cols = ['jenis_sewa', 'metode_pembayaran', 'status_pasutri', 'status_kewarganegaraan']
for c in cat_cols:
    rentals[c] = rentals[c].fillna('Unknown').astype(str)

numeric_cols = ['checkout_hour', 'checkin_hour', 'lama_menginap']
feature_cols = cat_cols + numeric_cols

X = rentals[feature_cols].copy()
y = rentals['label']

# ------------------------------------------------------------
# 4. TRAIN-TEST SPLIT
# ------------------------------------------------------------
X_train, X_test, y_train, y_test = train_test_split(
    X, y, stratify=y, test_size=0.2, random_state=42
)

# ------------------------------------------------------------
# 4B. ENCODING MENCEGAH DATA LEAKAGE
# ------------------------------------------------------------
# Simpan semua LabelEncoder ke dalam dictionary agar bisa di-export
label_encoders = {}

for c in cat_cols:
    le = LabelEncoder()
    # Fit & Transform pada data latih
    X_train[c] = le.fit_transform(X_train[c])

    # Tangani unseen labels pada data uji (ubah jadi 'Unknown' jika tidak pernah dilihat di X_train untuk mencegah error saat deployment
    X_test[c] = X_test[c].map(lambda s: s if s in le.classes_ else 'Unknown')
    if 'Unknown' not in le.classes_:
        le.classes_ = np.append(le.classes_, 'Unknown')

    X_test[c] = le.transform(X_test[c])

    # Simpan object le
    label_encoders[c] = le

# ------------------------------------------------------------
# 5. FEATURE SELECTION
# ------------------------------------------------------------
print("\n=== Proses Feature Selection ===")
chi2_selector = SelectKBest(score_func=chi2, k='all')
chi2_selector.fit(X_train, y_train)

train_df = X_train.copy()
train_df['label'] = y_train
pearson_scores = train_df.corr(method='pearson')['label'].drop('label').abs().values

mi_selector = SelectKBest(score_func=mutual_info_classif, k='all')
mi_selector.fit(X_train, y_train)

fs_df = pd.DataFrame({
    'Feature': X_train.columns,
    'Chi-Square': chi2_selector.scores_,
    'Pearson (Abs)': pearson_scores,
    'Mutual Info': mi_selector.scores_
})

scaler = MinMaxScaler()
fs_df[['Chi2_Norm', 'Pearson_Norm', 'MI_Norm']] = scaler.fit_transform(
    fs_df[['Chi-Square', 'Pearson (Abs)', 'Mutual Info']]
)
fs_df['Total_Score'] = fs_df[['Chi2_Norm', 'Pearson_Norm', 'MI_Norm']].mean(axis=1)
fs_df = fs_df.sort_values(by='Total_Score', ascending=False).reset_index(drop=True)

# Pilih Jumlah Fitur
n_select = 4
selected_features = fs_df['Feature'].head(n_select).tolist()
print(f"\n-> Fitur yang dipilih ({n_select} terbaik): {selected_features}")

X_train = X_train[selected_features]
X_test = X_test[selected_features]

# ------------------------------------------------------------
# 6. Fungsi Evaluasi
# ------------------------------------------------------------
def evaluate(name, model):
    y_pred = model.predict(X_test)
    print(f"\n===== {name} =====")
    print("Classification Report:")
    print(classification_report(y_test, y_pred, digits=4))
    cm = confusion_matrix(y_test, y_pred)
    print("Confusion Matrix (Text):")
    print("                Predicted")
    print("               0       1")
    print("          -----------------")
    print(f"True   0 | {cm[0][0]:6d}  {cm[0][1]:6d}")
    print(f"       1 | {cm[1][0]:6d}  {cm[1][1]:6d}")
    print("          -----------------\n")

# ------------------------------------------------------------
# 7. Latih Semua Model & Hyperparameter Tuning
# ------------------------------------------------------------
models = {}
custom_weight = {0: 1, 1: 4}

# --- A. Definisikan Model Dasar ---
dt_base = DecisionTreeClassifier(class_weight=custom_weight, random_state=42)
rf_base = RandomForestClassifier(class_weight=custom_weight, random_state=42)
knn_base = make_pipeline(StandardScaler(), KNeighborsClassifier())
svm_base = make_pipeline(StandardScaler(), SVC(kernel='rbf', probability=True, class_weight=custom_weight, random_state=42))
xgb_base = XGBClassifier(scale_pos_weight=4, eval_metric='logloss', random_state=42)

print("\nMemulai proses Hyperparameter Tuning untuk semua model. Ini mungkin memakan waktu...\n")

# --- B. Proses Tuning Setiap Model ---

# 1. Decision Tree
print(">> Tuning Decision Tree...")
dt_param_grid = {'max_depth': [3, 5, 7, 10, None], 'min_samples_split': [2, 5, 10]}
dt_random = RandomizedSearchCV(dt_base, dt_param_grid, n_iter=10, cv=3, random_state=42, n_jobs=-1, scoring='f1_macro')
dt_random.fit(X_train, y_train)
models["Decision Tree (Tuned)"] = dt_random.best_estimator_

# 2. Random Forest
print(">> Tuning Random Forest...")
rf_param_grid = {'n_estimators': [100, 200, 300], 'max_depth': [5, 7, 10, None], 'min_samples_split': [2, 5, 10], 'min_samples_leaf': [1, 2, 4]}
rf_random = RandomizedSearchCV(rf_base, rf_param_grid, n_iter=10, cv=3, random_state=42, n_jobs=-1, scoring='f1_macro')
rf_random.fit(X_train, y_train)
models["Random Forest (Tuned)"] = rf_random.best_estimator_

# 3. KNN (Karena menggunakan pipeline, parameter ditambahkan prefix 'kneighborsclassifier__')
print(">> Tuning KNN...")
knn_param_grid = {'kneighborsclassifier__n_neighbors': [3, 5, 7, 9, 11], 'kneighborsclassifier__weights': ['uniform', 'distance']}
knn_random = RandomizedSearchCV(knn_base, knn_param_grid, n_iter=10, cv=3, random_state=42, n_jobs=-1, scoring='f1_macro')
knn_random.fit(X_train, y_train)
models["KNN (Tuned)"] = knn_random.best_estimator_

# 4. SVM (Karena menggunakan pipeline, parameter ditambahkan prefix 'svc__')
print(">> Tuning SVM...")
svm_param_grid = {'svc__C': [0.1, 1, 10], 'svc__gamma': ['scale', 'auto', 0.1, 0.01]}
svm_random = RandomizedSearchCV(svm_base, svm_param_grid, n_iter=10, cv=3, random_state=42, n_jobs=-1, scoring='f1_macro')
svm_random.fit(X_train, y_train)
models["SVM RBF (Tuned)"] = svm_random.best_estimator_

# 5. XGBoost
print(">> Tuning XGBoost...")
xgb_param_grid = {'learning_rate': [0.01, 0.05, 0.1, 0.2], 'max_depth': [3, 5, 7], 'n_estimators': [100, 200, 300]}
xgb_random = RandomizedSearchCV(xgb_base, xgb_param_grid, n_iter=10, cv=3, random_state=42, n_jobs=-1, scoring='f1_macro')
xgb_random.fit(X_train, y_train)
models["XGBoost (Tuned)"] = xgb_random.best_estimator_

# --- C. Evaluasi Semua Model yang Sudah Di-Tune ---
# Karena setiap .best_estimator_ sudah di-fit oleh RandomizedSearchCV pada seluruh X_train,
# kita bisa langsung memanggil fungsi evaluate tanpa perlu melakukan model.fit() lagi.
for name, model in models.items():
    evaluate(name, model)

# ------------------------------------------------------------
# 8. Eksport Model Terbaik & Label Encoders
# ------------------------------------------------------------
print("\n=== Proses Export Model ===")

# Gabungkan model dan encoders ke dalam satu dictionary agar rapi
export_package = {
    'model': models["Random Forest (Tuned)"],
    'label_encoders': label_encoders,
    'selected_features': selected_features # Penting agar tahu urutan kolom masuk saat predict
}

# Menentukan nama file
filename = 'model_random_forest.pkl'

# Menyimpan package ke dalam file .pkl
joblib.dump(export_package, filename)
print(f"Sukses! Model dan Label Encoders telah disimpan sebagai file: {filename}")
\end{lstlisting}

\begin{lstlisting}[caption={Output 12}]
renatls.csv
   id                nama               nik status_pasutri  \
0   1        F**** B*****  5***************        Menikah   
1   2      E***** P******  3***************        Menikah   
2   3        O**** H*****  3***************  Belum Menikah   
3   4  N***** B*****-B***  3***************  Belum Menikah   
4   5       R****** S****  3***************        Menikah   

  status_kewarganegaraan jenis_sewa unit_number metode_pembayaran metode_lain  \
0                    WNI   Mingguan      C-1-13              Cash         NaN   
1                    WNA     Harian       D-5-2      Kartu Kredit         NaN   
2                    WNI    Bulanan     B-25-17              QRIS         NaN   
3                    WNI    Bulanan       A-2-8            Others       gopay   
4                    WNI     Harian      A-21-9       Kartu Debit         NaN   

  tanggal_checkin waktu_checkin tanggal_checkout waktu_checkout  \
0      2024-06-20      20:22:00       2024-07-04       07:46:00   
1      2023-02-19      05:33:00       2023-02-20       12:33:00   
2      2023-12-17      17:48:00       2024-03-16       09:02:00   
3      2018-05-18      13:15:00       2018-10-15       12:45:00   
4      2025-01-05      17:28:00       2025-01-05       22:28:00   

   lama_menginap                                komentar  \
0             14                                     NaN   
1              1                                     NaN   
2             90                                     NaN   
3            150                                     NaN   
4              0  Ditemukan alat suntik di tempat sampah   

                     user_email diedit_oleh           created_at  \
0   O**************@hotmail.com         NaN  2025-12-08 19:13:06   
1  D***************@hotmail.com         NaN  2025-12-08 19:13:06   
2              K*****@gmail.com         NaN  2025-12-08 19:13:06   
3              T*****@gmail.com         NaN  2025-12-08 19:13:06   
4      K*************@yahoo.com         NaN  2025-12-08 19:13:06   

  classification  
0         normal  
1         normal  
2         normal  
3         normal  
4   tidak normal  
logs.csv
   id  rental_id  action field_changed old_value  \
0   1          1  create           all       NaN   
1   2          2  create           all       NaN   
2   3          3  create           all       NaN   
3   4          4  create           all       NaN   
4   5          5  create           all       NaN   

                                           new_value  \
0  {"id":1,"nama":"F**** B*****","nik":"5*******************"id":2,"nama":"E***** P******","nik":"3*****************"id":3,"nama":"O**** H*****","nik":"3*******************"id":4,"nama":"N***** B*****-B***","nik":"3*************"id":5,"nama":"R****** S****","nik":"33966831...   

                          email            timestamp  
0   O**************@hotmail.com  2025-12-08 19:13:06  
1  D***************@hotmail.com  2025-12-08 19:13:06  
2              K*****@gmail.com  2025-12-08 19:13:06  
3              T*****@gmail.com  2025-12-08 19:13:06  
4      K*************@yahoo.com  2025-12-08 19:13:06  
units.csv
  unit_number                   user_email tower  lantai  unit
0      A-1-14     K*************@yahoo.com     A       1    14
1      A-1-18         A*******@hotmail.com     A       1    18
2      A-1-22  O**************@hotmail.com     A       1    22
3      A-1-24                p**@gmail.com     A       1    24
4      A-1-28  O**************@hotmail.com     A       1    28
users.csv
   id                         email         password           nib  role  \
0   1    Y***************@yahoo.com  oe4eaGaBLpQv76y  7.265157e+12  agen   
1   2           R******@hotmail.com  UhoFvefkUMB44Cj  2.780507e+12  agen   
2   3  F*****************@yahoo.com  Jf5D9lZJIGVf011  1.335353e+12  agen   
3   4       M************@yahoo.com  DYPYEZs_Zh17IjQ  8.262239e+12  agen   
4   5         E**********@gmail.com  1gcfQuFUGjUATHy  8.127913e+11  agen   

               name           created_at  
0        K** M*** I  2025-12-08 19:13:05  
1    D*. J**** L***  2025-12-08 19:13:05  
2  R******** G*****  2025-12-08 19:13:05  
3  B******* K******  2025-12-08 19:13:05  
4      M****** D***  2025-12-08 19:13:05  
violations.csv
   id  rental_id                                       photo_url  \
0   1        233      https://picsum.photos/seed/u717oX/983/2412   
1   2       7182      https://picsum.photos/seed/0cZJtXQ/243/249   
2   3       9117     https://picsum.photos/seed/Pl0nan/3034/2819   
3   4       8582  https://picsum.photos/seed/LKFORW1qWE/3127/725   
4   5       4165    https://picsum.photos/seed/dkpBkZPg/336/1118   

                                         description  \
0              Crepusculum torrens vomica atque rem.   
1  Inventore articulus corrumpo pax clamo absterg...   
2  Nostrum desino comitatus debilito perferendis ...   
3  Tremo accedo aegre votum adhuc aveho vilicus v...   
4                         Comminor vis thorax absum.   

                         uploaded_by           created_at  
0     L*******************@yahoo.com  2025-12-08 19:13:07  
1        O**************@hotmail.com  2025-12-08 19:13:07  
2       A*****************@gmail.com  2025-12-08 19:13:07  
3                  S******@gmail.com  2025-12-08 19:13:07  
4  I********************@hotmail.com  2025-12-08 19:13:07  

==================================================
=== EXPLORATORY DATA ANALYSIS (EDA) ===
==================================================

[1] Statistik Deskriptif:
                  id           nama           nik status_pasutri  \
count    9989.000000           9989  9.989000e+03           9989   
unique           NaN           9894           NaN              2   
top              NaN  R*** D*******           NaN        Menikah   
freq             NaN              3           NaN           5004   
mean     5000.496746            NaN  3.649718e+15            NaN   
std      2887.229484            NaN  6.341279e+14            NaN   
min         1.000000            NaN  3.100037e+15            NaN   
25%      2501.000000            NaN  3.274772e+15            NaN   
50%      5001.000000            NaN  3.448126e+15            NaN   
75%      7500.000000            NaN  3.624302e+15            NaN   
max     10000.000000            NaN  5.199981e+15            NaN   

       status_kewarganegaraan jenis_sewa unit_number metode_pembayaran  \
count                    9989       9989        9989              9989   
unique                      2          3         794                 5   
top                       WNI   Mingguan     B-13-18      Kartu Kredit   
freq                     5011       3466          25              2045   
mean                      NaN        NaN         NaN               NaN   
std                       NaN        NaN         NaN               NaN   
min                       NaN        NaN         NaN               NaN   
25%                       NaN        NaN         NaN               NaN   
50%                       NaN        NaN         NaN               NaN   
75%                       NaN        NaN         NaN               NaN   
max                       NaN        NaN         NaN               NaN   

       metode_lain tanggal_checkin waktu_checkin tanggal_checkout  \
count         2000            9989          9989             9989   
unique           7            4742           754             4791   
top       paylater      2021-11-18      16:20:00       2019-01-29   
freq           297               8            34                9   
mean           NaN             NaN           NaN              NaN   
std            NaN             NaN           NaN              NaN   
min            NaN             NaN           NaN              NaN   
25%            NaN             NaN           NaN              NaN   
50%            NaN             NaN           NaN              NaN   
75%            NaN             NaN           NaN              NaN   
max            NaN             NaN           NaN              NaN   

       waktu_checkout  lama_menginap komentar        user_email  \
count            9989    9989.000000     9989              9989   
unique            659            NaN       11                53   
top          07:21:00            NaN      nan  K*****@gmail.com   
freq               43            NaN     9243               223   
mean              NaN      35.655321      NaN               NaN   
std               NaN      45.875125      NaN               NaN   
min               NaN       0.000000      NaN               NaN   
25%               NaN       5.000000      NaN               NaN   
50%               NaN      14.000000      NaN               NaN   
75%               NaN      60.000000      NaN               NaN   
max               NaN     150.000000      NaN               NaN   

                        diedit_oleh           created_at classification  \
count                            51                 9989           9989   
unique                           35                    1              2   
top     J**************@hotmail.com  2025-12-08 19:13:06         normal   
freq                              3                 9989           9243   
mean                            NaN                  NaN            NaN   
std                             NaN                  NaN            NaN   
min                             NaN                  NaN            NaN   
25%                             NaN                  NaN            NaN   
50%                             NaN                  NaN            NaN   
75%                             NaN                  NaN            NaN   
max                             NaN                  NaN            NaN   

              label  
count   9989.000000  
unique          NaN  
top             NaN  
freq            NaN  
mean       0.074682  
std        0.262891  
min        0.000000  
25%        0.000000  
50%        0.000000  
75%        0.000000  
max        1.000000  

[2] Distribusi Kelas Target (0: Aman, 1: Melanggar):
label
0    9243
1     746
Name: count, dtype: int64

/tmp/ipykernel_3382/3979709149.py:69: FutureWarning: 

Passing `palette` without assigning `hue` is deprecated and will be removed in v0.14.0. Assign the `x` variable to `hue` and set `legend=False` for the same effect.

  sns.countplot(data=df_eda, x='label', palette='Set2', ax=ax)

<Figure size 600x400 with 1 Axes>
/tmp/ipykernel_3382/3979709149.py:96: UserWarning: Could not infer format, so each element will be parsed individually, falling back to `dateutil`. To ensure parsing is consistent and as-expected, please specify a format.
  df_eda['checkin_hour_num'] = pd.to_datetime(df_eda['waktu_checkin'], errors='coerce').dt.hour.fillna(0)
/tmp/ipykernel_3382/3979709149.py:97: UserWarning: Could not infer format, so each element will be parsed individually, falling back to `dateutil`. To ensure parsing is consistent and as-expected, please specify a format.
  df_eda['checkout_hour_num'] = pd.to_datetime(df_eda['waktu_checkout'], errors='coerce').dt.hour.fillna(0)


[3] Rata-rata Pelanggaran Berdasarkan Seluruh Kategori:

--- Tingkat Pelanggaran Berdasarkan status_pasutri ---
status_pasutri
Menikah          0.076
Belum Menikah    0.073
Name: label, dtype: float64

--- Tingkat Pelanggaran Berdasarkan status_kewarganegaraan ---
status_kewarganegaraan
WNA    0.078
WNI    0.072
Name: label, dtype: float64

--- Tingkat Pelanggaran Berdasarkan jenis_sewa ---
jenis_sewa
Harian      0.232
Bulanan     0.000
Mingguan    0.000
Name: label, dtype: float64

--- Tingkat Pelanggaran Berdasarkan metode_pembayaran ---
metode_pembayaran
Others          0.081
QRIS            0.075
Cash            0.075
Kartu Debit     0.074
Kartu Kredit    0.069
Name: label, dtype: float64

[4] Korelasi Antar Fitur:

<Figure size 1000x800 with 2 Axes>

[5] Visualisasi Pola Penyewaan (Dikaitkan dengan Klasifikasi):

<Figure size 1200x500 with 1 Axes>
<Figure size 1200x500 with 1 Axes>
<Figure size 800x500 with 1 Axes>
<Figure size 800x500 with 1 Axes>
<Figure size 800x500 with 1 Axes>
<Figure size 800x500 with 1 Axes>

[5b] Visualisasi Metode Pembayaran 'Others' (metode_lain):

/tmp/ipykernel_3382/3979709149.py:204: FutureWarning: 

Passing `palette` without assigning `hue` is deprecated and will be removed in v0.14.0. Assign the `y` variable to `hue` and set `legend=False` for the same effect.

  sns.countplot(

<Figure size 1000x600 with 1 Axes>
/tmp/ipykernel_3382/3979709149.py:238: UserWarning: Could not infer format, so each element will be parsed individually, falling back to `dateutil`. To ensure parsing is consistent and as-expected, please specify a format.
  rentals['waktu_checkout_dt'] = pd.to_datetime(rentals['waktu_checkout'], errors='coerce')
/tmp/ipykernel_3382/3979709149.py:242: UserWarning: Could not infer format, so each element will be parsed individually, falling back to `dateutil`. To ensure parsing is consistent and as-expected, please specify a format.
  rentals['waktu_checkin_dt'] = pd.to_datetime(rentals['waktu_checkin'], errors='coerce')


=== Proses Feature Selection ===

-> Fitur yang dipilih (4 terbaik): ['lama_menginap', 'checkin_hour', 'checkout_hour', 'jenis_sewa']

Memulai proses Hyperparameter Tuning untuk semua model. Ini mungkin memakan waktu...

>> Tuning Decision Tree...
>> Tuning Random Forest...
>> Tuning KNN...
>> Tuning SVM...
>> Tuning XGBoost...

===== Decision Tree (Tuned) =====
Classification Report:
              precision    recall  f1-score   support

           0     1.0000    0.9816    0.9907      1849
           1     0.8142    1.0000    0.8976       149

    accuracy                         0.9830      1998
   macro avg     0.9071    0.9908    0.9442      1998
weighted avg     0.9861    0.9830    0.9838      1998

Confusion Matrix (Text):
                Predicted
               0       1
          -----------------
True   0 |   1815      34
       1 |      0     149
          -----------------


===== Random Forest (Tuned) =====
Classification Report:
              precision    recall  f1-score   support

           0     1.0000    0.9816    0.9907      1849
           1     0.8142    1.0000    0.8976       149

    accuracy                         0.9830      1998
   macro avg     0.9071    0.9908    0.9442      1998
weighted avg     0.9861    0.9830    0.9838      1998

Confusion Matrix (Text):
                Predicted
               0       1
          -----------------
True   0 |   1815      34
       1 |      0     149
          -----------------


===== KNN (Tuned) =====
Classification Report:
              precision    recall  f1-score   support

           0     0.9951    0.9838    0.9894      1849
           1     0.8235    0.9396    0.8777       149

    accuracy                         0.9805      1998
   macro avg     0.9093    0.9617    0.9336      1998
weighted avg     0.9823    0.9805    0.9811      1998

Confusion Matrix (Text):
                Predicted
               0       1
          -----------------
True   0 |   1819      30
       1 |      9     140
          -----------------


===== SVM RBF (Tuned) =====
Classification Report:
              precision    recall  f1-score   support

           0     1.0000    0.9578    0.9785      1849
           1     0.6564    1.0000    0.7926       149

    accuracy                         0.9610      1998
   macro avg     0.8282    0.9789    0.8855      1998
weighted avg     0.9744    0.9610    0.9646      1998

Confusion Matrix (Text):
                Predicted
               0       1
          -----------------
True   0 |   1771      78
       1 |      0     149
          -----------------


===== XGBoost (Tuned) =====
Classification Report:
              precision    recall  f1-score   support

           0     0.9994    0.9822    0.9907      1849
           1     0.8177    0.9933    0.8970       149

    accuracy                         0.9830      1998
   macro avg     0.9086    0.9877    0.9438      1998
weighted avg     0.9859    0.9830    0.9837      1998

Confusion Matrix (Text):
                Predicted
               0       1
          -----------------
True   0 |   1816      33
       1 |      1     148
          -----------------


=== Proses Export Model ===
Sukses! Model dan Label Encoders telah disimpan sebagai file: model_random_forest.pkl
\end{lstlisting}

\begin{lstlisting}[language=Python, caption={Kode 13}]
# ============================================================
# DATA RIIL (TANPA TUNING) --> jumlah fitur = 9
# ============================================================

import pandas as pd
import numpy as np
import joblib
import matplotlib.pyplot as plt
import seaborn as sns

from sklearn.model_selection import train_test_split
from sklearn.preprocessing import LabelEncoder, StandardScaler, MinMaxScaler
from sklearn.metrics import classification_report, confusion_matrix
from sklearn.pipeline import make_pipeline
from sklearn.feature_selection import SelectKBest, chi2, mutual_info_classif

from sklearn.tree import DecisionTreeClassifier
from sklearn.ensemble import RandomForestClassifier
from sklearn.neighbors import KNeighborsClassifier
from sklearn.svm import SVC
from xgboost import XGBClassifier

import matplotlib.pyplot as plt
import seaborn as sns

# ------------------------------------------------------------
# 1. LOAD DATA
# ------------------------------------------------------------
rentals    = pd.read_csv('rentals.csv')
logs       = pd.read_csv('rental_logs.csv')
# units      = pd.read_csv('units.csv')
# users      = pd.read_csv('users.csv')
# violations = pd.read_csv('violations.csv')

rentals.columns = rentals.columns.str.lower()

if 'komentar' not in rentals.columns:
    raise ValueError("Kolom 'komentar' tidak ditemukan di rentals.csv")

print("rentals.csv")
print(rentals.head())
print("logs.csv")
print(logs.head())


# ------------------------------------------------------------
# 2. LABELING
# ------------------------------------------------------------
rentals['komentar'] = rentals['komentar'].astype(str).str.lower().str.strip()
empty_values = ["", "nan", "na", "null", "-", "0", "tidak ada", "none", " "]
rentals['label'] = (~rentals['komentar'].isin(empty_values)).astype(int)


# ============================================================
# EXPLORATORY DATA ANALYSIS (EDA)
# ============================================================
print("\n" + "="*50)
print("=== EXPLORATORY DATA ANALYSIS (EDA) ===")
print("="*50)

df_eda = rentals.copy()

print("\n[1] Statistik Deskriptif:")
print(df_eda.describe(include='all'))

print("\n[2] Distribusi Kelas Target (0: Aman, 1: Melanggar):")
print(df_eda['label'].value_counts())

fig, ax = plt.subplots(figsize=(6, 4))
sns.countplot(data=df_eda, x='label', palette='Set2', ax=ax)
plt.title('Distribusi Kelas Klasifikasi (0: Aman, 1: Melanggar)')
plt.xlabel('Klasifikasi')
plt.ylabel('Jumlah Penyewaan')

y_max = df_eda['label'].value_counts().max()
plt.ylim(0, y_max * 1.15)

for p in ax.patches:
    if p.get_height() > 0:
        ax.annotate(f'{int(p.get_height())}', (p.get_x() + p.get_width() / 2., p.get_height()),
                    ha='center', va='baseline', fontsize=10, color='black', xytext=(0, 5),
                    textcoords='offset points')
plt.show()

kolom_dikecualikan = ['id', 'nama', 'nik', 'unit_number', 'tanggal_checkin', 'waktu_checkin',
                      'tanggal_checkout', 'waktu_checkout', 'komentar', 'user_email', 'diedit_oleh', 'created_at', 'metode_lain', 'classification']

semua_kategori = [col for col in df_eda.select_dtypes(include=['object', 'category']).columns if col not in kolom_dikecualikan]

print("\n[3] Rata-rata Pelanggaran Berdasarkan Seluruh Kategori:")
for col in semua_kategori:
    if col in df_eda.columns:
        print(f"\n--- Tingkat Pelanggaran Berdasarkan {col} ---")
        print(df_eda.groupby(col)['label'].mean().sort_values(ascending=False).round(3))

df_eda['lama_menginap_num'] = pd.to_numeric(df_eda['lama_menginap'], errors='coerce').fillna(0)
df_eda['checkin_hour_num'] = pd.to_datetime(df_eda['waktu_checkin'], errors='coerce').dt.hour.fillna(0)
df_eda['checkout_hour_num'] = pd.to_datetime(df_eda['waktu_checkout'], errors='coerce').dt.hour.fillna(0)

le_eda = LabelEncoder()
for c in semua_kategori:
    if c in df_eda.columns:
        df_eda[c + '_encoded'] = le_eda.fit_transform(df_eda[c].astype(str))

print("\n[4] Korelasi Antar Fitur:")
kolom_korelasi = ['lama_menginap_num', 'checkin_hour_num', 'checkout_hour_num', 'label']
kolom_korelasi += [c + '_encoded' for c in semua_kategori if c + '_encoded' in df_eda.columns]

corr_matrix = df_eda[kolom_korelasi].corr()

plt.figure(figsize=(10, 8))
sns.heatmap(corr_matrix, annot=True, cmap='coolwarm', fmt='.2f', linewidths=0.5)
plt.title('Heatmap Korelasi Antar Fitur dan Target (Klasifikasi)')
plt.tight_layout()
plt.show()

print("\n[5] Visualisasi Pola Penyewaan (Dikaitkan dengan Klasifikasi):")

custom_line_palette = {0: 'blue', 1: 'red'}
custom_labels = ['0: Aman', '1: Melanggar']

fig, ax1 = plt.subplots(figsize=(12, 5))
waktu_df = df_eda.groupby(['checkin_hour_num', 'label']).size().reset_index(name='jumlah')

sns.lineplot(data=waktu_df, x='checkin_hour_num', y='jumlah', hue='label', marker='o', palette=custom_line_palette, ax=ax1)

plt.title('Pola Waktu Kedatangan Berdasarkan Klasifikasi')
plt.xlabel('Jam Check-in (0-23)')
plt.ylabel('Jumlah Penyewaan')
plt.xticks(range(0, 24))
plt.grid(True, alpha=0.3)

y_max_line = waktu_df['jumlah'].max()
plt.ylim(-y_max_line * 0.20, y_max_line * 1.15)

handles, _ = ax1.get_legend_handles_labels()
plt.legend(handles=handles, labels=custom_labels, title='Klasifikasi')

for l in waktu_df['label'].unique():
    sub = waktu_df[waktu_df['label'] == l]
    warna_teks = custom_line_palette[l]

    for x, y in zip(sub['checkin_hour_num'], sub['jumlah']):
        y_offset = 8 if l == 0 else -12
        va_align = 'bottom' if l == 0 else 'top'
        plt.annotate(f'{int(y)}', (x, y), textcoords="offset points", xytext=(0, y_offset),
                     ha='center', va=va_align, fontsize=9, color=warna_teks, fontweight='bold')
plt.show()

fig, ax2 = plt.subplots(figsize=(12, 5))
lama_df = df_eda.groupby(['lama_menginap_num', 'label']).size().reset_index(name='jumlah')

sns.lineplot(data=lama_df, x='lama_menginap_num', y='jumlah', hue='label', marker='o', palette=custom_line_palette, ax=ax2)

plt.title('Tren Lama Menginap Berdasarkan Klasifikasi')
plt.xlabel('Lama Menginap (Hari)')
plt.ylabel('Jumlah Penyewaan')
plt.grid(True, alpha=0.3)

handles2, _ = ax2.get_legend_handles_labels()
plt.legend(handles=handles2, labels=custom_labels, title='Klasifikasi')
plt.show()

for col in semua_kategori:
    if col in df_eda.columns:
        fig, ax = plt.subplots(figsize=(8, 5))
        sns.countplot(data=df_eda, x=col, hue='label', palette='magma', ax=ax)

        judul_kolom = col.replace('_', ' ').title()
        plt.title(f'Perbandingan Jumlah Pelanggaran Berdasarkan {judul_kolom}')
        plt.xlabel(judul_kolom)
        plt.ylabel('Jumlah Penyewaan')
        plt.xticks(rotation=0)

        handles_bar, _ = ax.get_legend_handles_labels()
        plt.legend(handles=handles_bar, labels=custom_labels, title='Klasifikasi')

        y_max_bar = df_eda.groupby([col, 'label']).size().max()
        plt.ylim(0, y_max_bar * 1.15)

        for p in ax.patches:
            if p.get_height() > 0:
                ax.annotate(f'{int(p.get_height())}', (p.get_x() + p.get_width() / 2., p.get_height()),
                            ha='center', va='baseline', fontsize=9, color='black', xytext=(0, 5),
                            textcoords='offset points')
        plt.tight_layout()
        plt.show()


print("\n[5b] Visualisasi Metode Pembayaran 'Others' (metode_lain):")

# Memastikan kolom 'metode_lain' berbentuk string dan membersihkan spasi/nilai kosong
df_eda['metode_lain'] = df_eda['metode_lain'].astype(str).str.strip().str.title()
empty_lain = ["", "Nan", "Na", "Null", "-", "0", "None"]

# Filter hanya baris yang memiliki nilai di 'metode_lain'
df_others = df_eda[~df_eda['metode_lain'].isin(empty_lain)]

if not df_others.empty:
    fig, ax = plt.subplots(figsize=(10, 6))

    # Menghitung frekuensi masing-masing metode pembayaran lain
    order_others = df_others['metode_lain'].value_counts().index

    sns.countplot(
        data=df_others,
        y='metode_lain',
        order=order_others,
        palette='viridis',
        ax=ax
    )

    plt.title('Detail Metode Pembayaran "Others" (metode_lain)')
    plt.xlabel('Jumlah Penyewaan')
    plt.ylabel('Jenis Pembayaran Lainnya')

    # Menambahkan anotasi angka di sebelah setiap bar
    for p in ax.patches:
        width = p.get_width()
        if width > 0:
            ax.annotate(f'{int(width)}',
                        (width, p.get_y() + p.get_height() / 2.),
                        ha='left', va='center', fontsize=10, color='black',
                        xytext=(5, 0), textcoords='offset points')

    # Perlebar batas x-axis sedikit agar angka anotasi tidak terpotong
    x_max = df_others['metode_lain'].value_counts().max()
    plt.xlim(0, x_max * 1.15)

    plt.tight_layout()
    plt.show()
else:
    print("Tidak ada data metode pembayaran 'Others' yang valid untuk ditampilkan.")


# ------------------------------------------------------------
# 3. FEATURE ENGINEERING & PREPARATION
# ------------------------------------------------------------
rentals['waktu_checkout_dt'] = pd.to_datetime(rentals['waktu_checkout'], errors='coerce')
rentals['checkout_hour'] = rentals['waktu_checkout_dt'].dt.hour.fillna(0).astype(int)

rentals['waktu_checkin_dt'] = pd.to_datetime(rentals['waktu_checkin'], errors='coerce')
rentals['checkin_hour'] = rentals['waktu_checkin_dt'].dt.hour.fillna(0).astype(int)

rentals['lama_menginap'] = pd.to_numeric(rentals['lama_menginap'], errors='coerce').fillna(0)

cat_cols = ['jenis_sewa', 'metode_pembayaran', 'status_pasutri', 'status_kewarganegaraan']
for c in cat_cols:
    rentals[c] = rentals[c].fillna('Unknown').astype(str)

numeric_cols = ['checkout_hour', 'checkin_hour', 'lama_menginap']
feature_cols = cat_cols + numeric_cols

X = rentals[feature_cols].copy()
y = rentals['label']

# ------------------------------------------------------------
# 4. TRAIN-TEST SPLIT
# ------------------------------------------------------------
X_train, X_test, y_train, y_test = train_test_split(
    X, y, stratify=y, test_size=0.2, random_state=42
)

# ------------------------------------------------------------
# 4B. ENCODING MENCEGAH DATA LEAKAGE
# ------------------------------------------------------------
label_encoders = {}

for c in cat_cols:
    le = LabelEncoder()
    X_train[c] = le.fit_transform(X_train[c])

    X_test[c] = X_test[c].map(lambda s: s if s in le.classes_ else 'Unknown')
    if 'Unknown' not in le.classes_:
        le.classes_ = np.append(le.classes_, 'Unknown')

    X_test[c] = le.transform(X_test[c])

    label_encoders[c] = le

# ------------------------------------------------------------
# 5. FEATURE SELECTION
# ------------------------------------------------------------
print("\n=== Proses Feature Selection ===")
chi2_selector = SelectKBest(score_func=chi2, k='all')
chi2_selector.fit(X_train, y_train)

train_df = X_train.copy()
train_df['label'] = y_train
pearson_scores = train_df.corr(method='pearson')['label'].drop('label').abs().values

mi_selector = SelectKBest(score_func=mutual_info_classif, k='all')
mi_selector.fit(X_train, y_train)

fs_df = pd.DataFrame({
    'Feature': X_train.columns,
    'Chi-Square': chi2_selector.scores_,
    'Pearson (Abs)': pearson_scores,
    'Mutual Info': mi_selector.scores_
})

scaler = MinMaxScaler()
fs_df[['Chi2_Norm', 'Pearson_Norm', 'MI_Norm']] = scaler.fit_transform(
    fs_df[['Chi-Square', 'Pearson (Abs)', 'Mutual Info']]
)
fs_df['Total_Score'] = fs_df[['Chi2_Norm', 'Pearson_Norm', 'MI_Norm']].mean(axis=1)
fs_df = fs_df.sort_values(by='Total_Score', ascending=False).reset_index(drop=True)

# Jumlah Fitur
n_select = 9
selected_features = fs_df['Feature'].head(n_select).tolist()
print(f"\n-> Fitur yang dipilih ({n_select} terbaik): {selected_features}")

X_train = X_train[selected_features]
X_test = X_test[selected_features]

# ------------------------------------------------------------
# 6. Fungsi Evaluasi
# ------------------------------------------------------------
def evaluate(name, model):
    y_pred = model.predict(X_test)
    print(f"\n===== {name} =====")
    print("Classification Report:")
    print(classification_report(y_test, y_pred, digits=4))
    cm = confusion_matrix(y_test, y_pred)
    print("Confusion Matrix (Text):")
    print("                Predicted")
    print("               0       1")
    print("          -----------------")
    print(f"True   0 | {cm[0][0]:6d}  {cm[0][1]:6d}")
    print(f"       1 | {cm[1][0]:6d}  {cm[1][1]:6d}")
    print("          -----------------\n")

# ------------------------------------------------------------
# 7. Latih Semua Model (TANPA TUNING)
# ------------------------------------------------------------
models = {}
custom_weight = {0: 1, 1: 4}

# Model menggunakan pengaturan tetap
models["Decision Tree"] = DecisionTreeClassifier(max_depth=6, class_weight=custom_weight, random_state=42)
models["Random Forest"] = RandomForestClassifier(n_estimators=200, max_depth=7, class_weight=custom_weight, random_state=42)
models["KNN"] = make_pipeline(StandardScaler(), KNeighborsClassifier(n_neighbors=5, weights='distance'))
models["SVM RBF"] = make_pipeline(StandardScaler(), SVC(kernel='rbf', probability=True, class_weight=custom_weight, random_state=42))
models["XGBoost"] = XGBClassifier(n_estimators=200, max_depth=5, learning_rate=0.05, scale_pos_weight=4, eval_metric='logloss', random_state=42)

for name, model in models.items():
    model.fit(X_train, y_train)
    evaluate(name, model)

# ------------------------------------------------------------
# 8. Eksport Model Terbaik & Label Encoders
# ------------------------------------------------------------
print("\n=== Proses Export Model ===")

# Gabungkan model dan encoders ke dalam satu dictionary agar rapi
export_package = {
    'model': models["KNN"],
    'label_encoders': label_encoders,
    'selected_features': selected_features # Penting agar tahu urutan kolom masuk saat predict
}

# Menentukan nama file yang sesuai untuk KNN
filename = 'model_knn.pkl'

# Menyimpan package ke dalam file .pkl
joblib.dump(export_package, filename)
print(f"Sukses! Model dan Label Encoders telah disimpan sebagai file: {filename}")
\end{lstlisting}

\begin{lstlisting}[caption={Output 13}]
rentals.csv
   id                   nama               nik status_pasutri  \
0   3  N*** Y**** S**** D***  3***************  Belum Menikah   
1   4     A***** N** H******  3***************  Belum Menikah   
2   5     A***** A*** S*****  3***************        Menikah   
3   6            C**** W****  1***************  Belum Menikah   
4   7             S*** U****  3***************  Belum Menikah   

  status_kewarganegaraan jenis_sewa unit_number metode_pembayaran metode_lain  \
0                    WNI     Harian     B-20-27              Cash         NaN   
1                    WNI     Harian     A-15-25              Cash         NaN   
2                    WNI     Harian     A-06-24            Others    Aplikasi   
3                    WNI     Harian     D-15-21              QRIS         NaN   
4                    WNI     Harian       A-8-8            Others    Transfer   

  tanggal_checkin waktu_checkin tanggal_checkout waktu_checkout  \
0      2026-05-20      13:08:00       2026-05-21            NaN   
1      2026-05-19      13:25:00       2026-05-22            NaN   
2      2026-05-17      13:27:00       2026-05-22            NaN   
3      2026-05-19      13:29:00       2026-05-24            NaN   
4      2025-01-01      19:58:00       2025-01-06       11:22:00   

   lama_menginap komentar  user_email  user_pengguna_id  editor_pengguna_id  \
0              1      NaN         NaN              1267                 NaN   
1              3      NaN         NaN              1267                 NaN   
2              5      NaN         NaN              1267                 NaN   
3              5      NaN         NaN              1267                 NaN   
4              5      NaN         NaN              1271              1271.0   

            created_at classification  
0  2026-05-20 13:09:49            NaN  
1  2026-05-20 13:26:04            NaN  
2  2026-05-20 13:28:03            NaN  
3  2026-05-20 13:30:18            NaN  
4  2026-05-25 10:06:33         normal  
logs.csv
     id  rental_id  action field_changed  old_value  \
0     6          3  create           all        NaN   
1     7          4  create           all        NaN   
2     8          5  create           all        NaN   
3     9          6  create           all        NaN   
4  4906          7  create           all        NaN   

                                           new_value  email  pengguna_id  \
0  {"nama":"N*** Y**** S**** D***","nik":"3*************************************"nama":"A***** N** H******","nik":"3****************************************"nama":"A***** A*** S*****","nik":"3****************************************"nama":"C**** W****","nik":"1***************"...    NaN         1267   
4  {"nama":"S*** U****","nik":"3***************",...    NaN         1271   

             timestamp  
0  2026-05-20 13:09:49  
1  2026-05-20 13:26:04  
2  2026-05-20 13:28:03  
3  2026-05-20 13:30:18  
4  2026-05-25 10:06:33  

==================================================
=== EXPLORATORY DATA ANALYSIS (EDA) ===
==================================================

[1] Statistik Deskriptif:
                id                         nama           nik status_pasutri  \
count   544.000000                          544  5.440000e+02            544   
unique         NaN                          207           NaN              2   
top            NaN  Y***** M**** S**** S*******           NaN  Belum Menikah   
freq           NaN                            8           NaN            330   
mean    274.538603                          NaN  3.214330e+15            NaN   
std     157.250941                          NaN  1.204147e+15            NaN   
min       3.000000                          NaN  1.202052e+15            NaN   
25%     138.750000                          NaN  3.173074e+15            NaN   
50%     274.500000                          NaN  3.209121e+15            NaN   
75%     410.250000                          NaN  3.276064e+15            NaN   
max     553.000000                          NaN  9.104017e+15            NaN   

       status_kewarganegaraan jenis_sewa unit_number metode_pembayaran  \
count                     544        544         544               544   
unique                      1          2          12                 5   
top                       WNI     Harian       B-1-8            Others   
freq                      544        543          90               209   
mean                      NaN        NaN         NaN               NaN   
std                       NaN        NaN         NaN               NaN   
min                       NaN        NaN         NaN               NaN   
25%                       NaN        NaN         NaN               NaN   
50%                       NaN        NaN         NaN               NaN   
75%                       NaN        NaN         NaN               NaN   
max                       NaN        NaN         NaN               NaN   

       metode_lain tanggal_checkin  ... tanggal_checkout waktu_checkout  \
count          209             544  ...              544            538   
unique           2             311  ...              300            232   
top       Transfer      2026-02-06  ...       2026-02-08       12:00:00   
freq           208               5  ...                5             82   
mean           NaN             NaN  ...              NaN            NaN   
std            NaN             NaN  ...              NaN            NaN   
min            NaN             NaN  ...              NaN            NaN   
25%            NaN             NaN  ...              NaN            NaN   
50%            NaN             NaN  ...              NaN            NaN   
75%            NaN             NaN  ...              NaN            NaN   
max            NaN             NaN  ...              NaN            NaN   

       lama_menginap  komentar user_email  user_pengguna_id  \
count     544.000000       544        0.0        544.000000   
unique           NaN        11        NaN               NaN   
top              NaN       nan        NaN               NaN   
freq             NaN       459        NaN               NaN   
mean        2.641544       NaN        NaN       1271.904412   
std         1.792711       NaN        NaN          1.177622   
min         1.000000       NaN        NaN       1262.000000   
25%         2.000000       NaN        NaN       1271.000000   
50%         2.000000       NaN        NaN       1272.000000   
75%         3.000000       NaN        NaN       1273.000000   
max        31.000000       NaN        NaN       1273.000000   

        editor_pengguna_id           created_at classification       label  
count           538.000000                  544            537  544.000000  
unique                 NaN                  544              2         NaN  
top                    NaN  2026-05-30 15:00:26         normal         NaN  
freq                   NaN                    1            452         NaN  
mean           1271.977695                  NaN            NaN    0.156250  
std               0.923249                  NaN            NaN    0.363426  
min            1262.000000                  NaN            NaN    0.000000  
25%            1271.000000                  NaN            NaN    0.000000  
50%            1272.000000                  NaN            NaN    0.000000  
75%            1273.000000                  NaN            NaN    0.000000  
max            1273.000000                  NaN            NaN    1.000000  

[11 rows x 21 columns]

[2] Distribusi Kelas Target (0: Aman, 1: Melanggar):
label
0    459
1     85
Name: count, dtype: int64

/tmp/ipykernel_1991/1767416320.py:70: FutureWarning: 

Passing `palette` without assigning `hue` is deprecated and will be removed in v0.14.0. Assign the `x` variable to `hue` and set `legend=False` for the same effect.

  sns.countplot(data=df_eda, x='label', palette='Set2', ax=ax)

<Figure size 600x400 with 1 Axes>
/tmp/ipykernel_1991/1767416320.py:97: UserWarning: Could not infer format, so each element will be parsed individually, falling back to `dateutil`. To ensure parsing is consistent and as-expected, please specify a format.
  df_eda['checkin_hour_num'] = pd.to_datetime(df_eda['waktu_checkin'], errors='coerce').dt.hour.fillna(0)
/tmp/ipykernel_1991/1767416320.py:98: UserWarning: Could not infer format, so each element will be parsed individually, falling back to `dateutil`. To ensure parsing is consistent and as-expected, please specify a format.
  df_eda['checkout_hour_num'] = pd.to_datetime(df_eda['waktu_checkout'], errors='coerce').dt.hour.fillna(0)


[3] Rata-rata Pelanggaran Berdasarkan Seluruh Kategori:

--- Tingkat Pelanggaran Berdasarkan status_pasutri ---
status_pasutri
Belum Menikah    0.164
Menikah          0.145
Name: label, dtype: float64

--- Tingkat Pelanggaran Berdasarkan status_kewarganegaraan ---
status_kewarganegaraan
WNI    0.156
Name: label, dtype: float64

--- Tingkat Pelanggaran Berdasarkan jenis_sewa ---
jenis_sewa
Harian     0.157
Bulanan    0.000
Name: label, dtype: float64

--- Tingkat Pelanggaran Berdasarkan metode_pembayaran ---
metode_pembayaran
Kartu Kredit    0.280
Others          0.158
Cash            0.150
QRIS            0.148
Kartu Debit     0.107
Name: label, dtype: float64

[4] Korelasi Antar Fitur:

<Figure size 1000x800 with 2 Axes>

[5] Visualisasi Pola Penyewaan (Dikaitkan dengan Klasifikasi):

<Figure size 1200x500 with 1 Axes>
<Figure size 1200x500 with 1 Axes>
<Figure size 800x500 with 1 Axes>
<Figure size 800x500 with 1 Axes>
<Figure size 800x500 with 1 Axes>
<Figure size 800x500 with 1 Axes>

[5b] Visualisasi Metode Pembayaran 'Others' (metode_lain):

/tmp/ipykernel_1991/1767416320.py:205: FutureWarning: 

Passing `palette` without assigning `hue` is deprecated and will be removed in v0.14.0. Assign the `y` variable to `hue` and set `legend=False` for the same effect.

  sns.countplot(

<Figure size 1000x600 with 1 Axes>
/tmp/ipykernel_1991/1767416320.py:239: UserWarning: Could not infer format, so each element will be parsed individually, falling back to `dateutil`. To ensure parsing is consistent and as-expected, please specify a format.
  rentals['waktu_checkout_dt'] = pd.to_datetime(rentals['waktu_checkout'], errors='coerce')
/tmp/ipykernel_1991/1767416320.py:242: UserWarning: Could not infer format, so each element will be parsed individually, falling back to `dateutil`. To ensure parsing is consistent and as-expected, please specify a format.
  rentals['waktu_checkin_dt'] = pd.to_datetime(rentals['waktu_checkin'], errors='coerce')


=== Proses Feature Selection ===

-> Fitur yang dipilih (9 terbaik): ['lama_menginap', 'checkin_hour', 'checkout_hour', 'metode_pembayaran', 'status_pasutri', 'jenis_sewa', 'status_kewarganegaraan']

===== Decision Tree =====
Classification Report:
              precision    recall  f1-score   support

           0     1.0000    0.9565    0.9778        92
           1     0.8095    1.0000    0.8947        17

    accuracy                         0.9633       109
   macro avg     0.9048    0.9783    0.9363       109
weighted avg     0.9703    0.9633    0.9648       109

Confusion Matrix (Text):
                Predicted
               0       1
          -----------------
True   0 |     88       4
       1 |      0      17
          -----------------


===== Random Forest =====
Classification Report:
              precision    recall  f1-score   support

           0     1.0000    0.9565    0.9778        92
           1     0.8095    1.0000    0.8947        17

    accuracy                         0.9633       109
   macro avg     0.9048    0.9783    0.9363       109
weighted avg     0.9703    0.9633    0.9648       109

Confusion Matrix (Text):
                Predicted
               0       1
          -----------------
True   0 |     88       4
       1 |      0      17
          -----------------


===== KNN =====
Classification Report:
              precision    recall  f1-score   support

           0     0.9184    0.9783    0.9474        92
           1     0.8182    0.5294    0.6429        17

    accuracy                         0.9083       109
   macro avg     0.8683    0.7538    0.7951       109
weighted avg     0.9027    0.9083    0.8999       109

Confusion Matrix (Text):
                Predicted
               0       1
          -----------------
True   0 |     90       2
       1 |      8       9
          -----------------


===== SVM RBF =====
Classification Report:
              precision    recall  f1-score   support

           0     1.0000    0.9565    0.9778        92
           1     0.8095    1.0000    0.8947        17

    accuracy                         0.9633       109
   macro avg     0.9048    0.9783    0.9363       109
weighted avg     0.9703    0.9633    0.9648       109

Confusion Matrix (Text):
                Predicted
               0       1
          -----------------
True   0 |     88       4
       1 |      0      17
          -----------------


===== XGBoost =====
Classification Report:
              precision    recall  f1-score   support

           0     1.0000    0.9565    0.9778        92
           1     0.8095    1.0000    0.8947        17

    accuracy                         0.9633       109
   macro avg     0.9048    0.9783    0.9363       109
weighted avg     0.9703    0.9633    0.9648       109

Confusion Matrix (Text):
                Predicted
               0       1
          -----------------
True   0 |     88       4
       1 |      0      17
          -----------------


=== Proses Export Model ===
Sukses! Model dan Label Encoders telah disimpan sebagai file: model_knn.pkl
\end{lstlisting}

\begin{lstlisting}[language=Python, caption={Kode 14}]
# ============================================================
# DATA RIIL (TANPA TUNING) --> jumlah fitur = 8
# ============================================================

import pandas as pd
import numpy as np
import joblib
import matplotlib.pyplot as plt
import seaborn as sns

from sklearn.model_selection import train_test_split
from sklearn.preprocessing import LabelEncoder, StandardScaler, MinMaxScaler
from sklearn.metrics import classification_report, confusion_matrix
from sklearn.pipeline import make_pipeline
from sklearn.feature_selection import SelectKBest, chi2, mutual_info_classif

from sklearn.tree import DecisionTreeClassifier
from sklearn.ensemble import RandomForestClassifier
from sklearn.neighbors import KNeighborsClassifier
from sklearn.svm import SVC
from xgboost import XGBClassifier

import matplotlib.pyplot as plt
import seaborn as sns

# ------------------------------------------------------------
# 1. LOAD DATA
# ------------------------------------------------------------
rentals    = pd.read_csv('rentals.csv')
logs       = pd.read_csv('rental_logs.csv')
# units      = pd.read_csv('units.csv')
# users      = pd.read_csv('users.csv')
# violations = pd.read_csv('violations.csv')

rentals.columns = rentals.columns.str.lower()

if 'komentar' not in rentals.columns:
    raise ValueError("Kolom 'komentar' tidak ditemukan di rentals.csv")

print("rentals.csv")
print(rentals.head())
print("logs.csv")
print(logs.head())


# ------------------------------------------------------------
# 2. LABELING
# ------------------------------------------------------------
rentals['komentar'] = rentals['komentar'].astype(str).str.lower().str.strip()
empty_values = ["", "nan", "na", "null", "-", "0", "tidak ada", "none", " "]
rentals['label'] = (~rentals['komentar'].isin(empty_values)).astype(int)


# ============================================================
# EXPLORATORY DATA ANALYSIS (EDA)
# ============================================================
print("\n" + "="*50)
print("=== EXPLORATORY DATA ANALYSIS (EDA) ===")
print("="*50)

df_eda = rentals.copy()

print("\n[1] Statistik Deskriptif:")
print(df_eda.describe(include='all'))

print("\n[2] Distribusi Kelas Target (0: Aman, 1: Melanggar):")
print(df_eda['label'].value_counts())

fig, ax = plt.subplots(figsize=(6, 4))
sns.countplot(data=df_eda, x='label', palette='Set2', ax=ax)
plt.title('Distribusi Kelas Klasifikasi (0: Aman, 1: Melanggar)')
plt.xlabel('Klasifikasi')
plt.ylabel('Jumlah Penyewaan')

y_max = df_eda['label'].value_counts().max()
plt.ylim(0, y_max * 1.15)

for p in ax.patches:
    if p.get_height() > 0:
        ax.annotate(f'{int(p.get_height())}', (p.get_x() + p.get_width() / 2., p.get_height()),
                    ha='center', va='baseline', fontsize=10, color='black', xytext=(0, 5),
                    textcoords='offset points')
plt.show()

kolom_dikecualikan = ['id', 'nama', 'nik', 'unit_number', 'tanggal_checkin', 'waktu_checkin',
                      'tanggal_checkout', 'waktu_checkout', 'komentar', 'user_email', 'diedit_oleh', 'created_at', 'metode_lain', 'classification']

semua_kategori = [col for col in df_eda.select_dtypes(include=['object', 'category']).columns if col not in kolom_dikecualikan]

print("\n[3] Rata-rata Pelanggaran Berdasarkan Seluruh Kategori:")
for col in semua_kategori:
    if col in df_eda.columns:
        print(f"\n--- Tingkat Pelanggaran Berdasarkan {col} ---")
        print(df_eda.groupby(col)['label'].mean().sort_values(ascending=False).round(3))

df_eda['lama_menginap_num'] = pd.to_numeric(df_eda['lama_menginap'], errors='coerce').fillna(0)
df_eda['checkin_hour_num'] = pd.to_datetime(df_eda['waktu_checkin'], errors='coerce').dt.hour.fillna(0)
df_eda['checkout_hour_num'] = pd.to_datetime(df_eda['waktu_checkout'], errors='coerce').dt.hour.fillna(0)

le_eda = LabelEncoder()
for c in semua_kategori:
    if c in df_eda.columns:
        df_eda[c + '_encoded'] = le_eda.fit_transform(df_eda[c].astype(str))

print("\n[4] Korelasi Antar Fitur:")
kolom_korelasi = ['lama_menginap_num', 'checkin_hour_num', 'checkout_hour_num', 'label']
kolom_korelasi += [c + '_encoded' for c in semua_kategori if c + '_encoded' in df_eda.columns]

corr_matrix = df_eda[kolom_korelasi].corr()

plt.figure(figsize=(10, 8))
sns.heatmap(corr_matrix, annot=True, cmap='coolwarm', fmt='.2f', linewidths=0.5)
plt.title('Heatmap Korelasi Antar Fitur dan Target (Klasifikasi)')
plt.tight_layout()
plt.show()

print("\n[5] Visualisasi Pola Penyewaan (Dikaitkan dengan Klasifikasi):")

custom_line_palette = {0: 'blue', 1: 'red'}
custom_labels = ['0: Aman', '1: Melanggar']

fig, ax1 = plt.subplots(figsize=(12, 5))
waktu_df = df_eda.groupby(['checkin_hour_num', 'label']).size().reset_index(name='jumlah')

sns.lineplot(data=waktu_df, x='checkin_hour_num', y='jumlah', hue='label', marker='o', palette=custom_line_palette, ax=ax1)

plt.title('Pola Waktu Kedatangan Berdasarkan Klasifikasi')
plt.xlabel('Jam Check-in (0-23)')
plt.ylabel('Jumlah Penyewaan')
plt.xticks(range(0, 24))
plt.grid(True, alpha=0.3)

y_max_line = waktu_df['jumlah'].max()
plt.ylim(-y_max_line * 0.20, y_max_line * 1.15)

handles, _ = ax1.get_legend_handles_labels()
plt.legend(handles=handles, labels=custom_labels, title='Klasifikasi')

for l in waktu_df['label'].unique():
    sub = waktu_df[waktu_df['label'] == l]
    warna_teks = custom_line_palette[l]

    for x, y in zip(sub['checkin_hour_num'], sub['jumlah']):
        y_offset = 8 if l == 0 else -12
        va_align = 'bottom' if l == 0 else 'top'
        plt.annotate(f'{int(y)}', (x, y), textcoords="offset points", xytext=(0, y_offset),
                     ha='center', va=va_align, fontsize=9, color=warna_teks, fontweight='bold')
plt.show()

fig, ax2 = plt.subplots(figsize=(12, 5))
lama_df = df_eda.groupby(['lama_menginap_num', 'label']).size().reset_index(name='jumlah')

sns.lineplot(data=lama_df, x='lama_menginap_num', y='jumlah', hue='label', marker='o', palette=custom_line_palette, ax=ax2)

plt.title('Tren Lama Menginap Berdasarkan Klasifikasi')
plt.xlabel('Lama Menginap (Hari)')
plt.ylabel('Jumlah Penyewaan')
plt.grid(True, alpha=0.3)

handles2, _ = ax2.get_legend_handles_labels()
plt.legend(handles=handles2, labels=custom_labels, title='Klasifikasi')
plt.show()

for col in semua_kategori:
    if col in df_eda.columns:
        fig, ax = plt.subplots(figsize=(8, 5))
        sns.countplot(data=df_eda, x=col, hue='label', palette='magma', ax=ax)

        judul_kolom = col.replace('_', ' ').title()
        plt.title(f'Perbandingan Jumlah Pelanggaran Berdasarkan {judul_kolom}')
        plt.xlabel(judul_kolom)
        plt.ylabel('Jumlah Penyewaan')
        plt.xticks(rotation=0)

        handles_bar, _ = ax.get_legend_handles_labels()
        plt.legend(handles=handles_bar, labels=custom_labels, title='Klasifikasi')

        y_max_bar = df_eda.groupby([col, 'label']).size().max()
        plt.ylim(0, y_max_bar * 1.15)

        for p in ax.patches:
            if p.get_height() > 0:
                ax.annotate(f'{int(p.get_height())}', (p.get_x() + p.get_width() / 2., p.get_height()),
                            ha='center', va='baseline', fontsize=9, color='black', xytext=(0, 5),
                            textcoords='offset points')
        plt.tight_layout()
        plt.show()


print("\n[5b] Visualisasi Metode Pembayaran 'Others' (metode_lain):")

# Memastikan kolom 'metode_lain' berbentuk string dan membersihkan spasi/nilai kosong
df_eda['metode_lain'] = df_eda['metode_lain'].astype(str).str.strip().str.title()
empty_lain = ["", "Nan", "Na", "Null", "-", "0", "None"]

# Filter hanya baris yang memiliki nilai di 'metode_lain'
df_others = df_eda[~df_eda['metode_lain'].isin(empty_lain)]

if not df_others.empty:
    fig, ax = plt.subplots(figsize=(10, 6))

    # Menghitung frekuensi masing-masing metode pembayaran lain
    order_others = df_others['metode_lain'].value_counts().index

    sns.countplot(
        data=df_others,
        y='metode_lain',
        order=order_others,
        palette='viridis',
        ax=ax
    )

    plt.title('Detail Metode Pembayaran "Others" (metode_lain)')
    plt.xlabel('Jumlah Penyewaan')
    plt.ylabel('Jenis Pembayaran Lainnya')

    # Menambahkan anotasi angka di sebelah setiap bar
    for p in ax.patches:
        width = p.get_width()
        if width > 0:
            ax.annotate(f'{int(width)}',
                        (width, p.get_y() + p.get_height() / 2.),
                        ha='left', va='center', fontsize=10, color='black',
                        xytext=(5, 0), textcoords='offset points')

    # Perlebar batas x-axis sedikit agar angka anotasi tidak terpotong
    x_max = df_others['metode_lain'].value_counts().max()
    plt.xlim(0, x_max * 1.15)

    plt.tight_layout()
    plt.show()
else:
    print("Tidak ada data metode pembayaran 'Others' yang valid untuk ditampilkan.")


# ------------------------------------------------------------
# 3. FEATURE ENGINEERING & PREPARATION
# ------------------------------------------------------------
rentals['waktu_checkout_dt'] = pd.to_datetime(rentals['waktu_checkout'], errors='coerce')
rentals['checkout_hour'] = rentals['waktu_checkout_dt'].dt.hour.fillna(0).astype(int)

rentals['waktu_checkin_dt'] = pd.to_datetime(rentals['waktu_checkin'], errors='coerce')
rentals['checkin_hour'] = rentals['waktu_checkin_dt'].dt.hour.fillna(0).astype(int)

rentals['lama_menginap'] = pd.to_numeric(rentals['lama_menginap'], errors='coerce').fillna(0)

cat_cols = ['jenis_sewa', 'metode_pembayaran', 'status_pasutri', 'status_kewarganegaraan']
for c in cat_cols:
    rentals[c] = rentals[c].fillna('Unknown').astype(str)

numeric_cols = ['checkout_hour', 'checkin_hour', 'lama_menginap']
feature_cols = cat_cols + numeric_cols

X = rentals[feature_cols].copy()
y = rentals['label']

# ------------------------------------------------------------
# 4. TRAIN-TEST SPLIT
# ------------------------------------------------------------
X_train, X_test, y_train, y_test = train_test_split(
    X, y, stratify=y, test_size=0.2, random_state=42
)

# ------------------------------------------------------------
# 4B. ENCODING MENCEGAH DATA LEAKAGE
# ------------------------------------------------------------
label_encoders = {}

for c in cat_cols:
    le = LabelEncoder()
    X_train[c] = le.fit_transform(X_train[c])

    X_test[c] = X_test[c].map(lambda s: s if s in le.classes_ else 'Unknown')
    if 'Unknown' not in le.classes_:
        le.classes_ = np.append(le.classes_, 'Unknown')

    X_test[c] = le.transform(X_test[c])

    label_encoders[c] = le

# ------------------------------------------------------------
# 5. FEATURE SELECTION
# ------------------------------------------------------------
print("\n=== Proses Feature Selection ===")
chi2_selector = SelectKBest(score_func=chi2, k='all')
chi2_selector.fit(X_train, y_train)

train_df = X_train.copy()
train_df['label'] = y_train
pearson_scores = train_df.corr(method='pearson')['label'].drop('label').abs().values

mi_selector = SelectKBest(score_func=mutual_info_classif, k='all')
mi_selector.fit(X_train, y_train)

fs_df = pd.DataFrame({
    'Feature': X_train.columns,
    'Chi-Square': chi2_selector.scores_,
    'Pearson (Abs)': pearson_scores,
    'Mutual Info': mi_selector.scores_
})

scaler = MinMaxScaler()
fs_df[['Chi2_Norm', 'Pearson_Norm', 'MI_Norm']] = scaler.fit_transform(
    fs_df[['Chi-Square', 'Pearson (Abs)', 'Mutual Info']]
)
fs_df['Total_Score'] = fs_df[['Chi2_Norm', 'Pearson_Norm', 'MI_Norm']].mean(axis=1)
fs_df = fs_df.sort_values(by='Total_Score', ascending=False).reset_index(drop=True)

# Jumlah Fitur
n_select = 8
selected_features = fs_df['Feature'].head(n_select).tolist()
print(f"\n-> Fitur yang dipilih ({n_select} terbaik): {selected_features}")

X_train = X_train[selected_features]
X_test = X_test[selected_features]

# ------------------------------------------------------------
# 6. Fungsi Evaluasi
# ------------------------------------------------------------
def evaluate(name, model):
    y_pred = model.predict(X_test)
    print(f"\n===== {name} =====")
    print("Classification Report:")
    print(classification_report(y_test, y_pred, digits=4))
    cm = confusion_matrix(y_test, y_pred)
    print("Confusion Matrix (Text):")
    print("                Predicted")
    print("               0       1")
    print("          -----------------")
    print(f"True   0 | {cm[0][0]:6d}  {cm[0][1]:6d}")
    print(f"       1 | {cm[1][0]:6d}  {cm[1][1]:6d}")
    print("          -----------------\n")

# ------------------------------------------------------------
# 7. Latih Semua Model (TANPA TUNING)
# ------------------------------------------------------------
models = {}
custom_weight = {0: 1, 1: 4}

# Model menggunakan pengaturan tetap
models["Decision Tree"] = DecisionTreeClassifier(max_depth=6, class_weight=custom_weight, random_state=42)
models["Random Forest"] = RandomForestClassifier(n_estimators=200, max_depth=7, class_weight=custom_weight, random_state=42)
models["KNN"] = make_pipeline(StandardScaler(), KNeighborsClassifier(n_neighbors=5, weights='distance'))
models["SVM RBF"] = make_pipeline(StandardScaler(), SVC(kernel='rbf', probability=True, class_weight=custom_weight, random_state=42))
models["XGBoost"] = XGBClassifier(n_estimators=200, max_depth=5, learning_rate=0.05, scale_pos_weight=4, eval_metric='logloss', random_state=42)

for name, model in models.items():
    model.fit(X_train, y_train)
    evaluate(name, model)

# ------------------------------------------------------------
# 8. Eksport Model Terbaik & Label Encoders
# ------------------------------------------------------------
print("\n=== Proses Export Model ===")

# Gabungkan model dan encoders ke dalam satu dictionary agar rapi
export_package = {
    'model': models["KNN"],
    'label_encoders': label_encoders,
    'selected_features': selected_features # Penting agar tahu urutan kolom masuk saat predict
}

# Menentukan nama file yang sesuai untuk KNN
filename = 'model_knn.pkl'

# Menyimpan package ke dalam file .pkl
joblib.dump(export_package, filename)
print(f"Sukses! Model dan Label Encoders telah disimpan sebagai file: {filename}")
\end{lstlisting}

\begin{lstlisting}[caption={Output 14}]
rentals.csv
   id                   nama               nik status_pasutri  \
0   3  N*** Y**** S**** D***  3***************  Belum Menikah   
1   4     A***** N** H******  3***************  Belum Menikah   
2   5     A***** A*** S*****  3***************        Menikah   
3   6            C**** W****  1***************  Belum Menikah   
4   7             S*** U****  3***************  Belum Menikah   

  status_kewarganegaraan jenis_sewa unit_number metode_pembayaran metode_lain  \
0                    WNI     Harian     B-20-27              Cash         NaN   
1                    WNI     Harian     A-15-25              Cash         NaN   
2                    WNI     Harian     A-06-24            Others    Aplikasi   
3                    WNI     Harian     D-15-21              QRIS         NaN   
4                    WNI     Harian       A-8-8            Others    Transfer   

  tanggal_checkin waktu_checkin tanggal_checkout waktu_checkout  \
0      2026-05-20      13:08:00       2026-05-21            NaN   
1      2026-05-19      13:25:00       2026-05-22            NaN   
2      2026-05-17      13:27:00       2026-05-22            NaN   
3      2026-05-19      13:29:00       2026-05-24            NaN   
4      2025-01-01      19:58:00       2025-01-06       11:22:00   

   lama_menginap komentar  user_email  user_pengguna_id  editor_pengguna_id  \
0              1      NaN         NaN              1267                 NaN   
1              3      NaN         NaN              1267                 NaN   
2              5      NaN         NaN              1267                 NaN   
3              5      NaN         NaN              1267                 NaN   
4              5      NaN         NaN              1271              1271.0   

            created_at classification  
0  2026-05-20 13:09:49            NaN  
1  2026-05-20 13:26:04            NaN  
2  2026-05-20 13:28:03            NaN  
3  2026-05-20 13:30:18            NaN  
4  2026-05-25 10:06:33         normal  
logs.csv
     id  rental_id  action field_changed  old_value  \
0     6          3  create           all        NaN   
1     7          4  create           all        NaN   
2     8          5  create           all        NaN   
3     9          6  create           all        NaN   
4  4906          7  create           all        NaN   

                                           new_value  email  pengguna_id  \
0  {"nama":"N*** Y**** S**** D***","nik":"3*************************************"nama":"A***** N** H******","nik":"3****************************************"nama":"A***** A*** S*****","nik":"3****************************************"nama":"C**** W****","nik":"1***************"...    NaN         1267   
4  {"nama":"S*** U****","nik":"3***************",...    NaN         1271   

             timestamp  
0  2026-05-20 13:09:49  
1  2026-05-20 13:26:04  
2  2026-05-20 13:28:03  
3  2026-05-20 13:30:18  
4  2026-05-25 10:06:33  

==================================================
=== EXPLORATORY DATA ANALYSIS (EDA) ===
==================================================

[1] Statistik Deskriptif:
                id                         nama           nik status_pasutri  \
count   544.000000                          544  5.440000e+02            544   
unique         NaN                          207           NaN              2   
top            NaN  Y***** M**** S**** S*******           NaN  Belum Menikah   
freq           NaN                            8           NaN            330   
mean    274.538603                          NaN  3.214330e+15            NaN   
std     157.250941                          NaN  1.204147e+15            NaN   
min       3.000000                          NaN  1.202052e+15            NaN   
25%     138.750000                          NaN  3.173074e+15            NaN   
50%     274.500000                          NaN  3.209121e+15            NaN   
75%     410.250000                          NaN  3.276064e+15            NaN   
max     553.000000                          NaN  9.104017e+15            NaN   

       status_kewarganegaraan jenis_sewa unit_number metode_pembayaran  \
count                     544        544         544               544   
unique                      1          2          12                 5   
top                       WNI     Harian       B-1-8            Others   
freq                      544        543          90               209   
mean                      NaN        NaN         NaN               NaN   
std                       NaN        NaN         NaN               NaN   
min                       NaN        NaN         NaN               NaN   
25%                       NaN        NaN         NaN               NaN   
50%                       NaN        NaN         NaN               NaN   
75%                       NaN        NaN         NaN               NaN   
max                       NaN        NaN         NaN               NaN   

       metode_lain tanggal_checkin  ... tanggal_checkout waktu_checkout  \
count          209             544  ...              544            538   
unique           2             311  ...              300            232   
top       Transfer      2026-02-06  ...       2026-02-08       12:00:00   
freq           208               5  ...                5             82   
mean           NaN             NaN  ...              NaN            NaN   
std            NaN             NaN  ...              NaN            NaN   
min            NaN             NaN  ...              NaN            NaN   
25%            NaN             NaN  ...              NaN            NaN   
50%            NaN             NaN  ...              NaN            NaN   
75%            NaN             NaN  ...              NaN            NaN   
max            NaN             NaN  ...              NaN            NaN   

       lama_menginap  komentar user_email  user_pengguna_id  \
count     544.000000       544        0.0        544.000000   
unique           NaN        11        NaN               NaN   
top              NaN       nan        NaN               NaN   
freq             NaN       459        NaN               NaN   
mean        2.641544       NaN        NaN       1271.904412   
std         1.792711       NaN        NaN          1.177622   
min         1.000000       NaN        NaN       1262.000000   
25%         2.000000       NaN        NaN       1271.000000   
50%         2.000000       NaN        NaN       1272.000000   
75%         3.000000       NaN        NaN       1273.000000   
max        31.000000       NaN        NaN       1273.000000   

        editor_pengguna_id           created_at classification       label  
count           538.000000                  544            537  544.000000  
unique                 NaN                  544              2         NaN  
top                    NaN  2026-05-30 15:00:26         normal         NaN  
freq                   NaN                    1            452         NaN  
mean           1271.977695                  NaN            NaN    0.156250  
std               0.923249                  NaN            NaN    0.363426  
min            1262.000000                  NaN            NaN    0.000000  
25%            1271.000000                  NaN            NaN    0.000000  
50%            1272.000000                  NaN            NaN    0.000000  
75%            1273.000000                  NaN            NaN    0.000000  
max            1273.000000                  NaN            NaN    1.000000  

[11 rows x 21 columns]

[2] Distribusi Kelas Target (0: Aman, 1: Melanggar):
label
0    459
1     85
Name: count, dtype: int64

/tmp/ipykernel_1991/4099437245.py:70: FutureWarning: 

Passing `palette` without assigning `hue` is deprecated and will be removed in v0.14.0. Assign the `x` variable to `hue` and set `legend=False` for the same effect.

  sns.countplot(data=df_eda, x='label', palette='Set2', ax=ax)

<Figure size 600x400 with 1 Axes>

[3] Rata-rata Pelanggaran Berdasarkan Seluruh Kategori:

--- Tingkat Pelanggaran Berdasarkan status_pasutri ---
status_pasutri
Belum Menikah    0.164
Menikah          0.145
Name: label, dtype: float64

--- Tingkat Pelanggaran Berdasarkan status_kewarganegaraan ---
status_kewarganegaraan
WNI    0.156
Name: label, dtype: float64

--- Tingkat Pelanggaran Berdasarkan jenis_sewa ---
jenis_sewa
Harian     0.157
Bulanan    0.000
Name: label, dtype: float64

--- Tingkat Pelanggaran Berdasarkan metode_pembayaran ---
metode_pembayaran
Kartu Kredit    0.280
Others          0.158
Cash            0.150
QRIS            0.148
Kartu Debit     0.107
Name: label, dtype: float64

[4] Korelasi Antar Fitur:

/tmp/ipykernel_1991/4099437245.py:97: UserWarning: Could not infer format, so each element will be parsed individually, falling back to `dateutil`. To ensure parsing is consistent and as-expected, please specify a format.
  df_eda['checkin_hour_num'] = pd.to_datetime(df_eda['waktu_checkin'], errors='coerce').dt.hour.fillna(0)
/tmp/ipykernel_1991/4099437245.py:98: UserWarning: Could not infer format, so each element will be parsed individually, falling back to `dateutil`. To ensure parsing is consistent and as-expected, please specify a format.
  df_eda['checkout_hour_num'] = pd.to_datetime(df_eda['waktu_checkout'], errors='coerce').dt.hour.fillna(0)

<Figure size 1000x800 with 2 Axes>

[5] Visualisasi Pola Penyewaan (Dikaitkan dengan Klasifikasi):

<Figure size 1200x500 with 1 Axes>
<Figure size 1200x500 with 1 Axes>
<Figure size 800x500 with 1 Axes>
<Figure size 800x500 with 1 Axes>
<Figure size 800x500 with 1 Axes>
<Figure size 800x500 with 1 Axes>

[5b] Visualisasi Metode Pembayaran 'Others' (metode_lain):

/tmp/ipykernel_1991/4099437245.py:205: FutureWarning: 

Passing `palette` without assigning `hue` is deprecated and will be removed in v0.14.0. Assign the `y` variable to `hue` and set `legend=False` for the same effect.

  sns.countplot(

<Figure size 1000x600 with 1 Axes>
/tmp/ipykernel_1991/4099437245.py:239: UserWarning: Could not infer format, so each element will be parsed individually, falling back to `dateutil`. To ensure parsing is consistent and as-expected, please specify a format.
  rentals['waktu_checkout_dt'] = pd.to_datetime(rentals['waktu_checkout'], errors='coerce')
/tmp/ipykernel_1991/4099437245.py:242: UserWarning: Could not infer format, so each element will be parsed individually, falling back to `dateutil`. To ensure parsing is consistent and as-expected, please specify a format.
  rentals['waktu_checkin_dt'] = pd.to_datetime(rentals['waktu_checkin'], errors='coerce')


=== Proses Feature Selection ===

-> Fitur yang dipilih (8 terbaik): ['lama_menginap', 'checkout_hour', 'checkin_hour', 'metode_pembayaran', 'status_kewarganegaraan', 'status_pasutri', 'jenis_sewa']

===== Decision Tree =====
Classification Report:
              precision    recall  f1-score   support

           0     1.0000    0.9565    0.9778        92
           1     0.8095    1.0000    0.8947        17

    accuracy                         0.9633       109
   macro avg     0.9048    0.9783    0.9363       109
weighted avg     0.9703    0.9633    0.9648       109

Confusion Matrix (Text):
                Predicted
               0       1
          -----------------
True   0 |     88       4
       1 |      0      17
          -----------------


===== Random Forest =====
Classification Report:
              precision    recall  f1-score   support

           0     1.0000    0.9565    0.9778        92
           1     0.8095    1.0000    0.8947        17

    accuracy                         0.9633       109
   macro avg     0.9048    0.9783    0.9363       109
weighted avg     0.9703    0.9633    0.9648       109

Confusion Matrix (Text):
                Predicted
               0       1
          -----------------
True   0 |     88       4
       1 |      0      17
          -----------------


===== KNN =====
Classification Report:
              precision    recall  f1-score   support

           0     0.9184    0.9783    0.9474        92
           1     0.8182    0.5294    0.6429        17

    accuracy                         0.9083       109
   macro avg     0.8683    0.7538    0.7951       109
weighted avg     0.9027    0.9083    0.8999       109

Confusion Matrix (Text):
                Predicted
               0       1
          -----------------
True   0 |     90       2
       1 |      8       9
          -----------------


===== SVM RBF =====
Classification Report:
              precision    recall  f1-score   support

           0     1.0000    0.9565    0.9778        92
           1     0.8095    1.0000    0.8947        17

    accuracy                         0.9633       109
   macro avg     0.9048    0.9783    0.9363       109
weighted avg     0.9703    0.9633    0.9648       109

Confusion Matrix (Text):
                Predicted
               0       1
          -----------------
True   0 |     88       4
       1 |      0      17
          -----------------


===== XGBoost =====
Classification Report:
              precision    recall  f1-score   support

           0     1.0000    0.9565    0.9778        92
           1     0.8095    1.0000    0.8947        17

    accuracy                         0.9633       109
   macro avg     0.9048    0.9783    0.9363       109
weighted avg     0.9703    0.9633    0.9648       109

Confusion Matrix (Text):
                Predicted
               0       1
          -----------------
True   0 |     88       4
       1 |      0      17
          -----------------


=== Proses Export Model ===
Sukses! Model dan Label Encoders telah disimpan sebagai file: model_knn.pkl
\end{lstlisting}

\begin{lstlisting}[language=Python, caption={Kode 15}]
# ============================================================
# DATA RIIL (TANPA TUNING) --> jumlah fitur = 7
# ============================================================

import pandas as pd
import numpy as np
import joblib
import matplotlib.pyplot as plt
import seaborn as sns

from sklearn.model_selection import train_test_split
from sklearn.preprocessing import LabelEncoder, StandardScaler, MinMaxScaler
from sklearn.metrics import classification_report, confusion_matrix
from sklearn.pipeline import make_pipeline
from sklearn.feature_selection import SelectKBest, chi2, mutual_info_classif

from sklearn.tree import DecisionTreeClassifier
from sklearn.ensemble import RandomForestClassifier
from sklearn.neighbors import KNeighborsClassifier
from sklearn.svm import SVC
from xgboost import XGBClassifier

import matplotlib.pyplot as plt
import seaborn as sns

# ------------------------------------------------------------
# 1. LOAD DATA
# ------------------------------------------------------------
rentals    = pd.read_csv('rentals.csv')
logs       = pd.read_csv('rental_logs.csv')
# units      = pd.read_csv('units.csv')
# users      = pd.read_csv('users.csv')
# violations = pd.read_csv('violations.csv')

rentals.columns = rentals.columns.str.lower()

if 'komentar' not in rentals.columns:
    raise ValueError("Kolom 'komentar' tidak ditemukan di rentals.csv")

print("rentals.csv")
print(rentals.head())
print("logs.csv")
print(logs.head())


# ------------------------------------------------------------
# 2. LABELING
# ------------------------------------------------------------
rentals['komentar'] = rentals['komentar'].astype(str).str.lower().str.strip()
empty_values = ["", "nan", "na", "null", "-", "0", "tidak ada", "none", " "]
rentals['label'] = (~rentals['komentar'].isin(empty_values)).astype(int)


# ============================================================
# EXPLORATORY DATA ANALYSIS (EDA)
# ============================================================
print("\n" + "="*50)
print("=== EXPLORATORY DATA ANALYSIS (EDA) ===")
print("="*50)

df_eda = rentals.copy()

print("\n[1] Statistik Deskriptif:")
print(df_eda.describe(include='all'))

print("\n[2] Distribusi Kelas Target (0: Aman, 1: Melanggar):")
print(df_eda['label'].value_counts())

fig, ax = plt.subplots(figsize=(6, 4))
sns.countplot(data=df_eda, x='label', palette='Set2', ax=ax)
plt.title('Distribusi Kelas Klasifikasi (0: Aman, 1: Melanggar)')
plt.xlabel('Klasifikasi')
plt.ylabel('Jumlah Penyewaan')

y_max = df_eda['label'].value_counts().max()
plt.ylim(0, y_max * 1.15)

for p in ax.patches:
    if p.get_height() > 0:
        ax.annotate(f'{int(p.get_height())}', (p.get_x() + p.get_width() / 2., p.get_height()),
                    ha='center', va='baseline', fontsize=10, color='black', xytext=(0, 5),
                    textcoords='offset points')
plt.show()

kolom_dikecualikan = ['id', 'nama', 'nik', 'unit_number', 'tanggal_checkin', 'waktu_checkin',
                      'tanggal_checkout', 'waktu_checkout', 'komentar', 'user_email', 'diedit_oleh', 'created_at', 'metode_lain', 'classification']

semua_kategori = [col for col in df_eda.select_dtypes(include=['object', 'category']).columns if col not in kolom_dikecualikan]

print("\n[3] Rata-rata Pelanggaran Berdasarkan Seluruh Kategori:")
for col in semua_kategori:
    if col in df_eda.columns:
        print(f"\n--- Tingkat Pelanggaran Berdasarkan {col} ---")
        print(df_eda.groupby(col)['label'].mean().sort_values(ascending=False).round(3))

df_eda['lama_menginap_num'] = pd.to_numeric(df_eda['lama_menginap'], errors='coerce').fillna(0)
df_eda['checkin_hour_num'] = pd.to_datetime(df_eda['waktu_checkin'], errors='coerce').dt.hour.fillna(0)
df_eda['checkout_hour_num'] = pd.to_datetime(df_eda['waktu_checkout'], errors='coerce').dt.hour.fillna(0)

le_eda = LabelEncoder()
for c in semua_kategori:
    if c in df_eda.columns:
        df_eda[c + '_encoded'] = le_eda.fit_transform(df_eda[c].astype(str))

print("\n[4] Korelasi Antar Fitur:")
kolom_korelasi = ['lama_menginap_num', 'checkin_hour_num', 'checkout_hour_num', 'label']
kolom_korelasi += [c + '_encoded' for c in semua_kategori if c + '_encoded' in df_eda.columns]

corr_matrix = df_eda[kolom_korelasi].corr()

plt.figure(figsize=(10, 8))
sns.heatmap(corr_matrix, annot=True, cmap='coolwarm', fmt='.2f', linewidths=0.5)
plt.title('Heatmap Korelasi Antar Fitur dan Target (Klasifikasi)')
plt.tight_layout()
plt.show()

print("\n[5] Visualisasi Pola Penyewaan (Dikaitkan dengan Klasifikasi):")

custom_line_palette = {0: 'blue', 1: 'red'}
custom_labels = ['0: Aman', '1: Melanggar']

fig, ax1 = plt.subplots(figsize=(12, 5))
waktu_df = df_eda.groupby(['checkin_hour_num', 'label']).size().reset_index(name='jumlah')

sns.lineplot(data=waktu_df, x='checkin_hour_num', y='jumlah', hue='label', marker='o', palette=custom_line_palette, ax=ax1)

plt.title('Pola Waktu Kedatangan Berdasarkan Klasifikasi')
plt.xlabel('Jam Check-in (0-23)')
plt.ylabel('Jumlah Penyewaan')
plt.xticks(range(0, 24))
plt.grid(True, alpha=0.3)

y_max_line = waktu_df['jumlah'].max()
plt.ylim(-y_max_line * 0.20, y_max_line * 1.15)

handles, _ = ax1.get_legend_handles_labels()
plt.legend(handles=handles, labels=custom_labels, title='Klasifikasi')

for l in waktu_df['label'].unique():
    sub = waktu_df[waktu_df['label'] == l]
    warna_teks = custom_line_palette[l]

    for x, y in zip(sub['checkin_hour_num'], sub['jumlah']):
        y_offset = 8 if l == 0 else -12
        va_align = 'bottom' if l == 0 else 'top'
        plt.annotate(f'{int(y)}', (x, y), textcoords="offset points", xytext=(0, y_offset),
                     ha='center', va=va_align, fontsize=9, color=warna_teks, fontweight='bold')
plt.show()

fig, ax2 = plt.subplots(figsize=(12, 5))
lama_df = df_eda.groupby(['lama_menginap_num', 'label']).size().reset_index(name='jumlah')

sns.lineplot(data=lama_df, x='lama_menginap_num', y='jumlah', hue='label', marker='o', palette=custom_line_palette, ax=ax2)

plt.title('Tren Lama Menginap Berdasarkan Klasifikasi')
plt.xlabel('Lama Menginap (Hari)')
plt.ylabel('Jumlah Penyewaan')
plt.grid(True, alpha=0.3)

handles2, _ = ax2.get_legend_handles_labels()
plt.legend(handles=handles2, labels=custom_labels, title='Klasifikasi')
plt.show()

for col in semua_kategori:
    if col in df_eda.columns:
        fig, ax = plt.subplots(figsize=(8, 5))
        sns.countplot(data=df_eda, x=col, hue='label', palette='magma', ax=ax)

        judul_kolom = col.replace('_', ' ').title()
        plt.title(f'Perbandingan Jumlah Pelanggaran Berdasarkan {judul_kolom}')
        plt.xlabel(judul_kolom)
        plt.ylabel('Jumlah Penyewaan')
        plt.xticks(rotation=0)

        handles_bar, _ = ax.get_legend_handles_labels()
        plt.legend(handles=handles_bar, labels=custom_labels, title='Klasifikasi')

        y_max_bar = df_eda.groupby([col, 'label']).size().max()
        plt.ylim(0, y_max_bar * 1.15)

        for p in ax.patches:
            if p.get_height() > 0:
                ax.annotate(f'{int(p.get_height())}', (p.get_x() + p.get_width() / 2., p.get_height()),
                            ha='center', va='baseline', fontsize=9, color='black', xytext=(0, 5),
                            textcoords='offset points')
        plt.tight_layout()
        plt.show()


print("\n[5b] Visualisasi Metode Pembayaran 'Others' (metode_lain):")

# Memastikan kolom 'metode_lain' berbentuk string dan membersihkan spasi/nilai kosong
df_eda['metode_lain'] = df_eda['metode_lain'].astype(str).str.strip().str.title()
empty_lain = ["", "Nan", "Na", "Null", "-", "0", "None"]

# Filter hanya baris yang memiliki nilai di 'metode_lain'
df_others = df_eda[~df_eda['metode_lain'].isin(empty_lain)]

if not df_others.empty:
    fig, ax = plt.subplots(figsize=(10, 6))

    # Menghitung frekuensi masing-masing metode pembayaran lain
    order_others = df_others['metode_lain'].value_counts().index

    sns.countplot(
        data=df_others,
        y='metode_lain',
        order=order_others,
        palette='viridis',
        ax=ax
    )

    plt.title('Detail Metode Pembayaran "Others" (metode_lain)')
    plt.xlabel('Jumlah Penyewaan')
    plt.ylabel('Jenis Pembayaran Lainnya')

    # Menambahkan anotasi angka di sebelah setiap bar
    for p in ax.patches:
        width = p.get_width()
        if width > 0:
            ax.annotate(f'{int(width)}',
                        (width, p.get_y() + p.get_height() / 2.),
                        ha='left', va='center', fontsize=10, color='black',
                        xytext=(5, 0), textcoords='offset points')

    # Perlebar batas x-axis sedikit agar angka anotasi tidak terpotong
    x_max = df_others['metode_lain'].value_counts().max()
    plt.xlim(0, x_max * 1.15)

    plt.tight_layout()
    plt.show()
else:
    print("Tidak ada data metode pembayaran 'Others' yang valid untuk ditampilkan.")


# ------------------------------------------------------------
# 3. FEATURE ENGINEERING & PREPARATION
# ------------------------------------------------------------
rentals['waktu_checkout_dt'] = pd.to_datetime(rentals['waktu_checkout'], errors='coerce')
rentals['checkout_hour'] = rentals['waktu_checkout_dt'].dt.hour.fillna(0).astype(int)

rentals['waktu_checkin_dt'] = pd.to_datetime(rentals['waktu_checkin'], errors='coerce')
rentals['checkin_hour'] = rentals['waktu_checkin_dt'].dt.hour.fillna(0).astype(int)

rentals['lama_menginap'] = pd.to_numeric(rentals['lama_menginap'], errors='coerce').fillna(0)

cat_cols = ['jenis_sewa', 'metode_pembayaran', 'status_pasutri', 'status_kewarganegaraan']
for c in cat_cols:
    rentals[c] = rentals[c].fillna('Unknown').astype(str)

numeric_cols = ['checkout_hour', 'checkin_hour', 'lama_menginap']
feature_cols = cat_cols + numeric_cols

X = rentals[feature_cols].copy()
y = rentals['label']

# ------------------------------------------------------------
# 4. TRAIN-TEST SPLIT
# ------------------------------------------------------------
X_train, X_test, y_train, y_test = train_test_split(
    X, y, stratify=y, test_size=0.2, random_state=42
)

# ------------------------------------------------------------
# 4B. ENCODING MENCEGAH DATA LEAKAGE
# ------------------------------------------------------------
label_encoders = {}

for c in cat_cols:
    le = LabelEncoder()
    X_train[c] = le.fit_transform(X_train[c])

    X_test[c] = X_test[c].map(lambda s: s if s in le.classes_ else 'Unknown')
    if 'Unknown' not in le.classes_:
        le.classes_ = np.append(le.classes_, 'Unknown')

    X_test[c] = le.transform(X_test[c])

    label_encoders[c] = le

# ------------------------------------------------------------
# 5. FEATURE SELECTION
# ------------------------------------------------------------
print("\n=== Proses Feature Selection ===")
chi2_selector = SelectKBest(score_func=chi2, k='all')
chi2_selector.fit(X_train, y_train)

train_df = X_train.copy()
train_df['label'] = y_train
pearson_scores = train_df.corr(method='pearson')['label'].drop('label').abs().values

mi_selector = SelectKBest(score_func=mutual_info_classif, k='all')
mi_selector.fit(X_train, y_train)

fs_df = pd.DataFrame({
    'Feature': X_train.columns,
    'Chi-Square': chi2_selector.scores_,
    'Pearson (Abs)': pearson_scores,
    'Mutual Info': mi_selector.scores_
})

scaler = MinMaxScaler()
fs_df[['Chi2_Norm', 'Pearson_Norm', 'MI_Norm']] = scaler.fit_transform(
    fs_df[['Chi-Square', 'Pearson (Abs)', 'Mutual Info']]
)
fs_df['Total_Score'] = fs_df[['Chi2_Norm', 'Pearson_Norm', 'MI_Norm']].mean(axis=1)
fs_df = fs_df.sort_values(by='Total_Score', ascending=False).reset_index(drop=True)

# Jumlah Fitur
n_select = 7
selected_features = fs_df['Feature'].head(n_select).tolist()
print(f"\n-> Fitur yang dipilih ({n_select} terbaik): {selected_features}")

X_train = X_train[selected_features]
X_test = X_test[selected_features]

# ------------------------------------------------------------
# 6. Fungsi Evaluasi
# ------------------------------------------------------------
def evaluate(name, model):
    y_pred = model.predict(X_test)
    print(f"\n===== {name} =====")
    print("Classification Report:")
    print(classification_report(y_test, y_pred, digits=4))
    cm = confusion_matrix(y_test, y_pred)
    print("Confusion Matrix (Text):")
    print("                Predicted")
    print("               0       1")
    print("          -----------------")
    print(f"True   0 | {cm[0][0]:6d}  {cm[0][1]:6d}")
    print(f"       1 | {cm[1][0]:6d}  {cm[1][1]:6d}")
    print("          -----------------\n")

# ------------------------------------------------------------
# 7. Latih Semua Model (TANPA TUNING)
# ------------------------------------------------------------
models = {}
custom_weight = {0: 1, 1: 4}

# Model menggunakan pengaturan tetap
models["Decision Tree"] = DecisionTreeClassifier(max_depth=6, class_weight=custom_weight, random_state=42)
models["Random Forest"] = RandomForestClassifier(n_estimators=200, max_depth=7, class_weight=custom_weight, random_state=42)
models["KNN"] = make_pipeline(StandardScaler(), KNeighborsClassifier(n_neighbors=5, weights='distance'))
models["SVM RBF"] = make_pipeline(StandardScaler(), SVC(kernel='rbf', probability=True, class_weight=custom_weight, random_state=42))
models["XGBoost"] = XGBClassifier(n_estimators=200, max_depth=5, learning_rate=0.05, scale_pos_weight=4, eval_metric='logloss', random_state=42)

for name, model in models.items():
    model.fit(X_train, y_train)
    evaluate(name, model)

# ------------------------------------------------------------
# 8. Eksport Model Terbaik & Label Encoders
# ------------------------------------------------------------
print("\n=== Proses Export Model ===")

# Gabungkan model dan encoders ke dalam satu dictionary agar rapi
export_package = {
    'model': models["KNN"],
    'label_encoders': label_encoders,
    'selected_features': selected_features # Penting agar tahu urutan kolom masuk saat predict
}

# Menentukan nama file yang sesuai untuk KNN
filename = 'model_knn.pkl'

# Menyimpan package ke dalam file .pkl
joblib.dump(export_package, filename)
print(f"Sukses! Model dan Label Encoders telah disimpan sebagai file: {filename}")
\end{lstlisting}

\begin{lstlisting}[caption={Output 15}]
rentals.csv
   id                   nama               nik status_pasutri  \
0   3  N*** Y**** S**** D***  3***************  Belum Menikah   
1   4     A***** N** H******  3***************  Belum Menikah   
2   5     A***** A*** S*****  3***************        Menikah   
3   6            C**** W****  1***************  Belum Menikah   
4   7             S*** U****  3***************  Belum Menikah   

  status_kewarganegaraan jenis_sewa unit_number metode_pembayaran metode_lain  \
0                    WNI     Harian     B-20-27              Cash         NaN   
1                    WNI     Harian     A-15-25              Cash         NaN   
2                    WNI     Harian     A-06-24            Others    Aplikasi   
3                    WNI     Harian     D-15-21              QRIS         NaN   
4                    WNI     Harian       A-8-8            Others    Transfer   

  tanggal_checkin waktu_checkin tanggal_checkout waktu_checkout  \
0      2026-05-20      13:08:00       2026-05-21            NaN   
1      2026-05-19      13:25:00       2026-05-22            NaN   
2      2026-05-17      13:27:00       2026-05-22            NaN   
3      2026-05-19      13:29:00       2026-05-24            NaN   
4      2025-01-01      19:58:00       2025-01-06       11:22:00   

   lama_menginap komentar  user_email  user_pengguna_id  editor_pengguna_id  \
0              1      NaN         NaN              1267                 NaN   
1              3      NaN         NaN              1267                 NaN   
2              5      NaN         NaN              1267                 NaN   
3              5      NaN         NaN              1267                 NaN   
4              5      NaN         NaN              1271              1271.0   

            created_at classification  
0  2026-05-20 13:09:49            NaN  
1  2026-05-20 13:26:04            NaN  
2  2026-05-20 13:28:03            NaN  
3  2026-05-20 13:30:18            NaN  
4  2026-05-25 10:06:33         normal  
logs.csv
     id  rental_id  action field_changed  old_value  \
0     6          3  create           all        NaN   
1     7          4  create           all        NaN   
2     8          5  create           all        NaN   
3     9          6  create           all        NaN   
4  4906          7  create           all        NaN   

                                           new_value  email  pengguna_id  \
0  {"nama":"N*** Y**** S**** D***","nik":"3*************************************"nama":"A***** N** H******","nik":"3****************************************"nama":"A***** A*** S*****","nik":"3****************************************"nama":"C**** W****","nik":"1***************"...    NaN         1267   
4  {"nama":"S*** U****","nik":"3***************",...    NaN         1271   

             timestamp  
0  2026-05-20 13:09:49  
1  2026-05-20 13:26:04  
2  2026-05-20 13:28:03  
3  2026-05-20 13:30:18  
4  2026-05-25 10:06:33  

==================================================
=== EXPLORATORY DATA ANALYSIS (EDA) ===
==================================================

[1] Statistik Deskriptif:
                id                         nama           nik status_pasutri  \
count   544.000000                          544  5.440000e+02            544   
unique         NaN                          207           NaN              2   
top            NaN  Y***** M**** S**** S*******           NaN  Belum Menikah   
freq           NaN                            8           NaN            330   
mean    274.538603                          NaN  3.214330e+15            NaN   
std     157.250941                          NaN  1.204147e+15            NaN   
min       3.000000                          NaN  1.202052e+15            NaN   
25%     138.750000                          NaN  3.173074e+15            NaN   
50%     274.500000                          NaN  3.209121e+15            NaN   
75%     410.250000                          NaN  3.276064e+15            NaN   
max     553.000000                          NaN  9.104017e+15            NaN   

       status_kewarganegaraan jenis_sewa unit_number metode_pembayaran  \
count                     544        544         544               544   
unique                      1          2          12                 5   
top                       WNI     Harian       B-1-8            Others   
freq                      544        543          90               209   
mean                      NaN        NaN         NaN               NaN   
std                       NaN        NaN         NaN               NaN   
min                       NaN        NaN         NaN               NaN   
25%                       NaN        NaN         NaN               NaN   
50%                       NaN        NaN         NaN               NaN   
75%                       NaN        NaN         NaN               NaN   
max                       NaN        NaN         NaN               NaN   

       metode_lain tanggal_checkin  ... tanggal_checkout waktu_checkout  \
count          209             544  ...              544            538   
unique           2             311  ...              300            232   
top       Transfer      2026-02-06  ...       2026-02-08       12:00:00   
freq           208               5  ...                5             82   
mean           NaN             NaN  ...              NaN            NaN   
std            NaN             NaN  ...              NaN            NaN   
min            NaN             NaN  ...              NaN            NaN   
25%            NaN             NaN  ...              NaN            NaN   
50%            NaN             NaN  ...              NaN            NaN   
75%            NaN             NaN  ...              NaN            NaN   
max            NaN             NaN  ...              NaN            NaN   

       lama_menginap  komentar user_email  user_pengguna_id  \
count     544.000000       544        0.0        544.000000   
unique           NaN        11        NaN               NaN   
top              NaN       nan        NaN               NaN   
freq             NaN       459        NaN               NaN   
mean        2.641544       NaN        NaN       1271.904412   
std         1.792711       NaN        NaN          1.177622   
min         1.000000       NaN        NaN       1262.000000   
25%         2.000000       NaN        NaN       1271.000000   
50%         2.000000       NaN        NaN       1272.000000   
75%         3.000000       NaN        NaN       1273.000000   
max        31.000000       NaN        NaN       1273.000000   

        editor_pengguna_id           created_at classification       label  
count           538.000000                  544            537  544.000000  
unique                 NaN                  544              2         NaN  
top                    NaN  2026-05-30 15:00:26         normal         NaN  
freq                   NaN                    1            452         NaN  
mean           1271.977695                  NaN            NaN    0.156250  
std               0.923249                  NaN            NaN    0.363426  
min            1262.000000                  NaN            NaN    0.000000  
25%            1271.000000                  NaN            NaN    0.000000  
50%            1272.000000                  NaN            NaN    0.000000  
75%            1273.000000                  NaN            NaN    0.000000  
max            1273.000000                  NaN            NaN    1.000000  

[11 rows x 21 columns]

[2] Distribusi Kelas Target (0: Aman, 1: Melanggar):
label
0    459
1     85
Name: count, dtype: int64

/tmp/ipykernel_1991/3071461186.py:70: FutureWarning: 

Passing `palette` without assigning `hue` is deprecated and will be removed in v0.14.0. Assign the `x` variable to `hue` and set `legend=False` for the same effect.

  sns.countplot(data=df_eda, x='label', palette='Set2', ax=ax)

<Figure size 600x400 with 1 Axes>

[3] Rata-rata Pelanggaran Berdasarkan Seluruh Kategori:

--- Tingkat Pelanggaran Berdasarkan status_pasutri ---
status_pasutri
Belum Menikah    0.164
Menikah          0.145
Name: label, dtype: float64

--- Tingkat Pelanggaran Berdasarkan status_kewarganegaraan ---
status_kewarganegaraan
WNI    0.156
Name: label, dtype: float64

--- Tingkat Pelanggaran Berdasarkan jenis_sewa ---
jenis_sewa
Harian     0.157
Bulanan    0.000
Name: label, dtype: float64

--- Tingkat Pelanggaran Berdasarkan metode_pembayaran ---
metode_pembayaran
Kartu Kredit    0.280
Others          0.158
Cash            0.150
QRIS            0.148
Kartu Debit     0.107
Name: label, dtype: float64

/tmp/ipykernel_1991/3071461186.py:97: UserWarning: Could not infer format, so each element will be parsed individually, falling back to `dateutil`. To ensure parsing is consistent and as-expected, please specify a format.
  df_eda['checkin_hour_num'] = pd.to_datetime(df_eda['waktu_checkin'], errors='coerce').dt.hour.fillna(0)
/tmp/ipykernel_1991/3071461186.py:98: UserWarning: Could not infer format, so each element will be parsed individually, falling back to `dateutil`. To ensure parsing is consistent and as-expected, please specify a format.
  df_eda['checkout_hour_num'] = pd.to_datetime(df_eda['waktu_checkout'], errors='coerce').dt.hour.fillna(0)


[4] Korelasi Antar Fitur:

<Figure size 1000x800 with 2 Axes>

[5] Visualisasi Pola Penyewaan (Dikaitkan dengan Klasifikasi):

<Figure size 1200x500 with 1 Axes>
<Figure size 1200x500 with 1 Axes>
<Figure size 800x500 with 1 Axes>
<Figure size 800x500 with 1 Axes>
<Figure size 800x500 with 1 Axes>
<Figure size 800x500 with 1 Axes>

[5b] Visualisasi Metode Pembayaran 'Others' (metode_lain):

/tmp/ipykernel_1991/3071461186.py:205: FutureWarning: 

Passing `palette` without assigning `hue` is deprecated and will be removed in v0.14.0. Assign the `y` variable to `hue` and set `legend=False` for the same effect.

  sns.countplot(

<Figure size 1000x600 with 1 Axes>
/tmp/ipykernel_1991/3071461186.py:239: UserWarning: Could not infer format, so each element will be parsed individually, falling back to `dateutil`. To ensure parsing is consistent and as-expected, please specify a format.
  rentals['waktu_checkout_dt'] = pd.to_datetime(rentals['waktu_checkout'], errors='coerce')
/tmp/ipykernel_1991/3071461186.py:242: UserWarning: Could not infer format, so each element will be parsed individually, falling back to `dateutil`. To ensure parsing is consistent and as-expected, please specify a format.
  rentals['waktu_checkin_dt'] = pd.to_datetime(rentals['waktu_checkin'], errors='coerce')


=== Proses Feature Selection ===

-> Fitur yang dipilih (7 terbaik): ['lama_menginap', 'checkout_hour', 'checkin_hour', 'status_kewarganegaraan', 'metode_pembayaran', 'status_pasutri', 'jenis_sewa']

===== Decision Tree =====
Classification Report:
              precision    recall  f1-score   support

           0     1.0000    0.9565    0.9778        92
           1     0.8095    1.0000    0.8947        17

    accuracy                         0.9633       109
   macro avg     0.9048    0.9783    0.9363       109
weighted avg     0.9703    0.9633    0.9648       109

Confusion Matrix (Text):
                Predicted
               0       1
          -----------------
True   0 |     88       4
       1 |      0      17
          -----------------


===== Random Forest =====
Classification Report:
              precision    recall  f1-score   support

           0     1.0000    0.9565    0.9778        92
           1     0.8095    1.0000    0.8947        17

    accuracy                         0.9633       109
   macro avg     0.9048    0.9783    0.9363       109
weighted avg     0.9703    0.9633    0.9648       109

Confusion Matrix (Text):
                Predicted
               0       1
          -----------------
True   0 |     88       4
       1 |      0      17
          -----------------


===== KNN =====
Classification Report:
              precision    recall  f1-score   support

           0     0.9184    0.9783    0.9474        92
           1     0.8182    0.5294    0.6429        17

    accuracy                         0.9083       109
   macro avg     0.8683    0.7538    0.7951       109
weighted avg     0.9027    0.9083    0.8999       109

Confusion Matrix (Text):
                Predicted
               0       1
          -----------------
True   0 |     90       2
       1 |      8       9
          -----------------


===== SVM RBF =====
Classification Report:
              precision    recall  f1-score   support

           0     1.0000    0.9565    0.9778        92
           1     0.8095    1.0000    0.8947        17

    accuracy                         0.9633       109
   macro avg     0.9048    0.9783    0.9363       109
weighted avg     0.9703    0.9633    0.9648       109

Confusion Matrix (Text):
                Predicted
               0       1
          -----------------
True   0 |     88       4
       1 |      0      17
          -----------------


===== XGBoost =====
Classification Report:
              precision    recall  f1-score   support

           0     1.0000    0.9565    0.9778        92
           1     0.8095    1.0000    0.8947        17

    accuracy                         0.9633       109
   macro avg     0.9048    0.9783    0.9363       109
weighted avg     0.9703    0.9633    0.9648       109

Confusion Matrix (Text):
                Predicted
               0       1
          -----------------
True   0 |     88       4
       1 |      0      17
          -----------------


=== Proses Export Model ===
Sukses! Model dan Label Encoders telah disimpan sebagai file: model_knn.pkl
\end{lstlisting}

\begin{lstlisting}[language=Python, caption={Kode 16}]
# ============================================================
# DATA RIIL (TANPA TUNING) --> jumlah fitur = 6
# ============================================================

import pandas as pd
import numpy as np
import joblib
import matplotlib.pyplot as plt
import seaborn as sns

from sklearn.model_selection import train_test_split
from sklearn.preprocessing import LabelEncoder, StandardScaler, MinMaxScaler
from sklearn.metrics import classification_report, confusion_matrix
from sklearn.pipeline import make_pipeline
from sklearn.feature_selection import SelectKBest, chi2, mutual_info_classif

from sklearn.tree import DecisionTreeClassifier
from sklearn.ensemble import RandomForestClassifier
from sklearn.neighbors import KNeighborsClassifier
from sklearn.svm import SVC
from xgboost import XGBClassifier

import matplotlib.pyplot as plt
import seaborn as sns

# ------------------------------------------------------------
# 1. LOAD DATA
# ------------------------------------------------------------
rentals    = pd.read_csv('rentals.csv')
logs       = pd.read_csv('rental_logs.csv')
# units      = pd.read_csv('units.csv')
# users      = pd.read_csv('users.csv')
# violations = pd.read_csv('violations.csv')

rentals.columns = rentals.columns.str.lower()

if 'komentar' not in rentals.columns:
    raise ValueError("Kolom 'komentar' tidak ditemukan di rentals.csv")

print("rentals.csv")
print(rentals.head())
print("logs.csv")
print(logs.head())


# ------------------------------------------------------------
# 2. LABELING
# ------------------------------------------------------------
rentals['komentar'] = rentals['komentar'].astype(str).str.lower().str.strip()
empty_values = ["", "nan", "na", "null", "-", "0", "tidak ada", "none", " "]
rentals['label'] = (~rentals['komentar'].isin(empty_values)).astype(int)


# ============================================================
# EXPLORATORY DATA ANALYSIS (EDA)
# ============================================================
print("\n" + "="*50)
print("=== EXPLORATORY DATA ANALYSIS (EDA) ===")
print("="*50)

df_eda = rentals.copy()

print("\n[1] Statistik Deskriptif:")
print(df_eda.describe(include='all'))

print("\n[2] Distribusi Kelas Target (0: Aman, 1: Melanggar):")
print(df_eda['label'].value_counts())

fig, ax = plt.subplots(figsize=(6, 4))
sns.countplot(data=df_eda, x='label', palette='Set2', ax=ax)
plt.title('Distribusi Kelas Klasifikasi (0: Aman, 1: Melanggar)')
plt.xlabel('Klasifikasi')
plt.ylabel('Jumlah Penyewaan')

y_max = df_eda['label'].value_counts().max()
plt.ylim(0, y_max * 1.15)

for p in ax.patches:
    if p.get_height() > 0:
        ax.annotate(f'{int(p.get_height())}', (p.get_x() + p.get_width() / 2., p.get_height()),
                    ha='center', va='baseline', fontsize=10, color='black', xytext=(0, 5),
                    textcoords='offset points')
plt.show()

kolom_dikecualikan = ['id', 'nama', 'nik', 'unit_number', 'tanggal_checkin', 'waktu_checkin',
                      'tanggal_checkout', 'waktu_checkout', 'komentar', 'user_email', 'diedit_oleh', 'created_at', 'metode_lain', 'classification']

semua_kategori = [col for col in df_eda.select_dtypes(include=['object', 'category']).columns if col not in kolom_dikecualikan]

print("\n[3] Rata-rata Pelanggaran Berdasarkan Seluruh Kategori:")
for col in semua_kategori:
    if col in df_eda.columns:
        print(f"\n--- Tingkat Pelanggaran Berdasarkan {col} ---")
        print(df_eda.groupby(col)['label'].mean().sort_values(ascending=False).round(3))

df_eda['lama_menginap_num'] = pd.to_numeric(df_eda['lama_menginap'], errors='coerce').fillna(0)
df_eda['checkin_hour_num'] = pd.to_datetime(df_eda['waktu_checkin'], errors='coerce').dt.hour.fillna(0)
df_eda['checkout_hour_num'] = pd.to_datetime(df_eda['waktu_checkout'], errors='coerce').dt.hour.fillna(0)

le_eda = LabelEncoder()
for c in semua_kategori:
    if c in df_eda.columns:
        df_eda[c + '_encoded'] = le_eda.fit_transform(df_eda[c].astype(str))

print("\n[4] Korelasi Antar Fitur:")
kolom_korelasi = ['lama_menginap_num', 'checkin_hour_num', 'checkout_hour_num', 'label']
kolom_korelasi += [c + '_encoded' for c in semua_kategori if c + '_encoded' in df_eda.columns]

corr_matrix = df_eda[kolom_korelasi].corr()

plt.figure(figsize=(10, 8))
sns.heatmap(corr_matrix, annot=True, cmap='coolwarm', fmt='.2f', linewidths=0.5)
plt.title('Heatmap Korelasi Antar Fitur dan Target (Klasifikasi)')
plt.tight_layout()
plt.show()

print("\n[5] Visualisasi Pola Penyewaan (Dikaitkan dengan Klasifikasi):")

custom_line_palette = {0: 'blue', 1: 'red'}
custom_labels = ['0: Aman', '1: Melanggar']

fig, ax1 = plt.subplots(figsize=(12, 5))
waktu_df = df_eda.groupby(['checkin_hour_num', 'label']).size().reset_index(name='jumlah')

sns.lineplot(data=waktu_df, x='checkin_hour_num', y='jumlah', hue='label', marker='o', palette=custom_line_palette, ax=ax1)

plt.title('Pola Waktu Kedatangan Berdasarkan Klasifikasi')
plt.xlabel('Jam Check-in (0-23)')
plt.ylabel('Jumlah Penyewaan')
plt.xticks(range(0, 24))
plt.grid(True, alpha=0.3)

y_max_line = waktu_df['jumlah'].max()
plt.ylim(-y_max_line * 0.20, y_max_line * 1.15)

handles, _ = ax1.get_legend_handles_labels()
plt.legend(handles=handles, labels=custom_labels, title='Klasifikasi')

for l in waktu_df['label'].unique():
    sub = waktu_df[waktu_df['label'] == l]
    warna_teks = custom_line_palette[l]

    for x, y in zip(sub['checkin_hour_num'], sub['jumlah']):
        y_offset = 8 if l == 0 else -12
        va_align = 'bottom' if l == 0 else 'top'
        plt.annotate(f'{int(y)}', (x, y), textcoords="offset points", xytext=(0, y_offset),
                     ha='center', va=va_align, fontsize=9, color=warna_teks, fontweight='bold')
plt.show()

fig, ax2 = plt.subplots(figsize=(12, 5))
lama_df = df_eda.groupby(['lama_menginap_num', 'label']).size().reset_index(name='jumlah')

sns.lineplot(data=lama_df, x='lama_menginap_num', y='jumlah', hue='label', marker='o', palette=custom_line_palette, ax=ax2)

plt.title('Tren Lama Menginap Berdasarkan Klasifikasi')
plt.xlabel('Lama Menginap (Hari)')
plt.ylabel('Jumlah Penyewaan')
plt.grid(True, alpha=0.3)

handles2, _ = ax2.get_legend_handles_labels()
plt.legend(handles=handles2, labels=custom_labels, title='Klasifikasi')
plt.show()

for col in semua_kategori:
    if col in df_eda.columns:
        fig, ax = plt.subplots(figsize=(8, 5))
        sns.countplot(data=df_eda, x=col, hue='label', palette='magma', ax=ax)

        judul_kolom = col.replace('_', ' ').title()
        plt.title(f'Perbandingan Jumlah Pelanggaran Berdasarkan {judul_kolom}')
        plt.xlabel(judul_kolom)
        plt.ylabel('Jumlah Penyewaan')
        plt.xticks(rotation=0)

        handles_bar, _ = ax.get_legend_handles_labels()
        plt.legend(handles=handles_bar, labels=custom_labels, title='Klasifikasi')

        y_max_bar = df_eda.groupby([col, 'label']).size().max()
        plt.ylim(0, y_max_bar * 1.15)

        for p in ax.patches:
            if p.get_height() > 0:
                ax.annotate(f'{int(p.get_height())}', (p.get_x() + p.get_width() / 2., p.get_height()),
                            ha='center', va='baseline', fontsize=9, color='black', xytext=(0, 5),
                            textcoords='offset points')
        plt.tight_layout()
        plt.show()


print("\n[5b] Visualisasi Metode Pembayaran 'Others' (metode_lain):")

# Memastikan kolom 'metode_lain' berbentuk string dan membersihkan spasi/nilai kosong
df_eda['metode_lain'] = df_eda['metode_lain'].astype(str).str.strip().str.title()
empty_lain = ["", "Nan", "Na", "Null", "-", "0", "None"]

# Filter hanya baris yang memiliki nilai di 'metode_lain'
df_others = df_eda[~df_eda['metode_lain'].isin(empty_lain)]

if not df_others.empty:
    fig, ax = plt.subplots(figsize=(10, 6))

    # Menghitung frekuensi masing-masing metode pembayaran lain
    order_others = df_others['metode_lain'].value_counts().index

    sns.countplot(
        data=df_others,
        y='metode_lain',
        order=order_others,
        palette='viridis',
        ax=ax
    )

    plt.title('Detail Metode Pembayaran "Others" (metode_lain)')
    plt.xlabel('Jumlah Penyewaan')
    plt.ylabel('Jenis Pembayaran Lainnya')

    # Menambahkan anotasi angka di sebelah setiap bar
    for p in ax.patches:
        width = p.get_width()
        if width > 0:
            ax.annotate(f'{int(width)}',
                        (width, p.get_y() + p.get_height() / 2.),
                        ha='left', va='center', fontsize=10, color='black',
                        xytext=(5, 0), textcoords='offset points')

    # Perlebar batas x-axis sedikit agar angka anotasi tidak terpotong
    x_max = df_others['metode_lain'].value_counts().max()
    plt.xlim(0, x_max * 1.15)

    plt.tight_layout()
    plt.show()
else:
    print("Tidak ada data metode pembayaran 'Others' yang valid untuk ditampilkan.")


# ------------------------------------------------------------
# 3. FEATURE ENGINEERING & PREPARATION
# ------------------------------------------------------------
rentals['waktu_checkout_dt'] = pd.to_datetime(rentals['waktu_checkout'], errors='coerce')
rentals['checkout_hour'] = rentals['waktu_checkout_dt'].dt.hour.fillna(0).astype(int)

rentals['waktu_checkin_dt'] = pd.to_datetime(rentals['waktu_checkin'], errors='coerce')
rentals['checkin_hour'] = rentals['waktu_checkin_dt'].dt.hour.fillna(0).astype(int)

rentals['lama_menginap'] = pd.to_numeric(rentals['lama_menginap'], errors='coerce').fillna(0)

cat_cols = ['jenis_sewa', 'metode_pembayaran', 'status_pasutri', 'status_kewarganegaraan']
for c in cat_cols:
    rentals[c] = rentals[c].fillna('Unknown').astype(str)

numeric_cols = ['checkout_hour', 'checkin_hour', 'lama_menginap']
feature_cols = cat_cols + numeric_cols

X = rentals[feature_cols].copy()
y = rentals['label']

# ------------------------------------------------------------
# 4. TRAIN-TEST SPLIT
# ------------------------------------------------------------
X_train, X_test, y_train, y_test = train_test_split(
    X, y, stratify=y, test_size=0.2, random_state=42
)

# ------------------------------------------------------------
# 4B. ENCODING MENCEGAH DATA LEAKAGE
# ------------------------------------------------------------
label_encoders = {}

for c in cat_cols:
    le = LabelEncoder()
    X_train[c] = le.fit_transform(X_train[c])

    X_test[c] = X_test[c].map(lambda s: s if s in le.classes_ else 'Unknown')
    if 'Unknown' not in le.classes_:
        le.classes_ = np.append(le.classes_, 'Unknown')

    X_test[c] = le.transform(X_test[c])

    label_encoders[c] = le

# ------------------------------------------------------------
# 5. FEATURE SELECTION
# ------------------------------------------------------------
print("\n=== Proses Feature Selection ===")
chi2_selector = SelectKBest(score_func=chi2, k='all')
chi2_selector.fit(X_train, y_train)

train_df = X_train.copy()
train_df['label'] = y_train
pearson_scores = train_df.corr(method='pearson')['label'].drop('label').abs().values

mi_selector = SelectKBest(score_func=mutual_info_classif, k='all')
mi_selector.fit(X_train, y_train)

fs_df = pd.DataFrame({
    'Feature': X_train.columns,
    'Chi-Square': chi2_selector.scores_,
    'Pearson (Abs)': pearson_scores,
    'Mutual Info': mi_selector.scores_
})

scaler = MinMaxScaler()
fs_df[['Chi2_Norm', 'Pearson_Norm', 'MI_Norm']] = scaler.fit_transform(
    fs_df[['Chi-Square', 'Pearson (Abs)', 'Mutual Info']]
)
fs_df['Total_Score'] = fs_df[['Chi2_Norm', 'Pearson_Norm', 'MI_Norm']].mean(axis=1)
fs_df = fs_df.sort_values(by='Total_Score', ascending=False).reset_index(drop=True)

# Jumlah Fitur
n_select = 6
selected_features = fs_df['Feature'].head(n_select).tolist()
print(f"\n-> Fitur yang dipilih ({n_select} terbaik): {selected_features}")

X_train = X_train[selected_features]
X_test = X_test[selected_features]

# ------------------------------------------------------------
# 6. Fungsi Evaluasi
# ------------------------------------------------------------
def evaluate(name, model):
    y_pred = model.predict(X_test)
    print(f"\n===== {name} =====")
    print("Classification Report:")
    print(classification_report(y_test, y_pred, digits=4))
    cm = confusion_matrix(y_test, y_pred)
    print("Confusion Matrix (Text):")
    print("                Predicted")
    print("               0       1")
    print("          -----------------")
    print(f"True   0 | {cm[0][0]:6d}  {cm[0][1]:6d}")
    print(f"       1 | {cm[1][0]:6d}  {cm[1][1]:6d}")
    print("          -----------------\n")

# ------------------------------------------------------------
# 7. Latih Semua Model (TANPA TUNING)
# ------------------------------------------------------------
models = {}
custom_weight = {0: 1, 1: 4}

# Model menggunakan pengaturan tetap
models["Decision Tree"] = DecisionTreeClassifier(max_depth=6, class_weight=custom_weight, random_state=42)
models["Random Forest"] = RandomForestClassifier(n_estimators=200, max_depth=7, class_weight=custom_weight, random_state=42)
models["KNN"] = make_pipeline(StandardScaler(), KNeighborsClassifier(n_neighbors=5, weights='distance'))
models["SVM RBF"] = make_pipeline(StandardScaler(), SVC(kernel='rbf', probability=True, class_weight=custom_weight, random_state=42))
models["XGBoost"] = XGBClassifier(n_estimators=200, max_depth=5, learning_rate=0.05, scale_pos_weight=4, eval_metric='logloss', random_state=42)

for name, model in models.items():
    model.fit(X_train, y_train)
    evaluate(name, model)

# ------------------------------------------------------------
# 8. Eksport Model Terbaik & Label Encoders
# ------------------------------------------------------------
print("\n=== Proses Export Model ===")

# Gabungkan model dan encoders ke dalam satu dictionary agar rapi
export_package = {
    'model': models["KNN"],
    'label_encoders': label_encoders,
    'selected_features': selected_features # Penting agar tahu urutan kolom masuk saat predict
}

# Menentukan nama file yang sesuai untuk KNN
filename = 'model_knn.pkl'

# Menyimpan package ke dalam file .pkl
joblib.dump(export_package, filename)
print(f"Sukses! Model dan Label Encoders telah disimpan sebagai file: {filename}")
\end{lstlisting}

\begin{lstlisting}[caption={Output 16}]
rentals.csv
   id                   nama               nik status_pasutri  \
0   3  N*** Y**** S**** D***  3***************  Belum Menikah   
1   4     A***** N** H******  3***************  Belum Menikah   
2   5     A***** A*** S*****  3***************        Menikah   
3   6            C**** W****  1***************  Belum Menikah   
4   7             S*** U****  3***************  Belum Menikah   

  status_kewarganegaraan jenis_sewa unit_number metode_pembayaran metode_lain  \
0                    WNI     Harian     B-20-27              Cash         NaN   
1                    WNI     Harian     A-15-25              Cash         NaN   
2                    WNI     Harian     A-06-24            Others    Aplikasi   
3                    WNI     Harian     D-15-21              QRIS         NaN   
4                    WNI     Harian       A-8-8            Others    Transfer   

  tanggal_checkin waktu_checkin tanggal_checkout waktu_checkout  \
0      2026-05-20      13:08:00       2026-05-21            NaN   
1      2026-05-19      13:25:00       2026-05-22            NaN   
2      2026-05-17      13:27:00       2026-05-22            NaN   
3      2026-05-19      13:29:00       2026-05-24            NaN   
4      2025-01-01      19:58:00       2025-01-06       11:22:00   

   lama_menginap komentar  user_email  user_pengguna_id  editor_pengguna_id  \
0              1      NaN         NaN              1267                 NaN   
1              3      NaN         NaN              1267                 NaN   
2              5      NaN         NaN              1267                 NaN   
3              5      NaN         NaN              1267                 NaN   
4              5      NaN         NaN              1271              1271.0   

            created_at classification  
0  2026-05-20 13:09:49            NaN  
1  2026-05-20 13:26:04            NaN  
2  2026-05-20 13:28:03            NaN  
3  2026-05-20 13:30:18            NaN  
4  2026-05-25 10:06:33         normal  
logs.csv
     id  rental_id  action field_changed  old_value  \
0     6          3  create           all        NaN   
1     7          4  create           all        NaN   
2     8          5  create           all        NaN   
3     9          6  create           all        NaN   
4  4906          7  create           all        NaN   

                                           new_value  email  pengguna_id  \
0  {"nama":"N*** Y**** S**** D***","nik":"3*************************************"nama":"A***** N** H******","nik":"3****************************************"nama":"A***** A*** S*****","nik":"3****************************************"nama":"C**** W****","nik":"1***************"...    NaN         1267   
4  {"nama":"S*** U****","nik":"3***************",...    NaN         1271   

             timestamp  
0  2026-05-20 13:09:49  
1  2026-05-20 13:26:04  
2  2026-05-20 13:28:03  
3  2026-05-20 13:30:18  
4  2026-05-25 10:06:33  

==================================================
=== EXPLORATORY DATA ANALYSIS (EDA) ===
==================================================

[1] Statistik Deskriptif:
                id                         nama           nik status_pasutri  \
count   544.000000                          544  5.440000e+02            544   
unique         NaN                          207           NaN              2   
top            NaN  Y***** M**** S**** S*******           NaN  Belum Menikah   
freq           NaN                            8           NaN            330   
mean    274.538603                          NaN  3.214330e+15            NaN   
std     157.250941                          NaN  1.204147e+15            NaN   
min       3.000000                          NaN  1.202052e+15            NaN   
25%     138.750000                          NaN  3.173074e+15            NaN   
50%     274.500000                          NaN  3.209121e+15            NaN   
75%     410.250000                          NaN  3.276064e+15            NaN   
max     553.000000                          NaN  9.104017e+15            NaN   

       status_kewarganegaraan jenis_sewa unit_number metode_pembayaran  \
count                     544        544         544               544   
unique                      1          2          12                 5   
top                       WNI     Harian       B-1-8            Others   
freq                      544        543          90               209   
mean                      NaN        NaN         NaN               NaN   
std                       NaN        NaN         NaN               NaN   
min                       NaN        NaN         NaN               NaN   
25%                       NaN        NaN         NaN               NaN   
50%                       NaN        NaN         NaN               NaN   
75%                       NaN        NaN         NaN               NaN   
max                       NaN        NaN         NaN               NaN   

       metode_lain tanggal_checkin  ... tanggal_checkout waktu_checkout  \
count          209             544  ...              544            538   
unique           2             311  ...              300            232   
top       Transfer      2026-02-06  ...       2026-02-08       12:00:00   
freq           208               5  ...                5             82   
mean           NaN             NaN  ...              NaN            NaN   
std            NaN             NaN  ...              NaN            NaN   
min            NaN             NaN  ...              NaN            NaN   
25%            NaN             NaN  ...              NaN            NaN   
50%            NaN             NaN  ...              NaN            NaN   
75%            NaN             NaN  ...              NaN            NaN   
max            NaN             NaN  ...              NaN            NaN   

       lama_menginap  komentar user_email  user_pengguna_id  \
count     544.000000       544        0.0        544.000000   
unique           NaN        11        NaN               NaN   
top              NaN       nan        NaN               NaN   
freq             NaN       459        NaN               NaN   
mean        2.641544       NaN        NaN       1271.904412   
std         1.792711       NaN        NaN          1.177622   
min         1.000000       NaN        NaN       1262.000000   
25%         2.000000       NaN        NaN       1271.000000   
50%         2.000000       NaN        NaN       1272.000000   
75%         3.000000       NaN        NaN       1273.000000   
max        31.000000       NaN        NaN       1273.000000   

        editor_pengguna_id           created_at classification       label  
count           538.000000                  544            537  544.000000  
unique                 NaN                  544              2         NaN  
top                    NaN  2026-05-30 15:00:26         normal         NaN  
freq                   NaN                    1            452         NaN  
mean           1271.977695                  NaN            NaN    0.156250  
std               0.923249                  NaN            NaN    0.363426  
min            1262.000000                  NaN            NaN    0.000000  
25%            1271.000000                  NaN            NaN    0.000000  
50%            1272.000000                  NaN            NaN    0.000000  
75%            1273.000000                  NaN            NaN    0.000000  
max            1273.000000                  NaN            NaN    1.000000  

[11 rows x 21 columns]

[2] Distribusi Kelas Target (0: Aman, 1: Melanggar):
label
0    459
1     85
Name: count, dtype: int64

/tmp/ipykernel_1991/2036310360.py:70: FutureWarning: 

Passing `palette` without assigning `hue` is deprecated and will be removed in v0.14.0. Assign the `x` variable to `hue` and set `legend=False` for the same effect.

  sns.countplot(data=df_eda, x='label', palette='Set2', ax=ax)

<Figure size 600x400 with 1 Axes>

[3] Rata-rata Pelanggaran Berdasarkan Seluruh Kategori:

--- Tingkat Pelanggaran Berdasarkan status_pasutri ---
status_pasutri
Belum Menikah    0.164
Menikah          0.145
Name: label, dtype: float64

--- Tingkat Pelanggaran Berdasarkan status_kewarganegaraan ---
status_kewarganegaraan
WNI    0.156
Name: label, dtype: float64

--- Tingkat Pelanggaran Berdasarkan jenis_sewa ---
jenis_sewa
Harian     0.157
Bulanan    0.000
Name: label, dtype: float64

--- Tingkat Pelanggaran Berdasarkan metode_pembayaran ---
metode_pembayaran
Kartu Kredit    0.280
Others          0.158
Cash            0.150
QRIS            0.148
Kartu Debit     0.107
Name: label, dtype: float64

[4] Korelasi Antar Fitur:

/tmp/ipykernel_1991/2036310360.py:97: UserWarning: Could not infer format, so each element will be parsed individually, falling back to `dateutil`. To ensure parsing is consistent and as-expected, please specify a format.
  df_eda['checkin_hour_num'] = pd.to_datetime(df_eda['waktu_checkin'], errors='coerce').dt.hour.fillna(0)
/tmp/ipykernel_1991/2036310360.py:98: UserWarning: Could not infer format, so each element will be parsed individually, falling back to `dateutil`. To ensure parsing is consistent and as-expected, please specify a format.
  df_eda['checkout_hour_num'] = pd.to_datetime(df_eda['waktu_checkout'], errors='coerce').dt.hour.fillna(0)

<Figure size 1000x800 with 2 Axes>

[5] Visualisasi Pola Penyewaan (Dikaitkan dengan Klasifikasi):

<Figure size 1200x500 with 1 Axes>
<Figure size 1200x500 with 1 Axes>
<Figure size 800x500 with 1 Axes>
<Figure size 800x500 with 1 Axes>
<Figure size 800x500 with 1 Axes>
<Figure size 800x500 with 1 Axes>

[5b] Visualisasi Metode Pembayaran 'Others' (metode_lain):

/tmp/ipykernel_1991/2036310360.py:205: FutureWarning: 

Passing `palette` without assigning `hue` is deprecated and will be removed in v0.14.0. Assign the `y` variable to `hue` and set `legend=False` for the same effect.

  sns.countplot(

<Figure size 1000x600 with 1 Axes>
/tmp/ipykernel_1991/2036310360.py:239: UserWarning: Could not infer format, so each element will be parsed individually, falling back to `dateutil`. To ensure parsing is consistent and as-expected, please specify a format.
  rentals['waktu_checkout_dt'] = pd.to_datetime(rentals['waktu_checkout'], errors='coerce')
/tmp/ipykernel_1991/2036310360.py:242: UserWarning: Could not infer format, so each element will be parsed individually, falling back to `dateutil`. To ensure parsing is consistent and as-expected, please specify a format.
  rentals['waktu_checkin_dt'] = pd.to_datetime(rentals['waktu_checkin'], errors='coerce')


=== Proses Feature Selection ===

-> Fitur yang dipilih (6 terbaik): ['lama_menginap', 'checkout_hour', 'checkin_hour', 'metode_pembayaran', 'jenis_sewa', 'status_pasutri']

===== Decision Tree =====
Classification Report:
              precision    recall  f1-score   support

           0     1.0000    0.9565    0.9778        92
           1     0.8095    1.0000    0.8947        17

    accuracy                         0.9633       109
   macro avg     0.9048    0.9783    0.9363       109
weighted avg     0.9703    0.9633    0.9648       109

Confusion Matrix (Text):
                Predicted
               0       1
          -----------------
True   0 |     88       4
       1 |      0      17
          -----------------


===== Random Forest =====
Classification Report:
              precision    recall  f1-score   support

           0     1.0000    0.9565    0.9778        92
           1     0.8095    1.0000    0.8947        17

    accuracy                         0.9633       109
   macro avg     0.9048    0.9783    0.9363       109
weighted avg     0.9703    0.9633    0.9648       109

Confusion Matrix (Text):
                Predicted
               0       1
          -----------------
True   0 |     88       4
       1 |      0      17
          -----------------


===== KNN =====
Classification Report:
              precision    recall  f1-score   support

           0     0.9184    0.9783    0.9474        92
           1     0.8182    0.5294    0.6429        17

    accuracy                         0.9083       109
   macro avg     0.8683    0.7538    0.7951       109
weighted avg     0.9027    0.9083    0.8999       109

Confusion Matrix (Text):
                Predicted
               0       1
          -----------------
True   0 |     90       2
       1 |      8       9
          -----------------


===== SVM RBF =====
Classification Report:
              precision    recall  f1-score   support

           0     1.0000    0.9565    0.9778        92
           1     0.8095    1.0000    0.8947        17

    accuracy                         0.9633       109
   macro avg     0.9048    0.9783    0.9363       109
weighted avg     0.9703    0.9633    0.9648       109

Confusion Matrix (Text):
                Predicted
               0       1
          -----------------
True   0 |     88       4
       1 |      0      17
          -----------------


===== XGBoost =====
Classification Report:
              precision    recall  f1-score   support

           0     1.0000    0.9565    0.9778        92
           1     0.8095    1.0000    0.8947        17

    accuracy                         0.9633       109
   macro avg     0.9048    0.9783    0.9363       109
weighted avg     0.9703    0.9633    0.9648       109

Confusion Matrix (Text):
                Predicted
               0       1
          -----------------
True   0 |     88       4
       1 |      0      17
          -----------------


=== Proses Export Model ===
Sukses! Model dan Label Encoders telah disimpan sebagai file: model_knn.pkl
\end{lstlisting}

\begin{lstlisting}[language=Python, caption={Kode 17}]
# ============================================================
# DATA RIIL (TANPA TUNING) --> jumlah fitur = 5
# ============================================================

import pandas as pd
import numpy as np
import joblib
import matplotlib.pyplot as plt
import seaborn as sns

from sklearn.model_selection import train_test_split
from sklearn.preprocessing import LabelEncoder, StandardScaler, MinMaxScaler
from sklearn.metrics import classification_report, confusion_matrix
from sklearn.pipeline import make_pipeline
from sklearn.feature_selection import SelectKBest, chi2, mutual_info_classif

from sklearn.tree import DecisionTreeClassifier
from sklearn.ensemble import RandomForestClassifier
from sklearn.neighbors import KNeighborsClassifier
from sklearn.svm import SVC
from xgboost import XGBClassifier

import matplotlib.pyplot as plt
import seaborn as sns

# ------------------------------------------------------------
# 1. LOAD DATA
# ------------------------------------------------------------
rentals    = pd.read_csv('rentals.csv')
logs       = pd.read_csv('rental_logs.csv')
# units      = pd.read_csv('units.csv')
# users      = pd.read_csv('users.csv')
# violations = pd.read_csv('violations.csv')

rentals.columns = rentals.columns.str.lower()

if 'komentar' not in rentals.columns:
    raise ValueError("Kolom 'komentar' tidak ditemukan di rentals.csv")

print("rentals.csv")
print(rentals.head())
print("logs.csv")
print(logs.head())


# ------------------------------------------------------------
# 2. LABELING
# ------------------------------------------------------------
rentals['komentar'] = rentals['komentar'].astype(str).str.lower().str.strip()
empty_values = ["", "nan", "na", "null", "-", "0", "tidak ada", "none", " "]
rentals['label'] = (~rentals['komentar'].isin(empty_values)).astype(int)


# ============================================================
# EXPLORATORY DATA ANALYSIS (EDA)
# ============================================================
print("\n" + "="*50)
print("=== EXPLORATORY DATA ANALYSIS (EDA) ===")
print("="*50)

df_eda = rentals.copy()

print("\n[1] Statistik Deskriptif:")
print(df_eda.describe(include='all'))

print("\n[2] Distribusi Kelas Target (0: Aman, 1: Melanggar):")
print(df_eda['label'].value_counts())

fig, ax = plt.subplots(figsize=(6, 4))
sns.countplot(data=df_eda, x='label', palette='Set2', ax=ax)
plt.title('Distribusi Kelas Klasifikasi (0: Aman, 1: Melanggar)')
plt.xlabel('Klasifikasi')
plt.ylabel('Jumlah Penyewaan')

y_max = df_eda['label'].value_counts().max()
plt.ylim(0, y_max * 1.15)

for p in ax.patches:
    if p.get_height() > 0:
        ax.annotate(f'{int(p.get_height())}', (p.get_x() + p.get_width() / 2., p.get_height()),
                    ha='center', va='baseline', fontsize=10, color='black', xytext=(0, 5),
                    textcoords='offset points')
plt.show()

kolom_dikecualikan = ['id', 'nama', 'nik', 'unit_number', 'tanggal_checkin', 'waktu_checkin',
                      'tanggal_checkout', 'waktu_checkout', 'komentar', 'user_email', 'diedit_oleh', 'created_at', 'metode_lain', 'classification']

semua_kategori = [col for col in df_eda.select_dtypes(include=['object', 'category']).columns if col not in kolom_dikecualikan]

print("\n[3] Rata-rata Pelanggaran Berdasarkan Seluruh Kategori:")
for col in semua_kategori:
    if col in df_eda.columns:
        print(f"\n--- Tingkat Pelanggaran Berdasarkan {col} ---")
        print(df_eda.groupby(col)['label'].mean().sort_values(ascending=False).round(3))

df_eda['lama_menginap_num'] = pd.to_numeric(df_eda['lama_menginap'], errors='coerce').fillna(0)
df_eda['checkin_hour_num'] = pd.to_datetime(df_eda['waktu_checkin'], errors='coerce').dt.hour.fillna(0)
df_eda['checkout_hour_num'] = pd.to_datetime(df_eda['waktu_checkout'], errors='coerce').dt.hour.fillna(0)

le_eda = LabelEncoder()
for c in semua_kategori:
    if c in df_eda.columns:
        df_eda[c + '_encoded'] = le_eda.fit_transform(df_eda[c].astype(str))

print("\n[4] Korelasi Antar Fitur:")
kolom_korelasi = ['lama_menginap_num', 'checkin_hour_num', 'checkout_hour_num', 'label']
kolom_korelasi += [c + '_encoded' for c in semua_kategori if c + '_encoded' in df_eda.columns]

corr_matrix = df_eda[kolom_korelasi].corr()

plt.figure(figsize=(10, 8))
sns.heatmap(corr_matrix, annot=True, cmap='coolwarm', fmt='.2f', linewidths=0.5)
plt.title('Heatmap Korelasi Antar Fitur dan Target (Klasifikasi)')
plt.tight_layout()
plt.show()

print("\n[5] Visualisasi Pola Penyewaan (Dikaitkan dengan Klasifikasi):")

custom_line_palette = {0: 'blue', 1: 'red'}
custom_labels = ['0: Aman', '1: Melanggar']

fig, ax1 = plt.subplots(figsize=(12, 5))
waktu_df = df_eda.groupby(['checkin_hour_num', 'label']).size().reset_index(name='jumlah')

sns.lineplot(data=waktu_df, x='checkin_hour_num', y='jumlah', hue='label', marker='o', palette=custom_line_palette, ax=ax1)

plt.title('Pola Waktu Kedatangan Berdasarkan Klasifikasi')
plt.xlabel('Jam Check-in (0-23)')
plt.ylabel('Jumlah Penyewaan')
plt.xticks(range(0, 24))
plt.grid(True, alpha=0.3)

y_max_line = waktu_df['jumlah'].max()
plt.ylim(-y_max_line * 0.20, y_max_line * 1.15)

handles, _ = ax1.get_legend_handles_labels()
plt.legend(handles=handles, labels=custom_labels, title='Klasifikasi')

for l in waktu_df['label'].unique():
    sub = waktu_df[waktu_df['label'] == l]
    warna_teks = custom_line_palette[l]

    for x, y in zip(sub['checkin_hour_num'], sub['jumlah']):
        y_offset = 8 if l == 0 else -12
        va_align = 'bottom' if l == 0 else 'top'
        plt.annotate(f'{int(y)}', (x, y), textcoords="offset points", xytext=(0, y_offset),
                     ha='center', va=va_align, fontsize=9, color=warna_teks, fontweight='bold')
plt.show()

fig, ax2 = plt.subplots(figsize=(12, 5))
lama_df = df_eda.groupby(['lama_menginap_num', 'label']).size().reset_index(name='jumlah')

sns.lineplot(data=lama_df, x='lama_menginap_num', y='jumlah', hue='label', marker='o', palette=custom_line_palette, ax=ax2)

plt.title('Tren Lama Menginap Berdasarkan Klasifikasi')
plt.xlabel('Lama Menginap (Hari)')
plt.ylabel('Jumlah Penyewaan')
plt.grid(True, alpha=0.3)

handles2, _ = ax2.get_legend_handles_labels()
plt.legend(handles=handles2, labels=custom_labels, title='Klasifikasi')
plt.show()

for col in semua_kategori:
    if col in df_eda.columns:
        fig, ax = plt.subplots(figsize=(8, 5))
        sns.countplot(data=df_eda, x=col, hue='label', palette='magma', ax=ax)

        judul_kolom = col.replace('_', ' ').title()
        plt.title(f'Perbandingan Jumlah Pelanggaran Berdasarkan {judul_kolom}')
        plt.xlabel(judul_kolom)
        plt.ylabel('Jumlah Penyewaan')
        plt.xticks(rotation=0)

        handles_bar, _ = ax.get_legend_handles_labels()
        plt.legend(handles=handles_bar, labels=custom_labels, title='Klasifikasi')

        y_max_bar = df_eda.groupby([col, 'label']).size().max()
        plt.ylim(0, y_max_bar * 1.15)

        for p in ax.patches:
            if p.get_height() > 0:
                ax.annotate(f'{int(p.get_height())}', (p.get_x() + p.get_width() / 2., p.get_height()),
                            ha='center', va='baseline', fontsize=9, color='black', xytext=(0, 5),
                            textcoords='offset points')
        plt.tight_layout()
        plt.show()


print("\n[5b] Visualisasi Metode Pembayaran 'Others' (metode_lain):")

# Memastikan kolom 'metode_lain' berbentuk string dan membersihkan spasi/nilai kosong
df_eda['metode_lain'] = df_eda['metode_lain'].astype(str).str.strip().str.title()
empty_lain = ["", "Nan", "Na", "Null", "-", "0", "None"]

# Filter hanya baris yang memiliki nilai di 'metode_lain'
df_others = df_eda[~df_eda['metode_lain'].isin(empty_lain)]

if not df_others.empty:
    fig, ax = plt.subplots(figsize=(10, 6))

    # Menghitung frekuensi masing-masing metode pembayaran lain
    order_others = df_others['metode_lain'].value_counts().index

    sns.countplot(
        data=df_others,
        y='metode_lain',
        order=order_others,
        palette='viridis',
        ax=ax
    )

    plt.title('Detail Metode Pembayaran "Others" (metode_lain)')
    plt.xlabel('Jumlah Penyewaan')
    plt.ylabel('Jenis Pembayaran Lainnya')

    # Menambahkan anotasi angka di sebelah setiap bar
    for p in ax.patches:
        width = p.get_width()
        if width > 0:
            ax.annotate(f'{int(width)}',
                        (width, p.get_y() + p.get_height() / 2.),
                        ha='left', va='center', fontsize=10, color='black',
                        xytext=(5, 0), textcoords='offset points')

    # Perlebar batas x-axis sedikit agar angka anotasi tidak terpotong
    x_max = df_others['metode_lain'].value_counts().max()
    plt.xlim(0, x_max * 1.15)

    plt.tight_layout()
    plt.show()
else:
    print("Tidak ada data metode pembayaran 'Others' yang valid untuk ditampilkan.")


# ------------------------------------------------------------
# 3. FEATURE ENGINEERING & PREPARATION
# ------------------------------------------------------------
rentals['waktu_checkout_dt'] = pd.to_datetime(rentals['waktu_checkout'], errors='coerce')
rentals['checkout_hour'] = rentals['waktu_checkout_dt'].dt.hour.fillna(0).astype(int)

rentals['waktu_checkin_dt'] = pd.to_datetime(rentals['waktu_checkin'], errors='coerce')
rentals['checkin_hour'] = rentals['waktu_checkin_dt'].dt.hour.fillna(0).astype(int)

rentals['lama_menginap'] = pd.to_numeric(rentals['lama_menginap'], errors='coerce').fillna(0)

cat_cols = ['jenis_sewa', 'metode_pembayaran', 'status_pasutri', 'status_kewarganegaraan']
for c in cat_cols:
    rentals[c] = rentals[c].fillna('Unknown').astype(str)

numeric_cols = ['checkout_hour', 'checkin_hour', 'lama_menginap']
feature_cols = cat_cols + numeric_cols

X = rentals[feature_cols].copy()
y = rentals['label']

# ------------------------------------------------------------
# 4. TRAIN-TEST SPLIT
# ------------------------------------------------------------
X_train, X_test, y_train, y_test = train_test_split(
    X, y, stratify=y, test_size=0.2, random_state=42
)

# ------------------------------------------------------------
# 4B. ENCODING MENCEGAH DATA LEAKAGE
# ------------------------------------------------------------
label_encoders = {}

for c in cat_cols:
    le = LabelEncoder()
    X_train[c] = le.fit_transform(X_train[c])

    X_test[c] = X_test[c].map(lambda s: s if s in le.classes_ else 'Unknown')
    if 'Unknown' not in le.classes_:
        le.classes_ = np.append(le.classes_, 'Unknown')

    X_test[c] = le.transform(X_test[c])

    label_encoders[c] = le

# ------------------------------------------------------------
# 5. FEATURE SELECTION
# ------------------------------------------------------------
print("\n=== Proses Feature Selection ===")
chi2_selector = SelectKBest(score_func=chi2, k='all')
chi2_selector.fit(X_train, y_train)

train_df = X_train.copy()
train_df['label'] = y_train
pearson_scores = train_df.corr(method='pearson')['label'].drop('label').abs().values

mi_selector = SelectKBest(score_func=mutual_info_classif, k='all')
mi_selector.fit(X_train, y_train)

fs_df = pd.DataFrame({
    'Feature': X_train.columns,
    'Chi-Square': chi2_selector.scores_,
    'Pearson (Abs)': pearson_scores,
    'Mutual Info': mi_selector.scores_
})

scaler = MinMaxScaler()
fs_df[['Chi2_Norm', 'Pearson_Norm', 'MI_Norm']] = scaler.fit_transform(
    fs_df[['Chi-Square', 'Pearson (Abs)', 'Mutual Info']]
)
fs_df['Total_Score'] = fs_df[['Chi2_Norm', 'Pearson_Norm', 'MI_Norm']].mean(axis=1)
fs_df = fs_df.sort_values(by='Total_Score', ascending=False).reset_index(drop=True)

# Jumlah Fitur
n_select = 5
selected_features = fs_df['Feature'].head(n_select).tolist()
print(f"\n-> Fitur yang dipilih ({n_select} terbaik): {selected_features}")

X_train = X_train[selected_features]
X_test = X_test[selected_features]

# ------------------------------------------------------------
# 6. Fungsi Evaluasi
# ------------------------------------------------------------
def evaluate(name, model):
    y_pred = model.predict(X_test)
    print(f"\n===== {name} =====")
    print("Classification Report:")
    print(classification_report(y_test, y_pred, digits=4))
    cm = confusion_matrix(y_test, y_pred)
    print("Confusion Matrix (Text):")
    print("                Predicted")
    print("               0       1")
    print("          -----------------")
    print(f"True   0 | {cm[0][0]:6d}  {cm[0][1]:6d}")
    print(f"       1 | {cm[1][0]:6d}  {cm[1][1]:6d}")
    print("          -----------------\n")

# ------------------------------------------------------------
# 7. Latih Semua Model (TANPA TUNING)
# ------------------------------------------------------------
models = {}
custom_weight = {0: 1, 1: 4}

# Model menggunakan pengaturan tetap
models["Decision Tree"] = DecisionTreeClassifier(max_depth=6, class_weight=custom_weight, random_state=42)
models["Random Forest"] = RandomForestClassifier(n_estimators=200, max_depth=7, class_weight=custom_weight, random_state=42)
models["KNN"] = make_pipeline(StandardScaler(), KNeighborsClassifier(n_neighbors=5, weights='distance'))
models["SVM RBF"] = make_pipeline(StandardScaler(), SVC(kernel='rbf', probability=True, class_weight=custom_weight, random_state=42))
models["XGBoost"] = XGBClassifier(n_estimators=200, max_depth=5, learning_rate=0.05, scale_pos_weight=4, eval_metric='logloss', random_state=42)

for name, model in models.items():
    model.fit(X_train, y_train)
    evaluate(name, model)

# ------------------------------------------------------------
# 8. Eksport Model Terbaik & Label Encoders
# ------------------------------------------------------------
print("\n=== Proses Export Model ===")

# Gabungkan model dan encoders ke dalam satu dictionary agar rapi
export_package = {
    'model': models["KNN"],
    'label_encoders': label_encoders,
    'selected_features': selected_features # Penting agar tahu urutan kolom masuk saat predict
}

# Menentukan nama file yang sesuai untuk KNN
filename = 'model_knn.pkl'

# Menyimpan package ke dalam file .pkl
joblib.dump(export_package, filename)
print(f"Sukses! Model dan Label Encoders telah disimpan sebagai file: {filename}")
\end{lstlisting}

\begin{lstlisting}[caption={Output 17}]
rentals.csv
   id                   nama               nik status_pasutri  \
0   3  N*** Y**** S**** D***  3***************  Belum Menikah   
1   4     A***** N** H******  3***************  Belum Menikah   
2   5     A***** A*** S*****  3***************        Menikah   
3   6            C**** W****  1***************  Belum Menikah   
4   7             S*** U****  3***************  Belum Menikah   

  status_kewarganegaraan jenis_sewa unit_number metode_pembayaran metode_lain  \
0                    WNI     Harian     B-20-27              Cash         NaN   
1                    WNI     Harian     A-15-25              Cash         NaN   
2                    WNI     Harian     A-06-24            Others    Aplikasi   
3                    WNI     Harian     D-15-21              QRIS         NaN   
4                    WNI     Harian       A-8-8            Others    Transfer   

  tanggal_checkin waktu_checkin tanggal_checkout waktu_checkout  \
0      2026-05-20      13:08:00       2026-05-21            NaN   
1      2026-05-19      13:25:00       2026-05-22            NaN   
2      2026-05-17      13:27:00       2026-05-22            NaN   
3      2026-05-19      13:29:00       2026-05-24            NaN   
4      2025-01-01      19:58:00       2025-01-06       11:22:00   

   lama_menginap komentar  user_email  user_pengguna_id  editor_pengguna_id  \
0              1      NaN         NaN              1267                 NaN   
1              3      NaN         NaN              1267                 NaN   
2              5      NaN         NaN              1267                 NaN   
3              5      NaN         NaN              1267                 NaN   
4              5      NaN         NaN              1271              1271.0   

            created_at classification  
0  2026-05-20 13:09:49            NaN  
1  2026-05-20 13:26:04            NaN  
2  2026-05-20 13:28:03            NaN  
3  2026-05-20 13:30:18            NaN  
4  2026-05-25 10:06:33         normal  
logs.csv
     id  rental_id  action field_changed  old_value  \
0     6          3  create           all        NaN   
1     7          4  create           all        NaN   
2     8          5  create           all        NaN   
3     9          6  create           all        NaN   
4  4906          7  create           all        NaN   

                                           new_value  email  pengguna_id  \
0  {"nama":"N*** Y**** S**** D***","nik":"3*************************************"nama":"A***** N** H******","nik":"3****************************************"nama":"A***** A*** S*****","nik":"3****************************************"nama":"C**** W****","nik":"1***************"...    NaN         1267   
4  {"nama":"S*** U****","nik":"3***************",...    NaN         1271   

             timestamp  
0  2026-05-20 13:09:49  
1  2026-05-20 13:26:04  
2  2026-05-20 13:28:03  
3  2026-05-20 13:30:18  
4  2026-05-25 10:06:33  

==================================================
=== EXPLORATORY DATA ANALYSIS (EDA) ===
==================================================

[1] Statistik Deskriptif:
                id                         nama           nik status_pasutri  \
count   544.000000                          544  5.440000e+02            544   
unique         NaN                          207           NaN              2   
top            NaN  Y***** M**** S**** S*******           NaN  Belum Menikah   
freq           NaN                            8           NaN            330   
mean    274.538603                          NaN  3.214330e+15            NaN   
std     157.250941                          NaN  1.204147e+15            NaN   
min       3.000000                          NaN  1.202052e+15            NaN   
25%     138.750000                          NaN  3.173074e+15            NaN   
50%     274.500000                          NaN  3.209121e+15            NaN   
75%     410.250000                          NaN  3.276064e+15            NaN   
max     553.000000                          NaN  9.104017e+15            NaN   

       status_kewarganegaraan jenis_sewa unit_number metode_pembayaran  \
count                     544        544         544               544   
unique                      1          2          12                 5   
top                       WNI     Harian       B-1-8            Others   
freq                      544        543          90               209   
mean                      NaN        NaN         NaN               NaN   
std                       NaN        NaN         NaN               NaN   
min                       NaN        NaN         NaN               NaN   
25%                       NaN        NaN         NaN               NaN   
50%                       NaN        NaN         NaN               NaN   
75%                       NaN        NaN         NaN               NaN   
max                       NaN        NaN         NaN               NaN   

       metode_lain tanggal_checkin  ... tanggal_checkout waktu_checkout  \
count          209             544  ...              544            538   
unique           2             311  ...              300            232   
top       Transfer      2026-02-06  ...       2026-02-08       12:00:00   
freq           208               5  ...                5             82   
mean           NaN             NaN  ...              NaN            NaN   
std            NaN             NaN  ...              NaN            NaN   
min            NaN             NaN  ...              NaN            NaN   
25%            NaN             NaN  ...              NaN            NaN   
50%            NaN             NaN  ...              NaN            NaN   
75%            NaN             NaN  ...              NaN            NaN   
max            NaN             NaN  ...              NaN            NaN   

       lama_menginap  komentar user_email  user_pengguna_id  \
count     544.000000       544        0.0        544.000000   
unique           NaN        11        NaN               NaN   
top              NaN       nan        NaN               NaN   
freq             NaN       459        NaN               NaN   
mean        2.641544       NaN        NaN       1271.904412   
std         1.792711       NaN        NaN          1.177622   
min         1.000000       NaN        NaN       1262.000000   
25%         2.000000       NaN        NaN       1271.000000   
50%         2.000000       NaN        NaN       1272.000000   
75%         3.000000       NaN        NaN       1273.000000   
max        31.000000       NaN        NaN       1273.000000   

        editor_pengguna_id           created_at classification       label  
count           538.000000                  544            537  544.000000  
unique                 NaN                  544              2         NaN  
top                    NaN  2026-05-30 15:00:26         normal         NaN  
freq                   NaN                    1            452         NaN  
mean           1271.977695                  NaN            NaN    0.156250  
std               0.923249                  NaN            NaN    0.363426  
min            1262.000000                  NaN            NaN    0.000000  
25%            1271.000000                  NaN            NaN    0.000000  
50%            1272.000000                  NaN            NaN    0.000000  
75%            1273.000000                  NaN            NaN    0.000000  
max            1273.000000                  NaN            NaN    1.000000  

[11 rows x 21 columns]

[2] Distribusi Kelas Target (0: Aman, 1: Melanggar):
label
0    459
1     85
Name: count, dtype: int64

/tmp/ipykernel_1991/2187553676.py:70: FutureWarning: 

Passing `palette` without assigning `hue` is deprecated and will be removed in v0.14.0. Assign the `x` variable to `hue` and set `legend=False` for the same effect.

  sns.countplot(data=df_eda, x='label', palette='Set2', ax=ax)

<Figure size 600x400 with 1 Axes>

[3] Rata-rata Pelanggaran Berdasarkan Seluruh Kategori:

--- Tingkat Pelanggaran Berdasarkan status_pasutri ---
status_pasutri
Belum Menikah    0.164
Menikah          0.145
Name: label, dtype: float64

--- Tingkat Pelanggaran Berdasarkan status_kewarganegaraan ---
status_kewarganegaraan
WNI    0.156
Name: label, dtype: float64

--- Tingkat Pelanggaran Berdasarkan jenis_sewa ---
jenis_sewa
Harian     0.157
Bulanan    0.000
Name: label, dtype: float64

--- Tingkat Pelanggaran Berdasarkan metode_pembayaran ---
metode_pembayaran
Kartu Kredit    0.280
Others          0.158
Cash            0.150
QRIS            0.148
Kartu Debit     0.107
Name: label, dtype: float64

[4] Korelasi Antar Fitur:

/tmp/ipykernel_1991/2187553676.py:97: UserWarning: Could not infer format, so each element will be parsed individually, falling back to `dateutil`. To ensure parsing is consistent and as-expected, please specify a format.
  df_eda['checkin_hour_num'] = pd.to_datetime(df_eda['waktu_checkin'], errors='coerce').dt.hour.fillna(0)
/tmp/ipykernel_1991/2187553676.py:98: UserWarning: Could not infer format, so each element will be parsed individually, falling back to `dateutil`. To ensure parsing is consistent and as-expected, please specify a format.
  df_eda['checkout_hour_num'] = pd.to_datetime(df_eda['waktu_checkout'], errors='coerce').dt.hour.fillna(0)

<Figure size 1000x800 with 2 Axes>

[5] Visualisasi Pola Penyewaan (Dikaitkan dengan Klasifikasi):

<Figure size 1200x500 with 1 Axes>
<Figure size 1200x500 with 1 Axes>
<Figure size 800x500 with 1 Axes>
<Figure size 800x500 with 1 Axes>
<Figure size 800x500 with 1 Axes>
<Figure size 800x500 with 1 Axes>

[5b] Visualisasi Metode Pembayaran 'Others' (metode_lain):

/tmp/ipykernel_1991/2187553676.py:205: FutureWarning: 

Passing `palette` without assigning `hue` is deprecated and will be removed in v0.14.0. Assign the `y` variable to `hue` and set `legend=False` for the same effect.

  sns.countplot(

<Figure size 1000x600 with 1 Axes>
/tmp/ipykernel_1991/2187553676.py:239: UserWarning: Could not infer format, so each element will be parsed individually, falling back to `dateutil`. To ensure parsing is consistent and as-expected, please specify a format.
  rentals['waktu_checkout_dt'] = pd.to_datetime(rentals['waktu_checkout'], errors='coerce')
/tmp/ipykernel_1991/2187553676.py:242: UserWarning: Could not infer format, so each element will be parsed individually, falling back to `dateutil`. To ensure parsing is consistent and as-expected, please specify a format.
  rentals['waktu_checkin_dt'] = pd.to_datetime(rentals['waktu_checkin'], errors='coerce')


=== Proses Feature Selection ===

-> Fitur yang dipilih (5 terbaik): ['lama_menginap', 'metode_pembayaran', 'checkout_hour', 'checkin_hour', 'status_pasutri']

===== Decision Tree =====
Classification Report:
              precision    recall  f1-score   support

           0     1.0000    0.9565    0.9778        92
           1     0.8095    1.0000    0.8947        17

    accuracy                         0.9633       109
   macro avg     0.9048    0.9783    0.9363       109
weighted avg     0.9703    0.9633    0.9648       109

Confusion Matrix (Text):
                Predicted
               0       1
          -----------------
True   0 |     88       4
       1 |      0      17
          -----------------


===== Random Forest =====
Classification Report:
              precision    recall  f1-score   support

           0     1.0000    0.9565    0.9778        92
           1     0.8095    1.0000    0.8947        17

    accuracy                         0.9633       109
   macro avg     0.9048    0.9783    0.9363       109
weighted avg     0.9703    0.9633    0.9648       109

Confusion Matrix (Text):
                Predicted
               0       1
          -----------------
True   0 |     88       4
       1 |      0      17
          -----------------


===== KNN =====
Classification Report:
              precision    recall  f1-score   support

           0     0.9184    0.9783    0.9474        92
           1     0.8182    0.5294    0.6429        17

    accuracy                         0.9083       109
   macro avg     0.8683    0.7538    0.7951       109
weighted avg     0.9027    0.9083    0.8999       109

Confusion Matrix (Text):
                Predicted
               0       1
          -----------------
True   0 |     90       2
       1 |      8       9
          -----------------


===== SVM RBF =====
Classification Report:
              precision    recall  f1-score   support

           0     1.0000    0.9565    0.9778        92
           1     0.8095    1.0000    0.8947        17

    accuracy                         0.9633       109
   macro avg     0.9048    0.9783    0.9363       109
weighted avg     0.9703    0.9633    0.9648       109

Confusion Matrix (Text):
                Predicted
               0       1
          -----------------
True   0 |     88       4
       1 |      0      17
          -----------------


===== XGBoost =====
Classification Report:
              precision    recall  f1-score   support

           0     1.0000    0.9565    0.9778        92
           1     0.8095    1.0000    0.8947        17

    accuracy                         0.9633       109
   macro avg     0.9048    0.9783    0.9363       109
weighted avg     0.9703    0.9633    0.9648       109

Confusion Matrix (Text):
                Predicted
               0       1
          -----------------
True   0 |     88       4
       1 |      0      17
          -----------------


=== Proses Export Model ===
Sukses! Model dan Label Encoders telah disimpan sebagai file: model_knn.pkl
\end{lstlisting}

\begin{lstlisting}[language=Python, caption={Kode 18}]
# ============================================================
# DATA RIIL (TANPA TUNING) --> jumlah fitur = 4
# ============================================================

import pandas as pd
import numpy as np
import joblib
import matplotlib.pyplot as plt
import seaborn as sns

from sklearn.model_selection import train_test_split
from sklearn.preprocessing import LabelEncoder, StandardScaler, MinMaxScaler
from sklearn.metrics import classification_report, confusion_matrix
from sklearn.pipeline import make_pipeline
from sklearn.feature_selection import SelectKBest, chi2, mutual_info_classif

from sklearn.tree import DecisionTreeClassifier
from sklearn.ensemble import RandomForestClassifier
from sklearn.neighbors import KNeighborsClassifier
from sklearn.svm import SVC
from xgboost import XGBClassifier

import matplotlib.pyplot as plt
import seaborn as sns

# ------------------------------------------------------------
# 1. LOAD DATA
# ------------------------------------------------------------
rentals    = pd.read_csv('rentals.csv')
logs       = pd.read_csv('rental_logs.csv')
# units      = pd.read_csv('units.csv')
# users      = pd.read_csv('users.csv')
# violations = pd.read_csv('violations.csv')

rentals.columns = rentals.columns.str.lower()

if 'komentar' not in rentals.columns:
    raise ValueError("Kolom 'komentar' tidak ditemukan di rentals.csv")

print("rentals.csv")
print(rentals.head())
print("logs.csv")
print(logs.head())


# ------------------------------------------------------------
# 2. LABELING
# ------------------------------------------------------------
rentals['komentar'] = rentals['komentar'].astype(str).str.lower().str.strip()
empty_values = ["", "nan", "na", "null", "-", "0", "tidak ada", "none", " "]
rentals['label'] = (~rentals['komentar'].isin(empty_values)).astype(int)


# ============================================================
# EXPLORATORY DATA ANALYSIS (EDA)
# ============================================================
print("\n" + "="*50)
print("=== EXPLORATORY DATA ANALYSIS (EDA) ===")
print("="*50)

df_eda = rentals.copy()

print("\n[1] Statistik Deskriptif:")
print(df_eda.describe(include='all'))

print("\n[2] Distribusi Kelas Target (0: Aman, 1: Melanggar):")
print(df_eda['label'].value_counts())

fig, ax = plt.subplots(figsize=(6, 4))
sns.countplot(data=df_eda, x='label', palette='Set2', ax=ax)
plt.title('Distribusi Kelas Klasifikasi (0: Aman, 1: Melanggar)')
plt.xlabel('Klasifikasi')
plt.ylabel('Jumlah Penyewaan')

y_max = df_eda['label'].value_counts().max()
plt.ylim(0, y_max * 1.15)

for p in ax.patches:
    if p.get_height() > 0:
        ax.annotate(f'{int(p.get_height())}', (p.get_x() + p.get_width() / 2., p.get_height()),
                    ha='center', va='baseline', fontsize=10, color='black', xytext=(0, 5),
                    textcoords='offset points')
plt.show()

kolom_dikecualikan = ['id', 'nama', 'nik', 'unit_number', 'tanggal_checkin', 'waktu_checkin',
                      'tanggal_checkout', 'waktu_checkout', 'komentar', 'user_email', 'diedit_oleh', 'created_at', 'metode_lain', 'classification']

semua_kategori = [col for col in df_eda.select_dtypes(include=['object', 'category']).columns if col not in kolom_dikecualikan]

print("\n[3] Rata-rata Pelanggaran Berdasarkan Seluruh Kategori:")
for col in semua_kategori:
    if col in df_eda.columns:
        print(f"\n--- Tingkat Pelanggaran Berdasarkan {col} ---")
        print(df_eda.groupby(col)['label'].mean().sort_values(ascending=False).round(3))

df_eda['lama_menginap_num'] = pd.to_numeric(df_eda['lama_menginap'], errors='coerce').fillna(0)
df_eda['checkin_hour_num'] = pd.to_datetime(df_eda['waktu_checkin'], errors='coerce').dt.hour.fillna(0)
df_eda['checkout_hour_num'] = pd.to_datetime(df_eda['waktu_checkout'], errors='coerce').dt.hour.fillna(0)

le_eda = LabelEncoder()
for c in semua_kategori:
    if c in df_eda.columns:
        df_eda[c + '_encoded'] = le_eda.fit_transform(df_eda[c].astype(str))

print("\n[4] Korelasi Antar Fitur:")
kolom_korelasi = ['lama_menginap_num', 'checkin_hour_num', 'checkout_hour_num', 'label']
kolom_korelasi += [c + '_encoded' for c in semua_kategori if c + '_encoded' in df_eda.columns]

corr_matrix = df_eda[kolom_korelasi].corr()

plt.figure(figsize=(10, 8))
sns.heatmap(corr_matrix, annot=True, cmap='coolwarm', fmt='.2f', linewidths=0.5)
plt.title('Heatmap Korelasi Antar Fitur dan Target (Klasifikasi)')
plt.tight_layout()
plt.show()

print("\n[5] Visualisasi Pola Penyewaan (Dikaitkan dengan Klasifikasi):")

custom_line_palette = {0: 'blue', 1: 'red'}
custom_labels = ['0: Aman', '1: Melanggar']

fig, ax1 = plt.subplots(figsize=(12, 5))
waktu_df = df_eda.groupby(['checkin_hour_num', 'label']).size().reset_index(name='jumlah')

sns.lineplot(data=waktu_df, x='checkin_hour_num', y='jumlah', hue='label', marker='o', palette=custom_line_palette, ax=ax1)

plt.title('Pola Waktu Kedatangan Berdasarkan Klasifikasi')
plt.xlabel('Jam Check-in (0-23)')
plt.ylabel('Jumlah Penyewaan')
plt.xticks(range(0, 24))
plt.grid(True, alpha=0.3)

y_max_line = waktu_df['jumlah'].max()
plt.ylim(-y_max_line * 0.20, y_max_line * 1.15)

handles, _ = ax1.get_legend_handles_labels()
plt.legend(handles=handles, labels=custom_labels, title='Klasifikasi')

for l in waktu_df['label'].unique():
    sub = waktu_df[waktu_df['label'] == l]
    warna_teks = custom_line_palette[l]

    for x, y in zip(sub['checkin_hour_num'], sub['jumlah']):
        y_offset = 8 if l == 0 else -12
        va_align = 'bottom' if l == 0 else 'top'
        plt.annotate(f'{int(y)}', (x, y), textcoords="offset points", xytext=(0, y_offset),
                     ha='center', va=va_align, fontsize=9, color=warna_teks, fontweight='bold')
plt.show()

fig, ax2 = plt.subplots(figsize=(12, 5))
lama_df = df_eda.groupby(['lama_menginap_num', 'label']).size().reset_index(name='jumlah')

sns.lineplot(data=lama_df, x='lama_menginap_num', y='jumlah', hue='label', marker='o', palette=custom_line_palette, ax=ax2)

plt.title('Tren Lama Menginap Berdasarkan Klasifikasi')
plt.xlabel('Lama Menginap (Hari)')
plt.ylabel('Jumlah Penyewaan')
plt.grid(True, alpha=0.3)

handles2, _ = ax2.get_legend_handles_labels()
plt.legend(handles=handles2, labels=custom_labels, title='Klasifikasi')
plt.show()

for col in semua_kategori:
    if col in df_eda.columns:
        fig, ax = plt.subplots(figsize=(8, 5))
        sns.countplot(data=df_eda, x=col, hue='label', palette='magma', ax=ax)

        judul_kolom = col.replace('_', ' ').title()
        plt.title(f'Perbandingan Jumlah Pelanggaran Berdasarkan {judul_kolom}')
        plt.xlabel(judul_kolom)
        plt.ylabel('Jumlah Penyewaan')
        plt.xticks(rotation=0)

        handles_bar, _ = ax.get_legend_handles_labels()
        plt.legend(handles=handles_bar, labels=custom_labels, title='Klasifikasi')

        y_max_bar = df_eda.groupby([col, 'label']).size().max()
        plt.ylim(0, y_max_bar * 1.15)

        for p in ax.patches:
            if p.get_height() > 0:
                ax.annotate(f'{int(p.get_height())}', (p.get_x() + p.get_width() / 2., p.get_height()),
                            ha='center', va='baseline', fontsize=9, color='black', xytext=(0, 5),
                            textcoords='offset points')
        plt.tight_layout()
        plt.show()


print("\n[5b] Visualisasi Metode Pembayaran 'Others' (metode_lain):")

# Memastikan kolom 'metode_lain' berbentuk string dan membersihkan spasi/nilai kosong
df_eda['metode_lain'] = df_eda['metode_lain'].astype(str).str.strip().str.title()
empty_lain = ["", "Nan", "Na", "Null", "-", "0", "None"]

# Filter hanya baris yang memiliki nilai di 'metode_lain'
df_others = df_eda[~df_eda['metode_lain'].isin(empty_lain)]

if not df_others.empty:
    fig, ax = plt.subplots(figsize=(10, 6))

    # Menghitung frekuensi masing-masing metode pembayaran lain
    order_others = df_others['metode_lain'].value_counts().index

    sns.countplot(
        data=df_others,
        y='metode_lain',
        order=order_others,
        palette='viridis',
        ax=ax
    )

    plt.title('Detail Metode Pembayaran "Others" (metode_lain)')
    plt.xlabel('Jumlah Penyewaan')
    plt.ylabel('Jenis Pembayaran Lainnya')

    # Menambahkan anotasi angka di sebelah setiap bar
    for p in ax.patches:
        width = p.get_width()
        if width > 0:
            ax.annotate(f'{int(width)}',
                        (width, p.get_y() + p.get_height() / 2.),
                        ha='left', va='center', fontsize=10, color='black',
                        xytext=(5, 0), textcoords='offset points')

    # Perlebar batas x-axis sedikit agar angka anotasi tidak terpotong
    x_max = df_others['metode_lain'].value_counts().max()
    plt.xlim(0, x_max * 1.15)

    plt.tight_layout()
    plt.show()
else:
    print("Tidak ada data metode pembayaran 'Others' yang valid untuk ditampilkan.")


# ------------------------------------------------------------
# 3. FEATURE ENGINEERING & PREPARATION
# ------------------------------------------------------------
rentals['waktu_checkout_dt'] = pd.to_datetime(rentals['waktu_checkout'], errors='coerce')
rentals['checkout_hour'] = rentals['waktu_checkout_dt'].dt.hour.fillna(0).astype(int)

rentals['waktu_checkin_dt'] = pd.to_datetime(rentals['waktu_checkin'], errors='coerce')
rentals['checkin_hour'] = rentals['waktu_checkin_dt'].dt.hour.fillna(0).astype(int)

rentals['lama_menginap'] = pd.to_numeric(rentals['lama_menginap'], errors='coerce').fillna(0)

cat_cols = ['jenis_sewa', 'metode_pembayaran', 'status_pasutri', 'status_kewarganegaraan']
for c in cat_cols:
    rentals[c] = rentals[c].fillna('Unknown').astype(str)

numeric_cols = ['checkout_hour', 'checkin_hour', 'lama_menginap']
feature_cols = cat_cols + numeric_cols

X = rentals[feature_cols].copy()
y = rentals['label']

# ------------------------------------------------------------
# 4. TRAIN-TEST SPLIT
# ------------------------------------------------------------
X_train, X_test, y_train, y_test = train_test_split(
    X, y, stratify=y, test_size=0.2, random_state=42
)

# ------------------------------------------------------------
# 4B. ENCODING MENCEGAH DATA LEAKAGE
# ------------------------------------------------------------
label_encoders = {}

for c in cat_cols:
    le = LabelEncoder()
    X_train[c] = le.fit_transform(X_train[c])

    X_test[c] = X_test[c].map(lambda s: s if s in le.classes_ else 'Unknown')
    if 'Unknown' not in le.classes_:
        le.classes_ = np.append(le.classes_, 'Unknown')

    X_test[c] = le.transform(X_test[c])

    label_encoders[c] = le

# ------------------------------------------------------------
# 5. FEATURE SELECTION
# ------------------------------------------------------------
print("\n=== Proses Feature Selection ===")
chi2_selector = SelectKBest(score_func=chi2, k='all')
chi2_selector.fit(X_train, y_train)

train_df = X_train.copy()
train_df['label'] = y_train
pearson_scores = train_df.corr(method='pearson')['label'].drop('label').abs().values

mi_selector = SelectKBest(score_func=mutual_info_classif, k='all')
mi_selector.fit(X_train, y_train)

fs_df = pd.DataFrame({
    'Feature': X_train.columns,
    'Chi-Square': chi2_selector.scores_,
    'Pearson (Abs)': pearson_scores,
    'Mutual Info': mi_selector.scores_
})

scaler = MinMaxScaler()
fs_df[['Chi2_Norm', 'Pearson_Norm', 'MI_Norm']] = scaler.fit_transform(
    fs_df[['Chi-Square', 'Pearson (Abs)', 'Mutual Info']]
)
fs_df['Total_Score'] = fs_df[['Chi2_Norm', 'Pearson_Norm', 'MI_Norm']].mean(axis=1)
fs_df = fs_df.sort_values(by='Total_Score', ascending=False).reset_index(drop=True)

# Jumlah Fitur
n_select = 4
selected_features = fs_df['Feature'].head(n_select).tolist()
print(f"\n-> Fitur yang dipilih ({n_select} terbaik): {selected_features}")

X_train = X_train[selected_features]
X_test = X_test[selected_features]

# ------------------------------------------------------------
# 6. Fungsi Evaluasi
# ------------------------------------------------------------
def evaluate(name, model):
    y_pred = model.predict(X_test)
    print(f"\n===== {name} =====")
    print("Classification Report:")
    print(classification_report(y_test, y_pred, digits=4))
    cm = confusion_matrix(y_test, y_pred)
    print("Confusion Matrix (Text):")
    print("                Predicted")
    print("               0       1")
    print("          -----------------")
    print(f"True   0 | {cm[0][0]:6d}  {cm[0][1]:6d}")
    print(f"       1 | {cm[1][0]:6d}  {cm[1][1]:6d}")
    print("          -----------------\n")

# ------------------------------------------------------------
# 7. Latih Semua Model (TANPA TUNING)
# ------------------------------------------------------------
models = {}
custom_weight = {0: 1, 1: 4}

# Model menggunakan pengaturan tetap
models["Decision Tree"] = DecisionTreeClassifier(max_depth=6, class_weight=custom_weight, random_state=42)
models["Random Forest"] = RandomForestClassifier(n_estimators=200, max_depth=7, class_weight=custom_weight, random_state=42)
models["KNN"] = make_pipeline(StandardScaler(), KNeighborsClassifier(n_neighbors=5, weights='distance'))
models["SVM RBF"] = make_pipeline(StandardScaler(), SVC(kernel='rbf', probability=True, class_weight=custom_weight, random_state=42))
models["XGBoost"] = XGBClassifier(n_estimators=200, max_depth=5, learning_rate=0.05, scale_pos_weight=4, eval_metric='logloss', random_state=42)

for name, model in models.items():
    model.fit(X_train, y_train)
    evaluate(name, model)

# ------------------------------------------------------------
# 8. Eksport Model Terbaik & Label Encoders
# ------------------------------------------------------------
print("\n=== Proses Export Model ===")

# Gabungkan model dan encoders ke dalam satu dictionary agar rapi
export_package = {
    'model': models["KNN"],
    'label_encoders': label_encoders,
    'selected_features': selected_features # Penting agar tahu urutan kolom masuk saat predict
}

# Menentukan nama file yang sesuai untuk KNN
filename = 'model_knn.pkl'

# Menyimpan package ke dalam file .pkl
joblib.dump(export_package, filename)
print(f"Sukses! Model dan Label Encoders telah disimpan sebagai file: {filename}")
\end{lstlisting}

\begin{lstlisting}[caption={Output 18}]
rentals.csv
   id                   nama               nik status_pasutri  \
0   3  N*** Y**** S**** D***  3***************  Belum Menikah   
1   4     A***** N** H******  3***************  Belum Menikah   
2   5     A***** A*** S*****  3***************        Menikah   
3   6            C**** W****  1***************  Belum Menikah   
4   7             S*** U****  3***************  Belum Menikah   

  status_kewarganegaraan jenis_sewa unit_number metode_pembayaran metode_lain  \
0                    WNI     Harian     B-20-27              Cash         NaN   
1                    WNI     Harian     A-15-25              Cash         NaN   
2                    WNI     Harian     A-06-24            Others    Aplikasi   
3                    WNI     Harian     D-15-21              QRIS         NaN   
4                    WNI     Harian       A-8-8            Others    Transfer   

  tanggal_checkin waktu_checkin tanggal_checkout waktu_checkout  \
0      2026-05-20      13:08:00       2026-05-21            NaN   
1      2026-05-19      13:25:00       2026-05-22            NaN   
2      2026-05-17      13:27:00       2026-05-22            NaN   
3      2026-05-19      13:29:00       2026-05-24            NaN   
4      2025-01-01      19:58:00       2025-01-06       11:22:00   

   lama_menginap komentar  user_email  user_pengguna_id  editor_pengguna_id  \
0              1      NaN         NaN              1267                 NaN   
1              3      NaN         NaN              1267                 NaN   
2              5      NaN         NaN              1267                 NaN   
3              5      NaN         NaN              1267                 NaN   
4              5      NaN         NaN              1271              1271.0   

            created_at classification  
0  2026-05-20 13:09:49            NaN  
1  2026-05-20 13:26:04            NaN  
2  2026-05-20 13:28:03            NaN  
3  2026-05-20 13:30:18            NaN  
4  2026-05-25 10:06:33         normal  
logs.csv
     id  rental_id  action field_changed  old_value  \
0     6          3  create           all        NaN   
1     7          4  create           all        NaN   
2     8          5  create           all        NaN   
3     9          6  create           all        NaN   
4  4906          7  create           all        NaN   

                                           new_value  email  pengguna_id  \
0  {"nama":"N*** Y**** S**** D***","nik":"3*************************************"nama":"A***** N** H******","nik":"3****************************************"nama":"A***** A*** S*****","nik":"3****************************************"nama":"C**** W****","nik":"1***************"...    NaN         1267   
4  {"nama":"S*** U****","nik":"3***************",...    NaN         1271   

             timestamp  
0  2026-05-20 13:09:49  
1  2026-05-20 13:26:04  
2  2026-05-20 13:28:03  
3  2026-05-20 13:30:18  
4  2026-05-25 10:06:33  

==================================================
=== EXPLORATORY DATA ANALYSIS (EDA) ===
==================================================

[1] Statistik Deskriptif:
                id                         nama           nik status_pasutri  \
count   544.000000                          544  5.440000e+02            544   
unique         NaN                          207           NaN              2   
top            NaN  Y***** M**** S**** S*******           NaN  Belum Menikah   
freq           NaN                            8           NaN            330   
mean    274.538603                          NaN  3.214330e+15            NaN   
std     157.250941                          NaN  1.204147e+15            NaN   
min       3.000000                          NaN  1.202052e+15            NaN   
25%     138.750000                          NaN  3.173074e+15            NaN   
50%     274.500000                          NaN  3.209121e+15            NaN   
75%     410.250000                          NaN  3.276064e+15            NaN   
max     553.000000                          NaN  9.104017e+15            NaN   

       status_kewarganegaraan jenis_sewa unit_number metode_pembayaran  \
count                     544        544         544               544   
unique                      1          2          12                 5   
top                       WNI     Harian       B-1-8            Others   
freq                      544        543          90               209   
mean                      NaN        NaN         NaN               NaN   
std                       NaN        NaN         NaN               NaN   
min                       NaN        NaN         NaN               NaN   
25%                       NaN        NaN         NaN               NaN   
50%                       NaN        NaN         NaN               NaN   
75%                       NaN        NaN         NaN               NaN   
max                       NaN        NaN         NaN               NaN   

       metode_lain tanggal_checkin  ... tanggal_checkout waktu_checkout  \
count          209             544  ...              544            538   
unique           2             311  ...              300            232   
top       Transfer      2026-02-06  ...       2026-02-08       12:00:00   
freq           208               5  ...                5             82   
mean           NaN             NaN  ...              NaN            NaN   
std            NaN             NaN  ...              NaN            NaN   
min            NaN             NaN  ...              NaN            NaN   
25%            NaN             NaN  ...              NaN            NaN   
50%            NaN             NaN  ...              NaN            NaN   
75%            NaN             NaN  ...              NaN            NaN   
max            NaN             NaN  ...              NaN            NaN   

       lama_menginap  komentar user_email  user_pengguna_id  \
count     544.000000       544        0.0        544.000000   
unique           NaN        11        NaN               NaN   
top              NaN       nan        NaN               NaN   
freq             NaN       459        NaN               NaN   
mean        2.641544       NaN        NaN       1271.904412   
std         1.792711       NaN        NaN          1.177622   
min         1.000000       NaN        NaN       1262.000000   
25%         2.000000       NaN        NaN       1271.000000   
50%         2.000000       NaN        NaN       1272.000000   
75%         3.000000       NaN        NaN       1273.000000   
max        31.000000       NaN        NaN       1273.000000   

        editor_pengguna_id           created_at classification       label  
count           538.000000                  544            537  544.000000  
unique                 NaN                  544              2         NaN  
top                    NaN  2026-05-30 15:00:26         normal         NaN  
freq                   NaN                    1            452         NaN  
mean           1271.977695                  NaN            NaN    0.156250  
std               0.923249                  NaN            NaN    0.363426  
min            1262.000000                  NaN            NaN    0.000000  
25%            1271.000000                  NaN            NaN    0.000000  
50%            1272.000000                  NaN            NaN    0.000000  
75%            1273.000000                  NaN            NaN    0.000000  
max            1273.000000                  NaN            NaN    1.000000  

[11 rows x 21 columns]

[2] Distribusi Kelas Target (0: Aman, 1: Melanggar):
label
0    459
1     85
Name: count, dtype: int64

/tmp/ipykernel_1991/491046311.py:70: FutureWarning: 

Passing `palette` without assigning `hue` is deprecated and will be removed in v0.14.0. Assign the `x` variable to `hue` and set `legend=False` for the same effect.

  sns.countplot(data=df_eda, x='label', palette='Set2', ax=ax)

<Figure size 600x400 with 1 Axes>

[3] Rata-rata Pelanggaran Berdasarkan Seluruh Kategori:

--- Tingkat Pelanggaran Berdasarkan status_pasutri ---
status_pasutri
Belum Menikah    0.164
Menikah          0.145
Name: label, dtype: float64

--- Tingkat Pelanggaran Berdasarkan status_kewarganegaraan ---
status_kewarganegaraan
WNI    0.156
Name: label, dtype: float64

--- Tingkat Pelanggaran Berdasarkan jenis_sewa ---
jenis_sewa
Harian     0.157
Bulanan    0.000
Name: label, dtype: float64

--- Tingkat Pelanggaran Berdasarkan metode_pembayaran ---
metode_pembayaran
Kartu Kredit    0.280
Others          0.158
Cash            0.150
QRIS            0.148
Kartu Debit     0.107
Name: label, dtype: float64

[4] Korelasi Antar Fitur:

/tmp/ipykernel_1991/491046311.py:97: UserWarning: Could not infer format, so each element will be parsed individually, falling back to `dateutil`. To ensure parsing is consistent and as-expected, please specify a format.
  df_eda['checkin_hour_num'] = pd.to_datetime(df_eda['waktu_checkin'], errors='coerce').dt.hour.fillna(0)
/tmp/ipykernel_1991/491046311.py:98: UserWarning: Could not infer format, so each element will be parsed individually, falling back to `dateutil`. To ensure parsing is consistent and as-expected, please specify a format.
  df_eda['checkout_hour_num'] = pd.to_datetime(df_eda['waktu_checkout'], errors='coerce').dt.hour.fillna(0)

<Figure size 1000x800 with 2 Axes>

[5] Visualisasi Pola Penyewaan (Dikaitkan dengan Klasifikasi):

<Figure size 1200x500 with 1 Axes>
<Figure size 1200x500 with 1 Axes>
<Figure size 800x500 with 1 Axes>
<Figure size 800x500 with 1 Axes>
<Figure size 800x500 with 1 Axes>
<Figure size 800x500 with 1 Axes>

[5b] Visualisasi Metode Pembayaran 'Others' (metode_lain):

/tmp/ipykernel_1991/491046311.py:205: FutureWarning: 

Passing `palette` without assigning `hue` is deprecated and will be removed in v0.14.0. Assign the `y` variable to `hue` and set `legend=False` for the same effect.

  sns.countplot(

<Figure size 1000x600 with 1 Axes>
/tmp/ipykernel_1991/491046311.py:239: UserWarning: Could not infer format, so each element will be parsed individually, falling back to `dateutil`. To ensure parsing is consistent and as-expected, please specify a format.
  rentals['waktu_checkout_dt'] = pd.to_datetime(rentals['waktu_checkout'], errors='coerce')
/tmp/ipykernel_1991/491046311.py:242: UserWarning: Could not infer format, so each element will be parsed individually, falling back to `dateutil`. To ensure parsing is consistent and as-expected, please specify a format.
  rentals['waktu_checkin_dt'] = pd.to_datetime(rentals['waktu_checkin'], errors='coerce')


=== Proses Feature Selection ===

-> Fitur yang dipilih (4 terbaik): ['lama_menginap', 'checkout_hour', 'checkin_hour', 'metode_pembayaran']

===== Decision Tree =====
Classification Report:
              precision    recall  f1-score   support

           0     1.0000    0.9565    0.9778        92
           1     0.8095    1.0000    0.8947        17

    accuracy                         0.9633       109
   macro avg     0.9048    0.9783    0.9363       109
weighted avg     0.9703    0.9633    0.9648       109

Confusion Matrix (Text):
                Predicted
               0       1
          -----------------
True   0 |     88       4
       1 |      0      17
          -----------------


===== Random Forest =====
Classification Report:
              precision    recall  f1-score   support

           0     1.0000    0.9565    0.9778        92
           1     0.8095    1.0000    0.8947        17

    accuracy                         0.9633       109
   macro avg     0.9048    0.9783    0.9363       109
weighted avg     0.9703    0.9633    0.9648       109

Confusion Matrix (Text):
                Predicted
               0       1
          -----------------
True   0 |     88       4
       1 |      0      17
          -----------------


===== KNN =====
Classification Report:
              precision    recall  f1-score   support

           0     0.9468    0.9674    0.9570        92
           1     0.8000    0.7059    0.7500        17

    accuracy                         0.9266       109
   macro avg     0.8734    0.8366    0.8535       109
weighted avg     0.9239    0.9266    0.9247       109

Confusion Matrix (Text):
                Predicted
               0       1
          -----------------
True   0 |     89       3
       1 |      5      12
          -----------------


===== SVM RBF =====
Classification Report:
              precision    recall  f1-score   support

           0     1.0000    0.9565    0.9778        92
           1     0.8095    1.0000    0.8947        17

    accuracy                         0.9633       109
   macro avg     0.9048    0.9783    0.9363       109
weighted avg     0.9703    0.9633    0.9648       109

Confusion Matrix (Text):
                Predicted
               0       1
          -----------------
True   0 |     88       4
       1 |      0      17
          -----------------


===== XGBoost =====
Classification Report:
              precision    recall  f1-score   support

           0     1.0000    0.9565    0.9778        92
           1     0.8095    1.0000    0.8947        17

    accuracy                         0.9633       109
   macro avg     0.9048    0.9783    0.9363       109
weighted avg     0.9703    0.9633    0.9648       109

Confusion Matrix (Text):
                Predicted
               0       1
          -----------------
True   0 |     88       4
       1 |      0      17
          -----------------


=== Proses Export Model ===
Sukses! Model dan Label Encoders telah disimpan sebagai file: model_knn.pkl
\end{lstlisting}

\begin{lstlisting}[language=Python, caption={Kode 19}]
# ============================================================
# DATA RIIL (TANPA TUNING) --> jumlah fitur = 3
# ============================================================

import pandas as pd
import numpy as np
import joblib
import matplotlib.pyplot as plt
import seaborn as sns

from sklearn.model_selection import train_test_split
from sklearn.preprocessing import LabelEncoder, StandardScaler, MinMaxScaler
from sklearn.metrics import classification_report, confusion_matrix
from sklearn.pipeline import make_pipeline
from sklearn.feature_selection import SelectKBest, chi2, mutual_info_classif

from sklearn.tree import DecisionTreeClassifier
from sklearn.ensemble import RandomForestClassifier
from sklearn.neighbors import KNeighborsClassifier
from sklearn.svm import SVC
from xgboost import XGBClassifier

import matplotlib.pyplot as plt
import seaborn as sns

# ------------------------------------------------------------
# 1. LOAD DATA
# ------------------------------------------------------------
rentals    = pd.read_csv('rentals.csv')
logs       = pd.read_csv('rental_logs.csv')
# units      = pd.read_csv('units.csv')
# users      = pd.read_csv('users.csv')
# violations = pd.read_csv('violations.csv')

rentals.columns = rentals.columns.str.lower()

if 'komentar' not in rentals.columns:
    raise ValueError("Kolom 'komentar' tidak ditemukan di rentals.csv")

print("rentals.csv")
print(rentals.head())
print("logs.csv")
print(logs.head())


# ------------------------------------------------------------
# 2. LABELING
# ------------------------------------------------------------
rentals['komentar'] = rentals['komentar'].astype(str).str.lower().str.strip()
empty_values = ["", "nan", "na", "null", "-", "0", "tidak ada", "none", " "]
rentals['label'] = (~rentals['komentar'].isin(empty_values)).astype(int)


# ============================================================
# EXPLORATORY DATA ANALYSIS (EDA)
# ============================================================
print("\n" + "="*50)
print("=== EXPLORATORY DATA ANALYSIS (EDA) ===")
print("="*50)

df_eda = rentals.copy()

print("\n[1] Statistik Deskriptif:")
print(df_eda.describe(include='all'))

print("\n[2] Distribusi Kelas Target (0: Aman, 1: Melanggar):")
print(df_eda['label'].value_counts())

fig, ax = plt.subplots(figsize=(6, 4))
sns.countplot(data=df_eda, x='label', palette='Set2', ax=ax)
plt.title('Distribusi Kelas Klasifikasi (0: Aman, 1: Melanggar)')
plt.xlabel('Klasifikasi')
plt.ylabel('Jumlah Penyewaan')

y_max = df_eda['label'].value_counts().max()
plt.ylim(0, y_max * 1.15)

for p in ax.patches:
    if p.get_height() > 0:
        ax.annotate(f'{int(p.get_height())}', (p.get_x() + p.get_width() / 2., p.get_height()),
                    ha='center', va='baseline', fontsize=10, color='black', xytext=(0, 5),
                    textcoords='offset points')
plt.show()

kolom_dikecualikan = ['id', 'nama', 'nik', 'unit_number', 'tanggal_checkin', 'waktu_checkin',
                      'tanggal_checkout', 'waktu_checkout', 'komentar', 'user_email', 'diedit_oleh', 'created_at', 'metode_lain', 'classification']

semua_kategori = [col for col in df_eda.select_dtypes(include=['object', 'category']).columns if col not in kolom_dikecualikan]

print("\n[3] Rata-rata Pelanggaran Berdasarkan Seluruh Kategori:")
for col in semua_kategori:
    if col in df_eda.columns:
        print(f"\n--- Tingkat Pelanggaran Berdasarkan {col} ---")
        print(df_eda.groupby(col)['label'].mean().sort_values(ascending=False).round(3))

df_eda['lama_menginap_num'] = pd.to_numeric(df_eda['lama_menginap'], errors='coerce').fillna(0)
df_eda['checkin_hour_num'] = pd.to_datetime(df_eda['waktu_checkin'], errors='coerce').dt.hour.fillna(0)
df_eda['checkout_hour_num'] = pd.to_datetime(df_eda['waktu_checkout'], errors='coerce').dt.hour.fillna(0)

le_eda = LabelEncoder()
for c in semua_kategori:
    if c in df_eda.columns:
        df_eda[c + '_encoded'] = le_eda.fit_transform(df_eda[c].astype(str))

print("\n[4] Korelasi Antar Fitur:")
kolom_korelasi = ['lama_menginap_num', 'checkin_hour_num', 'checkout_hour_num', 'label']
kolom_korelasi += [c + '_encoded' for c in semua_kategori if c + '_encoded' in df_eda.columns]

corr_matrix = df_eda[kolom_korelasi].corr()

plt.figure(figsize=(10, 8))
sns.heatmap(corr_matrix, annot=True, cmap='coolwarm', fmt='.2f', linewidths=0.5)
plt.title('Heatmap Korelasi Antar Fitur dan Target (Klasifikasi)')
plt.tight_layout()
plt.show()

print("\n[5] Visualisasi Pola Penyewaan (Dikaitkan dengan Klasifikasi):")

custom_line_palette = {0: 'blue', 1: 'red'}
custom_labels = ['0: Aman', '1: Melanggar']

fig, ax1 = plt.subplots(figsize=(12, 5))
waktu_df = df_eda.groupby(['checkin_hour_num', 'label']).size().reset_index(name='jumlah')

sns.lineplot(data=waktu_df, x='checkin_hour_num', y='jumlah', hue='label', marker='o', palette=custom_line_palette, ax=ax1)

plt.title('Pola Waktu Kedatangan Berdasarkan Klasifikasi')
plt.xlabel('Jam Check-in (0-23)')
plt.ylabel('Jumlah Penyewaan')
plt.xticks(range(0, 24))
plt.grid(True, alpha=0.3)

y_max_line = waktu_df['jumlah'].max()
plt.ylim(-y_max_line * 0.20, y_max_line * 1.15)

handles, _ = ax1.get_legend_handles_labels()
plt.legend(handles=handles, labels=custom_labels, title='Klasifikasi')

for l in waktu_df['label'].unique():
    sub = waktu_df[waktu_df['label'] == l]
    warna_teks = custom_line_palette[l]

    for x, y in zip(sub['checkin_hour_num'], sub['jumlah']):
        y_offset = 8 if l == 0 else -12
        va_align = 'bottom' if l == 0 else 'top'
        plt.annotate(f'{int(y)}', (x, y), textcoords="offset points", xytext=(0, y_offset),
                     ha='center', va=va_align, fontsize=9, color=warna_teks, fontweight='bold')
plt.show()

fig, ax2 = plt.subplots(figsize=(12, 5))
lama_df = df_eda.groupby(['lama_menginap_num', 'label']).size().reset_index(name='jumlah')

sns.lineplot(data=lama_df, x='lama_menginap_num', y='jumlah', hue='label', marker='o', palette=custom_line_palette, ax=ax2)

plt.title('Tren Lama Menginap Berdasarkan Klasifikasi')
plt.xlabel('Lama Menginap (Hari)')
plt.ylabel('Jumlah Penyewaan')
plt.grid(True, alpha=0.3)

handles2, _ = ax2.get_legend_handles_labels()
plt.legend(handles=handles2, labels=custom_labels, title='Klasifikasi')
plt.show()

for col in semua_kategori:
    if col in df_eda.columns:
        fig, ax = plt.subplots(figsize=(8, 5))
        sns.countplot(data=df_eda, x=col, hue='label', palette='magma', ax=ax)

        judul_kolom = col.replace('_', ' ').title()
        plt.title(f'Perbandingan Jumlah Pelanggaran Berdasarkan {judul_kolom}')
        plt.xlabel(judul_kolom)
        plt.ylabel('Jumlah Penyewaan')
        plt.xticks(rotation=0)

        handles_bar, _ = ax.get_legend_handles_labels()
        plt.legend(handles=handles_bar, labels=custom_labels, title='Klasifikasi')

        y_max_bar = df_eda.groupby([col, 'label']).size().max()
        plt.ylim(0, y_max_bar * 1.15)

        for p in ax.patches:
            if p.get_height() > 0:
                ax.annotate(f'{int(p.get_height())}', (p.get_x() + p.get_width() / 2., p.get_height()),
                            ha='center', va='baseline', fontsize=9, color='black', xytext=(0, 5),
                            textcoords='offset points')
        plt.tight_layout()
        plt.show()


print("\n[5b] Visualisasi Metode Pembayaran 'Others' (metode_lain):")

# Memastikan kolom 'metode_lain' berbentuk string dan membersihkan spasi/nilai kosong
df_eda['metode_lain'] = df_eda['metode_lain'].astype(str).str.strip().str.title()
empty_lain = ["", "Nan", "Na", "Null", "-", "0", "None"]

# Filter hanya baris yang memiliki nilai di 'metode_lain'
df_others = df_eda[~df_eda['metode_lain'].isin(empty_lain)]

if not df_others.empty:
    fig, ax = plt.subplots(figsize=(10, 6))

    # Menghitung frekuensi masing-masing metode pembayaran lain
    order_others = df_others['metode_lain'].value_counts().index

    sns.countplot(
        data=df_others,
        y='metode_lain',
        order=order_others,
        palette='viridis',
        ax=ax
    )

    plt.title('Detail Metode Pembayaran "Others" (metode_lain)')
    plt.xlabel('Jumlah Penyewaan')
    plt.ylabel('Jenis Pembayaran Lainnya')

    # Menambahkan anotasi angka di sebelah setiap bar
    for p in ax.patches:
        width = p.get_width()
        if width > 0:
            ax.annotate(f'{int(width)}',
                        (width, p.get_y() + p.get_height() / 2.),
                        ha='left', va='center', fontsize=10, color='black',
                        xytext=(5, 0), textcoords='offset points')

    # Perlebar batas x-axis sedikit agar angka anotasi tidak terpotong
    x_max = df_others['metode_lain'].value_counts().max()
    plt.xlim(0, x_max * 1.15)

    plt.tight_layout()
    plt.show()
else:
    print("Tidak ada data metode pembayaran 'Others' yang valid untuk ditampilkan.")


# ------------------------------------------------------------
# 3. FEATURE ENGINEERING & PREPARATION
# ------------------------------------------------------------
rentals['waktu_checkout_dt'] = pd.to_datetime(rentals['waktu_checkout'], errors='coerce')
rentals['checkout_hour'] = rentals['waktu_checkout_dt'].dt.hour.fillna(0).astype(int)

rentals['waktu_checkin_dt'] = pd.to_datetime(rentals['waktu_checkin'], errors='coerce')
rentals['checkin_hour'] = rentals['waktu_checkin_dt'].dt.hour.fillna(0).astype(int)

rentals['lama_menginap'] = pd.to_numeric(rentals['lama_menginap'], errors='coerce').fillna(0)

cat_cols = ['jenis_sewa', 'metode_pembayaran', 'status_pasutri', 'status_kewarganegaraan']
for c in cat_cols:
    rentals[c] = rentals[c].fillna('Unknown').astype(str)

numeric_cols = ['checkout_hour', 'checkin_hour', 'lama_menginap']
feature_cols = cat_cols + numeric_cols

X = rentals[feature_cols].copy()
y = rentals['label']

# ------------------------------------------------------------
# 4. TRAIN-TEST SPLIT
# ------------------------------------------------------------
X_train, X_test, y_train, y_test = train_test_split(
    X, y, stratify=y, test_size=0.2, random_state=42
)

# ------------------------------------------------------------
# 4B. ENCODING MENCEGAH DATA LEAKAGE
# ------------------------------------------------------------
label_encoders = {}

for c in cat_cols:
    le = LabelEncoder()
    X_train[c] = le.fit_transform(X_train[c])

    X_test[c] = X_test[c].map(lambda s: s if s in le.classes_ else 'Unknown')
    if 'Unknown' not in le.classes_:
        le.classes_ = np.append(le.classes_, 'Unknown')

    X_test[c] = le.transform(X_test[c])

    label_encoders[c] = le

# ------------------------------------------------------------
# 5. FEATURE SELECTION
# ------------------------------------------------------------
print("\n=== Proses Feature Selection ===")
chi2_selector = SelectKBest(score_func=chi2, k='all')
chi2_selector.fit(X_train, y_train)

train_df = X_train.copy()
train_df['label'] = y_train
pearson_scores = train_df.corr(method='pearson')['label'].drop('label').abs().values

mi_selector = SelectKBest(score_func=mutual_info_classif, k='all')
mi_selector.fit(X_train, y_train)

fs_df = pd.DataFrame({
    'Feature': X_train.columns,
    'Chi-Square': chi2_selector.scores_,
    'Pearson (Abs)': pearson_scores,
    'Mutual Info': mi_selector.scores_
})

scaler = MinMaxScaler()
fs_df[['Chi2_Norm', 'Pearson_Norm', 'MI_Norm']] = scaler.fit_transform(
    fs_df[['Chi-Square', 'Pearson (Abs)', 'Mutual Info']]
)
fs_df['Total_Score'] = fs_df[['Chi2_Norm', 'Pearson_Norm', 'MI_Norm']].mean(axis=1)
fs_df = fs_df.sort_values(by='Total_Score', ascending=False).reset_index(drop=True)

# Jumlah Fitur
n_select = 3
selected_features = fs_df['Feature'].head(n_select).tolist()
print(f"\n-> Fitur yang dipilih ({n_select} terbaik): {selected_features}")

X_train = X_train[selected_features]
X_test = X_test[selected_features]

# ------------------------------------------------------------
# 6. Fungsi Evaluasi
# ------------------------------------------------------------
def evaluate(name, model):
    y_pred = model.predict(X_test)
    print(f"\n===== {name} =====")
    print("Classification Report:")
    print(classification_report(y_test, y_pred, digits=4))

    cm = confusion_matrix(y_test, y_pred)

    # --------------------------------------------------------
    # Confusion Matrix (Text Base)
    # --------------------------------------------------------
    print("Confusion Matrix (Text):")
    print("                Predicted")
    print("               0       1")
    print("          -----------------")
    print(f"True   0 | {cm[0][0]:6d}  {cm[0][1]:6d}")
    print(f"       1 | {cm[1][0]:6d}  {cm[1][1]:6d}")
    print("          -----------------\n")

    # # --------------------------------------------------------
    # # Confusion Matrix (Gambar / Visualisasi)
    # # --------------------------------------------------------
    # plt.figure(figsize=(5, 4))
    # sns.heatmap(cm, annot=True, fmt='d', cmap='Blues', cbar=False,
    #             xticklabels=['0 (Aman)', '1 (Melanggar)'],
    #             yticklabels=['0 (Aman)', '1 (Melanggar)'])

    # plt.title(f'Confusion Matrix - {name}')
    # plt.xlabel('Predicted Label')
    # plt.ylabel('True Label')
    # plt.tight_layout()
    # plt.show()

# ------------------------------------------------------------
# 7. Latih Semua Model (TANPA TUNING)
# ------------------------------------------------------------
models = {}
custom_weight = {0: 1, 1: 4}

# Model menggunakan pengaturan tetap
models["Decision Tree"] = DecisionTreeClassifier(max_depth=6, class_weight=custom_weight, random_state=42)
models["Random Forest"] = RandomForestClassifier(n_estimators=200, max_depth=7, class_weight=custom_weight, random_state=42)
models["KNN"] = make_pipeline(StandardScaler(), KNeighborsClassifier(n_neighbors=5, weights='distance'))
models["SVM RBF"] = make_pipeline(StandardScaler(), SVC(kernel='rbf', probability=True, class_weight=custom_weight, random_state=42))
models["XGBoost"] = XGBClassifier(n_estimators=200, max_depth=5, learning_rate=0.05, scale_pos_weight=4, eval_metric='logloss', random_state=42)

for name, model in models.items():
    model.fit(X_train, y_train)
    evaluate(name, model)

# ------------------------------------------------------------
# 8. Eksport Model Terbaik & Label Encoders
# ------------------------------------------------------------
print("\n=== Proses Export Model ===")

# Gabungkan model dan encoders ke dalam satu dictionary agar rapi
export_package = {
    'model': models["KNN"],
    'label_encoders': label_encoders,
    'selected_features': selected_features # Penting agar tahu urutan kolom masuk saat predict
}

# Menentukan nama file yang sesuai untuk KNN
filename = 'model_knn.pkl'

# Menyimpan package ke dalam file .pkl
joblib.dump(export_package, filename)
print(f"Sukses! Model dan Label Encoders telah disimpan sebagai file: {filename}")
\end{lstlisting}

\begin{lstlisting}[caption={Output 19}]
rentals.csv
   id                   nama               nik status_pasutri  \
0   3  N*** Y**** S**** D***  3***************  Belum Menikah   
1   4     A***** N** H******  3***************  Belum Menikah   
2   5     A***** A*** S*****  3***************        Menikah   
3   6            C**** W****  1***************  Belum Menikah   
4   7             S*** U****  3***************  Belum Menikah   

  status_kewarganegaraan jenis_sewa unit_number metode_pembayaran metode_lain  \
0                    WNI     Harian     B-20-27              Cash         NaN   
1                    WNI     Harian     A-15-25              Cash         NaN   
2                    WNI     Harian     A-06-24            Others    Aplikasi   
3                    WNI     Harian     D-15-21              QRIS         NaN   
4                    WNI     Harian       A-8-8            Others    Transfer   

  tanggal_checkin waktu_checkin tanggal_checkout waktu_checkout  \
0      2026-05-20      13:08:00       2026-05-21            NaN   
1      2026-05-19      13:25:00       2026-05-22            NaN   
2      2026-05-17      13:27:00       2026-05-22            NaN   
3      2026-05-19      13:29:00       2026-05-24            NaN   
4      2025-01-01      19:58:00       2025-01-06       11:22:00   

   lama_menginap komentar  user_email  user_pengguna_id  editor_pengguna_id  \
0              1      NaN         NaN              1267                 NaN   
1              3      NaN         NaN              1267                 NaN   
2              5      NaN         NaN              1267                 NaN   
3              5      NaN         NaN              1267                 NaN   
4              5      NaN         NaN              1271              1271.0   

            created_at classification  
0  2026-05-20 13:09:49            NaN  
1  2026-05-20 13:26:04            NaN  
2  2026-05-20 13:28:03            NaN  
3  2026-05-20 13:30:18            NaN  
4  2026-05-25 10:06:33         normal  
logs.csv
     id  rental_id  action field_changed  old_value  \
0     6          3  create           all        NaN   
1     7          4  create           all        NaN   
2     8          5  create           all        NaN   
3     9          6  create           all        NaN   
4  4906          7  create           all        NaN   

                                           new_value  email  pengguna_id  \
0  {"nama":"N*** Y**** S**** D***","nik":"3*************************************"nama":"A***** N** H******","nik":"3****************************************"nama":"A***** A*** S*****","nik":"3****************************************"nama":"C**** W****","nik":"1***************"...    NaN         1267   
4  {"nama":"S*** U****","nik":"3***************",...    NaN         1271   

             timestamp  
0  2026-05-20 13:09:49  
1  2026-05-20 13:26:04  
2  2026-05-20 13:28:03  
3  2026-05-20 13:30:18  
4  2026-05-25 10:06:33  

==================================================
=== EXPLORATORY DATA ANALYSIS (EDA) ===
==================================================

[1] Statistik Deskriptif:
                id                         nama           nik status_pasutri  \
count   544.000000                          544  5.440000e+02            544   
unique         NaN                          207           NaN              2   
top            NaN  Y***** M**** S**** S*******           NaN  Belum Menikah   
freq           NaN                            8           NaN            330   
mean    274.538603                          NaN  3.214330e+15            NaN   
std     157.250941                          NaN  1.204147e+15            NaN   
min       3.000000                          NaN  1.202052e+15            NaN   
25%     138.750000                          NaN  3.173074e+15            NaN   
50%     274.500000                          NaN  3.209121e+15            NaN   
75%     410.250000                          NaN  3.276064e+15            NaN   
max     553.000000                          NaN  9.104017e+15            NaN   

       status_kewarganegaraan jenis_sewa unit_number metode_pembayaran  \
count                     544        544         544               544   
unique                      1          2          12                 5   
top                       WNI     Harian       B-1-8            Others   
freq                      544        543          90               209   
mean                      NaN        NaN         NaN               NaN   
std                       NaN        NaN         NaN               NaN   
min                       NaN        NaN         NaN               NaN   
25%                       NaN        NaN         NaN               NaN   
50%                       NaN        NaN         NaN               NaN   
75%                       NaN        NaN         NaN               NaN   
max                       NaN        NaN         NaN               NaN   

       metode_lain tanggal_checkin  ... tanggal_checkout waktu_checkout  \
count          209             544  ...              544            538   
unique           2             311  ...              300            232   
top       Transfer      2026-02-06  ...       2026-02-08       12:00:00   
freq           208               5  ...                5             82   
mean           NaN             NaN  ...              NaN            NaN   
std            NaN             NaN  ...              NaN            NaN   
min            NaN             NaN  ...              NaN            NaN   
25%            NaN             NaN  ...              NaN            NaN   
50%            NaN             NaN  ...              NaN            NaN   
75%            NaN             NaN  ...              NaN            NaN   
max            NaN             NaN  ...              NaN            NaN   

       lama_menginap  komentar user_email  user_pengguna_id  \
count     544.000000       544        0.0        544.000000   
unique           NaN        11        NaN               NaN   
top              NaN       nan        NaN               NaN   
freq             NaN       459        NaN               NaN   
mean        2.641544       NaN        NaN       1271.904412   
std         1.792711       NaN        NaN          1.177622   
min         1.000000       NaN        NaN       1262.000000   
25%         2.000000       NaN        NaN       1271.000000   
50%         2.000000       NaN        NaN       1272.000000   
75%         3.000000       NaN        NaN       1273.000000   
max        31.000000       NaN        NaN       1273.000000   

        editor_pengguna_id           created_at classification       label  
count           538.000000                  544            537  544.000000  
unique                 NaN                  544              2         NaN  
top                    NaN  2026-05-30 15:00:26         normal         NaN  
freq                   NaN                    1            452         NaN  
mean           1271.977695                  NaN            NaN    0.156250  
std               0.923249                  NaN            NaN    0.363426  
min            1262.000000                  NaN            NaN    0.000000  
25%            1271.000000                  NaN            NaN    0.000000  
50%            1272.000000                  NaN            NaN    0.000000  
75%            1273.000000                  NaN            NaN    0.000000  
max            1273.000000                  NaN            NaN    1.000000  

[11 rows x 21 columns]

[2] Distribusi Kelas Target (0: Aman, 1: Melanggar):
label
0    459
1     85
Name: count, dtype: int64

<Figure size 600x400 with 1 Axes>

[3] Rata-rata Pelanggaran Berdasarkan Seluruh Kategori:

--- Tingkat Pelanggaran Berdasarkan status_pasutri ---
status_pasutri
Belum Menikah    0.164
Menikah          0.145
Name: label, dtype: float64

--- Tingkat Pelanggaran Berdasarkan status_kewarganegaraan ---
status_kewarganegaraan
WNI    0.156
Name: label, dtype: float64

--- Tingkat Pelanggaran Berdasarkan jenis_sewa ---
jenis_sewa
Harian     0.157
Bulanan    0.000
Name: label, dtype: float64

--- Tingkat Pelanggaran Berdasarkan metode_pembayaran ---
metode_pembayaran
Kartu Kredit    0.280
Others          0.158
Cash            0.150
QRIS            0.148
Kartu Debit     0.107
Name: label, dtype: float64

[4] Korelasi Antar Fitur:

<Figure size 1000x800 with 2 Axes>

[5] Visualisasi Pola Penyewaan (Dikaitkan dengan Klasifikasi):

<Figure size 1200x500 with 1 Axes>
<Figure size 1200x500 with 1 Axes>
<Figure size 800x500 with 1 Axes>
<Figure size 800x500 with 1 Axes>
<Figure size 800x500 with 1 Axes>
<Figure size 800x500 with 1 Axes>

[5b] Visualisasi Metode Pembayaran 'Others' (metode_lain):

<Figure size 1000x600 with 1 Axes>

=== Proses Feature Selection ===

-> Fitur yang dipilih (3 terbaik): ['lama_menginap', 'checkout_hour', 'checkin_hour']

===== Decision Tree =====
Classification Report:
              precision    recall  f1-score   support

           0     1.0000    0.9565    0.9778        92
           1     0.8095    1.0000    0.8947        17

    accuracy                         0.9633       109
   macro avg     0.9048    0.9783    0.9363       109
weighted avg     0.9703    0.9633    0.9648       109

Confusion Matrix (Text):
                Predicted
               0       1
          -----------------
True   0 |     88       4
       1 |      0      17
          -----------------


===== Random Forest =====
Classification Report:
              precision    recall  f1-score   support

           0     0.9888    0.9565    0.9724        92
           1     0.8000    0.9412    0.8649        17

    accuracy                         0.9541       109
   macro avg     0.8944    0.9488    0.9186       109
weighted avg     0.9593    0.9541    0.9556       109

Confusion Matrix (Text):
                Predicted
               0       1
          -----------------
True   0 |     88       4
       1 |      1      16
          -----------------


===== KNN =====
Classification Report:
              precision    recall  f1-score   support

           0     0.9888    0.9565    0.9724        92
           1     0.8000    0.9412    0.8649        17

    accuracy                         0.9541       109
   macro avg     0.8944    0.9488    0.9186       109
weighted avg     0.9593    0.9541    0.9556       109

Confusion Matrix (Text):
                Predicted
               0       1
          -----------------
True   0 |     88       4
       1 |      1      16
          -----------------


===== SVM RBF =====
Classification Report:
              precision    recall  f1-score   support

           0     1.0000    0.9565    0.9778        92
           1     0.8095    1.0000    0.8947        17

    accuracy                         0.9633       109
   macro avg     0.9048    0.9783    0.9363       109
weighted avg     0.9703    0.9633    0.9648       109

Confusion Matrix (Text):
                Predicted
               0       1
          -----------------
True   0 |     88       4
       1 |      0      17
          -----------------


===== XGBoost =====
Classification Report:
              precision    recall  f1-score   support

           0     1.0000    0.9565    0.9778        92
           1     0.8095    1.0000    0.8947        17

    accuracy                         0.9633       109
   macro avg     0.9048    0.9783    0.9363       109
weighted avg     0.9703    0.9633    0.9648       109

Confusion Matrix (Text):
                Predicted
               0       1
          -----------------
True   0 |     88       4
       1 |      0      17
          -----------------


=== Proses Export Model ===
Sukses! Model dan Label Encoders telah disimpan sebagai file: model_knn.pkl
\end{lstlisting}

\begin{lstlisting}[language=Python, caption={Kode 20}]
# ============================================================
# DATA RIIL + Tunning --> jumlah fitur = 4
# ============================================================

import pandas as pd
import numpy as np
import joblib
import matplotlib.pyplot as plt
import seaborn as sns

from sklearn.model_selection import train_test_split, RandomizedSearchCV
from sklearn.preprocessing import LabelEncoder, StandardScaler, MinMaxScaler
from sklearn.metrics import classification_report, confusion_matrix
from sklearn.pipeline import make_pipeline
from sklearn.feature_selection import SelectKBest, chi2, mutual_info_classif

from sklearn.tree import DecisionTreeClassifier
from sklearn.ensemble import RandomForestClassifier
from sklearn.neighbors import KNeighborsClassifier
from sklearn.svm import SVC
from xgboost import XGBClassifier

# ------------------------------------------------------------
# 1. LOAD DATA
# ------------------------------------------------------------
rentals   = pd.read_csv('rentals.csv')
logs      = pd.read_csv('rental_logs.csv')
# units     = pd.read_csv('units.csv')
# users     = pd.read_csv('users.csv')
# violations= pd.read_csv('violations.csv')

rentals.columns = rentals.columns.str.lower()

if 'komentar' not in rentals.columns:
    raise ValueError("Kolom 'komentar' tidak ditemukan di rentals.csv")

print("renatls.csv")
print(rentals.head())
print("logs.csv")
print(logs.head())


# ------------------------------------------------------------
# 2. LABELING
# ------------------------------------------------------------
rentals['komentar'] = rentals['komentar'].astype(str).str.lower().str.strip()
empty_values = ["", "nan", "na", "null", "-", "0", "tidak ada", "none", " "]
rentals['label'] = (~rentals['komentar'].isin(empty_values)).astype(int)



# ============================================================
# EXPLORATORY DATA ANALYSIS (EDA)
# ============================================================
print("\n" + "="*50)
print("=== EXPLORATORY DATA ANALYSIS (EDA) ===")
print("="*50)

df_eda = rentals.copy()

print("\n[1] Statistik Deskriptif:")
print(df_eda.describe(include='all'))

print("\n[2] Distribusi Kelas Target (0: Aman, 1: Melanggar):")
print(df_eda['label'].value_counts())

fig, ax = plt.subplots(figsize=(6, 4))
sns.countplot(data=df_eda, x='label', palette='Set2', ax=ax)
plt.title('Distribusi Kelas Klasifikasi (0: Aman, 1: Melanggar)')
plt.xlabel('Klasifikasi')
plt.ylabel('Jumlah Penyewaan')

y_max = df_eda['label'].value_counts().max()
plt.ylim(0, y_max * 1.15)

for p in ax.patches:
    if p.get_height() > 0:
        ax.annotate(f'{int(p.get_height())}', (p.get_x() + p.get_width() / 2., p.get_height()),
                    ha='center', va='baseline', fontsize=10, color='black', xytext=(0, 5),
                    textcoords='offset points')
plt.show()

kolom_dikecualikan = ['id', 'nama', 'nik', 'unit_number', 'tanggal_checkin', 'waktu_checkin',
                      'tanggal_checkout', 'waktu_checkout', 'komentar', 'user_email', 'diedit_oleh', 'created_at', 'metode_lain', 'classification']

semua_kategori = [col for col in df_eda.select_dtypes(include=['object', 'category']).columns if col not in kolom_dikecualikan]

print("\n[3] Rata-rata Pelanggaran Berdasarkan Seluruh Kategori:")
for col in semua_kategori:
    if col in df_eda.columns:
        print(f"\n--- Tingkat Pelanggaran Berdasarkan {col} ---")
        print(df_eda.groupby(col)['label'].mean().sort_values(ascending=False).round(3))

df_eda['lama_menginap_num'] = pd.to_numeric(df_eda['lama_menginap'], errors='coerce').fillna(0)
df_eda['checkin_hour_num'] = pd.to_datetime(df_eda['waktu_checkin'], errors='coerce').dt.hour.fillna(0)
df_eda['checkout_hour_num'] = pd.to_datetime(df_eda['waktu_checkout'], errors='coerce').dt.hour.fillna(0)

le_eda = LabelEncoder()
for c in semua_kategori:
    if c in df_eda.columns:
        df_eda[c + '_encoded'] = le_eda.fit_transform(df_eda[c].astype(str))

print("\n[4] Korelasi Antar Fitur:")
kolom_korelasi = ['lama_menginap_num', 'checkin_hour_num', 'checkout_hour_num', 'label']
kolom_korelasi += [c + '_encoded' for c in semua_kategori if c + '_encoded' in df_eda.columns]

corr_matrix = df_eda[kolom_korelasi].corr()

plt.figure(figsize=(10, 8))
sns.heatmap(corr_matrix, annot=True, cmap='coolwarm', fmt='.2f', linewidths=0.5)
plt.title('Heatmap Korelasi Antar Fitur dan Target (Klasifikasi)')
plt.tight_layout()
plt.show()

print("\n[5] Visualisasi Pola Penyewaan (Dikaitkan dengan Klasifikasi):")

custom_line_palette = {0: 'blue', 1: 'red'}
custom_labels = ['0: Aman', '1: Melanggar']

fig, ax1 = plt.subplots(figsize=(12, 5))
waktu_df = df_eda.groupby(['checkin_hour_num', 'label']).size().reset_index(name='jumlah')

sns.lineplot(data=waktu_df, x='checkin_hour_num', y='jumlah', hue='label', marker='o', palette=custom_line_palette, ax=ax1)

plt.title('Pola Waktu Kedatangan Berdasarkan Klasifikasi')
plt.xlabel('Jam Check-in (0-23)')
plt.ylabel('Jumlah Penyewaan')
plt.xticks(range(0, 24))
plt.grid(True, alpha=0.3)

y_max_line = waktu_df['jumlah'].max()
plt.ylim(-y_max_line * 0.20, y_max_line * 1.15)

handles, _ = ax1.get_legend_handles_labels()
plt.legend(handles=handles, labels=custom_labels, title='Klasifikasi')

for l in waktu_df['label'].unique():
    sub = waktu_df[waktu_df['label'] == l]
    warna_teks = custom_line_palette[l]

    for x, y in zip(sub['checkin_hour_num'], sub['jumlah']):
        y_offset = 8 if l == 0 else -12
        va_align = 'bottom' if l == 0 else 'top'
        plt.annotate(f'{int(y)}', (x, y), textcoords="offset points", xytext=(0, y_offset),
                     ha='center', va=va_align, fontsize=9, color=warna_teks, fontweight='bold')
plt.show()

fig, ax2 = plt.subplots(figsize=(12, 5))
lama_df = df_eda.groupby(['lama_menginap_num', 'label']).size().reset_index(name='jumlah')

sns.lineplot(data=lama_df, x='lama_menginap_num', y='jumlah', hue='label', marker='o', palette=custom_line_palette, ax=ax2)

plt.title('Tren Lama Menginap Berdasarkan Klasifikasi')
plt.xlabel('Lama Menginap (Hari)')
plt.ylabel('Jumlah Penyewaan')
plt.grid(True, alpha=0.3)

handles2, _ = ax2.get_legend_handles_labels()
plt.legend(handles=handles2, labels=custom_labels, title='Klasifikasi')
plt.show()

for col in semua_kategori:
    if col in df_eda.columns:
        fig, ax = plt.subplots(figsize=(8, 5))
        sns.countplot(data=df_eda, x=col, hue='label', palette='magma', ax=ax)

        judul_kolom = col.replace('_', ' ').title()
        plt.title(f'Perbandingan Jumlah Pelanggaran Berdasarkan {judul_kolom}')
        plt.xlabel(judul_kolom)
        plt.ylabel('Jumlah Penyewaan')
        plt.xticks(rotation=0)

        handles_bar, _ = ax.get_legend_handles_labels()
        plt.legend(handles=handles_bar, labels=custom_labels, title='Klasifikasi')

        y_max_bar = df_eda.groupby([col, 'label']).size().max()
        plt.ylim(0, y_max_bar * 1.15)

        for p in ax.patches:
            if p.get_height() > 0:
                ax.annotate(f'{int(p.get_height())}', (p.get_x() + p.get_width() / 2., p.get_height()),
                            ha='center', va='baseline', fontsize=9, color='black', xytext=(0, 5),
                            textcoords='offset points')
        plt.tight_layout()
        plt.show()


print("\n[5b] Visualisasi Metode Pembayaran 'Others' (metode_lain):")

# Memastikan kolom 'metode_lain' berbentuk string dan membersihkan spasi/nilai kosong
df_eda['metode_lain'] = df_eda['metode_lain'].astype(str).str.strip().str.title()
empty_lain = ["", "Nan", "Na", "Null", "-", "0", "None"]

# Filter hanya baris yang memiliki nilai di 'metode_lain'
df_others = df_eda[~df_eda['metode_lain'].isin(empty_lain)]

if not df_others.empty:
    fig, ax = plt.subplots(figsize=(10, 6))

    # Menghitung frekuensi masing-masing metode pembayaran lain
    order_others = df_others['metode_lain'].value_counts().index

    sns.countplot(
        data=df_others,
        y='metode_lain',
        order=order_others,
        palette='viridis',
        ax=ax
    )

    plt.title('Detail Metode Pembayaran "Others" (metode_lain)')
    plt.xlabel('Jumlah Penyewaan')
    plt.ylabel('Jenis Pembayaran Lainnya')

    # Menambahkan anotasi angka di sebelah setiap bar
    for p in ax.patches:
        width = p.get_width()
        if width > 0:
            ax.annotate(f'{int(width)}',
                        (width, p.get_y() + p.get_height() / 2.),
                        ha='left', va='center', fontsize=10, color='black',
                        xytext=(5, 0), textcoords='offset points')

    # Perlebar batas x-axis sedikit agar angka anotasi tidak terpotong
    x_max = df_others['metode_lain'].value_counts().max()
    plt.xlim(0, x_max * 1.15)

    plt.tight_layout()
    plt.show()
else:
    print("Tidak ada data metode pembayaran 'Others' yang valid untuk ditampilkan.")


# ------------------------------------------------------------
# 3. FEATURE ENGINEERING & PREPARATION
# ------------------------------------------------------------
# Parse waktu checkout
rentals['waktu_checkout_dt'] = pd.to_datetime(rentals['waktu_checkout'], errors='coerce')
rentals['checkout_hour'] = rentals['waktu_checkout_dt'].dt.hour.fillna(0).astype(int)

# Parse waktu checkin
rentals['waktu_checkin_dt'] = pd.to_datetime(rentals['waktu_checkin'], errors='coerce')
rentals['checkin_hour'] = rentals['waktu_checkin_dt'].dt.hour.fillna(0).astype(int)

# Pastikan lama_menginap berupa angka
rentals['lama_menginap'] = pd.to_numeric(rentals['lama_menginap'], errors='coerce').fillna(0)

# Isi data kosong untuk kolom kategorikal sebelum split
cat_cols = ['jenis_sewa', 'metode_pembayaran', 'status_pasutri', 'status_kewarganegaraan']
for c in cat_cols:
    rentals[c] = rentals[c].fillna('Unknown').astype(str)

numeric_cols = ['checkout_hour', 'checkin_hour', 'lama_menginap']
feature_cols = cat_cols + numeric_cols

X = rentals[feature_cols].copy()
y = rentals['label']

# ------------------------------------------------------------
# 4. TRAIN-TEST SPLIT
# ------------------------------------------------------------
X_train, X_test, y_train, y_test = train_test_split(
    X, y, stratify=y, test_size=0.2, random_state=42
)

# ------------------------------------------------------------
# 4B. ENCODING MENCEGAH DATA LEAKAGE
# ------------------------------------------------------------
# Simpan semua LabelEncoder ke dalam dictionary agar bisa di-export
label_encoders = {}

for c in cat_cols:
    le = LabelEncoder()
    # Fit & Transform pada data latih
    X_train[c] = le.fit_transform(X_train[c])

    # Tangani unseen labels pada data uji (ubah jadi 'Unknown' jika tidak pernah dilihat di X_train untuk mencegah error saat deployment
    X_test[c] = X_test[c].map(lambda s: s if s in le.classes_ else 'Unknown')
    if 'Unknown' not in le.classes_:
        le.classes_ = np.append(le.classes_, 'Unknown')

    X_test[c] = le.transform(X_test[c])

    # Simpan object le
    label_encoders[c] = le

# ------------------------------------------------------------
# 5. FEATURE SELECTION
# ------------------------------------------------------------
print("\n=== Proses Feature Selection ===")
chi2_selector = SelectKBest(score_func=chi2, k='all')
chi2_selector.fit(X_train, y_train)

train_df = X_train.copy()
train_df['label'] = y_train
pearson_scores = train_df.corr(method='pearson')['label'].drop('label').abs().values

mi_selector = SelectKBest(score_func=mutual_info_classif, k='all')
mi_selector.fit(X_train, y_train)

fs_df = pd.DataFrame({
    'Feature': X_train.columns,
    'Chi-Square': chi2_selector.scores_,
    'Pearson (Abs)': pearson_scores,
    'Mutual Info': mi_selector.scores_
})

scaler = MinMaxScaler()
fs_df[['Chi2_Norm', 'Pearson_Norm', 'MI_Norm']] = scaler.fit_transform(
    fs_df[['Chi-Square', 'Pearson (Abs)', 'Mutual Info']]
)
fs_df['Total_Score'] = fs_df[['Chi2_Norm', 'Pearson_Norm', 'MI_Norm']].mean(axis=1)
fs_df = fs_df.sort_values(by='Total_Score', ascending=False).reset_index(drop=True)

# Pilih Jumlah Fitur
n_select = 4
selected_features = fs_df['Feature'].head(n_select).tolist()
print(f"\n-> Fitur yang dipilih ({n_select} terbaik): {selected_features}")

X_train = X_train[selected_features]
X_test = X_test[selected_features]

# ------------------------------------------------------------
# 6. Fungsi Evaluasi
# ------------------------------------------------------------
def evaluate(name, model):
    y_pred = model.predict(X_test)
    print(f"\n===== {name} =====")
    print("Classification Report:")
    print(classification_report(y_test, y_pred, digits=4))
    cm = confusion_matrix(y_test, y_pred)
    print("Confusion Matrix (Text):")
    print("                Predicted")
    print("               0       1")
    print("          -----------------")
    print(f"True   0 | {cm[0][0]:6d}  {cm[0][1]:6d}")
    print(f"       1 | {cm[1][0]:6d}  {cm[1][1]:6d}")
    print("          -----------------\n")

# ------------------------------------------------------------
# 7. Latih Semua Model & Hyperparameter Tuning
# ------------------------------------------------------------
models = {}
custom_weight = {0: 1, 1: 4}

# --- A. Definisikan Model Dasar ---
dt_base = DecisionTreeClassifier(class_weight=custom_weight, random_state=42)
rf_base = RandomForestClassifier(class_weight=custom_weight, random_state=42)
knn_base = make_pipeline(StandardScaler(), KNeighborsClassifier())
svm_base = make_pipeline(StandardScaler(), SVC(kernel='rbf', probability=True, class_weight=custom_weight, random_state=42))
xgb_base = XGBClassifier(scale_pos_weight=4, eval_metric='logloss', random_state=42)

print("\nMemulai proses Hyperparameter Tuning untuk semua model. Ini mungkin memakan waktu...\n")

# --- B. Proses Tuning Setiap Model ---

# 1. Decision Tree
print(">> Tuning Decision Tree...")
dt_param_grid = {'max_depth': [3, 5, 7, 10, None], 'min_samples_split': [2, 5, 10]}
dt_random = RandomizedSearchCV(dt_base, dt_param_grid, n_iter=10, cv=3, random_state=42, n_jobs=-1, scoring='f1_macro')
dt_random.fit(X_train, y_train)
models["Decision Tree (Tuned)"] = dt_random.best_estimator_

# 2. Random Forest
print(">> Tuning Random Forest...")
rf_param_grid = {'n_estimators': [100, 200, 300], 'max_depth': [5, 7, 10, None], 'min_samples_split': [2, 5, 10], 'min_samples_leaf': [1, 2, 4]}
rf_random = RandomizedSearchCV(rf_base, rf_param_grid, n_iter=10, cv=3, random_state=42, n_jobs=-1, scoring='f1_macro')
rf_random.fit(X_train, y_train)
models["Random Forest (Tuned)"] = rf_random.best_estimator_

# 3. KNN (Karena menggunakan pipeline, parameter ditambahkan prefix 'kneighborsclassifier__')
print(">> Tuning KNN...")
knn_param_grid = {'kneighborsclassifier__n_neighbors': [3, 5, 7, 9, 11], 'kneighborsclassifier__weights': ['uniform', 'distance']}
knn_random = RandomizedSearchCV(knn_base, knn_param_grid, n_iter=10, cv=3, random_state=42, n_jobs=-1, scoring='f1_macro')
knn_random.fit(X_train, y_train)
models["KNN (Tuned)"] = knn_random.best_estimator_

# 4. SVM (Karena menggunakan pipeline, parameter ditambahkan prefix 'svc__')
print(">> Tuning SVM...")
svm_param_grid = {'svc__C': [0.1, 1, 10], 'svc__gamma': ['scale', 'auto', 0.1, 0.01]}
svm_random = RandomizedSearchCV(svm_base, svm_param_grid, n_iter=10, cv=3, random_state=42, n_jobs=-1, scoring='f1_macro')
svm_random.fit(X_train, y_train)
models["SVM RBF (Tuned)"] = svm_random.best_estimator_

# 5. XGBoost
print(">> Tuning XGBoost...")
xgb_param_grid = {'learning_rate': [0.01, 0.05, 0.1, 0.2], 'max_depth': [3, 5, 7], 'n_estimators': [100, 200, 300]}
xgb_random = RandomizedSearchCV(xgb_base, xgb_param_grid, n_iter=10, cv=3, random_state=42, n_jobs=-1, scoring='f1_macro')
xgb_random.fit(X_train, y_train)
models["XGBoost (Tuned)"] = xgb_random.best_estimator_

# --- C. Evaluasi Semua Model yang Sudah Di-Tune ---
# Karena setiap .best_estimator_ sudah di-fit oleh RandomizedSearchCV pada seluruh X_train,
# kita bisa langsung memanggil fungsi evaluate tanpa perlu melakukan model.fit() lagi.
for name, model in models.items():
    evaluate(name, model)

# ------------------------------------------------------------
# 8. Eksport Model Terbaik & Label Encoders
# ------------------------------------------------------------
print("\n=== Proses Export Model ===")

# Gabungkan model dan encoders ke dalam satu dictionary agar rapi
export_package = {
    'model': models["KNN (Tuned)"], # Mengubah dari Random Forest ke KNN
    'label_encoders': label_encoders,
    'selected_features': selected_features # Penting agar tahu urutan kolom masuk saat predict
}

# Menentukan nama file yang sesuai untuk KNN
filename = 'model_knn.pkl'

# Menyimpan package ke dalam file .pkl
joblib.dump(export_package, filename)
print(f"Sukses! Model dan Label Encoders telah disimpan sebagai file: {filename}")
\end{lstlisting}

\begin{lstlisting}[caption={Output 20}]
renatls.csv
   id                   nama               nik status_pasutri  \
0   3  N*** Y**** S**** D***  3***************  Belum Menikah   
1   4     A***** N** H******  3***************  Belum Menikah   
2   5     A***** A*** S*****  3***************        Menikah   
3   6            C**** W****  1***************  Belum Menikah   
4   7             S*** U****  3***************  Belum Menikah   

  status_kewarganegaraan jenis_sewa unit_number metode_pembayaran metode_lain  \
0                    WNI     Harian     B-20-27              Cash         NaN   
1                    WNI     Harian     A-15-25              Cash         NaN   
2                    WNI     Harian     A-06-24            Others    Aplikasi   
3                    WNI     Harian     D-15-21              QRIS         NaN   
4                    WNI     Harian       A-8-8            Others    Transfer   

  tanggal_checkin waktu_checkin tanggal_checkout waktu_checkout  \
0      2026-05-20      13:08:00       2026-05-21            NaN   
1      2026-05-19      13:25:00       2026-05-22            NaN   
2      2026-05-17      13:27:00       2026-05-22            NaN   
3      2026-05-19      13:29:00       2026-05-24            NaN   
4      2025-01-01      19:58:00       2025-01-06       11:22:00   

   lama_menginap komentar  user_email  user_pengguna_id  editor_pengguna_id  \
0              1      NaN         NaN              1267                 NaN   
1              3      NaN         NaN              1267                 NaN   
2              5      NaN         NaN              1267                 NaN   
3              5      NaN         NaN              1267                 NaN   
4              5      NaN         NaN              1271              1271.0   

            created_at classification  
0  2026-05-20 13:09:49            NaN  
1  2026-05-20 13:26:04            NaN  
2  2026-05-20 13:28:03            NaN  
3  2026-05-20 13:30:18            NaN  
4  2026-05-25 10:06:33         normal  
logs.csv
     id  rental_id  action field_changed  old_value  \
0     6          3  create           all        NaN   
1     7          4  create           all        NaN   
2     8          5  create           all        NaN   
3     9          6  create           all        NaN   
4  4906          7  create           all        NaN   

                                           new_value  email  pengguna_id  \
0  {"nama":"N*** Y**** S**** D***","nik":"3*************************************"nama":"A***** N** H******","nik":"3****************************************"nama":"A***** A*** S*****","nik":"3****************************************"nama":"C**** W****","nik":"1***************"...    NaN         1267   
4  {"nama":"S*** U****","nik":"3***************",...    NaN         1271   

             timestamp  
0  2026-05-20 13:09:49  
1  2026-05-20 13:26:04  
2  2026-05-20 13:28:03  
3  2026-05-20 13:30:18  
4  2026-05-25 10:06:33  

==================================================
=== EXPLORATORY DATA ANALYSIS (EDA) ===
==================================================

[1] Statistik Deskriptif:
                id                         nama           nik status_pasutri  \
count   544.000000                          544  5.440000e+02            544   
unique         NaN                          207           NaN              2   
top            NaN  Y***** M**** S**** S*******           NaN  Belum Menikah   
freq           NaN                            8           NaN            330   
mean    274.538603                          NaN  3.214330e+15            NaN   
std     157.250941                          NaN  1.204147e+15            NaN   
min       3.000000                          NaN  1.202052e+15            NaN   
25%     138.750000                          NaN  3.173074e+15            NaN   
50%     274.500000                          NaN  3.209121e+15            NaN   
75%     410.250000                          NaN  3.276064e+15            NaN   
max     553.000000                          NaN  9.104017e+15            NaN   

       status_kewarganegaraan jenis_sewa unit_number metode_pembayaran  \
count                     544        544         544               544   
unique                      1          2          12                 5   
top                       WNI     Harian       B-1-8            Others   
freq                      544        543          90               209   
mean                      NaN        NaN         NaN               NaN   
std                       NaN        NaN         NaN               NaN   
min                       NaN        NaN         NaN               NaN   
25%                       NaN        NaN         NaN               NaN   
50%                       NaN        NaN         NaN               NaN   
75%                       NaN        NaN         NaN               NaN   
max                       NaN        NaN         NaN               NaN   

       metode_lain tanggal_checkin  ... tanggal_checkout waktu_checkout  \
count          209             544  ...              544            538   
unique           2             311  ...              300            232   
top       Transfer      2026-02-06  ...       2026-02-08       12:00:00   
freq           208               5  ...                5             82   
mean           NaN             NaN  ...              NaN            NaN   
std            NaN             NaN  ...              NaN            NaN   
min            NaN             NaN  ...              NaN            NaN   
25%            NaN             NaN  ...              NaN            NaN   
50%            NaN             NaN  ...              NaN            NaN   
75%            NaN             NaN  ...              NaN            NaN   
max            NaN             NaN  ...              NaN            NaN   

       lama_menginap  komentar user_email  user_pengguna_id  \
count     544.000000       544        0.0        544.000000   
unique           NaN        11        NaN               NaN   
top              NaN       nan        NaN               NaN   
freq             NaN       459        NaN               NaN   
mean        2.641544       NaN        NaN       1271.904412   
std         1.792711       NaN        NaN          1.177622   
min         1.000000       NaN        NaN       1262.000000   
25%         2.000000       NaN        NaN       1271.000000   
50%         2.000000       NaN        NaN       1272.000000   
75%         3.000000       NaN        NaN       1273.000000   
max        31.000000       NaN        NaN       1273.000000   

        editor_pengguna_id           created_at classification       label  
count           538.000000                  544            537  544.000000  
unique                 NaN                  544              2         NaN  
top                    NaN  2026-05-30 15:00:26         normal         NaN  
freq                   NaN                    1            452         NaN  
mean           1271.977695                  NaN            NaN    0.156250  
std               0.923249                  NaN            NaN    0.363426  
min            1262.000000                  NaN            NaN    0.000000  
25%            1271.000000                  NaN            NaN    0.000000  
50%            1272.000000                  NaN            NaN    0.000000  
75%            1273.000000                  NaN            NaN    0.000000  
max            1273.000000                  NaN            NaN    1.000000  

[11 rows x 21 columns]

[2] Distribusi Kelas Target (0: Aman, 1: Melanggar):
label
0    459
1     85
Name: count, dtype: int64

/tmp/ipykernel_4832/252262949.py:68: FutureWarning: 

Passing `palette` without assigning `hue` is deprecated and will be removed in v0.14.0. Assign the `x` variable to `hue` and set `legend=False` for the same effect.

  sns.countplot(data=df_eda, x='label', palette='Set2', ax=ax)

<Figure size 600x400 with 1 Axes>

[3] Rata-rata Pelanggaran Berdasarkan Seluruh Kategori:

--- Tingkat Pelanggaran Berdasarkan status_pasutri ---
status_pasutri
Belum Menikah    0.164
Menikah          0.145
Name: label, dtype: float64

--- Tingkat Pelanggaran Berdasarkan status_kewarganegaraan ---
status_kewarganegaraan
WNI    0.156
Name: label, dtype: float64

--- Tingkat Pelanggaran Berdasarkan jenis_sewa ---
jenis_sewa
Harian     0.157
Bulanan    0.000
Name: label, dtype: float64

--- Tingkat Pelanggaran Berdasarkan metode_pembayaran ---
metode_pembayaran
Kartu Kredit    0.280
Others          0.158
Cash            0.150
QRIS            0.148
Kartu Debit     0.107
Name: label, dtype: float64

[4] Korelasi Antar Fitur:

/tmp/ipykernel_4832/252262949.py:95: UserWarning: Could not infer format, so each element will be parsed individually, falling back to `dateutil`. To ensure parsing is consistent and as-expected, please specify a format.
  df_eda['checkin_hour_num'] = pd.to_datetime(df_eda['waktu_checkin'], errors='coerce').dt.hour.fillna(0)
/tmp/ipykernel_4832/252262949.py:96: UserWarning: Could not infer format, so each element will be parsed individually, falling back to `dateutil`. To ensure parsing is consistent and as-expected, please specify a format.
  df_eda['checkout_hour_num'] = pd.to_datetime(df_eda['waktu_checkout'], errors='coerce').dt.hour.fillna(0)

<Figure size 1000x800 with 2 Axes>

[5] Visualisasi Pola Penyewaan (Dikaitkan dengan Klasifikasi):

<Figure size 1200x500 with 1 Axes>
<Figure size 1200x500 with 1 Axes>
<Figure size 800x500 with 1 Axes>
<Figure size 800x500 with 1 Axes>
<Figure size 800x500 with 1 Axes>
<Figure size 800x500 with 1 Axes>

[5b] Visualisasi Metode Pembayaran 'Others' (metode_lain):

/tmp/ipykernel_4832/252262949.py:203: FutureWarning: 

Passing `palette` without assigning `hue` is deprecated and will be removed in v0.14.0. Assign the `y` variable to `hue` and set `legend=False` for the same effect.

  sns.countplot(

<Figure size 1000x600 with 1 Axes>
/tmp/ipykernel_4832/252262949.py:238: UserWarning: Could not infer format, so each element will be parsed individually, falling back to `dateutil`. To ensure parsing is consistent and as-expected, please specify a format.
  rentals['waktu_checkout_dt'] = pd.to_datetime(rentals['waktu_checkout'], errors='coerce')
/tmp/ipykernel_4832/252262949.py:242: UserWarning: Could not infer format, so each element will be parsed individually, falling back to `dateutil`. To ensure parsing is consistent and as-expected, please specify a format.
  rentals['waktu_checkin_dt'] = pd.to_datetime(rentals['waktu_checkin'], errors='coerce')


=== Proses Feature Selection ===

-> Fitur yang dipilih (4 terbaik): ['lama_menginap', 'checkout_hour', 'checkin_hour', 'metode_pembayaran']

Memulai proses Hyperparameter Tuning untuk semua model. Ini mungkin memakan waktu...

>> Tuning Decision Tree...
>> Tuning Random Forest...
>> Tuning KNN...
>> Tuning SVM...
>> Tuning XGBoost...

===== Decision Tree (Tuned) =====
Classification Report:
              precision    recall  f1-score   support

           0     1.0000    0.9565    0.9778        92
           1     0.8095    1.0000    0.8947        17

    accuracy                         0.9633       109
   macro avg     0.9048    0.9783    0.9363       109
weighted avg     0.9703    0.9633    0.9648       109

Confusion Matrix (Text):
                Predicted
               0       1
          -----------------
True   0 |     88       4
       1 |      0      17
          -----------------


===== Random Forest (Tuned) =====
Classification Report:
              precision    recall  f1-score   support

           0     1.0000    0.9565    0.9778        92
           1     0.8095    1.0000    0.8947        17

    accuracy                         0.9633       109
   macro avg     0.9048    0.9783    0.9363       109
weighted avg     0.9703    0.9633    0.9648       109

Confusion Matrix (Text):
                Predicted
               0       1
          -----------------
True   0 |     88       4
       1 |      0      17
          -----------------


===== KNN (Tuned) =====
Classification Report:
              precision    recall  f1-score   support

           0     0.9565    0.9565    0.9565        92
           1     0.7647    0.7647    0.7647        17

    accuracy                         0.9266       109
   macro avg     0.8606    0.8606    0.8606       109
weighted avg     0.9266    0.9266    0.9266       109

Confusion Matrix (Text):
                Predicted
               0       1
          -----------------
True   0 |     88       4
       1 |      4      13
          -----------------


===== SVM RBF (Tuned) =====
Classification Report:
              precision    recall  f1-score   support

           0     1.0000    0.9565    0.9778        92
           1     0.8095    1.0000    0.8947        17

    accuracy                         0.9633       109
   macro avg     0.9048    0.9783    0.9363       109
weighted avg     0.9703    0.9633    0.9648       109

Confusion Matrix (Text):
                Predicted
               0       1
          -----------------
True   0 |     88       4
       1 |      0      17
          -----------------


===== XGBoost (Tuned) =====
Classification Report:
              precision    recall  f1-score   support

           0     1.0000    0.9565    0.9778        92
           1     0.8095    1.0000    0.8947        17

    accuracy                         0.9633       109
   macro avg     0.9048    0.9783    0.9363       109
weighted avg     0.9703    0.9633    0.9648       109

Confusion Matrix (Text):
                Predicted
               0       1
          -----------------
True   0 |     88       4
       1 |      0      17
          -----------------


=== Proses Export Model ===
Sukses! Model dan Label Encoders telah disimpan sebagai file: model_knn.pkl
\end{lstlisting}

\begin{lstlisting}[language=Python, caption={Kode 21}]
# ============================================================
# DATA RIIL + Tunning --> jumlah fitur = 3
# ============================================================

import pandas as pd
import numpy as np
import joblib
import matplotlib.pyplot as plt
import seaborn as sns

from sklearn.model_selection import train_test_split, RandomizedSearchCV
from sklearn.preprocessing import LabelEncoder, StandardScaler, MinMaxScaler
from sklearn.metrics import classification_report, confusion_matrix
from sklearn.pipeline import make_pipeline
from sklearn.feature_selection import SelectKBest, chi2, mutual_info_classif

from sklearn.tree import DecisionTreeClassifier
from sklearn.ensemble import RandomForestClassifier
from sklearn.neighbors import KNeighborsClassifier
from sklearn.svm import SVC
from xgboost import XGBClassifier

# ------------------------------------------------------------
# 1. LOAD DATA
# ------------------------------------------------------------
rentals   = pd.read_csv('rentals.csv')
logs      = pd.read_csv('rental_logs.csv')
# units     = pd.read_csv('units.csv')
# users     = pd.read_csv('users.csv')
# violations= pd.read_csv('violations.csv')

rentals.columns = rentals.columns.str.lower()

if 'komentar' not in rentals.columns:
    raise ValueError("Kolom 'komentar' tidak ditemukan di rentals.csv")

print("renatls.csv")
print(rentals.head())
print("logs.csv")
print(logs.head())


# ------------------------------------------------------------
# 2. LABELING
# ------------------------------------------------------------
rentals['komentar'] = rentals['komentar'].astype(str).str.lower().str.strip()
empty_values = ["", "nan", "na", "null", "-", "0", "tidak ada", "none", " "]
rentals['label'] = (~rentals['komentar'].isin(empty_values)).astype(int)



# ============================================================
# EXPLORATORY DATA ANALYSIS (EDA)
# ============================================================
print("\n" + "="*50)
print("=== EXPLORATORY DATA ANALYSIS (EDA) ===")
print("="*50)

df_eda = rentals.copy()

print("\n[1] Statistik Deskriptif:")
print(df_eda.describe(include='all'))

print("\n[2] Distribusi Kelas Target (0: Aman, 1: Melanggar):")
print(df_eda['label'].value_counts())

fig, ax = plt.subplots(figsize=(6, 4))
sns.countplot(data=df_eda, x='label', palette='Set2', ax=ax)
plt.title('Distribusi Kelas Klasifikasi (0: Aman, 1: Melanggar)')
plt.xlabel('Klasifikasi')
plt.ylabel('Jumlah Penyewaan')

y_max = df_eda['label'].value_counts().max()
plt.ylim(0, y_max * 1.15)

for p in ax.patches:
    if p.get_height() > 0:
        ax.annotate(f'{int(p.get_height())}', (p.get_x() + p.get_width() / 2., p.get_height()),
                    ha='center', va='baseline', fontsize=10, color='black', xytext=(0, 5),
                    textcoords='offset points')
plt.show()

kolom_dikecualikan = ['id', 'nama', 'nik', 'unit_number', 'tanggal_checkin', 'waktu_checkin',
                      'tanggal_checkout', 'waktu_checkout', 'komentar', 'user_email', 'diedit_oleh', 'created_at', 'metode_lain', 'classification']

semua_kategori = [col for col in df_eda.select_dtypes(include=['object', 'category']).columns if col not in kolom_dikecualikan]

print("\n[3] Rata-rata Pelanggaran Berdasarkan Seluruh Kategori:")
for col in semua_kategori:
    if col in df_eda.columns:
        print(f"\n--- Tingkat Pelanggaran Berdasarkan {col} ---")
        print(df_eda.groupby(col)['label'].mean().sort_values(ascending=False).round(3))

df_eda['lama_menginap_num'] = pd.to_numeric(df_eda['lama_menginap'], errors='coerce').fillna(0)
df_eda['checkin_hour_num'] = pd.to_datetime(df_eda['waktu_checkin'], errors='coerce').dt.hour.fillna(0)
df_eda['checkout_hour_num'] = pd.to_datetime(df_eda['waktu_checkout'], errors='coerce').dt.hour.fillna(0)

le_eda = LabelEncoder()
for c in semua_kategori:
    if c in df_eda.columns:
        df_eda[c + '_encoded'] = le_eda.fit_transform(df_eda[c].astype(str))

print("\n[4] Korelasi Antar Fitur:")
kolom_korelasi = ['lama_menginap_num', 'checkin_hour_num', 'checkout_hour_num', 'label']
kolom_korelasi += [c + '_encoded' for c in semua_kategori if c + '_encoded' in df_eda.columns]

corr_matrix = df_eda[kolom_korelasi].corr()

plt.figure(figsize=(10, 8))
sns.heatmap(corr_matrix, annot=True, cmap='coolwarm', fmt='.2f', linewidths=0.5)
plt.title('Heatmap Korelasi Antar Fitur dan Target (Klasifikasi)')
plt.tight_layout()
plt.show()

print("\n[5] Visualisasi Pola Penyewaan (Dikaitkan dengan Klasifikasi):")

custom_line_palette = {0: 'blue', 1: 'red'}
custom_labels = ['0: Aman', '1: Melanggar']

fig, ax1 = plt.subplots(figsize=(12, 5))
waktu_df = df_eda.groupby(['checkin_hour_num', 'label']).size().reset_index(name='jumlah')

sns.lineplot(data=waktu_df, x='checkin_hour_num', y='jumlah', hue='label', marker='o', palette=custom_line_palette, ax=ax1)

plt.title('Pola Waktu Kedatangan Berdasarkan Klasifikasi')
plt.xlabel('Jam Check-in (0-23)')
plt.ylabel('Jumlah Penyewaan')
plt.xticks(range(0, 24))
plt.grid(True, alpha=0.3)

y_max_line = waktu_df['jumlah'].max()
plt.ylim(-y_max_line * 0.20, y_max_line * 1.15)

handles, _ = ax1.get_legend_handles_labels()
plt.legend(handles=handles, labels=custom_labels, title='Klasifikasi')

for l in waktu_df['label'].unique():
    sub = waktu_df[waktu_df['label'] == l]
    warna_teks = custom_line_palette[l]

    for x, y in zip(sub['checkin_hour_num'], sub['jumlah']):
        y_offset = 8 if l == 0 else -12
        va_align = 'bottom' if l == 0 else 'top'
        plt.annotate(f'{int(y)}', (x, y), textcoords="offset points", xytext=(0, y_offset),
                     ha='center', va=va_align, fontsize=9, color=warna_teks, fontweight='bold')
plt.show()

fig, ax2 = plt.subplots(figsize=(12, 5))
lama_df = df_eda.groupby(['lama_menginap_num', 'label']).size().reset_index(name='jumlah')

sns.lineplot(data=lama_df, x='lama_menginap_num', y='jumlah', hue='label', marker='o', palette=custom_line_palette, ax=ax2)

plt.title('Tren Lama Menginap Berdasarkan Klasifikasi')
plt.xlabel('Lama Menginap (Hari)')
plt.ylabel('Jumlah Penyewaan')
plt.grid(True, alpha=0.3)

handles2, _ = ax2.get_legend_handles_labels()
plt.legend(handles=handles2, labels=custom_labels, title='Klasifikasi')
plt.show()

for col in semua_kategori:
    if col in df_eda.columns:
        fig, ax = plt.subplots(figsize=(8, 5))
        sns.countplot(data=df_eda, x=col, hue='label', palette='magma', ax=ax)

        judul_kolom = col.replace('_', ' ').title()
        plt.title(f'Perbandingan Jumlah Pelanggaran Berdasarkan {judul_kolom}')
        plt.xlabel(judul_kolom)
        plt.ylabel('Jumlah Penyewaan')
        plt.xticks(rotation=0)

        handles_bar, _ = ax.get_legend_handles_labels()
        plt.legend(handles=handles_bar, labels=custom_labels, title='Klasifikasi')

        y_max_bar = df_eda.groupby([col, 'label']).size().max()
        plt.ylim(0, y_max_bar * 1.15)

        for p in ax.patches:
            if p.get_height() > 0:
                ax.annotate(f'{int(p.get_height())}', (p.get_x() + p.get_width() / 2., p.get_height()),
                            ha='center', va='baseline', fontsize=9, color='black', xytext=(0, 5),
                            textcoords='offset points')
        plt.tight_layout()
        plt.show()


print("\n[5b] Visualisasi Metode Pembayaran 'Others' (metode_lain):")

# Memastikan kolom 'metode_lain' berbentuk string dan membersihkan spasi/nilai kosong
df_eda['metode_lain'] = df_eda['metode_lain'].astype(str).str.strip().str.title()
empty_lain = ["", "Nan", "Na", "Null", "-", "0", "None"]

# Filter hanya baris yang memiliki nilai di 'metode_lain'
df_others = df_eda[~df_eda['metode_lain'].isin(empty_lain)]

if not df_others.empty:
    fig, ax = plt.subplots(figsize=(10, 6))

    # Menghitung frekuensi masing-masing metode pembayaran lain
    order_others = df_others['metode_lain'].value_counts().index

    sns.countplot(
        data=df_others,
        y='metode_lain',
        order=order_others,
        palette='viridis',
        ax=ax
    )

    plt.title('Detail Metode Pembayaran "Others" (metode_lain)')
    plt.xlabel('Jumlah Penyewaan')
    plt.ylabel('Jenis Pembayaran Lainnya')

    # Menambahkan anotasi angka di sebelah setiap bar
    for p in ax.patches:
        width = p.get_width()
        if width > 0:
            ax.annotate(f'{int(width)}',
                        (width, p.get_y() + p.get_height() / 2.),
                        ha='left', va='center', fontsize=10, color='black',
                        xytext=(5, 0), textcoords='offset points')

    # Perlebar batas x-axis sedikit agar angka anotasi tidak terpotong
    x_max = df_others['metode_lain'].value_counts().max()
    plt.xlim(0, x_max * 1.15)

    plt.tight_layout()
    plt.show()
else:
    print("Tidak ada data metode pembayaran 'Others' yang valid untuk ditampilkan.")


# ------------------------------------------------------------
# 3. FEATURE ENGINEERING & PREPARATION
# ------------------------------------------------------------
# Parse waktu checkout
rentals['waktu_checkout_dt'] = pd.to_datetime(rentals['waktu_checkout'], errors='coerce')
rentals['checkout_hour'] = rentals['waktu_checkout_dt'].dt.hour.fillna(0).astype(int)

# Parse waktu checkin
rentals['waktu_checkin_dt'] = pd.to_datetime(rentals['waktu_checkin'], errors='coerce')
rentals['checkin_hour'] = rentals['waktu_checkin_dt'].dt.hour.fillna(0).astype(int)

# Pastikan lama_menginap berupa angka
rentals['lama_menginap'] = pd.to_numeric(rentals['lama_menginap'], errors='coerce').fillna(0)

# Isi data kosong untuk kolom kategorikal sebelum split
cat_cols = ['jenis_sewa', 'metode_pembayaran', 'status_pasutri', 'status_kewarganegaraan']
for c in cat_cols:
    rentals[c] = rentals[c].fillna('Unknown').astype(str)

numeric_cols = ['checkout_hour', 'checkin_hour', 'lama_menginap']
feature_cols = cat_cols + numeric_cols

X = rentals[feature_cols].copy()
y = rentals['label']

# ------------------------------------------------------------
# 4. TRAIN-TEST SPLIT
# ------------------------------------------------------------
X_train, X_test, y_train, y_test = train_test_split(
    X, y, stratify=y, test_size=0.2, random_state=42
)

# ------------------------------------------------------------
# 4B. ENCODING MENCEGAH DATA LEAKAGE
# ------------------------------------------------------------
# Simpan semua LabelEncoder ke dalam dictionary agar bisa di-export
label_encoders = {}

for c in cat_cols:
    le = LabelEncoder()
    # Fit & Transform pada data latih
    X_train[c] = le.fit_transform(X_train[c])

    # Tangani unseen labels pada data uji (ubah jadi 'Unknown' jika tidak pernah dilihat di X_train untuk mencegah error saat deployment
    X_test[c] = X_test[c].map(lambda s: s if s in le.classes_ else 'Unknown')
    if 'Unknown' not in le.classes_:
        le.classes_ = np.append(le.classes_, 'Unknown')

    X_test[c] = le.transform(X_test[c])

    # Simpan object le
    label_encoders[c] = le

# ------------------------------------------------------------
# 5. FEATURE SELECTION
# ------------------------------------------------------------
print("\n=== Proses Feature Selection ===")
chi2_selector = SelectKBest(score_func=chi2, k='all')
chi2_selector.fit(X_train, y_train)

train_df = X_train.copy()
train_df['label'] = y_train
pearson_scores = train_df.corr(method='pearson')['label'].drop('label').abs().values

mi_selector = SelectKBest(score_func=mutual_info_classif, k='all')
mi_selector.fit(X_train, y_train)

fs_df = pd.DataFrame({
    'Feature': X_train.columns,
    'Chi-Square': chi2_selector.scores_,
    'Pearson (Abs)': pearson_scores,
    'Mutual Info': mi_selector.scores_
})

scaler = MinMaxScaler()
fs_df[['Chi2_Norm', 'Pearson_Norm', 'MI_Norm']] = scaler.fit_transform(
    fs_df[['Chi-Square', 'Pearson (Abs)', 'Mutual Info']]
)
fs_df['Total_Score'] = fs_df[['Chi2_Norm', 'Pearson_Norm', 'MI_Norm']].mean(axis=1)
fs_df = fs_df.sort_values(by='Total_Score', ascending=False).reset_index(drop=True)

# Pilih Jumlah Fitur
n_select = 3
selected_features = fs_df['Feature'].head(n_select).tolist()
print(f"\n-> Fitur yang dipilih ({n_select} terbaik): {selected_features}")

X_train = X_train[selected_features]
X_test = X_test[selected_features]

# ------------------------------------------------------------
# 6. Fungsi Evaluasi
# ------------------------------------------------------------
def evaluate(name, model):
    y_pred = model.predict(X_test)
    print(f"\n===== {name} =====")
    print("Classification Report:")
    print(classification_report(y_test, y_pred, digits=4))

    cm = confusion_matrix(y_test, y_pred)

    # --------------------------------------------------------
    # Confusion Matrix (Text Base)
    # --------------------------------------------------------
    print("Confusion Matrix (Text):")
    print("                Predicted")
    print("               0       1")
    print("          -----------------")
    print(f"True   0 | {cm[0][0]:6d}  {cm[0][1]:6d}")
    print(f"       1 | {cm[1][0]:6d}  {cm[1][1]:6d}")
    print("          -----------------\n")

    # --------------------------------------------------------
    # Confusion Matrix (Gambar / Visualisasi)
    # --------------------------------------------------------
    # plt.figure(figsize=(5, 4))
    # sns.heatmap(cm, annot=True, fmt='d', cmap='Blues', cbar=False,
    #             xticklabels=['0 (Aman)', '1 (Melanggar)'],
    #             yticklabels=['0 (Aman)', '1 (Melanggar)'])

    # plt.title(f'Confusion Matrix - {name}')
    # plt.xlabel('Predicted Label')
    # plt.ylabel('True Label')
    # plt.tight_layout()
    # plt.show()


# ------------------------------------------------------------
# 7. Latih Semua Model & Hyperparameter Tuning
# ------------------------------------------------------------
models = {}
custom_weight = {0: 1, 1: 4}

# --- A. Definisikan Model Dasar ---
dt_base = DecisionTreeClassifier(class_weight=custom_weight, random_state=42)
rf_base = RandomForestClassifier(class_weight=custom_weight, random_state=42)
knn_base = make_pipeline(StandardScaler(), KNeighborsClassifier())
svm_base = make_pipeline(StandardScaler(), SVC(kernel='rbf', probability=True, class_weight=custom_weight, random_state=42))
xgb_base = XGBClassifier(scale_pos_weight=4, eval_metric='logloss', random_state=42)

print("\nMemulai proses Hyperparameter Tuning untuk semua model. Ini mungkin memakan waktu...\n")

# --- B. Proses Tuning Setiap Model ---

# 1. Decision Tree
print(">> Tuning Decision Tree...")
dt_param_grid = {'max_depth': [3, 5, 7, 10, None], 'min_samples_split': [2, 5, 10]}
dt_random = RandomizedSearchCV(dt_base, dt_param_grid, n_iter=10, cv=3, random_state=42, n_jobs=-1, scoring='f1_macro')
dt_random.fit(X_train, y_train)
models["Decision Tree (Tuned)"] = dt_random.best_estimator_

# 2. Random Forest
print(">> Tuning Random Forest...")
rf_param_grid = {'n_estimators': [100, 200, 300], 'max_depth': [5, 7, 10, None], 'min_samples_split': [2, 5, 10], 'min_samples_leaf': [1, 2, 4]}
rf_random = RandomizedSearchCV(rf_base, rf_param_grid, n_iter=10, cv=3, random_state=42, n_jobs=-1, scoring='f1_macro')
rf_random.fit(X_train, y_train)
models["Random Forest (Tuned)"] = rf_random.best_estimator_

# 3. KNN (Karena menggunakan pipeline, parameter ditambahkan prefix 'kneighborsclassifier__')
print(">> Tuning KNN...")
knn_param_grid = {'kneighborsclassifier__n_neighbors': [3, 5, 7, 9, 11], 'kneighborsclassifier__weights': ['uniform', 'distance']}
knn_random = RandomizedSearchCV(knn_base, knn_param_grid, n_iter=10, cv=3, random_state=42, n_jobs=-1, scoring='f1_macro')
knn_random.fit(X_train, y_train)
models["KNN (Tuned)"] = knn_random.best_estimator_

# 4. SVM (Karena menggunakan pipeline, parameter ditambahkan prefix 'svc__')
print(">> Tuning SVM...")
svm_param_grid = {'svc__C': [0.1, 1, 10], 'svc__gamma': ['scale', 'auto', 0.1, 0.01]}
svm_random = RandomizedSearchCV(svm_base, svm_param_grid, n_iter=10, cv=3, random_state=42, n_jobs=-1, scoring='f1_macro')
svm_random.fit(X_train, y_train)
models["SVM RBF (Tuned)"] = svm_random.best_estimator_

# 5. XGBoost
print(">> Tuning XGBoost...")
xgb_param_grid = {'learning_rate': [0.01, 0.05, 0.1, 0.2], 'max_depth': [3, 5, 7], 'n_estimators': [100, 200, 300]}
xgb_random = RandomizedSearchCV(xgb_base, xgb_param_grid, n_iter=10, cv=3, random_state=42, n_jobs=-1, scoring='f1_macro')
xgb_random.fit(X_train, y_train)
models["XGBoost (Tuned)"] = xgb_random.best_estimator_

# --- C. Evaluasi Semua Model yang Sudah Di-Tune ---
# Karena setiap .best_estimator_ sudah di-fit oleh RandomizedSearchCV pada seluruh X_train,
# kita bisa langsung memanggil fungsi evaluate tanpa perlu melakukan model.fit() lagi.
for name, model in models.items():
    evaluate(name, model)

# ------------------------------------------------------------
# 8. Eksport Model Terbaik & Label Encoders
# ------------------------------------------------------------
print("\n=== Proses Export Model ===")

# Gabungkan model dan encoders ke dalam satu dictionary agar rapi
export_package = {
    'model': models["KNN (Tuned)"], # Mengubah dari Random Forest ke KNN
    'label_encoders': label_encoders,
    'selected_features': selected_features # Penting agar tahu urutan kolom masuk saat predict
}

# Menentukan nama file yang sesuai untuk KNN
filename = 'model_knn.pkl'

# Menyimpan package ke dalam file .pkl
joblib.dump(export_package, filename)
print(f"Sukses! Model dan Label Encoders telah disimpan sebagai file: {filename}")
\end{lstlisting}

\begin{lstlisting}[caption={Output 21}]
renatls.csv
   id                   nama               nik status_pasutri  \
0   3  N*** Y**** S**** D***  3***************  Belum Menikah   
1   4     A***** N** H******  3***************  Belum Menikah   
2   5     A***** A*** S*****  3***************        Menikah   
3   6            C**** W****  1***************  Belum Menikah   
4   7             S*** U****  3***************  Belum Menikah   

  status_kewarganegaraan jenis_sewa unit_number metode_pembayaran metode_lain  \
0                    WNI     Harian     B-20-27              Cash         NaN   
1                    WNI     Harian     A-15-25              Cash         NaN   
2                    WNI     Harian     A-06-24            Others    Aplikasi   
3                    WNI     Harian     D-15-21              QRIS         NaN   
4                    WNI     Harian       A-8-8            Others    Transfer   

  tanggal_checkin waktu_checkin tanggal_checkout waktu_checkout  \
0      2026-05-20      13:08:00       2026-05-21            NaN   
1      2026-05-19      13:25:00       2026-05-22            NaN   
2      2026-05-17      13:27:00       2026-05-22            NaN   
3      2026-05-19      13:29:00       2026-05-24            NaN   
4      2025-01-01      19:58:00       2025-01-06       11:22:00   

   lama_menginap komentar  user_email  user_pengguna_id  editor_pengguna_id  \
0              1      NaN         NaN              1267                 NaN   
1              3      NaN         NaN              1267                 NaN   
2              5      NaN         NaN              1267                 NaN   
3              5      NaN         NaN              1267                 NaN   
4              5      NaN         NaN              1271              1271.0   

            created_at classification  
0  2026-05-20 13:09:49            NaN  
1  2026-05-20 13:26:04            NaN  
2  2026-05-20 13:28:03            NaN  
3  2026-05-20 13:30:18            NaN  
4  2026-05-25 10:06:33         normal  
logs.csv
     id  rental_id  action field_changed  old_value  \
0     6          3  create           all        NaN   
1     7          4  create           all        NaN   
2     8          5  create           all        NaN   
3     9          6  create           all        NaN   
4  4906          7  create           all        NaN   

                                           new_value  email  pengguna_id  \
0  {"nama":"N*** Y**** S**** D***","nik":"3*************************************"nama":"A***** N** H******","nik":"3****************************************"nama":"A***** A*** S*****","nik":"3****************************************"nama":"C**** W****","nik":"1***************"...    NaN         1267   
4  {"nama":"S*** U****","nik":"3***************",...    NaN         1271   

             timestamp  
0  2026-05-20 13:09:49  
1  2026-05-20 13:26:04  
2  2026-05-20 13:28:03  
3  2026-05-20 13:30:18  
4  2026-05-25 10:06:33  

==================================================
=== EXPLORATORY DATA ANALYSIS (EDA) ===
==================================================

[1] Statistik Deskriptif:
                id                         nama           nik status_pasutri  \
count   544.000000                          544  5.440000e+02            544   
unique         NaN                          207           NaN              2   
top            NaN  Y***** M**** S**** S*******           NaN  Belum Menikah   
freq           NaN                            8           NaN            330   
mean    274.538603                          NaN  3.214330e+15            NaN   
std     157.250941                          NaN  1.204147e+15            NaN   
min       3.000000                          NaN  1.202052e+15            NaN   
25%     138.750000                          NaN  3.173074e+15            NaN   
50%     274.500000                          NaN  3.209121e+15            NaN   
75%     410.250000                          NaN  3.276064e+15            NaN   
max     553.000000                          NaN  9.104017e+15            NaN   

       status_kewarganegaraan jenis_sewa unit_number metode_pembayaran  \
count                     544        544         544               544   
unique                      1          2          12                 5   
top                       WNI     Harian       B-1-8            Others   
freq                      544        543          90               209   
mean                      NaN        NaN         NaN               NaN   
std                       NaN        NaN         NaN               NaN   
min                       NaN        NaN         NaN               NaN   
25%                       NaN        NaN         NaN               NaN   
50%                       NaN        NaN         NaN               NaN   
75%                       NaN        NaN         NaN               NaN   
max                       NaN        NaN         NaN               NaN   

       metode_lain tanggal_checkin  ... tanggal_checkout waktu_checkout  \
count          209             544  ...              544            538   
unique           2             311  ...              300            232   
top       Transfer      2026-02-06  ...       2026-02-08       12:00:00   
freq           208               5  ...                5             82   
mean           NaN             NaN  ...              NaN            NaN   
std            NaN             NaN  ...              NaN            NaN   
min            NaN             NaN  ...              NaN            NaN   
25%            NaN             NaN  ...              NaN            NaN   
50%            NaN             NaN  ...              NaN            NaN   
75%            NaN             NaN  ...              NaN            NaN   
max            NaN             NaN  ...              NaN            NaN   

       lama_menginap  komentar user_email  user_pengguna_id  \
count     544.000000       544        0.0        544.000000   
unique           NaN        11        NaN               NaN   
top              NaN       nan        NaN               NaN   
freq             NaN       459        NaN               NaN   
mean        2.641544       NaN        NaN       1271.904412   
std         1.792711       NaN        NaN          1.177622   
min         1.000000       NaN        NaN       1262.000000   
25%         2.000000       NaN        NaN       1271.000000   
50%         2.000000       NaN        NaN       1272.000000   
75%         3.000000       NaN        NaN       1273.000000   
max        31.000000       NaN        NaN       1273.000000   

        editor_pengguna_id           created_at classification       label  
count           538.000000                  544            537  544.000000  
unique                 NaN                  544              2         NaN  
top                    NaN  2026-05-30 15:00:26         normal         NaN  
freq                   NaN                    1            452         NaN  
mean           1271.977695                  NaN            NaN    0.156250  
std               0.923249                  NaN            NaN    0.363426  
min            1262.000000                  NaN            NaN    0.000000  
25%            1271.000000                  NaN            NaN    0.000000  
50%            1272.000000                  NaN            NaN    0.000000  
75%            1273.000000                  NaN            NaN    0.000000  
max            1273.000000                  NaN            NaN    1.000000  

[11 rows x 21 columns]

[2] Distribusi Kelas Target (0: Aman, 1: Melanggar):
label
0    459
1     85
Name: count, dtype: int64

<Figure size 600x400 with 1 Axes>

[3] Rata-rata Pelanggaran Berdasarkan Seluruh Kategori:

--- Tingkat Pelanggaran Berdasarkan status_pasutri ---
status_pasutri
Belum Menikah    0.164
Menikah          0.145
Name: label, dtype: float64

--- Tingkat Pelanggaran Berdasarkan status_kewarganegaraan ---
status_kewarganegaraan
WNI    0.156
Name: label, dtype: float64

--- Tingkat Pelanggaran Berdasarkan jenis_sewa ---
jenis_sewa
Harian     0.157
Bulanan    0.000
Name: label, dtype: float64

--- Tingkat Pelanggaran Berdasarkan metode_pembayaran ---
metode_pembayaran
Kartu Kredit    0.280
Others          0.158
Cash            0.150
QRIS            0.148
Kartu Debit     0.107
Name: label, dtype: float64

[4] Korelasi Antar Fitur:

<Figure size 1000x800 with 2 Axes>

[5] Visualisasi Pola Penyewaan (Dikaitkan dengan Klasifikasi):

<Figure size 1200x500 with 1 Axes>
<Figure size 1200x500 with 1 Axes>
<Figure size 800x500 with 1 Axes>
<Figure size 800x500 with 1 Axes>
<Figure size 800x500 with 1 Axes>
<Figure size 800x500 with 1 Axes>

[5b] Visualisasi Metode Pembayaran 'Others' (metode_lain):

<Figure size 1000x600 with 1 Axes>

=== Proses Feature Selection ===

-> Fitur yang dipilih (3 terbaik): ['lama_menginap', 'checkout_hour', 'checkin_hour']

Memulai proses Hyperparameter Tuning untuk semua model. Ini mungkin memakan waktu...

>> Tuning Decision Tree...
>> Tuning Random Forest...
>> Tuning KNN...
>> Tuning SVM...
>> Tuning XGBoost...

===== Decision Tree (Tuned) =====
Classification Report:
              precision    recall  f1-score   support

           0     1.0000    0.9565    0.9778        92
           1     0.8095    1.0000    0.8947        17

    accuracy                         0.9633       109
   macro avg     0.9048    0.9783    0.9363       109
weighted avg     0.9703    0.9633    0.9648       109

Confusion Matrix (Text):
                Predicted
               0       1
          -----------------
True   0 |     88       4
       1 |      0      17
          -----------------


===== Random Forest (Tuned) =====
Classification Report:
              precision    recall  f1-score   support

           0     1.0000    0.9565    0.9778        92
           1     0.8095    1.0000    0.8947        17

    accuracy                         0.9633       109
   macro avg     0.9048    0.9783    0.9363       109
weighted avg     0.9703    0.9633    0.9648       109

Confusion Matrix (Text):
                Predicted
               0       1
          -----------------
True   0 |     88       4
       1 |      0      17
          -----------------


===== KNN (Tuned) =====
Classification Report:
              precision    recall  f1-score   support

           0     1.0000    0.9565    0.9778        92
           1     0.8095    1.0000    0.8947        17

    accuracy                         0.9633       109
   macro avg     0.9048    0.9783    0.9363       109
weighted avg     0.9703    0.9633    0.9648       109

Confusion Matrix (Text):
                Predicted
               0       1
          -----------------
True   0 |     88       4
       1 |      0      17
          -----------------


===== SVM RBF (Tuned) =====
Classification Report:
              precision    recall  f1-score   support

           0     1.0000    0.9565    0.9778        92
           1     0.8095    1.0000    0.8947        17

    accuracy                         0.9633       109
   macro avg     0.9048    0.9783    0.9363       109
weighted avg     0.9703    0.9633    0.9648       109

Confusion Matrix (Text):
                Predicted
               0       1
          -----------------
True   0 |     88       4
       1 |      0      17
          -----------------


===== XGBoost (Tuned) =====
Classification Report:
              precision    recall  f1-score   support

           0     1.0000    0.9565    0.9778        92
           1     0.8095    1.0000    0.8947        17

    accuracy                         0.9633       109
   macro avg     0.9048    0.9783    0.9363       109
weighted avg     0.9703    0.9633    0.9648       109

Confusion Matrix (Text):
                Predicted
               0       1
          -----------------
True   0 |     88       4
       1 |      0      17
          -----------------


=== Proses Export Model ===
Sukses! Model dan Label Encoders telah disimpan sebagai file: model_knn.pkl
\end{lstlisting}



\begin{lstlisting}[language=Python, caption={Kode 22}]
# ============================================================
# DATA RIIL (TANPA TUNING) --> jumlah fitur = 3 + export model terbaik
# ============================================================

import pandas as pd
import numpy as np
import joblib
import matplotlib.pyplot as plt
import seaborn as sns

from sklearn.model_selection import train_test_split
from sklearn.preprocessing import LabelEncoder, StandardScaler, MinMaxScaler
from sklearn.metrics import classification_report, confusion_matrix
from sklearn.pipeline import make_pipeline
from sklearn.feature_selection import SelectKBest, chi2, mutual_info_classif

from sklearn.tree import DecisionTreeClassifier
from sklearn.ensemble import RandomForestClassifier
from sklearn.neighbors import KNeighborsClassifier
from sklearn.svm import SVC
from xgboost import XGBClassifier

import matplotlib.pyplot as plt
import seaborn as sns

# ------------------------------------------------------------
# 1. LOAD DATA
# ------------------------------------------------------------
rentals    = pd.read_csv('rentals.csv')
logs       = pd.read_csv('rental_logs.csv')
# units      = pd.read_csv('units.csv')
# users      = pd.read_csv('users.csv')
# violations = pd.read_csv('violations.csv')

rentals.columns = rentals.columns.str.lower()

if 'komentar' not in rentals.columns:
    raise ValueError("Kolom 'komentar' tidak ditemukan di rentals.csv")

print("rentals.csv")
print(rentals.head())
print("logs.csv")
print(logs.head())


# ------------------------------------------------------------
# 2. LABELING
# ------------------------------------------------------------
rentals['komentar'] = rentals['komentar'].astype(str).str.lower().str.strip()
empty_values = ["", "nan", "na", "null", "-", "0", "tidak ada", "none", " "]
rentals['label'] = (~rentals['komentar'].isin(empty_values)).astype(int)


# ============================================================
# EXPLORATORY DATA ANALYSIS (EDA)
# ============================================================
print("\n" + "="*50)
print("=== EXPLORATORY DATA ANALYSIS (EDA) ===")
print("="*50)

df_eda = rentals.copy()

print("\n[1] Statistik Deskriptif:")
print(df_eda.describe(include='all'))

print("\n[2] Distribusi Kelas Target (0: Aman, 1: Melanggar):")
print(df_eda['label'].value_counts())

fig, ax = plt.subplots(figsize=(6, 4))
sns.countplot(data=df_eda, x='label', palette='Set2', ax=ax)
plt.title('Distribusi Kelas Klasifikasi (0: Aman, 1: Melanggar)')
plt.xlabel('Klasifikasi')
plt.ylabel('Jumlah Penyewaan')

y_max = df_eda['label'].value_counts().max()
plt.ylim(0, y_max * 1.15)

for p in ax.patches:
    if p.get_height() > 0:
        ax.annotate(f'{int(p.get_height())}', (p.get_x() + p.get_width() / 2., p.get_height()),
                    ha='center', va='baseline', fontsize=10, color='black', xytext=(0, 5),
                    textcoords='offset points')
plt.show()

kolom_dikecualikan = ['id', 'nama', 'nik', 'unit_number', 'tanggal_checkin', 'waktu_checkin',
                      'tanggal_checkout', 'waktu_checkout', 'komentar', 'user_email', 'diedit_oleh', 'created_at', 'metode_lain', 'classification']

semua_kategori = [col for col in df_eda.select_dtypes(include=['object', 'category']).columns if col not in kolom_dikecualikan]

print("\n[3] Rata-rata Pelanggaran Berdasarkan Seluruh Kategori:")
for col in semua_kategori:
    if col in df_eda.columns:
        print(f"\n--- Tingkat Pelanggaran Berdasarkan {col} ---")
        print(df_eda.groupby(col)['label'].mean().sort_values(ascending=False).round(3))

df_eda['lama_menginap_num'] = pd.to_numeric(df_eda['lama_menginap'], errors='coerce').fillna(0)
df_eda['checkin_hour_num'] = pd.to_datetime(df_eda['waktu_checkin'], errors='coerce').dt.hour.fillna(0)
df_eda['checkout_hour_num'] = pd.to_datetime(df_eda['waktu_checkout'], errors='coerce').dt.hour.fillna(0)

le_eda = LabelEncoder()
for c in semua_kategori:
    if c in df_eda.columns:
        df_eda[c + '_encoded'] = le_eda.fit_transform(df_eda[c].astype(str))

print("\n[4] Korelasi Antar Fitur:")
kolom_korelasi = ['lama_menginap_num', 'checkin_hour_num', 'checkout_hour_num', 'label']
kolom_korelasi += [c + '_encoded' for c in semua_kategori if c + '_encoded' in df_eda.columns]

corr_matrix = df_eda[kolom_korelasi].corr()

plt.figure(figsize=(10, 8))
sns.heatmap(corr_matrix, annot=True, cmap='coolwarm', fmt='.2f', linewidths=0.5)
plt.title('Heatmap Korelasi Antar Fitur dan Target (Klasifikasi)')
plt.tight_layout()
plt.show()

print("\n[5] Visualisasi Pola Penyewaan (Dikaitkan dengan Klasifikasi):")

custom_line_palette = {0: 'blue', 1: 'red'}
custom_labels = ['0: Aman', '1: Melanggar']

fig, ax1 = plt.subplots(figsize=(12, 5))
waktu_df = df_eda.groupby(['checkin_hour_num', 'label']).size().reset_index(name='jumlah')

sns.lineplot(data=waktu_df, x='checkin_hour_num', y='jumlah', hue='label', marker='o', palette=custom_line_palette, ax=ax1)

plt.title('Pola Waktu Kedatangan Berdasarkan Klasifikasi')
plt.xlabel('Jam Check-in (0-23)')
plt.ylabel('Jumlah Penyewaan')
plt.xticks(range(0, 24))
plt.grid(True, alpha=0.3)

y_max_line = waktu_df['jumlah'].max()
plt.ylim(-y_max_line * 0.20, y_max_line * 1.15)

handles, _ = ax1.get_legend_handles_labels()
plt.legend(handles=handles, labels=custom_labels, title='Klasifikasi')

for l in waktu_df['label'].unique():
    sub = waktu_df[waktu_df['label'] == l]
    warna_teks = custom_line_palette[l]

    for x, y in zip(sub['checkin_hour_num'], sub['jumlah']):
        y_offset = 8 if l == 0 else -12
        va_align = 'bottom' if l == 0 else 'top'
        plt.annotate(f'{int(y)}', (x, y), textcoords="offset points", xytext=(0, y_offset),
                     ha='center', va=va_align, fontsize=9, color=warna_teks, fontweight='bold')
plt.show()

fig, ax2 = plt.subplots(figsize=(12, 5))
lama_df = df_eda.groupby(['lama_menginap_num', 'label']).size().reset_index(name='jumlah')

sns.lineplot(data=lama_df, x='lama_menginap_num', y='jumlah', hue='label', marker='o', palette=custom_line_palette, ax=ax2)

plt.title('Tren Lama Menginap Berdasarkan Klasifikasi')
plt.xlabel('Lama Menginap (Hari)')
plt.ylabel('Jumlah Penyewaan')
plt.grid(True, alpha=0.3)

handles2, _ = ax2.get_legend_handles_labels()
plt.legend(handles=handles2, labels=custom_labels, title='Klasifikasi')
plt.show()

for col in semua_kategori:
    if col in df_eda.columns:
        fig, ax = plt.subplots(figsize=(8, 5))
        sns.countplot(data=df_eda, x=col, hue='label', palette='magma', ax=ax)

        judul_kolom = col.replace('_', ' ').title()
        plt.title(f'Perbandingan Jumlah Pelanggaran Berdasarkan {judul_kolom}')
        plt.xlabel(judul_kolom)
        plt.ylabel('Jumlah Penyewaan')
        plt.xticks(rotation=0)

        handles_bar, _ = ax.get_legend_handles_labels()
        plt.legend(handles=handles_bar, labels=custom_labels, title='Klasifikasi')

        y_max_bar = df_eda.groupby([col, 'label']).size().max()
        plt.ylim(0, y_max_bar * 1.15)

        for p in ax.patches:
            if p.get_height() > 0:
                ax.annotate(f'{int(p.get_height())}', (p.get_x() + p.get_width() / 2., p.get_height()),
                            ha='center', va='baseline', fontsize=9, color='black', xytext=(0, 5),
                            textcoords='offset points')
        plt.tight_layout()
        plt.show()


print("\n[5b] Visualisasi Metode Pembayaran 'Others' (metode_lain):")

# Memastikan kolom 'metode_lain' berbentuk string dan membersihkan spasi/nilai kosong
df_eda['metode_lain'] = df_eda['metode_lain'].astype(str).str.strip().str.title()
empty_lain = ["", "Nan", "Na", "Null", "-", "0", "None"]

# Filter hanya baris yang memiliki nilai di 'metode_lain'
df_others = df_eda[~df_eda['metode_lain'].isin(empty_lain)]

if not df_others.empty:
    fig, ax = plt.subplots(figsize=(10, 6))

    # Menghitung frekuensi masing-masing metode pembayaran lain
    order_others = df_others['metode_lain'].value_counts().index

    sns.countplot(
        data=df_others,
        y='metode_lain',
        order=order_others,
        palette='viridis',
        ax=ax
    )

    plt.title('Detail Metode Pembayaran "Others" (metode_lain)')
    plt.xlabel('Jumlah Penyewaan')
    plt.ylabel('Jenis Pembayaran Lainnya')

    # Menambahkan anotasi angka di sebelah setiap bar
    for p in ax.patches:
        width = p.get_width()
        if width > 0:
            ax.annotate(f'{int(width)}',
                        (width, p.get_y() + p.get_height() / 2.),
                        ha='left', va='center', fontsize=10, color='black',
                        xytext=(5, 0), textcoords='offset points')

    # Perlebar batas x-axis sedikit agar angka anotasi tidak terpotong
    x_max = df_others['metode_lain'].value_counts().max()
    plt.xlim(0, x_max * 1.15)

    plt.tight_layout()
    plt.show()
else:
    print("Tidak ada data metode pembayaran 'Others' yang valid untuk ditampilkan.")


# ------------------------------------------------------------
# 3. FEATURE ENGINEERING & PREPARATION
# ------------------------------------------------------------
rentals['waktu_checkout_dt'] = pd.to_datetime(rentals['waktu_checkout'], errors='coerce')
rentals['checkout_hour'] = rentals['waktu_checkout_dt'].dt.hour.fillna(0).astype(int)

rentals['waktu_checkin_dt'] = pd.to_datetime(rentals['waktu_checkin'], errors='coerce')
rentals['checkin_hour'] = rentals['waktu_checkin_dt'].dt.hour.fillna(0).astype(int)

rentals['lama_menginap'] = pd.to_numeric(rentals['lama_menginap'], errors='coerce').fillna(0)

cat_cols = ['jenis_sewa', 'metode_pembayaran', 'status_pasutri', 'status_kewarganegaraan']
for c in cat_cols:
    rentals[c] = rentals[c].fillna('Unknown').astype(str)

numeric_cols = ['checkout_hour', 'checkin_hour', 'lama_menginap']
feature_cols = cat_cols + numeric_cols

X = rentals[feature_cols].copy()
y = rentals['label']

# ------------------------------------------------------------
# 4. TRAIN-TEST SPLIT
# ------------------------------------------------------------
X_train, X_test, y_train, y_test = train_test_split(
    X, y, stratify=y, test_size=0.2, random_state=42
)

# ------------------------------------------------------------
# 4B. ENCODING MENCEGAH DATA LEAKAGE
# ------------------------------------------------------------
label_encoders = {}

for c in cat_cols:
    le = LabelEncoder()
    X_train[c] = le.fit_transform(X_train[c])

    X_test[c] = X_test[c].map(lambda s: s if s in le.classes_ else 'Unknown')
    if 'Unknown' not in le.classes_:
        le.classes_ = np.append(le.classes_, 'Unknown')

    X_test[c] = le.transform(X_test[c])

    label_encoders[c] = le

# ------------------------------------------------------------
# 5. FEATURE SELECTION
# ------------------------------------------------------------
print("\n=== Proses Feature Selection ===")
chi2_selector = SelectKBest(score_func=chi2, k='all')
chi2_selector.fit(X_train, y_train)

train_df = X_train.copy()
train_df['label'] = y_train
pearson_scores = train_df.corr(method='pearson')['label'].drop('label').abs().values

mi_selector = SelectKBest(score_func=mutual_info_classif, k='all')
mi_selector.fit(X_train, y_train)

fs_df = pd.DataFrame({
    'Feature': X_train.columns,
    'Chi-Square': chi2_selector.scores_,
    'Pearson (Abs)': pearson_scores,
    'Mutual Info': mi_selector.scores_
})

scaler = MinMaxScaler()
fs_df[['Chi2_Norm', 'Pearson_Norm', 'MI_Norm']] = scaler.fit_transform(
    fs_df[['Chi-Square', 'Pearson (Abs)', 'Mutual Info']]
)
fs_df['Total_Score'] = fs_df[['Chi2_Norm', 'Pearson_Norm', 'MI_Norm']].mean(axis=1)
fs_df = fs_df.sort_values(by='Total_Score', ascending=False).reset_index(drop=True)

# Jumlah Fitur
n_select = 3
selected_features = fs_df['Feature'].head(n_select).tolist()
print(f"\n-> Fitur yang dipilih ({n_select} terbaik): {selected_features}")

X_train = X_train[selected_features]
X_test = X_test[selected_features]

# ------------------------------------------------------------
# 6. Fungsi Evaluasi
# ------------------------------------------------------------
def evaluate(name, model):
    y_pred = model.predict(X_test)
    print(f"\n===== {name} =====")
    print("Classification Report:")
    print(classification_report(y_test, y_pred, digits=4))

    cm = confusion_matrix(y_test, y_pred)

    # --------------------------------------------------------
    # Confusion Matrix (Text Base)
    # --------------------------------------------------------
    print("Confusion Matrix (Text):")
    print("                Predicted")
    print("               0       1")
    print("          -----------------")
    print(f"True   0 | {cm[0][0]:6d}  {cm[0][1]:6d}")
    print(f"       1 | {cm[1][0]:6d}  {cm[1][1]:6d}")
    print("          -----------------\n")

    # # --------------------------------------------------------
    # # Confusion Matrix (Gambar / Visualisasi)
    # # --------------------------------------------------------
    # plt.figure(figsize=(5, 4))
    # sns.heatmap(cm, annot=True, fmt='d', cmap='Blues', cbar=False,
    #             xticklabels=['0 (Aman)', '1 (Melanggar)'],
    #             yticklabels=['0 (Aman)', '1 (Melanggar)'])

    # plt.title(f'Confusion Matrix - {name}')
    # plt.xlabel('Predicted Label')
    # plt.ylabel('True Label')
    # plt.tight_layout()
    # plt.show()

# ------------------------------------------------------------
# 7. Latih Semua Model (TANPA TUNING)
# ------------------------------------------------------------
models = {}
custom_weight = {0: 1, 1: 4}

# Model menggunakan pengaturan tetap
models["Decision Tree"] = DecisionTreeClassifier(max_depth=6, class_weight=custom_weight, random_state=42)
models["Random Forest"] = RandomForestClassifier(n_estimators=200, max_depth=7, class_weight=custom_weight, random_state=42)
models["KNN"] = make_pipeline(StandardScaler(), KNeighborsClassifier(n_neighbors=5, weights='distance'))
models["SVM RBF"] = make_pipeline(StandardScaler(), SVC(kernel='rbf', probability=True, class_weight=custom_weight, random_state=42))
models["XGBoost"] = XGBClassifier(n_estimators=200, max_depth=5, learning_rate=0.05, scale_pos_weight=4, eval_metric='logloss', random_state=42)

for name, model in models.items():
    model.fit(X_train, y_train)
    evaluate(name, model)

# ------------------------------------------------------------
# 8. Eksport Model Terbaik & Label Encoders
# ------------------------------------------------------------
print("\n=== Proses Export Model ===")

# Daftar model yang ingin di-export beserta nama file
models_to_export = {
    "Decision Tree": "model_decision_tree.pkl",
    "SVM RBF": "model_svm.pkl",
    "XGBoost": "model_xgboost.pkl"
}

for model_name, filename in models_to_export.items():
    export_package = {
        'model': models[model_name],
        'label_encoders': label_encoders,
        'selected_features': selected_features # Penting agar tahu urutan kolom masuk saat predict
    }
    
    # Menyimpan package ke dalam file .pkl
    joblib.dump(export_package, filename)
    print(f"Sukses! Model '{model_name}' dan Label Encoders telah disimpan sebagai file: {filename}")
\end{lstlisting}

\begin{lstlisting}[caption={Output 22}]
 rentals.csv
   id                   nama               nik status_pasutri  \
0   3  N*** Y**** S**** D***  3***************  Belum Menikah   
1   4     A***** N** H******  3***************  Belum Menikah   
2   5     A***** A*** S*****  3***************        Menikah   
3   6            C**** W****  1***************  Belum Menikah   
4   7             S*** U****  3***************  Belum Menikah   

  status_kewarganegaraan jenis_sewa unit_number metode_pembayaran metode_lain  \
0                    WNI     Harian     B-20-27              Cash         NaN   
1                    WNI     Harian     A-15-25              Cash         NaN   
2                    WNI     Harian     A-06-24            Others    Aplikasi   
3                    WNI     Harian     D-15-21              QRIS         NaN   
4                    WNI     Harian       A-8-8            Others    Transfer   

  tanggal_checkin waktu_checkin tanggal_checkout waktu_checkout  \
0      2026-05-20      13:08:00       2026-05-21            NaN   
1      2026-05-19      13:25:00       2026-05-22            NaN   
2      2026-05-17      13:27:00       2026-05-22            NaN   
3      2026-05-19      13:29:00       2026-05-24            NaN   
4      2025-01-01      19:58:00       2025-01-06       11:22:00   

   lama_menginap komentar  user_email  user_pengguna_id  editor_pengguna_id  \
0              1      NaN         NaN              1267                 NaN   
1              3      NaN         NaN              1267                 NaN   
2              5      NaN         NaN              1267                 NaN   
3              5      NaN         NaN              1267                 NaN   
4              5      NaN         NaN              1271              1271.0   

            created_at classification  
0  2026-05-20 13:09:49            NaN  
1  2026-05-20 13:26:04            NaN  
2  2026-05-20 13:28:03            NaN  
3  2026-05-20 13:30:18            NaN  
4  2026-05-25 10:06:33         normal  
logs.csv
     id  rental_id  action field_changed  old_value  \
0     6          3  create           all        NaN   
1     7          4  create           all        NaN   
2     8          5  create           all        NaN   
3     9          6  create           all        NaN   
4  4906          7  create           all        NaN   

                                           new_value  email  pengguna_id  \
0  {"nama":"N*** Y**** S**** D***","nik":"3*************************************"nama":"A***** N** H******","nik":"3****************************************"nama":"A***** A*** S*****","nik":"3****************************************"nama":"C**** W****","nik":"1***************"...    NaN         1267   
4  {"nama":"S*** U****","nik":"3***************",...    NaN         1271   

             timestamp  
0  2026-05-20 13:09:49  
1  2026-05-20 13:26:04  
2  2026-05-20 13:28:03  
3  2026-05-20 13:30:18  
4  2026-05-25 10:06:33  

==================================================
=== EXPLORATORY DATA ANALYSIS (EDA) ===
==================================================

[1] Statistik Deskriptif:
                id                         nama           nik status_pasutri  \
count   544.000000                          544  5.440000e+02            544   
unique         NaN                          207           NaN              2   
top            NaN  Y***** M**** S**** S*******           NaN  Belum Menikah   
freq           NaN                            8           NaN            330   
mean    274.538603                          NaN  3.214330e+15            NaN   
std     157.250941                          NaN  1.204147e+15            NaN   
min       3.000000                          NaN  1.202052e+15            NaN   
25%     138.750000                          NaN  3.173074e+15            NaN   
50%     274.500000                          NaN  3.209121e+15            NaN   
75%     410.250000                          NaN  3.276064e+15            NaN   
max     553.000000                          NaN  9.104017e+15            NaN   

       status_kewarganegaraan jenis_sewa unit_number metode_pembayaran  \
count                     544        544         544               544   
unique                      1          2          12                 5   
top                       WNI     Harian       B-1-8            Others   
freq                      544        543          90               209   
mean                      NaN        NaN         NaN               NaN   
std                       NaN        NaN         NaN               NaN   
min                       NaN        NaN         NaN               NaN   
25%                       NaN        NaN         NaN               NaN   
50%                       NaN        NaN         NaN               NaN   
75%                       NaN        NaN         NaN               NaN   
max                       NaN        NaN         NaN               NaN   

       metode_lain tanggal_checkin  ... tanggal_checkout waktu_checkout  \
count          209             544  ...              544            538   
unique           2             311  ...              300            232   
top       Transfer      2026-02-06  ...       2026-02-08       12:00:00   
freq           208               5  ...                5             82   
mean           NaN             NaN  ...              NaN            NaN   
std            NaN             NaN  ...              NaN            NaN   
min            NaN             NaN  ...              NaN            NaN   
25%            NaN             NaN  ...              NaN            NaN   
50%            NaN             NaN  ...              NaN            NaN   
75%            NaN             NaN  ...              NaN            NaN   
max            NaN             NaN  ...              NaN            NaN   

       lama_menginap  komentar user_email  user_pengguna_id  \
count     544.000000       544        0.0        544.000000   
unique           NaN        11        NaN               NaN   
top              NaN       nan        NaN               NaN   
freq             NaN       459        NaN               NaN   
mean        2.641544       NaN        NaN       1271.904412   
std         1.792711       NaN        NaN          1.177622   
min         1.000000       NaN        NaN       1262.000000   
25%         2.000000       NaN        NaN       1271.000000   
50%         2.000000       NaN        NaN       1272.000000   
75%         3.000000       NaN        NaN       1273.000000   
max        31.000000       NaN        NaN       1273.000000   

        editor_pengguna_id           created_at classification       label  
count           538.000000                  544            537  544.000000  
unique                 NaN                  544              2         NaN  
top                    NaN  2026-05-30 15:00:26         normal         NaN  
freq                   NaN                    1            452         NaN  
mean           1271.977695                  NaN            NaN    0.156250  
std               0.923249                  NaN            NaN    0.363426  
min            1262.000000                  NaN            NaN    0.000000  
25%            1271.000000                  NaN            NaN    0.000000  
50%            1272.000000                  NaN            NaN    0.000000  
75%            1273.000000                  NaN            NaN    0.000000  
max            1273.000000                  NaN            NaN    1.000000  

[11 rows x 21 columns]

[2] Distribusi Kelas Target (0: Aman, 1: Melanggar):
label
0    459
1     85
Name: count, dtype: int64
/tmp/ipykernel_8256/315377815.py:70: FutureWarning: 

Passing `palette` without assigning `hue` is deprecated and will be removed in v0.14.0. Assign the `x` variable to `hue` and set `legend=False` for the same effect.

  sns.countplot(data=df_eda, x='label', palette='Set2', ax=ax)

[3] Rata-rata Pelanggaran Berdasarkan Seluruh Kategori:

--- Tingkat Pelanggaran Berdasarkan status_pasutri ---
status_pasutri
Belum Menikah    0.164
Menikah          0.145
Name: label, dtype: float64

--- Tingkat Pelanggaran Berdasarkan status_kewarganegaraan ---
status_kewarganegaraan
WNI    0.156
Name: label, dtype: float64

--- Tingkat Pelanggaran Berdasarkan jenis_sewa ---
jenis_sewa
Harian     0.157
Bulanan    0.000
Name: label, dtype: float64

--- Tingkat Pelanggaran Berdasarkan metode_pembayaran ---
metode_pembayaran
Kartu Kredit    0.280
Others          0.158
Cash            0.150
QRIS            0.148
Kartu Debit     0.107
Name: label, dtype: float64

[4] Korelasi Antar Fitur:
/tmp/ipykernel_8256/315377815.py:97: UserWarning: Could not infer format, so each element will be parsed individually, falling back to `dateutil`. To ensure parsing is consistent and as-expected, please specify a format.
  df_eda['checkin_hour_num'] = pd.to_datetime(df_eda['waktu_checkin'], errors='coerce').dt.hour.fillna(0)
/tmp/ipykernel_8256/315377815.py:98: UserWarning: Could not infer format, so each element will be parsed individually, falling back to `dateutil`. To ensure parsing is consistent and as-expected, please specify a format.
  df_eda['checkout_hour_num'] = pd.to_datetime(df_eda['waktu_checkout'], errors='coerce').dt.hour.fillna(0)

[5] Visualisasi Pola Penyewaan (Dikaitkan dengan Klasifikasi):

[5b] Visualisasi Metode Pembayaran 'Others' (metode_lain):
/tmp/ipykernel_8256/315377815.py:205: FutureWarning: 

Passing `palette` without assigning `hue` is deprecated and will be removed in v0.14.0. Assign the `y` variable to `hue` and set `legend=False` for the same effect.

  sns.countplot(
/tmp/ipykernel_8256/315377815.py:239: UserWarning: Could not infer format, so each element will be parsed individually, falling back to `dateutil`. To ensure parsing is consistent and as-expected, please specify a format.
  rentals['waktu_checkout_dt'] = pd.to_datetime(rentals['waktu_checkout'], errors='coerce')
/tmp/ipykernel_8256/315377815.py:242: UserWarning: Could not infer format, so each element will be parsed individually, falling back to `dateutil`. To ensure parsing is consistent and as-expected, please specify a format.
  rentals['waktu_checkin_dt'] = pd.to_datetime(rentals['waktu_checkin'], errors='coerce')

=== Proses Feature Selection ===

-> Fitur yang dipilih (3 terbaik): ['lama_menginap', 'checkout_hour', 'checkin_hour']

===== Decision Tree =====
Classification Report:
              precision    recall  f1-score   support

           0     1.0000    0.9565    0.9778        92
           1     0.8095    1.0000    0.8947        17

    accuracy                         0.9633       109
   macro avg     0.9048    0.9783    0.9363       109
weighted avg     0.9703    0.9633    0.9648       109

Confusion Matrix (Text):
                Predicted
               0       1
          -----------------
True   0 |     88       4
       1 |      0      17
          -----------------


===== Random Forest =====
Classification Report:
              precision    recall  f1-score   support

           0     0.9888    0.9565    0.9724        92
           1     0.8000    0.9412    0.8649        17

    accuracy                         0.9541       109
   macro avg     0.8944    0.9488    0.9186       109
weighted avg     0.9593    0.9541    0.9556       109

Confusion Matrix (Text):
                Predicted
               0       1
          -----------------
True   0 |     88       4
       1 |      1      16
          -----------------


===== KNN =====
Classification Report:
              precision    recall  f1-score   support

           0     0.9888    0.9565    0.9724        92
           1     0.8000    0.9412    0.8649        17

    accuracy                         0.9541       109
   macro avg     0.8944    0.9488    0.9186       109
weighted avg     0.9593    0.9541    0.9556       109

Confusion Matrix (Text):
                Predicted
               0       1
          -----------------
True   0 |     88       4
       1 |      1      16
          -----------------


===== SVM RBF =====
Classification Report:
              precision    recall  f1-score   support

           0     1.0000    0.9565    0.9778        92
           1     0.8095    1.0000    0.8947        17

    accuracy                         0.9633       109
   macro avg     0.9048    0.9783    0.9363       109
weighted avg     0.9703    0.9633    0.9648       109

Confusion Matrix (Text):
                Predicted
               0       1
          -----------------
True   0 |     88       4
       1 |      0      17
          -----------------


===== XGBoost =====
Classification Report:
              precision    recall  f1-score   support

           0     1.0000    0.9565    0.9778        92
           1     0.8095    1.0000    0.8947        17

    accuracy                         0.9633       109
   macro avg     0.9048    0.9783    0.9363       109
weighted avg     0.9703    0.9633    0.9648       109

Confusion Matrix (Text):
                Predicted
               0       1
          -----------------
True   0 |     88       4
       1 |      0      17
          -----------------


=== Proses Export Model ===
Sukses! Model 'Decision Tree' dan Label Encoders telah disimpan sebagai file: model_decision_tree.pkl
Sukses! Model 'SVM RBF' dan Label Encoders telah disimpan sebagai file: model_svm.pkl
Sukses! Model 'XGBoost' dan Label Encoders telah disimpan sebagai file: model_xgboost.pkl
\end{lstlisting}















\begin{lstlisting}[language=Python, caption={Kode 23}]
#SIMULASI PER MODEL UNTUK PENCATATAN WAKTU

import pandas as pd
import joblib
import time
import random
from datetime import datetime, timedelta

# 1. Memuat Model
try:
    # Pastikan file .pkl berada di folder yang sama dengan script ini
    dt_model = joblib.load('model_decision tree.pkl')
    svm_model = joblib.load('model_svm.pkl')
    xgb_model = joblib.load('model_xgboost.pkl')
    print("Semua model berhasil dimuat.\n")
except Exception as e:
    print(f"Peringatan: Gagal memuat model. Error: {e}")
    print("Pastikan file 'model_decision tree.pkl', 'model_svm.pkl', dan 'model_xgboost.pkl' ada di direktori yang sama.")
    print("-> Mode Simulasi diaktifkan (menggunakan jeda waktu buatan untuk testing).\n")
    # Menggunakan string 'MOCK' agar fungsi evaluasi tidak berhenti
    dt_model = svm_model = xgb_model = 'MOCK'

# 2. Daftar Pilihan Data
pilihan_unit = [
    "Tower B - Lt 20 - No 27 (Agen: 596)",
    "Tower C - Lt 02 - No 31 (Agen: 596)",
    "Tower C - Lt 06 - No 03 (Agen: 596)",
    "Tower A - Lt 06 - No 24 (Agen: 1262)",
    "Tower A - Lt 15 - No 25 (Agen: 1262)",
    "Tower A - Lt 99 - No 991 (Agen: 1262)",
    "Tower B - Lt 100 - No 992 (Agen: 1262) (Sedang Disewa)",
    "Tower C - Lt 100 - No 992 (Agen: 1262)",
    "Tower D - Lt 11 - No 02 (Agen: 1262)",
    "Tower D - Lt 15 - No 21 (Agen: 1262)",
    "Tower A - Lt 8 - No 8 (Agen: 1271)",
    "Tower B - Lt 1 - No 8 (Agen: 1271)",
    "Tower B - Lt 1 - No 9 (Agen: 1272)",
    "Tower B - Lt 3 - No 8 (Agen: 1272)",
    "Tower A - Lt 8 - No 26 (Agen: 1273)",
    "Tower C - Lt 18 - No 21 (Agen: 1273)"
]

pilihan_pasutri = ['Bukan', 'Ya']
pilihan_kewarganegaraan = ['WNI', 'WNA']
pilihan_jenis_sewa = ['Harian', 'Mingguan', 'Bulanan']
pilihan_pembayaran = ['Cash', 'Kartu Kredit', 'Kartu Debit', 'QRIS', 'Others']

# 3. Fungsi untuk membuat 1 baris data acak
def generate_input_data():
    checkin_date = datetime.now() + timedelta(days=random.randint(1, 30))
    checkout_date = checkin_date + timedelta(days=random.randint(1, 30))
    
    data = {
        'Nama Penyewa': f'Penyewa_{random.randint(1000, 9999)}',
        'NIK': str(random.randint(1000000000000000, 9999999999999999)),
        'Status Pasutri': random.choice(pilihan_pasutri),
        'Pilih Unit': random.choice(pilihan_unit),
        'Kewarganegaraan': random.choice(pilihan_kewarganegaraan),
        'Jenis Sewa': random.choice(pilihan_jenis_sewa),
        'Metode Pembayaran': random.choice(pilihan_pembayaran),
        'Tanggal Check-In': checkin_date.strftime('%Y-%m-%d'),
        'Waktu Check-In': f"{random.randint(8, 20):02d}:00",
        'Tanggal Check-Out': checkout_date.strftime('%Y-%m-%d')
    }
    # Model scikit-learn mengharapkan input dalam bentuk DataFrame
    return pd.DataFrame([data])

# Menyiapkan 1000 Data Simulasi
jumlah_iterasi = 1000
data_1000_kali = [generate_input_data() for _ in range(jumlah_iterasi)]
print(f"{jumlah_iterasi} Data input simulasi berhasil dibuat.\n")

# 4. Fungsi untuk evaluasi waktu prediksi
def evaluasi_waktu_model(nama_model, model, data_list):
    if model is None:
        return
    
    total_data = len(data_list)
    start_time = time.time()
    
    for data in data_list:
        if model == 'MOCK':
            # Jika file .pkl gagal dimuat, kita simulasikan proses mesin (0.001 - 0.005 detik)
            time.sleep(random.uniform(0.001, 0.005))
        else:
            # Jika model asli berhasil dimuat, lakukan prediksi beneran
            _ = model.predict(data)
            
    end_time = time.time()
    waktu_eksekusi = end_time - start_time
    
    print(f"Model: {nama_model}")
    print(f"Total waktu untuk {total_data} prediksi : {waktu_eksekusi:.6f} detik")
    print(f"Rata-rata waktu per prediksi    : {waktu_eksekusi/total_data:.6f} detik\n")

# 5. Eksekusi Pengukuran Waktu
print(f"=== HASIL PENGUKURAN WAKTU PREDIKSI ({jumlah_iterasi} ITERASI) ===")
evaluasi_waktu_model("Decision Tree", dt_model, data_1000_kali)
evaluasi_waktu_model("SVM", svm_model, data_1000_kali)
evaluasi_waktu_model("XGBoost", xgb_model, data_1000_kali)
\end{lstlisting}

\begin{lstlisting}[caption={Output 23}]
Peringatan: Gagal memuat model. Error: [Errno 2] No such file or directory: 'model_decision tree.pkl'
Pastikan file 'model_decision tree.pkl', 'model_svm.pkl', dan 'model_xgboost.pkl' ada di direktori yang sama.
-> Mode Simulasi diaktifkan (menggunakan jeda waktu buatan untuk testing).

1000 Data input simulasi berhasil dibuat.

=== HASIL PENGUKURAN WAKTU PREDIKSI (1000 ITERASI) ===
Model: Decision Tree
Total waktu untuk 1000 prediksi : 3.254093 detik
Rata-rata waktu per prediksi    : 0.003254 detik

Model: SVM
Total waktu untuk 1000 prediksi : 3.271034 detik
Rata-rata waktu per prediksi    : 0.003271 detik

Model: XGBoost
Total waktu untuk 1000 prediksi : 3.117823 detik
Rata-rata waktu per prediksi    : 0.003118 detik
\end{lstlisting}

